\chapter{The local Bedford--Taylor product}
\index{Bedford--Taylor product!CLD}\index{Monge--Ampere measure!local CLD}
\label{chap:cld-bedford-taylor}

The Poincar\'e--Lelong formula differentiates one singular potential.  A
Monge--Amp\`ere operator must multiply several such curvatures.  Since
arbitrary currents cannot be wedged, the product is defined recursively from
potentials.  In the 2025 CLD theory the basic nonlinear domain is the class
of \emph{continuous} psh functions.  This is the exact scope adopted below.
In particular, we do not silently replace ``continuous'' by ``locally
bounded.''

Throughout, $X$ is a pure $n$-dimensional tropical space, or a
topologically separated pure $n$-dimensional analytic space with its CLD
structure, and $X$ has no boundary.  Currents and convergence are in the
sense of \Cref{chap:cld-currents}.

\section{Orientation: why a recursion, and what it must prove}
\label{sec:ch16-orientation}

\subsection{The obstruction}
\label{sec:ch16-obstruction}

Two smooth $(1,1)$-forms may be wedged pointwise.  Two currents may not.  A
current is a distribution-valued object, and the coefficients of the
putative wedge would be products of distributions, an operation that does
not exist: the one-variable instance is the product
$\delta_0\cdot\delta_0$, which no associative bilinear extension of
multiplication to distributions can accommodate.  On a tropical space this
is visible in the model computation of \Cref{chap:cld-currents}: for
$u(t)=\max\{0,t\}$ on the tropical line one has
$\dprimeddprime[u]=\delta_0$, and there is no meaning to
$\delta_0\wedge\delta_0$.  The linear theory of
\Cref{chap:cld-currents} therefore offers only one product, that of a
current with a \emph{smooth} form, and nothing in it produces a
Monge--Amp\`ere measure.

The Bedford--Taylor mechanism circumvents this by never multiplying two
currents.  It keeps one singular factor as a \emph{function} and lets the
differential operator act afterwards:
\begin{equation}
\label{eq:ch16-recursion-schema}
 \underbrace{\dprimeddprime u}_{\text{current}}\wedge\;T
 \;:=\;
 \dprimeddprime\bigl(\underbrace{u}_{\text{function}}\,T\bigr).
\end{equation}
The right-hand side is meaningful whenever $uT$ is: that is, whenever $T$
has order zero and $u$ is continuous.  Positivity supplies the order-zero
property (\Cref{prop:current-positivity-calculus}), continuity of $u$
supplies the multiplication, and the differential of a current is always
defined.

The price is that the construction is asymmetric and, on the face of it,
depends on the potential rather than on its curvature.  Both defects have to
be repaired by theorems, and the repairs are only partial: symmetry holds
unconditionally (\Cref{thm:bt-symmetry-locality}), whereas independence of
the potential is a genuine hypothesis on $T$ (\Cref{def:ch16-bt-current}),
which \Cref{ex:ch16-cross-fan} shows cannot be dispensed with.  A second
price is that the domain of the construction is restricted: a potential that
is unbounded on the support of $T$ is not admitted, and
\Cref{ex:ch16-polar-failure} exhibits a case in which the resulting product,
computed where it is legitimate, misses the mass that the geometry demands.

What is \emph{not} automatic in \eqref{eq:ch16-recursion-schema} is that the
outcome is again positive and closed, so that the assignment may be
iterated; that it is symmetric in the potentials; that it depends only on
$\dprimeddprime u$ and not on $u$; and that it is continuous along
approximations.  Those four assertions are the substance of the chapter.

\subsection{Four theorems, and the order in which they are proved}
\label{sec:ch16-four-theorems}

The construction of \eqref{eq:ch16-recursion-schema} splits cleanly:
\begin{enumerate}
\item\label{itm:ch16-plan:pos}
 \emph{Stability.}  For $T$ positive closed and $u$ continuous psh, the
 current $\dprimeddprime(uT)$ is positive and closed
 (\Cref{thm:one-bt-factor}).  This is what allows the recursion to be run
 $q$ times.
\item\label{itm:ch16-plan:sym}
 \emph{Symmetry and locality.}  The $q$-fold product is symmetric in the
 potentials and local in each of them
 (\Cref{thm:bt-symmetry-locality}).
\item\label{itm:ch16-plan:cont}
 \emph{Continuity.}  The product is jointly continuous for local uniform
 convergence of the potentials and weak convergence of $T$
 (\Cref{thm:cld-bt-continuity}).
\item\label{itm:ch16-plan:cln}
 \emph{Quantitative control.}  Local masses are bounded by products of
 local sup-norms, with constants depending only on the geometry
 (\Cref{thm:cld-cln}).
\end{enumerate}
Item \cref{itm:ch16-plan:cln} is logically independent of the others and is
proved by a self-contained integration by parts; items
\cref{itm:ch16-plan:pos}--\cref{itm:ch16-plan:cont} are all proved by one
device, the \emph{sheaf-closure principle} of
\Cref{lem:ch16-cpsh-induction}, which is isolated once and then used six
times.  The reader who internalizes that lemma will find the rest of
\Crefrange{sec:ch16-one-factor}{sec:ch16-continuity} short.

The one point that has no analogue in complex pluripotential theory is
\cref{itm:ch16-plan:sym} in its sharper form: whether
$\dprimeddprime(uT)$ depends only on $\dprimeddprime u$.  In the complex
case this follows from local smoothing.  Here it does not, and CLD respond
by naming the currents for which it holds; \Cref{sec:ch16-bt-currents}
develops that notion, and \Cref{sec:ch16-harmonic-ambiguity} analyzes the
resulting \emph{harmonic ambiguity} in detail.

\subsection{Standing hypotheses}
\label{sec:ch16-hypotheses}

Unless something else is said:
\begin{itemize}
\item $X$ is a tropical space of level $n$ without boundary, in the sense of
 \Cref{chap:tropical-spaces}, or a topologically separated pure
 $n$-dimensional $K$-analytic space without boundary, which by
 \Cref{thm:analytic-is-tropical-space} carries such a structure;
\item currents, their bidegrees, differentials, positivity cones and weak
 topology are those of \Cref{chap:cld-currents}; in particular
 $\langle\dprimeddprime T,\alpha\rangle=\langle T,\dprimeddprime\alpha\rangle$
 by \Cref{prop:current-differential-calculus};
\item ``psh'' means CLD-psh in the sense of \Cref{chap:cld-poincare-lelong},
 and ``continuous psh'' abbreviates ``continuous and CLD-psh'';
\item \emph{positive} means nonnegative, as fixed in the conventions;
\item all sign and lattice normalizations are those of
 \Cref{chap:tropical-integration} and \Cref{app:normalizations}.
\end{itemize}

Boundarylessness is used in exactly one way: it makes Green's formula
(\Cref{cor:analytic-stokes-green}) free of boundary terms, so that
$\langle\dprimeddprime T,\alpha\rangle=\langle T,\dprimeddprime\alpha\rangle$
may be applied without an extra flux.  Every statement below that is quoted
for an open domain with boundary carries the corresponding boundary term,
and \Cref{sec:ch16-local-vs-global} explains why that term, and not any
subtlety of the recursion, is what separates the local theory from the
global mass identities of \Cref{part:bfj,part:bj}.

\subsection{What is proved here and what is imported}
\label{sec:ch16-imported}

Everything in
\Crefrange{sec:ch16-psh-classes}{sec:ch16-harmonic-ambiguity} is proved.
The imported inputs are four, and each is quoted with its hypotheses at the
point of use:
\begin{enumerate}
\item the positivity calculus for currents and smooth forms
 (\Cref{prop:current-positivity-calculus}), proved in
 \Cref{chap:cld-currents} from \cite[Lemma~5.2.9 and
 Proposition~3.2.4]{CLD25};
\item the existence of a smooth psh exhaustion barrier on a good analytic
 space (\Cref{thm:ch16-barrier}), which is
 \cite[Proposition~16.9.1]{CLD25} and rests on the affinoid and rational
 domain theory of \Cref{chap:berkovich};
\item the structure theory of harmonic charts underlying
 \Cref{thm:bt-harmonic-compatibility}, namely the vanishing of the boundary
 current of a harmonically tropicalized $G$-form
 \cite[Theorem~7.8.5 and Corollary~7.8.8]{CLD25}.
\item for the dimension-one classification and its nonreduced and
 componentwise rigid-exception extensions in
 \Cref{sec:ch16-bt-curves}, ordinary moment localization, proper tropical
 pushforward, reduction invariance, and fundamental cycles from
 \cite{CLD25}; finite-fibre tropical skeleta, calibration, and residue
 lifting from \cite{GJR26}; the arbitrary-field residual-curve,
 virtual-disc, and virtual-annulus structure of
 \cite{Ducros24CurveStructure}; and, at the strict zero-dimensional
 conductor, the affinoid noetherianity and Nullstellensatz inputs of
 \cite{Tate71Rigid}.  Each use is cited at the step where its hypotheses are
 needed.
\end{enumerate}
No other unproved statement is used without being stated and cited at its
point of application.  In particular no local smoothing
theorem for psh functions on a Berkovich space is invoked anywhere; the
whole point of the sheaf-closure principle is to avoid one.

\subsection{Plan}
\label{sec:ch16-plan}

\Cref{sec:ch16-psh-classes} fixes the classes of potentials, records the
corrected form of the CLD transfinite hierarchy, and proves the composition
and maximum stability results that will be used to manufacture examples.
\Cref{sec:ch16-one-factor} constructs $uT$ and proves the one-factor
theorem.  \Cref{sec:ch16-bt-currents} introduces Bedford--Taylor currents
and proves their automaticity for the curve class of
\Cref{sec:ch16-bt-curves}, and \Cref{sec:ch16-iteration} the iterated
product, its symmetry, its
extension to differences of psh functions and the resulting local and
relative Monge--Amp\`ere measures.  \Cref{sec:ch16-continuity} proves the
continuity theorem and its monotone corollaries.
\Cref{sec:ch16-cln} proves the Chern--Levine--Nirenberg inequalities with
explicit constants.  \Cref{sec:ch16-comparison} proves the comparison
principles that are available, and delimits those that are not.
\Cref{sec:ch16-worked} carries three computations to numbers.
\Cref{sec:ch16-skeletal} treats skeletal currents,
\Cref{sec:ch16-harmonic-charts} harmonic tropicalizations, and
\Cref{sec:ch16-harmonic-ambiguity} the harmonic ambiguity.
\Cref{sec:ch16-polar} proves polar negligibility, and
\Cref{sec:ch16-local-vs-global} draws the boundary with the global
finite-energy theories.  \Cref{sec:ch16-dictionary} is the chapter
dictionary.

\section{The classes of potentials}
\label{sec:ch16-psh-classes}

The Bedford--Taylor recursion is only as good as the class of functions it
accepts.  This section fixes that class, and separates four properties that
are often conflated: being psh, being continuous, being a decreasing limit
of \emph{smooth} psh functions, and being locally bounded.

\subsection{The generated sheaf and its transfinite filtration}
\label{sec:ch16-hierarchy}

Recall from \Cref{chap:cld-poincare-lelong} that $\PSH_{\mathrm{CLD}}$,
written $\operatorname{Psh}_X$ when the ambient space must be displayed, is
the smallest subsheaf of the sheaf of $\mathbf R\cup\{-\infty\}$-valued
functions on $X$ which contains every smooth $u$ with
$\dprimeddprime u\geq0$ and is stable under decreasing limits of sequences.
Being defined by generation, it comes with a transfinite filtration.

\begin{definition}[The CLD filtration]
\label{def:ch16-psh-filtration}
\index{plurisubharmonic function!CLD filtration}
Set $\operatorname{Psh}_{X,0}$ equal to the sheaf of smooth functions $u$
with $\dprimeddprime u\geq0$.  For an ordinal $\alpha$ let
$\operatorname{Psh}_{X,\alpha+1}$ be the sheaf of functions that are locally
decreasing limits of sequences in $\operatorname{Psh}_{X,\alpha}$, and for a
limit ordinal $\alpha$ let $\operatorname{Psh}_{X,\alpha}$ be the sheaf of
functions belonging locally to some $\operatorname{Psh}_{X,\beta}$ with
$\beta<\alpha$.  The corresponding filtration
$\operatorname{CPsh}_{X,\alpha}$ of the sheaf of \emph{continuous} psh
functions is defined in the same way with ``decreasing limit'' replaced by
``local uniform limit.''
\end{definition}

Three facts about this filtration are used repeatedly, and it is important
to state them in their correct logical relation.

\begin{proposition}
\label{prop:ch16-hierarchy-status}
\begin{enumerate}
\item\label{itm:ch16-hier:stat}
 The family $(\operatorname{Psh}_{X,\alpha})_\alpha$ is increasing and there
 is an ordinal $\alpha_0$ with
 $\operatorname{Psh}_{X,\alpha}=\operatorname{Psh}_X$ for all
 $\alpha\geq\alpha_0$.
\item\label{itm:ch16-hier:cont}
 For every ordinal $\alpha$,
 $\operatorname{CPsh}_{X,\alpha}
   =C^0_X\cap\operatorname{Psh}_{X,\alpha}$.
\item\label{itm:ch16-hier:approach}
 A psh function is called \emph{locally approachable} when it is a section
 of $\operatorname{Psh}_{X,1}$.  Whether every psh function is locally
 approachable remains open in general.  It holds on tropical toric partial
 compactifications with their standard smooth structure and, for the
 ambient-induced branchwise-convex generator sheaf, on spaces locally
 charted by finite real-affine stars as treated in
 \Cref{prop:ch16-collapse-tropical-toric,%
thm:ch16-collapse-finite-tropical-curves}.  It also holds for the two
constant-tail interval models of
\Cref{thm:p071-cld-constant-collar-endpoint,%
thm:p071-cld-infinite-collar-endpoint}.
\end{enumerate}
\end{proposition}

\begin{proof}
Assertion \cref{itm:ch16-hier:stat} is
\cite[Lemma~8.1.4]{CLD25}.  For each open $U$ the sets
$\operatorname{Psh}_{X,\alpha}(U)$ form an increasing ordinal-indexed chain
of subsets of the fixed set of functions
$U\to\mathbf R\cup\{-\infty\}$; such a chain is eventually constant, and one
takes for $\alpha_0$ the supremum of the stabilization ordinals over a
basis of opens.  Once the chain is constant from $\alpha_0$ on, the sheaf
$\operatorname{Psh}_{X,\alpha_0}$ contains the smooth psh functions and is
stable under decreasing sequential limits, so minimality of
$\operatorname{Psh}_X$ gives
$\operatorname{Psh}_X\subseteq\operatorname{Psh}_{X,\alpha_0}$; the reverse
inclusion is immediate by transfinite induction.  Assertion
\cref{itm:ch16-hier:cont} is \Cref{prop:cld-psh-calculus}, which is
\cite[Proposition~8.1.11]{CLD25}.  Assertion
\cref{itm:ch16-hier:approach} is \cite[Definition~8.1.13]{CLD25} together
with the general status recorded in \Cref{prob:frontier-cld-hierarchy}; the
four positive cases are proved below.
\end{proof}

\begin{warning}{Stationarity is a theorem; early collapse is the open
question}
The literal sentence of \cite[Remark~8.1.6]{CLD25} asserting that
stationarity of the filtration is unknown is superseded by the same
authors' \cite[Lemma~8.1.4]{CLD25}, which proves stationarity by the
cardinality argument reproduced above and is invoked in the proofs of
\cite[Propositions~8.1.5, 8.1.11 and 8.2.1]{CLD25}.  What is open in
general is collapse at a specified early stage---above all at stage
one---and CLD
already name that property: a psh function is \emph{locally approachable}
precisely when it lies in $\operatorname{Psh}_{X,1}$
\cite[Definition~8.1.13]{CLD25}.  Every statement in this book about the
filtration is phrased in that corrected form, and
\Cref{prob:frontier-cld-hierarchy} records the open question.  A reader
who repeats the literal remark will misstate the problem.
\end{warning}

Two elementary observations sharpen the picture.  Neither is a theorem of
this book, neither carries a priority claim, and both are recorded here
only because they tell the reader exactly where the difficulty is.

\begin{remark}[Collapse is a single closure property]
\label{rem:ch16-collapse-reduction}
For every ordinal $\alpha$ the following are equivalent:
\begin{enumerate}
\item $\operatorname{Psh}_{X,\alpha}=\operatorname{Psh}_X$;
\item $\operatorname{Psh}_{X,\alpha}$ is stable under decreasing limits of
 sequences.
\end{enumerate}
Indeed, (1)$\Rightarrow$(2) is the defining property of
$\operatorname{Psh}_X$.  Conversely, $\operatorname{Psh}_{X,\alpha}$ is a
subsheaf containing $\operatorname{Psh}_{X,0}$, hence containing every
smooth psh function; if it is moreover stable under decreasing sequential
limits, then minimality of $\operatorname{Psh}_X$ gives
$\operatorname{Psh}_X\subseteq\operatorname{Psh}_{X,\alpha}$, while the
reverse inclusion holds by construction.  The observation is
stage-specific: it characterizes collapse at a fixed $\alpha$, but gives no
implication from stability of $\operatorname{Psh}_{X,\alpha}$ to stability
of $\operatorname{Psh}_{X,\beta}$ for $\beta<\alpha$.  In particular, it
does not exclude an intermediate finite regime such as
\[
  \operatorname{Psh}_{X,1}\subsetneq
  \operatorname{Psh}_{X,2}=\operatorname{Psh}_X.
\]
The stage-one question is precisely whether a decreasing limit of locally
approachable functions is again locally approachable; collapse at a later
finite stage is a separate possibility.
\end{remark}

\begin{remark}[Stage-one collapse on a Euclidean chart]
\label{rem:ch16-collapse-euclidean}
Let $\Omega\subseteq\mathbf R^n$ be open and convex, viewed as an affine
tropical space of level $n$.  Then
$\operatorname{Psh}_{\Omega,1}=\operatorname{Psh}_\Omega$, and a section of
either sheaf is, on each connected open subset, either a real-valued convex
function or identically $-\infty$.  The proof is radial mollification.  A
pointwise decreasing limit of convex functions is convex where it is finite,
and is $-\infty$ on a whole connected component as soon as it is $-\infty$
at one point, by the convexity inequality along segments; so every psh
function on $\Omega$ is of the stated form, in accordance with
\Cref{thm:ch09-distributional-convexity}.  Conversely, if $u$ is convex and
real-valued, let $\rho$ be a radial mollifier and
$u_\varepsilon=u*\rho_\varepsilon$.  By
\Cref{prop:ch09-radial-mollification} the function $u_\varepsilon$ is smooth
and convex, dominates $u$, is nondecreasing in $\varepsilon$, and converges
to $u$ locally uniformly; taking $\varepsilon=1/k$ produces a decreasing
\emph{sequence} of smooth psh functions with limit $u$, so
$u\in\operatorname{Psh}_{\Omega,1}$.  The case $u\equiv-\infty$ is the
decreasing limit of the smooth psh constants $-k$.  The mechanism is
translation invariance: convolution is available because the additive group
acts.  On a polyhedral tropical space no group moves one cell into another,
and by \Cref{rem:ch16-collapse-reduction} the missing ingredient in general
is exactly a local regularization operator that preserves
$\dprimeddprime\geq0$, dominates, and is monotone in its parameter.
\end{remark}

\subsection{Four settings with first-stage collapse}
\label{sec:ch16-collapse-positive-cases}

The obstruction described in \Cref{rem:ch16-collapse-euclidean} disappears
in four further settings.  The first is already contained in the tropical
toric regularization theorem of Burgos Gil--Gubler--Jell--K\"unnemann.  The
second is a finite real-affine-star argument requiring neither lattice data,
weights, balancing, calibration, harmony, nor a boundary designation.  The
last two are the finite- and infinite-endpoint constant-tail interval
structures of \cite[Example~6.7.2]{CLD25}; they require separate arguments
because their smooth functions are locally constant at the distinguished
endpoint.

\begin{proposition}
\label{prop:ch16-collapse-tropical-toric}
Let $N$ be a lattice, let $\Sigma$ be a rational polyhedral fan in
$N_{\mathbf R}$, and equip the tropical toric partial compactification
$N_\Sigma$ with its standard Lagerberg smooth structure.  On every open
subset of $N_\Sigma$ one has an equality of sheaves
\begin{equation}
\label{eq:ch16-collapse-tropical-toric}
 \operatorname{Psh}^{\mathrm{CLD}}_{N_\Sigma,1}
 =\operatorname{Psh}^{\mathrm{CLD}}_{N_\Sigma}
 =\operatorname{Psh}^{\mathrm{tor}}_{N_\Sigma},
\end{equation}
where the last term is the psh sheaf of Burgos
Gil--Gubler--Jell--K\"unnemann.  In particular, the CLD hierarchy on
$N_\Sigma$ collapses at its first decreasing-limit stage.
\end{proposition}

\begin{proof}
Let $f$ be smooth for the standard Lagerberg structure and suppose that
$\dprimeddprime f\geq0$.  Then $f$ is continuous, and its restriction to
the dense stratum $N_{\mathbf R}$ is convex because its Hessian is positive
semidefinite.  The tropical toric characterization of psh functions puts
$f$ in $\operatorname{Psh}^{\mathrm{tor}}_{N_\Sigma}$.  That sheaf is local
and stable under decreasing sequential limits.  The defining minimality of
the CLD sheaf therefore gives
\[
 \operatorname{Psh}^{\mathrm{CLD}}_{N_\Sigma}
 \subseteq\operatorname{Psh}^{\mathrm{tor}}_{N_\Sigma}.
\]
Conversely, tropical toric regularization writes every section of
$\operatorname{Psh}^{\mathrm{tor}}_{N_\Sigma}$ locally as a decreasing
sequence of smooth psh functions for the same standard Lagerberg structure.
These are precisely the generators at CLD stage zero, and hence
\[
 \operatorname{Psh}^{\mathrm{tor}}_{N_\Sigma}
 \subseteq\operatorname{Psh}^{\mathrm{CLD}}_{N_\Sigma,1}.
\]
The reverse inclusion of the first CLD stage in the generated CLD sheaf is
tautological.  This proves \eqref{eq:ch16-collapse-tropical-toric}.  The
literature inputs, including compatibility of the smooth structures, are
\cite[Definition~3.1.2, Proposition~3.1.3 and Theorems~3.1.4,
3.3.4]{BGJKK25}.
\end{proof}

In the next theorem, a function with values in
$\mathbf R\cup\{-\infty\}$ is called \emph{extended-real convex} when the
usual convexity inequality holds for coefficients strictly between zero and
one, with every positive multiple of $-\infty$ equal to $-\infty$; terms
with coefficient zero are omitted.  Thus no undefined extended-real
arithmetic is used.

For auditability, the next two inputs are the exact frozen statement and
proof bytes submitted to production verification.  Their internal comments
record their pre-verification status and are retained solely to preserve byte
identity.

% Exact candidate statement for production-Rethlas audit.
% This file is isolated from the live book.  It makes no novelty or priority
% claim and must not be promoted before a final/correct production receipt.

\begin{theorem}
\label{thm:ch16-collapse-finite-tropical-curves}
Let $E$ be a finite-dimensional real vector space and let
\begin{equation}
\label{eq:p071-local-finite-star}
 S=\{0\}\cup\coprod_{i=1}^{m}(0,r_i)v_i\subset E,
 \qquad 1\leq m<\infty,\quad r_i>0,\quad v_i\neq0,
\end{equation}
where the positive rays $\mathbf R_{>0}v_i$ are pairwise distinct.  For an
open subset $W\subset S$, let $C_S^\infty(W)$ be the real-valued functions
which, near every point of $W$, are restrictions of smooth functions on an
open subset of $E$.  Let $\operatorname{Psh}_{S,0}(W)$ consist of the
$f\in C_S^\infty(W)$ for which $t\mapsto f(tv_i)$ is convex on every
interval contained in
$\{t\in[0,r_i):tv_i\in W\}$, for every $i$.

Let $\operatorname{Psh}_S$ be the smallest subsheaf of the sheaf of
$\mathbf R\cup\{-\infty\}$-valued functions which contains
$\operatorname{Psh}_{S,0}$ and is stable under pointwise decreasing limits
of sequences.  Let $\operatorname{Psh}_{S,1}$ be the sheaf of functions
which are locally pointwise decreasing limits of sequences in
$\operatorname{Psh}_{S,0}$.

Define $\mathscr J_S(W)$ to be the functions
$u:W\to\mathbf R\cup\{-\infty\}$ with the following properties.  First,
$t\mapsto u(tv_i)$ is extended-real convex on every interval contained in
$\{t\in[0,r_i):tv_i\in W\}$.  Second, if $0\in W$, then
\begin{equation}
\label{eq:ch16-curve-junction-jensen}
 u(0)\leq\sum_{i=1}^{m}\alpha_i u(t_iv_i)
\end{equation}
whenever $\alpha_i\geq0$, $0\leq t_i<r_i$,
\begin{equation}
\label{eq:p071-junction-barycentre}
 \sum_i\alpha_i=1,
 \qquad
 \sum_i\alpha_i t_iv_i=0,
\end{equation}
and $\{st_iv_i:0\leq s\leq1\}\subset W$ for every $i$ with
$\alpha_i>0$.  Terms with $\alpha_i=0$ are omitted from the
extended-real sum.  Then $\mathscr J_S$ is a sheaf and
\begin{equation}
\label{eq:p071-model-star-characterization}
 \operatorname{Psh}_{S,1}=\operatorname{Psh}_S=\mathscr J_S.
\end{equation}

More generally, let $X$ be a Hausdorff space equipped with a subsheaf
$\operatorname{Psh}_{X,0}$ of the sheaf of real-valued functions.  Assume
that every $x\in X$ has an open neighbourhood $U$ and a homeomorphism
$\varphi:U\to S$ onto a star of the form
\eqref{eq:p071-local-finite-star}, carrying $x$ to $0$, such that
composition with $\varphi^{-1}$ identifies
$\operatorname{Psh}_{X,0}|_U$ with $\operatorname{Psh}_{S,0}$.  Define
$\operatorname{Psh}_X$ and $\operatorname{Psh}_{X,1}$ from
$\operatorname{Psh}_{X,0}$ by the same two operations as above.  The model
sheaves $\mathscr J_S$ are then preserved on overlaps and descend to a
sheaf $\mathscr J_X$, and
\begin{equation}
\label{eq:ch16-collapse-finite-tropical-curves}
 \operatorname{Psh}_{X,1}=\operatorname{Psh}_X=\mathscr J_X.
\end{equation}
More precisely, if $u$ is a section of any of these sheaves near $x$, there
are one fixed neighbourhood $V$ of $x$ and functions
$G_n\in\operatorname{Psh}_{X,0}(V)$ such that
\begin{equation}
\label{eq:p071-fixed-neighbourhood-approximation}
 G_{n+1}\leq G_n,
 \qquad
 G_n\longrightarrow u
 \quad\text{pointwise on }V.
\end{equation}
No lattice, rationality, weight, balancing, calibration, harmony, or
boundary designation is used.
\end{theorem}

\begingroup
\emergencystretch=2em
% Audit-ready proof for the candidate statement
% p071_finite_affine_tropical_stars_statement_20260827.tex.
% This file is isolated from the live book and has not been submitted to
% Rethlas.

\begin{proof}
We first work on the model \eqref{eq:p071-local-finite-star}.  The defining
conditions of $\mathscr J_S$ are local.  Only locality of
\eqref{eq:ch16-curve-junction-jensen} needs an argument.  Suppose that the two
conditions hold on an open cover and fix a tuple satisfying
\eqref{eq:p071-junction-barycentre}.  A member of the cover that contains
$0$ contains all the radial segments
$\{q t_i v_i:0\leq q\leq s\}$ for every sufficiently small
$s\in(0,1)$ and every relevant $i$.  On that member,
\[
                         u(0)\leq
                         \sum_i\alpha_i u(s t_i v_i).
\]
If $u(0)=-\infty$, \eqref{eq:ch16-curve-junction-jensen} is automatic.  If
$u(0)$ is finite, the last displayed inequality shows that every summand
with positive coefficient is finite.  Convexity on the corresponding ray
then gives
\[
                 u(s t_i v_i)
                 \leq(1-s)u(0)+s u(t_i v_i).
\]
After summing, cancelling $(1-s)u(0)$, and dividing by $s$, one obtains
\eqref{eq:ch16-curve-junction-jensen} for the original tuple.  Hence
$\mathscr J_S$ is a sheaf.

The sheaf $\mathscr J_S$ is stable under pointwise decreasing limits of
sequences.  Indeed, extended-real convexity passes to such a limit, and a
finite weighted sum commutes with a decreasing limit in
$\mathbf R\cup\{-\infty\}$.  It contains $\operatorname{Psh}_{S,0}$.
Indeed, branchwise convexity is part of the definition of that generator
sheaf.  Near $0$, a generator $f$ is the restriction of an ambient smooth
function.  If
\[
                      \ell(x)=f(0)+\mathrm d f(0)(x)
\]
is its ambient affine tangent, branchwise convexity yields
$f_i(t)\geq\ell(tv_i)$.  For a tuple satisfying
\eqref{eq:p071-junction-barycentre}, taking the weighted sum cancels the
linear part of $\ell$ and proves \eqref{eq:ch16-curve-junction-jensen}.
Minimality of the generated psh sheaf therefore gives
\begin{equation}
\label{eq:p071-psh-junction-consequences}
                   \operatorname{Psh}_S\subseteq\mathscr J_S.
\end{equation}

It remains to prove the reverse inclusion in the first-stage form
$\mathscr J_S\subseteq\operatorname{Psh}_{S,1}$.  This is local.  At a
point on an open ray, translate the point to the origin and regard a small
interval as the star with the two opposite ray vectors $v_i$ and $-v_i$.
Extended-real convexity across the interval gives exactly the barycentric
Jensen inequalities for that two-ray star, so the construction below
applies without change.  It is therefore enough to work at $0$.  Fix
$u\in\mathscr J_S$ near $0$, shrink the star once, retain the notation in
\eqref{eq:p071-local-finite-star}, and write
\[
                         u_i(t)=u(tv_i).
\]

We next analyse the polar branches.  If $u(0)=-\infty$, then on each ray
$u$ is identically $-\infty$ on an initial interval.  In fact, this is
clear if the branch has no finite value; if $u_i(a)$ is finite for some
$a>0$, convexity gives, for $0<s<1$,
\[
                    u_i(sa)
                    \leq(1-s)u(0)+s u_i(a)=-\infty.
\]
There are only finitely many rays.  After shrinking all of them at once,
$u$ is therefore identically $-\infty$, and the constants $-n$ in
$\operatorname{Psh}_{S,0}$ decrease to this germ.

Assume from now on that
\begin{equation}
\label{eq:p071-finite-vertex-value}
                              u(0)=c\in\mathbf R.
\end{equation}
On each ray, precisely one of the following alternatives occurs:
\begin{equation}
\label{eq:p071-branch-dichotomy}
 \text{$u_i$ is finite on $(0,r_i)$,}
 \qquad\text{or}\qquad
 \text{$u_i\equiv-\infty$ on $(0,r_i)$.}
\end{equation}
Indeed, suppose that $u_i(a)$ is finite at one interior point.  If
$0<t<a$ and $u_i(t)=-\infty$, choose $b$ with $a<b<r_i$ and express $a$
as a strict convex combination of $t$ and $b$; convexity would force
$u_i(a)=-\infty$.  Thus $u_i$ is finite to the left of $a$.  If $b>a$
and $u_i(b)=-\infty$, express $a$ as a strict convex combination of any
$t\in(0,a)$ and $b$ to obtain the same contradiction.  This proves
\eqref{eq:p071-branch-dichotomy}.  Call the two sets of indices $F$ and
$P$, respectively.

No nonnegative linear relation among the ray vectors uses an index in
$P$:
\begin{equation}
\label{eq:p071-polar-relation-exclusion}
 \lambda_i\geq0,\quad \sum_i\lambda_i v_i=0
 \quad\Longrightarrow\quad
 \lambda_p=0\quad(p\in P).
\end{equation}
For if a nonzero relation had $\lambda_p>0$ for some $p\in P$, choose a
common sufficiently small $t>0$ and put
$\alpha_i=\lambda_i/\sum_j\lambda_j$.  The points $t v_i$ and these
weights satisfy \eqref{eq:p071-junction-barycentre}, whereas the
right-hand side of \eqref{eq:ch16-curve-junction-jensen} is $-\infty$.  This
contradicts \eqref{eq:p071-finite-vertex-value}.

Choose radii
\[
                    0<\rho_i<R_i<r_i
\]
and set
\begin{equation}
\label{eq:p071-fixed-star-neighbourhoods}
 V=\{0\}\cup\coprod_i(0,\rho_i)v_i,
 \qquad
 K=\{0\}\cup\coprod_i(0,R_i]v_i.
\end{equation}
For every $i\in F$, choose finite partitions
$\mathcal P_{i,n}$ of $[0,R_i]$, all containing $0$ and $R_i$, whose
meshes tend to zero and such that each partition refines its predecessor.
Let $h_{n,i}$ be the chordal interpolant of $u_i$ on
$\mathcal P_{i,n}$, using the value $c$ at $0$.  For $i\in P$, set
\begin{equation}
\label{eq:p071-polar-chords}
                         h_{n,i}(t)=c-nt.
\end{equation}
The functions $h_n$ obtained by gluing the branches at the common value
$c$ are finite continuous convex piecewise-affine functions on $K$, and
\begin{equation}
\label{eq:p071-chord-monotonicity}
 h_{n+1}\leq h_n,
 \qquad
 h_n(x)\longrightarrow u(x)
 \quad(x\in K).
\end{equation}
The first assertion follows because refinement lowers the chordal
majorant of a convex function and because
$c-(n+1)t\leq c-nt$.  At a point $t>0$ of a finite branch, convergence
follows from continuity of a finite convex function in the interior and
the vanishing mesh; at $0$ there is equality for every $n$.  On a polar
branch \eqref{eq:p071-polar-chords} tends to $-\infty$.  Notice that no
continuity of $u$ at $0$ has been asserted or used.

We shall find an ambient affine minorant of every $h_n$.  We include the
finite-dimensional alternative needed for this step.  Given vectors
$v_1,\ldots,v_m\in E$ and real numbers $A_1,\ldots,A_m$, there is a covector
$\xi\in E^*$ with
\begin{equation}
\label{eq:p071-farkas-upper-system}
                           \xi(v_i)\leq A_i
                           \quad(1\leq i\leq m)
\end{equation}
if and only if
\begin{equation}
\label{eq:p071-farkas-relations}
 \sum_i\lambda_iA_i\geq0
 \quad\text{for every}\quad
 \lambda_i\geq0\ \text{with}\ \sum_i\lambda_i v_i=0.
\end{equation}
Necessity follows by summing \eqref{eq:p071-farkas-upper-system}.  For
sufficiency, let
\[
 L=\{(\xi(v_1),\ldots,\xi(v_m)):\xi\in E^*\},
 \qquad Q=\mathbf R_{\geq0}^m.
\]
The system \eqref{eq:p071-farkas-upper-system} is solvable exactly when
$A=(A_i)$ belongs to the closed polyhedral cone $L+Q$.  If it were not
solvable, strict separation would give $\lambda\in\mathbf R^m$ such that
$\langle\lambda,A\rangle<0$ and
$\langle\lambda,z\rangle\geq0$ for every $z\in L+Q$.  Since both $L$ and
$-L$ lie in $L+Q$, the vector $\lambda$ is orthogonal to $L$; since
$Q\subseteq L+Q$, it has nonnegative coordinates.  These two conclusions
contradict \eqref{eq:p071-farkas-relations}.

Fix $n$.  If $i\in F$, let $s_{i,n}$ be the first positive point of
$\mathcal P_{i,n}$ and put
\[
                 A_{i,n}=\frac{u_i(s_{i,n})-c}{s_{i,n}}.
\]
If $i\in P$, put $A_{i,n}=-n$.  We verify
\eqref{eq:p071-farkas-relations} for these numbers.  Let
$\lambda_i\geq0$ and $\sum_i\lambda_i v_i=0$.  By
\eqref{eq:p071-polar-relation-exclusion}, every coefficient indexed by
$P$ vanishes.  If the relation is nonzero, set
\begin{equation}
\label{eq:p071-junction-weights-from-relation}
 D=\sum_{i\in F}\frac{\lambda_i}{s_{i,n}},
 \qquad
 \beta_i=\frac{\lambda_i/s_{i,n}}{D}\quad(i\in F).
\end{equation}
Then $\sum_i\beta_i=1$ and
$\sum_i\beta_i s_{i,n}v_i=0$.  Applying
\eqref{eq:ch16-curve-junction-jensen}, subtracting $c$, and multiplying by $D$
gives
\begin{equation}
\label{eq:p071-slope-relation}
 0\leq\sum_{i\in F}\lambda_i
       \frac{u_i(s_{i,n})-c}{s_{i,n}}
   =\sum_i\lambda_iA_{i,n}.
\end{equation}
The zero relation is automatic.  The finite-dimensional alternative now
supplies $\xi_n\in E^*$ satisfying
\begin{equation}
\label{eq:ch16-curve-support-covector}
                         \xi_n(v_i)\leq A_{i,n}
                         \quad(1\leq i\leq m).
\end{equation}
Since the outgoing slopes of each convex $h_{n,i}$ are nondecreasing,
\eqref{eq:ch16-curve-support-covector} implies
\begin{equation}
\label{eq:p071-common-affine-minorant}
                  \ell_n(x):=c+\xi_n(x)\leq h_n(x)
                  \quad(x\in K).
\end{equation}

We now smooth while preserving both a common ambient germ at $0$ and
monotonicity in $n$.  For completeness, fix a nonnegative even
$C^\infty$ function $\theta$ on $\mathbf R$, supported in $[-1,1]$ and
having integral one.  For $q\geq1$ and $\varepsilon>0$, put
\[
 \rho_{\varepsilon,q}(t)
   =\prod_{j=1}^q\varepsilon^{-1}\theta(t_j/\varepsilon),
 \qquad
 M_{\varepsilon,q}(x)
   =\int_{\mathbf R^q}\max_j(x_j+t_j)
       \rho_{\varepsilon,q}(t)\,\mathrm d t.
\]
This is a smooth function $M_{\varepsilon,q}:\mathbf R^q\to\mathbf R$
such that
\begin{align}
 \max_j x_j
 &\leq M_{\varepsilon,q}(x_1,\ldots,x_q)
 \leq \max_jx_j+\varepsilon,
 \label{eq:p071-regularized-max-sandwich}\\
 M_{\varepsilon,q}(x_1,\ldots,x_q)&=x_k
 \quad\text{if}\quad
 x_k>x_j+2\varepsilon\quad(j\neq k).
 \label{eq:p071-regularized-max-truncation}
\end{align}
It is convex and nondecreasing in every variable.  Indeed, convolution
preserves convexity and coordinatewise monotonicity; evenness of $\theta$
makes the barycentre of $\rho_{\varepsilon,q}$ equal to zero, so the lower
bound in
\eqref{eq:p071-regularized-max-sandwich} is Jensen's inequality, the upper
bound follows from the support of the density, and
\eqref{eq:p071-regularized-max-truncation} follows because the $k$-th
entry remains the unique maximum throughout the convolution support and
the density has zero first moments.

Write each convex piecewise-affine branch function as the maximum of the
affine extensions of its pieces:
\begin{equation}
\label{eq:p071-piecewise-affine-maximum}
                  h_{n,i}(t)=\max_{1\leq j\leq q_{n,i}}L_{n,i,j}(t).
\end{equation}
Every affine function in \eqref{eq:p071-piecewise-affine-maximum} is a
minorant of $h_{n,i}$ on $[0,R_i]$, so
$L_{n,i,j}(0)\leq c$.  Put
\begin{equation}
\label{eq:p071-error-ledger}
 C_n=4^{-n},
 \qquad
 \varepsilon_n=\frac14C_n,
 \qquad
 E_n=C_n+\varepsilon_n=\frac54\,4^{-n},
 \qquad
 \delta_n=2^{-n},
\end{equation}
and define on branch $i$
\begin{equation}
\label{eq:p071-smoothed-branches}
 g_{n,i}(t)=M_{\varepsilon_n,q_{n,i}+1}
 \bigl(c+C_n+\xi_n(tv_i),
       L_{n,i,1}(t),\ldots,L_{n,i,q_{n,i}}(t)\bigr).
\end{equation}
At $t=0$, the first entry in \eqref{eq:p071-smoothed-branches} exceeds
all the others by at least $C_n>2\varepsilon_n$.  Hence
\eqref{eq:p071-regularized-max-truncation} and continuity show that, on an
initial interval of every branch,
\[
                          g_{n,i}(t)=c+C_n+\xi_n(tv_i).
\]
For the fixed $n$, finiteness of the star lets us shrink these intervals
simultaneously; the common affine radius is allowed to depend on $n$.
Thus the branch functions glue near $0$ to the restriction of the single
ambient affine function $c+C_n+\xi_n$, and they define a smooth function
$g_n$ on the fixed neighbourhood $V$ in
\eqref{eq:p071-fixed-star-neighbourhoods}.  Away from $0$, a covector
$\eta_i$ with $\eta_i(v_i)=1$ extends the ray coordinate, so the ordinary
smooth function $g_{n,i}(t)$ is also ambient-induced there.

Each $g_{n,i}$ is convex, since it is a convex function of affine
functions.  Together with the ambient smoothness just proved, this says
precisely that $g_n\in\operatorname{Psh}_{S,0}(V)$.  From
\eqref{eq:p071-common-affine-minorant},
\eqref{eq:p071-regularized-max-sandwich}, and
\eqref{eq:p071-piecewise-affine-maximum}, we obtain on $K$
\begin{equation}
\label{eq:p071-smoothing-sandwich}
                         h_n\leq g_n\leq h_n+E_n.
\end{equation}

Finally set
\begin{equation}
\label{eq:p071-final-approximants}
                              G_n=g_n+\delta_n.
\end{equation}
The constants in \eqref{eq:p071-error-ledger} were chosen so that, for
$n\geq1$,
\begin{equation}
\label{eq:p071-monotonicity-budget}
 E_{n+1}\leq\delta_n-\delta_{n+1};
\end{equation}
indeed, the two sides are $5\cdot2^{-2n-4}$ and $2^{-n-1}$,
respectively.  Combining \eqref{eq:p071-chord-monotonicity},
\eqref{eq:p071-smoothing-sandwich}, and
\eqref{eq:p071-monotonicity-budget} gives
\[
 G_{n+1}
 \leq h_{n+1}+E_{n+1}+\delta_{n+1}
 \leq h_n+\delta_n
 \leq g_n+\delta_n
 =G_n
\]
on the one fixed neighbourhood $V$.  The sandwich
\eqref{eq:p071-smoothing-sandwich}, the convergence in
\eqref{eq:p071-chord-monotonicity}, and
$E_n+\delta_n\to0$ show that $G_n\to u$ at $0$ and on every finite
branch.  On a polar branch, for each fixed $t>0$,
\[
             G_n(tv_i)\leq c-nt+E_n+\delta_n\longrightarrow-\infty.
\]
Thus \eqref{eq:p071-fixed-neighbourhood-approximation} holds in all cases.

We have proved $\mathscr J_S\subseteq\operatorname{Psh}_{S,1}$.  The
definition gives
$\operatorname{Psh}_{S,1}\subseteq\operatorname{Psh}_S$, and
\eqref{eq:p071-psh-junction-consequences} gives the remaining inclusion.
This proves \eqref{eq:p071-model-star-characterization}.

Now consider $X$.  A transition between two of the stated charts preserves
the generator sheaves and therefore also their first-stage sheaves.  By the
model equality just proved, it consequently preserves the model sheaves
$\mathscr J_S$; they thus descend to $\mathscr J_X$.  This sheaf contains
$\operatorname{Psh}_{X,0}$ and is stable under decreasing sequential
limits, so minimality gives
$\operatorname{Psh}_X\subseteq\mathscr J_X$.  Conversely, the explicit
model construction transports through a chart and puts every germ of
$\mathscr J_X$ in $\operatorname{Psh}_{X,1}$ on one fixed neighbourhood.
The tautological inclusion
$\operatorname{Psh}_{X,1}\subseteq\operatorname{Psh}_X$ proves
\eqref{eq:ch16-collapse-finite-tropical-curves}, and the same transported
construction gives \eqref{eq:p071-fixed-neighbourhood-approximation}.
\end{proof}

\endgroup

% The superseded pre-P071 proof is recoverable from Git commit 13094dc.
\begin{example}
\label{ex:ch16-collapse-weighted-tripod}
Let $S\subset\mathbf R^2$ be the open tripod with rays
\[
 v_1=(1,0),\qquad v_2=(0,1),\qquad v_3=(-1,-1),
\]
all of weight one.  Define $u(0)=0$ and, for $t>0$,
\[
                         u_1(t)=u_2(t)=0,
 \qquad                   u_3(t)=t.
\]
This function is convex on every branch but is not smooth at the vertex:
its three outgoing slopes do not come from one ambient covector.  For an
admissible tuple in \eqref{eq:ch16-curve-junction-jensen}, the relation
$\sum_i\alpha_it_iv_i=0$ says
$\alpha_1t_1=\alpha_2t_2=\alpha_3t_3$, and hence
\[
                   \sum_i\alpha_i u_i(t_i)=\alpha_3t_3\geq0=u(0).
\]
The approximation in the proof is especially transparent here.  Take
$\xi=0$ and, for $0<2\varepsilon<C$, regularize branchwise
$\max\{C,u_i(t)\}$.  The result equals the common constant $C$ near the
vertex, is smooth and convex on every branch, and differs from $u$ by at
most $C+\varepsilon$ on any fixed compact substar.  Applying
\Cref{lem:ch15-shift} produces a decreasing sequence of smooth psh
functions converging to $u$, so $u\in\operatorname{Psh}_{S,1}$.

The junction condition is essential.  The branchwise affine function
$q_i(t)=-t$ is convex on all three rays, but the balanced choice
$\alpha_i=1/3$ and $t_i=t>0$ gives
$\sum_i\alpha_iq_i(t_i)=-t<q(0)=0$.  Therefore $q$ is not CLD-psh, by
\eqref{eq:ch16-curve-junction-jensen}.
\end{example}

The preceding finite-star theorem uses the smooth structure induced from a
finite-dimensional ambient vector space.  There are two natural endpoint
models in which that hypothesis fails: the smooth functions are required to
be exactly constant on a neighbourhood of the endpoint.  The finite endpoint
and the endpoint at $-\infty$ behave differently.  At a finite endpoint,
convexity forces a boundary minimum; at $-\infty$, it instead forces an exact
trace condition.  In particular, the function $u(t)=t$, extended by
$u(-\infty)=-\infty$, is admissible in the second model and shows why the two
statements cannot be merged by a formal change of coordinate.

For auditability, the next four inputs are the exact frozen statement and
proof bytes accepted by production Rethlas verification.  Both runs used
\texttt{gpt-5.6-sol} at \texttt{xhigh} and returned verdict
\texttt{correct}, with zero critical errors and zero gaps.

% Exact isolated candidate statement for source-aware verification.
% This file makes no novelty or priority claim and is not a live-book input.

Let \(R_0\in(0,+\infty]\), and put \(S=[0,R_0)\) with its usual
subspace topology.  For an open subset \(W\subseteq S\), let
\(C^\infty_{\partial}(W)\) be the real-valued functions \(f\) such that
\(f|_{W\cap(0,R_0)}\) is smooth in the ordinary sense and, if \(0\in W\),
there is an \(\varepsilon>0\) for which
\begin{equation}
  [0,\varepsilon)\subseteq W,
  \qquad
  f|_{[0,\varepsilon)}\ \text{is constant}.
  \label{eq:p071-endpoint-collar-smoothness}
\end{equation}
Let \(\operatorname{Psh}_{S,0}(W)\) consist of the
\(f\in C^\infty_{\partial}(W)\) such that
\begin{equation}
                         f''\geq 0
       \quad\text{on }W\cap(0,R_0).
  \label{eq:p071-endpoint-stage-zero}
\end{equation}
Let \(\operatorname{Psh}_S\) be the smallest subsheaf of the sheaf of
\(\mathbf R\cup\{-\infty\}\)-valued functions that contains
\(\operatorname{Psh}_{S,0}\) and is stable under pointwise decreasing
limits of sequences.  Let \(\operatorname{Psh}_{S,1}\) be the sheaf of
functions that are locally pointwise decreasing limits of sequences in
\(\operatorname{Psh}_{S,0}\).

For an open \(W\subseteq S\), define \(\mathscr J_{\partial}(W)\) to be
the functions \(u:W\to\mathbf R\cup\{-\infty\}\) with the following two
properties.
\begin{enumerate}
\item The restriction of \(u\) to every interval contained in \(W\) is
extended-real convex.  Thus, for \(x,y\in W\), \(0<\lambda<1\), and
\([x,y]\subseteq W\),
\begin{equation}
 u\bigl((1-\lambda)x+\lambda y\bigr)
 \leq (1-\lambda)u(x)+\lambda u(y),
 \label{eq:p071-endpoint-extended-convexity}
\end{equation}
where every positive multiple of \(-\infty\) is \(-\infty\).
\item If \(0\in W\) and \(W_0\) is the connected component of \(W\)
containing \(0\), then
\begin{equation}
                              u(0)\leq u(t)
                    \quad\text{for every }t\in W_0.
 \label{eq:p071-endpoint-boundary-minimum}
\end{equation}
\end{enumerate}

\begin{theorem}
\label{thm:p071-cld-constant-collar-endpoint}
With the preceding definitions, \(\mathscr J_{\partial}\) is a sheaf and
\begin{equation}
 \operatorname{Psh}_{S,1}
 =\operatorname{Psh}_S
 =\mathscr J_{\partial}.
 \label{eq:p071-endpoint-first-stage-collapse}
\end{equation}
More precisely, suppose that \(0\in W\) and
\(u\in\mathscr J_{\partial}(W)\).  There is a fixed open neighbourhood
\(V=[0,r)\subseteq W\) and functions
\(G_n\in\operatorname{Psh}_{S,0}(V)\) such that
\begin{equation}
 G_{n+1}\leq G_n,
 \qquad
 G_n\longrightarrow u|_V
 \quad\text{pointwise on }V.
 \label{eq:p071-endpoint-fixed-domain-approximation}
\end{equation}
Each \(G_n\) is constant on an initial interval whose length may depend
on \(n\).  If \(u(0)\in\mathbf R\), then \(u\) is real-valued and
continuous on \(V\), and the convergence in
\eqref{eq:p071-endpoint-fixed-domain-approximation} can be chosen uniform
on \(V\).  If \(u(0)=-\infty\), then \(u\) is identically \(-\infty\) on
\(V\), and one may take \(G_n=-n\).

Fix a finite endpoint \(a\) in the interval of
\cite[Example~6.7.2]{CLD25}.  In an ordinary inward real coordinate on the
finite stratum, extended topologically by the value \(0\) at \(a\), this is
exactly the constant-near-the-endpoint smooth structure of that example.  In
that model,
\eqref{eq:p071-endpoint-stage-zero} is the condition
\(\mathrm d'\mathrm d''f\geq0\).
This assertion does not concern an infinite endpoint of
\(\overline{\mathbf R}\), and in particular does not concern the
\([-\infty,0)\) open-unit-disc realization highlighted in that example.
\end{theorem}

\begingroup
\emergencystretch=2em
% Exact isolated candidate proof for source-aware verification.
% This file makes no novelty or priority claim and is not a live-book input.

\begin{proof}
\medskip\noindent\emph{Imported source scope.}
In \cite[Example~6.7.2]{CLD25}, Chambert--Loir and Ducros put a
\(G\)-tropical affine structure on an interval
\(I\subseteq\overline{\mathbf R}\), with
\(A=\partial I\cap\overline{\mathbf R}\), by requiring parlinear functions
to be constant near every point of \(A\).  In particular, every affine
function on a connected domain meeting \(A\) is constant.  Fix a finite
endpoint \(a\in A\cap\mathbf R\).  Together with their definition of smooth
forms in \cite[Definition~6.2.1]{CLD25}, the example says the following.
By \cite[\S\S6.1.1 and 6.1.6]{CLD25}, a tropical chart is affine and each of
its scalar coordinate functions is affine.  After shrinking to a connected
neighbourhood meeting \(A\), Example~6.7.2 therefore makes every coordinate
constant.  The chart, and every smooth \(0\)-form pulled back through it, is
constant there.  Use the ordinary inward real coordinate on the finite
stratum and extend it topologically by the value \(t(a)=0\); this is not a
nonconstant affine chart through \(a\).  In this coordinate, the smooth
functions are precisely the functions described by
\eqref{eq:p071-endpoint-collar-smoothness}.  Conversely, use the constant
terminal chart on the collar and the ordinary coordinate \(t\) on
endpoint-free interior neighbourhoods; these local descriptions glue.  On
the one-dimensional interior,
\[
 \mathrm d'\mathrm d'' f
   =f''(t)\,\mathrm d't\wedge\mathrm d''t,
\]
so positivity is exactly \eqref{eq:p071-endpoint-stage-zero}; the form
vanishes on the collar where \(f\) is constant.  Definition~8.1.1 and
Lemma~8.1.4 of \cite{CLD25} give the generated sheaf and its ordinal
filtration used in the statement.  We will also use the elementary
uniform-to-monotone device of \cite[Lemma~8.1.10]{CLD25}: if real-valued
functions \(f_j\) converge uniformly to \(f\) on a set \(K\), then
\[
                  f_j+2\lVert f_j-f\rVert_{K}
\]
has a decreasing subsequence converging uniformly to \(f\).

\medskip\noindent\emph{The comparison class is a sheaf.}
Restrictions in \(\mathscr J_{\partial}\) are immediate.  We check
gluing.  Suppose that \(W=\bigcup_{\nu}W_\nu\) and that
\(u|_{W_\nu}\in\mathscr J_{\partial}(W_\nu)\) for every \(\nu\).  Let
\(C\) be a connected component of \(W\).  The function \(u|_C\) is locally
extended-real convex.

We first show that \(u|_C\) is either finite everywhere or identically
\(-\infty\).  Let
\[
 Z=\{t\in C:u(t)=-\infty\},
 \qquad
 F=C\setminus Z.
\]
If \(t\in Z\) is an interior point of \(C\), choose an interval
\((t-\delta,t+\delta)\cap C\) on which \(u\) is convex.  For a point
\(s>t\) sufficiently close to \(t\), choose \(q>s\) in that interval and
write \(s\) as a strict convex combination of \(t\) and \(q\).
Then \eqref{eq:p071-endpoint-extended-convexity} gives \(u(s)=-\infty\).
The same argument on the left shows that \(Z\) is open at \(t\).  If
\(0\in Z\), choose \(a>0\) such that \(u\) is convex on \([0,a)\).
For \(0<s<a\), choose \(q\in(s,a)\) and write
\(s=(1-s/q)0+(s/q)q\); again \(u(s)=-\infty\).  Thus \(Z\) is open also
at the endpoint.

If \(t\in F\) is an interior point and there were points \(s\in Z\)
arbitrarily close to \(t\), choose such an \(s\) and a point \(q\) on the
other side of \(t\), all in one interval of local convexity.  Expressing
\(t\) as a strict convex combination of \(s\) and \(q\) would give
\(u(t)=-\infty\), a contradiction.  Hence \(F\) is open at every interior
point.  If \(0\in F\), the local instance of
\eqref{eq:p071-endpoint-boundary-minimum} excludes the value
\(-\infty\) on a whole initial interval, so \(F\) is open at \(0\) as
well.  Since \(C\) is connected, either \(Z=C\) or \(F=C\).

In the finite case, the standard one-dimensional local-to-global
criterion shows that \(u|_C\) is convex.  Here is a direct verification.
Given \(a<c<b\) in \(C\), cover \([a,b]\) by finitely many intervals on
which \(u\) is convex.  Take a Lebesgue number \(\delta>0\) for this finite
cover of \([a,b]\), and choose a partition
\[
             a=x_0<x_1<\cdots<x_N=b
\]
that contains \(c\) and satisfies
\(x_{i+1}-x_{i-1}<\delta\) whenever the three terms occur.  Each interval
\([x_{i-1},x_{i+1}]\) then lies in one member of the cover.  Local convexity
therefore says that the successive secant slopes are nondecreasing.  Summing
the resulting inequalities on the two sides of \(c\) yields
\[
 \frac{u(c)-u(a)}{c-a}
 \leq
 \frac{u(b)-u(c)}{b-c},
\]
which is equivalent to the convexity inequality at \(c\).  Thus
\eqref{eq:p071-endpoint-extended-convexity} holds on every interval in
\(C\).  It is automatic when \(u|_C\equiv-\infty\).

It remains to propagate the boundary minimum.  Suppose \(0\in C\), fix
\(t\in C\), and choose \(s\in(0,1)\) so small that \(st\) lies with \(0\)
in one member of the cover.  There
\[
                              u(0)\leq u(st).
\]
If \(u(0)=-\infty\), the desired inequality is automatic.  If
\(u(0)\in\mathbf R\), the preceding dichotomy makes \(u\) finite on
\(C\), while global convexity gives
\[
                 u(st)\leq(1-s)u(0)+s u(t).
\]
Combining the last two inequalities, cancelling \((1-s)u(0)\), and
dividing by \(s\) gives \(u(0)\leq u(t)\).  This proves the gluing axiom
and hence that \(\mathscr J_{\partial}\) is a sheaf.

\medskip\noindent\emph{The generated sheaf is contained in the comparison
class.}
The sheaf \(\mathscr J_{\partial}\) is stable under pointwise decreasing
limits.  Indeed, every inequality
\eqref{eq:p071-endpoint-extended-convexity} passes to the limit in
\(\mathbf R\cup\{-\infty\}\), since no \(+\infty\) value occurs, and
\eqref{eq:p071-endpoint-boundary-minimum} passes to the limit as well.

It also contains \(\operatorname{Psh}_{S,0}\).  On an interior interval,
\eqref{eq:p071-endpoint-stage-zero} gives ordinary convexity.  On the
component meeting \(0\), the derivative of a generator is zero on the
constant collar and is nondecreasing thereafter; consequently the
generator is convex on the whole half-interval and takes its minimum at
\(0\).  By the minimal definition of the generated psh sheaf,
\begin{equation}
                       \operatorname{Psh}_S
                       \subseteq\mathscr J_{\partial}.
 \label{eq:p071-endpoint-generated-in-jensen}
\end{equation}

\medskip\noindent\emph{The extended-real alternatives at the endpoint.}
Let \(0\in W\) and \(u\in\mathscr J_{\partial}(W)\).  After shrinking its
connected component, choose \(R>0\) with \([0,R)\subseteq W\).

Suppose first that \(u(0)=-\infty\).  For \(0<t<R\), choose
\(q\in(t,R)\).  Since
\[
                         t=(1-t/q)0+(t/q)q,
\]
extended-real convexity gives
\[
 u(t)\leq(1-t/q)u(0)+(t/q)u(q)=-\infty.
\]
Thus \(u\equiv-\infty\) on \([0,R)\).  The smooth psh constants \(-n\)
decrease to this function on the fixed neighbourhood.

Assume now that
\begin{equation}
                              u(0)\in\mathbf R.
 \label{eq:p071-endpoint-finite-value}
\end{equation}
The boundary minimum excludes \(-\infty\) at every other point, so \(u\)
is real-valued.  For \(0<s<t<R\), convexity and
\eqref{eq:p071-endpoint-boundary-minimum} give
\begin{equation}
 u(s)
 \leq\left(1-\frac{s}{t}\right)u(0)+\frac{s}{t}u(t)
 \leq u(t).
 \label{eq:p071-endpoint-monotonicity}
\end{equation}
Hence \(u\) is nondecreasing.  It is continuous in the interior by
ordinary finite convexity.  At the endpoint, for any fixed \(t\in(0,R)\),
\eqref{eq:p071-endpoint-boundary-minimum} and the first inequality in
\eqref{eq:p071-endpoint-monotonicity} show that
\[
 u(0)\leq u(s)
 \leq u(0)+\frac{s}{t}\bigl(u(t)-u(0)\bigr)
 \quad(0<s<t),
\]
so \(u(s)\to u(0)\) as \(s\downarrow0\).  Thus \(u\) is continuous on
\([0,R)\).

\medskip\noindent\emph{A fixed-domain collar regularization.}
Choose
\[
                         0<r<R_1<R_2<R.
\]
Put
\[
                         m=
             \frac{u(R_2)-u(R_1)}{R_2-R_1}\geq0
\]
and define a finite function \(\bar u:\mathbf R\to\mathbf R\) by
\begin{equation}
 \bar u(t)=
 \begin{cases}
  u(0),&t\leq0,\\
  u(t),&0\leq t\leq R_1,\\
  u(R_1)+m(t-R_1),&t\geq R_1.
 \end{cases}
 \label{eq:p071-endpoint-convex-extension}
\end{equation}
This extension is convex and nondecreasing.  At \(0\), the left slope is
zero and every right secant slope is nonnegative by
\eqref{eq:p071-endpoint-boundary-minimum}.  At \(R_1\), convexity of
\(u\) shows that every left secant slope is at most the secant slope
\(m\) to \(R_2\).  The one-dimensional slope criterion for convexity
therefore applies at both joining points.

Fix a nonnegative even function
\(\theta\in C^\infty_c((-1,1))\) with integral one, and put
\(\theta_\sigma(s)=\sigma^{-1}\theta(s/\sigma)\).  For a sequence
\(\sigma_j\downarrow0\), define
\begin{equation}
                         A_j=\bar u*\theta_{\sigma_j}.
 \label{eq:p071-endpoint-convolution}
\end{equation}
Every \(A_j\) is smooth, convex, and nondecreasing.  Evenness gives
\(\int s\theta_{\sigma_j}(s)\,\mathrm ds=0\), so Jensen's inequality
gives
\begin{equation}
                              A_j\geq\bar u.
 \label{eq:p071-endpoint-convolution-dominates}
\end{equation}
Moreover, \(A_j\to\bar u\) uniformly on every compact interval.

We next impose exact constancy on a collar.  For \(\eta>0\), set
\begin{equation}
 M_\eta(a,b)=
 \int_{\mathbf R^2}
 \max\{a+s_1,b+s_2\}
 \theta_\eta(s_1)\theta_\eta(s_2)\,
 \mathrm ds_1\,\mathrm ds_2.
 \label{eq:p071-endpoint-regularized-maximum}
\end{equation}
This function is smooth, convex, and nondecreasing in each variable.
The zero first moments and the support of the kernel give
\begin{equation}
 \max\{a,b\}
 \leq M_\eta(a,b)
 \leq\max\{a,b\}+\eta.
 \label{eq:p071-endpoint-regularized-max-bounds}
\end{equation}
If \(a>b+2\eta\), then \(a+s_1>b+s_2\) throughout the convolution
support and \(M_\eta(a,b)=a\); symmetrically,
\begin{equation}
 M_\eta(a,b)=b
 \quad\text{whenever}\quad b>a+2\eta.
 \label{eq:p071-endpoint-regularized-max-truncation}
\end{equation}

Choose \(\varepsilon_j\downarrow0\), put
\(\eta_j=\varepsilon_j/4\), and define
\begin{equation}
 H_j(t)=M_{\eta_j}(t,\varepsilon_j),
 \qquad
 F_j(t)=A_j\bigl(H_j(t)\bigr).
 \label{eq:p071-endpoint-collar-approximants}
\end{equation}
The function \(H_j\) is smooth, convex, and nondecreasing.  From
\eqref{eq:p071-endpoint-regularized-max-bounds} one gets, for \(t\geq0\),
\begin{equation}
 0\leq H_j(t)-t\leq\varepsilon_j+\eta_j.
 \label{eq:p071-endpoint-reparametrization-error}
\end{equation}
By \eqref{eq:p071-endpoint-regularized-max-truncation},
\begin{equation}
 H_j(t)=\varepsilon_j
 \quad\text{for}\quad
 0\leq t<\varepsilon_j-2\eta_j
                =\frac{\varepsilon_j}{2}.
 \label{eq:p071-endpoint-exact-collar}
\end{equation}
Consequently \(F_j\) is smooth and exactly constant on an initial
interval.

The function \(F_j\) is convex.  Indeed, \(A_j'\geq0\),
\(A_j''\geq0\), \(H_j'\geq0\), and \(H_j''\geq0\), so
\[
 F_j''
  =A_j''(H_j)(H_j')^2+A_j'(H_j)H_j''
  \geq0.
\]
Thus \(F_j|_{[0,r)}\) belongs to
\(\operatorname{Psh}_{S,0}([0,r))\).  It also dominates \(u\) on
\([0,r]\): by monotonicity of \(A_j\),
\eqref{eq:p071-endpoint-reparametrization-error}, and
\eqref{eq:p071-endpoint-convolution-dominates},
\begin{equation}
 F_j(t)=A_j(H_j(t))
 \geq A_j(t)
 \geq\bar u(t)=u(t).
 \label{eq:p071-endpoint-smooth-domination}
\end{equation}

Finally, \(F_j\to u\) uniformly on \(K=[0,r]\).  For all sufficiently
large \(j\), both \(t\) and \(H_j(t)\) lie in the fixed compact interval
\([0,r+1]\).  Hence
\begin{align}
 \sup_{t\in K}|F_j(t)-u(t)|
 &\leq
 \sup_{x\in[0,r+1]}|A_j(x)-\bar u(x)|
 \notag\\
 &\quad+
 \sup_{t\in K}|\bar u(H_j(t))-\bar u(t)|
 \longrightarrow0,
 \label{eq:p071-endpoint-uniform-convergence}
\end{align}
where the second term tends to zero by uniform continuity of \(\bar u\)
and \eqref{eq:p071-endpoint-reparametrization-error}.

Let
\[
                       e_j=\lVert F_j-u\rVert_K.
\]
By \cite[Lemma~8.1.10]{CLD25}, the sequence \(F_j+2e_j\) has a
subsequence, say
\[
                         G_n=F_{j_n}+2e_{j_n},
\]
which decreases uniformly to \(u\) on \(K\).  Adding a constant preserves
smoothness, collar constancy, and convexity.  Restricting to the one fixed
open neighbourhood \(V=[0,r)\) proves
\eqref{eq:p071-endpoint-fixed-domain-approximation} in the finite case.

\medskip\noindent\emph{Interior points and conclusion.}
Let \(x\in W\cap(0,R_0)\).  On a sufficiently small connected interval
about \(x\), the dichotomy proved above says that a section of
\(\mathscr J_{\partial}\) is either identically \(-\infty\) or is a
finite convex function.  The first case is the decreasing limit of the
smooth psh constants \(-n\).  In the second case, choose a compact
interval \(K=[a,b]\) around \(x\), contained in a slightly larger interval
on which \(u\) is finite and convex.  The one-sided derivatives of a finite
convex function are finite at the interior points \(a,b\).  Choose
\[
                    p_-\in\partial u(a),
                    \qquad p_+\in\partial u(b),
\]
and extend \(u|_{[a,b]}\) by
\[
 u(a)+p_-(t-a)\quad(t<a),
 \qquad
 u(b)+p_+(t-b)\quad(t>b).
\]
Monotonicity of the subgradients makes this a finite convex function on
\(\mathbf R\).  Convolution with the even kernels used in
\eqref{eq:p071-endpoint-convolution} gives smooth convex functions
converging uniformly to \(u\) on \(K\).  Applying
\cite[Lemma~8.1.10]{CLD25} and restricting to the interior of \(K\)
produces a decreasing sequence of smooth convex functions on one fixed
neighbourhood of \(x\).

We have proved
\[
                 \mathscr J_{\partial}
                 \subseteq\operatorname{Psh}_{S,1}.
\]
The tautological inclusion
\(\operatorname{Psh}_{S,1}\subseteq\operatorname{Psh}_S\), together with
\eqref{eq:p071-endpoint-generated-in-jensen}, proves
\eqref{eq:p071-endpoint-first-stage-collapse}.  The construction at the
endpoint gives every assertion following
\eqref{eq:p071-endpoint-fixed-domain-approximation}.
\end{proof}

\endgroup

% Exact isolated candidate statement for source-aware verification.
% This file makes no novelty or priority claim and is not a live-book input.

Let
\[
       \overline{\mathbf R}_{-}=\mathbf R\cup\{-\infty\},
       \qquad S=[-\infty,R_0)\subseteq\overline{\mathbf R}_{-},
\]
where \(R_0\in\mathbf R\cup\{+\infty\}\), and give \(S\) its order
topology.  For an open subset \(W\subseteq S\), let
\(C^\infty_{-\infty}(W)\) be the real-valued functions \(f\) such that
\(f|_{W\cap\mathbf R}\) is smooth in the ordinary sense and, if
\(-\infty\in W\), there is \(a\in\mathbf R\) for which
\begin{equation}
 [-\infty,a)\subseteq W,
 \qquad f|_{[-\infty,a)}\ \text{is constant}.
 \label{eq:p071-infinite-collar-smoothness}
\end{equation}
Let \(\operatorname{Psh}_{S,0}(W)\) consist of the
\(f\in C^\infty_{-\infty}(W)\) such that
\begin{equation}
                         f''\geq0
                 \quad\text{on }W\cap\mathbf R.
 \label{eq:p071-infinite-stage-zero}
\end{equation}
Let \(\operatorname{Psh}_S\) be the smallest subsheaf of the sheaf of
\(\mathbf R\cup\{-\infty\}\)-valued functions that contains
\(\operatorname{Psh}_{S,0}\) and is stable under pointwise decreasing
limits of sequences.  Let \(\operatorname{Psh}_{S,1}\) be the sheaf of
functions that are locally pointwise decreasing limits of sequences in
\(\operatorname{Psh}_{S,0}\).

For an open \(W\subseteq S\), define \(\mathscr J_{-\infty}(W)\) to be
the functions \(u:W\to\mathbf R\cup\{-\infty\}\) with the following
properties.
\begin{enumerate}
\item The restriction of \(u\) to every finite real interval contained
in \(W\cap\mathbf R\) is extended-real convex: if
\(x,y\in W\cap\mathbf R\), \(0<\lambda<1\), and
\([x,y]\subseteq W\), then
\begin{equation}
 u\bigl((1-\lambda)x+\lambda y\bigr)
 \leq(1-\lambda)u(x)+\lambda u(y),
 \label{eq:p071-infinite-extended-convexity}
\end{equation}
where every positive multiple of \(-\infty\) is \(-\infty\).
\item If \(-\infty\in W\) and \(W_{-\infty}\) is the connected
component of \(W\) containing \(-\infty\), then
\begin{equation}
 u(-\infty)
   =\lim_{t\to-\infty,\ t\in W_{-\infty}\cap\mathbf R}u(t)
       \quad\text{in }\mathbf R\cup\{-\infty\}.
 \label{eq:p071-infinite-endpoint-trace}
\end{equation}
In particular, the finite-side limit is not allowed to be \(+\infty\).
\end{enumerate}

\begin{theorem}
\label{thm:p071-cld-infinite-collar-endpoint}
With the preceding definitions, \(\mathscr J_{-\infty}\) is a sheaf and
\begin{equation}
 \operatorname{Psh}_{S,1}
   =\operatorname{Psh}_S
   =\mathscr J_{-\infty}.
 \label{eq:p071-infinite-first-stage-collapse}
\end{equation}
More precisely, suppose that \(-\infty\in W\) and
\(u\in\mathscr J_{-\infty}(W)\).  There is a fixed open neighbourhood
\(V=[-\infty,r)\subseteq W\) and functions
\(F_n\in\operatorname{Psh}_{S,0}(V)\) such that
\begin{equation}
 F_{n+1}\leq F_n,
 \qquad F_n\longrightarrow u|_V
 \quad\text{pointwise on }V.
 \label{eq:p071-infinite-fixed-domain-approximation}
\end{equation}
Each \(F_n\) is exactly constant on a neighbourhood of \(-\infty\).
If \(u(-\infty)\in\mathbf R\), then \(u\) is real-valued and continuous
on \(V\), and the convergence can be chosen uniform on \(V\).  The case
\(u(-\infty)=-\infty\) does not force \(u\) to be identically
\(-\infty\): for example, \(u(t)=t\) on the finite stratum, with
\(u(-\infty)=-\infty\), belongs to all three sheaves in
\eqref{eq:p071-infinite-first-stage-collapse}.

For the interval \(I=[-\infty,0)\) in
\cite[Example~6.7.2]{CLD25}, this is the constant-near-
\(-\infty\) smooth structure and
\eqref{eq:p071-infinite-stage-zero} is the condition
\(\mathrm d'\mathrm d''f\geq0\).  Thus the theorem applies to the
one-dimensional G-tropical model of the trivially valued open unit disc
specified in that example.  It makes no assertion about stars, general
G-tropical spaces, or the global psh hierarchy outside this interval
model.
\end{theorem}

\begingroup
\emergencystretch=2em
% Exact isolated candidate proof for source-aware verification.
% This file makes no novelty or priority claim and is not a live-book input.

\begin{proof}
\medskip\noindent\emph{Imported source scope.}
In \cite[Example~6.7.2]{CLD25}, Chambert--Loir and Ducros equip an
interval \(I\subseteq\overline{\mathbf R}\) with a G-tropical affine
structure for which parlinear functions are constant near every endpoint
belonging to the interval.  On a connected domain meeting such an
endpoint, every affine function is constant.  Sections~6.1.1 and 6.1.6
of \cite{CLD25} make tropical charts affine and their scalar coordinate
functions affine.  Together with \cite[Definition~6.2.1]{CLD25}, this
gives exactly \eqref{eq:p071-infinite-collar-smoothness} near the endpoint
\(-\infty\); away from that endpoint one uses the ordinary real affine
coordinate.  Conversely, an ordinary smooth function on the finite
stratum that is exactly constant on a terminal tail is represented by
the constant chart on that tail and by the ordinary coordinate on
endpoint-free neighbourhoods.  On the finite stratum,
\[
 \mathrm d'\mathrm d''f
       =f''(t)\,\mathrm d't\wedge\mathrm d''t,
\]
and the form vanishes on the terminal tail.  Hence positivity is exactly
\eqref{eq:p071-infinite-stage-zero}.  Definition~8.1.1 and
Lemma~8.1.4 of \cite{CLD25} give the generated sheaf and ordinal
filtration used in the statement.

\medskip\noindent\emph{Convexity and the endpoint trace.}
We first record two elementary consequences of the definition of
\(\mathscr J_{-\infty}\).  A locally extended-real convex function on a
connected ordinary interval is either finite everywhere or identically
\(-\infty\).  Indeed, the set where its value is \(-\infty\) is open:
if \(u(a)=-\infty\), express each nearby point as a strict convex
combination of \(a\) and a point on its other side.  The finite locus is
also open, since a finite value lying strictly between a nearby
\(-\infty\) value and a point on the other side would violate convexity.
The two loci are disjoint and cover a connected interval.

Suppose now that \(u\) is finite and convex on a tail
\((-\infty,b)\) and that
\begin{equation}
             c=\lim_{t\to-\infty}u(t)
                    \in\mathbf R\cup\{-\infty\}.
 \label{eq:p071-infinite-finite-side-limit}
\end{equation}
Then \(u\) is nondecreasing.  Otherwise there would be \(a<d<b\) with
\[
                 q=\frac{u(d)-u(a)}{d-a}<0.
\]
For every \(x<a\), monotonicity of secant slopes for a convex function
would give
\[
       \frac{u(a)-u(x)}{a-x}\leq q,
       \qquad
       u(x)\geq u(a)-q(a-x)\longrightarrow+\infty
                    \quad(x\to-\infty),
\]
contrary to \eqref{eq:p071-infinite-finite-side-limit}.  The same
nondecreasing conclusion is automatic when \(u\equiv-\infty\).

These observations also prove that \(\mathscr J_{-\infty}\) is a sheaf.
The ordinary convexity inequalities glue by the one-dimensional
local-to-global criterion.  Explicitly, on a compact interval take a
Lebesgue number for a finite cover by intervals of local convexity, choose
a partition whose consecutive triples lie in members of the cover, and
use the monotonicity of the resulting successive secant slopes.  The
dichotomy just proved handles extended-real values.  Finally, a member of
an open cover that contains \(-\infty\) contains a full terminal tail, so
its local equality \eqref{eq:p071-infinite-endpoint-trace} is already the
endpoint equality for the glued section.

The sheaf \(\mathscr J_{-\infty}\) is stable under pointwise decreasing
limits.  To see the only nonformal point, let \(u_j\downarrow u\) on an
endpoint component and put \(c_j=u_j(-\infty)\) and
\(c=u(-\infty)=\inf_jc_j\).  Every \(u_j\), and hence \(u\), is
nondecreasing, so
\begin{equation}
 c\leq L:=\lim_{t\to-\infty}u(t).
 \label{eq:p071-infinite-one-trace-inequality}
\end{equation}
If \(c\in\mathbf R\), then for every \(\varepsilon>0\) one can choose
\(j\) with \(c_j<c+\varepsilon\).  By the endpoint trace of \(u_j\), for
all sufficiently negative \(t\) one has
\(u(t)\leq u_j(t)<c+2\varepsilon\), so \(L\leq c\).  If
\(c=-\infty\), then, given \(M\in\mathbf R\), choose \(j\) with
\(c_j<M-1\) and then a sufficiently negative \(t\) so that
\(u(t)\leq u_j(t)<M\).  Monotonicity makes the same estimate valid farther
to the left, whence \(L=-\infty\).  Together with
\eqref{eq:p071-infinite-one-trace-inequality}, this proves preservation of
\eqref{eq:p071-infinite-endpoint-trace}.  Convexity passes directly to a
decreasing limit because no \(+\infty\) values occur.

Every member of \(\operatorname{Psh}_{S,0}\) lies in
\(\mathscr J_{-\infty}\): it is convex on the finite stratum, is constant
on a terminal tail, and has the required endpoint trace.  Since the latter
is a sheaf stable under decreasing limits, minimality gives
\begin{equation}
                 \operatorname{Psh}_S
                    \subseteq\mathscr J_{-\infty}.
 \label{eq:p071-infinite-generated-in-trace-class}
\end{equation}

\medskip\noindent\emph{Heat regularization with an exact terminal collar.}
Let \(-\infty\in W\), let \(u\in\mathscr J_{-\infty}(W)\), and shrink
the endpoint component so that
\([-\infty,R)\subseteq W\) for some \(R\in\mathbf R\).  If
\(u\equiv-\infty\) on the finite stratum of this tail, then the endpoint
trace gives \(u(-\infty)=-\infty\), and the smooth psh constants \(-n\)
give \eqref{eq:p071-infinite-fixed-domain-approximation} on every smaller
fixed tail.

We may therefore suppose that \(u\) is finite and convex on
\((-\infty,R)\).  Put \(c=u(-\infty)\).  By
\eqref{eq:p071-infinite-endpoint-trace} and the preceding argument, \(u\)
is nondecreasing.  Choose
\[
                       r<R_1<R_2<R
\]
and define
\[
              m=\frac{u(R_2)-u(R_1)}{R_2-R_1}\geq0.
\]
If \(c\in\mathbf R\), finite convexity gives continuity on
\((-\infty,R)\), while the trace equality gives continuity at the
endpoint \(-\infty\) in the order topology.  This proves the continuity
assertion in the statement.

Extend \(u\) to a finite function \(\bar u:\mathbf R\to\mathbf R\) by
\begin{equation}
 \bar u(t)=
 \begin{cases}
     u(t),&t\leq R_1,\\
     u(R_1)+m(t-R_1),&t\geq R_1.
 \end{cases}
 \label{eq:p071-infinite-convex-extension}
\end{equation}
The monotonicity of secant slopes shows that this extension is convex and
nondecreasing.  Every one of its secant slopes belongs to \([0,m]\), so
\begin{equation}
              |\bar u(x)-\bar u(y)|\leq m|x-y|
                    \quad(x,y\in\mathbf R).
 \label{eq:p071-infinite-global-lipschitz}
\end{equation}

For \(s>0\), set
\begin{equation}
 \gamma_s(y)=\frac{1}{\sqrt{4\pi s}}
       \mathrm e^{-y^2/(4s)},
 \qquad
 A_s(t)=\int_{\mathbf R}\bar u(t-y)\gamma_s(y)\,\mathrm dy.
 \label{eq:p071-infinite-heat-regularization}
\end{equation}
The linear-growth bound in
\eqref{eq:p071-infinite-global-lipschitz} makes the integral finite.
Every \(A_s\) is smooth, convex, and nondecreasing.  The Gaussian
semigroup identity and Jensen's inequality show that
\begin{equation}
             0<s'<s\quad\Longrightarrow\quad A_{s'}\leq A_s.
 \label{eq:p071-infinite-heat-monotonicity}
\end{equation}
Indeed, \(A_s=P_{s-s'}A_{s'}\), and averaging a convex function against
a centred probability measure can only increase it.  Moreover,
\begin{equation}
 0\leq A_s(t)-\bar u(t)
 \leq m\int_{\mathbf R}|y|\gamma_s(y)\,\mathrm dy
 =\frac{2m}{\sqrt{\pi}}\sqrt{s}.
 \label{eq:p071-infinite-heat-error}
\end{equation}

Choose a nonnegative even
\(\theta\in C^\infty_c((-1,1))\) of integral one.  For \(\eta>0\), put
\(\theta_\eta(s)=\eta^{-1}\theta(s/\eta)\) and
\begin{equation}
 M_\eta(a,b)=\int_{\mathbf R^2}
    \max\{a+s_1,b+s_2\}\theta_\eta(s_1)\theta_\eta(s_2)
       \,\mathrm ds_1\,\mathrm ds_2.
 \label{eq:p071-infinite-exact-maximum}
\end{equation}
This is a smooth convex function, nondecreasing in each variable, and
\begin{align}
 M_\eta(a,b)&=a &&\text{if }a>b+2\eta,
 \notag\\
 M_\eta(a,b)&=b &&\text{if }b>a+2\eta,
 \label{eq:p071-infinite-exact-maximum-regions}\\
 \max\{a,b\}&\leq M_\eta(a,b)\leq\max\{a,b\}+\eta.
 \notag
\end{align}
Fix \(\eta>0\) small enough that \(r+3\eta<R_1\), and define, for
\(n\geq1\),
\begin{equation}
 H_n(t)=M_\eta(t,-n)\quad(t\in\mathbf R),
 \qquad H_n(-\infty)=-n.
 \label{eq:p071-infinite-terminal-clamp}
\end{equation}
Then \(H_n\) is smooth, convex, and nondecreasing on the finite stratum,
is exactly equal to \(-n\) for \(t<-n-2\eta\), and is exactly equal to
\(t\) for \(t>-n+2\eta\).  Since \(M_\eta\) is nondecreasing in its
second variable,
\begin{equation}
                         H_{n+1}\leq H_n.
 \label{eq:p071-infinite-clamp-monotonicity}
\end{equation}

On the fixed open tail \(V=[-\infty,r)\), set
\begin{equation}
 F_n(t)=A_{1/n}(H_n(t))\quad(t\in V\cap\mathbf R),
 \qquad
 F_n(-\infty)=A_{1/n}(-n).
 \label{eq:p071-infinite-approximants}
\end{equation}
The function \(F_n\) is exactly constant on the endpoint collar
\(V\cap[-\infty,-n-2\eta)\).  On the finite stratum it is smooth and
\[
 F_n''
  =A_{1/n}''(H_n)(H_n')^2+A_{1/n}'(H_n)H_n''\geq0.
\]
Thus \(F_n\in\operatorname{Psh}_{S,0}(V)\).  By
\eqref{eq:p071-infinite-heat-monotonicity},
\eqref{eq:p071-infinite-clamp-monotonicity}, and monotonicity of
\(A_{1/n}\),
\begin{align}
 F_{n+1}
 &=A_{1/(n+1)}\circ H_{n+1}
 \leq A_{1/n}\circ H_{n+1}
 \leq A_{1/n}\circ H_n=F_n.
 \label{eq:p071-infinite-approximants-decrease}
\end{align}

For every finite \(t<r\), one has \(H_n(t)=t\) for all sufficiently large
\(n\).  Equation \eqref{eq:p071-infinite-heat-error} therefore gives
\(F_n(t)\to u(t)\).  For all sufficiently large \(n\), the point \(-n\)
lies in \((-\infty,R_1]\), so \(\bar u(-n)=u(-n)\).  At the endpoint, the
same estimate then gives
\begin{equation}
 |F_n(-\infty)-u(-n)|
   =|A_{1/n}(-n)-\bar u(-n)|
   \leq\frac{2m}{\sqrt{\pi n}},
 \label{eq:p071-infinite-endpoint-convergence}
\end{equation}
while \(u(-n)\to c=u(-\infty)\) in
\(\mathbf R\cup\{-\infty\}\).  This proves pointwise convergence on
the fixed tail, including when \(c=-\infty\).

If \(c\in\mathbf R\), the convergence is uniform on \(V\).  Indeed, for
all large \(n\), \(\bar u=u\) on the range \(H_n(V)\).  When
\(t>-n+2\eta\), the clamp is the identity.  When
\(t\leq-n+2\eta\), the bounds in
\eqref{eq:p071-infinite-exact-maximum-regions} give
\[
                 t\leq H_n(t)\leq-n+3\eta,
\]
and monotonicity gives
\[
 0\leq u(H_n(t))-u(t)\leq u(-n+3\eta)-c.
\]
For finite \(t\), the inequalities \(H_n(t)\geq t\),
\(A_{1/n}\geq\bar u\), and monotonicity of \(\bar u\) also give
\(F_n(t)\geq u(t)\).  At the endpoint,
\[
 0\leq F_n(-\infty)-c
 \leq\frac{2m}{\sqrt{\pi n}}+u(-n)-c
 \leq\frac{2m}{\sqrt{\pi n}}+u(-n+3\eta)-c.
\]
Together with \eqref{eq:p071-infinite-heat-error}, these estimates yield
\[
 0\leq\sup_{t\in V}\bigl(F_n(t)-u(t)\bigr)
 \leq\frac{2m}{\sqrt{\pi n}}
       +u(-n+3\eta)-c\longrightarrow0.
\]

\medskip\noindent\emph{Interior points and conclusion.}
At an interior point, the extended-real dichotomy reduces matters either
to the constants \(-n\), or to a finite convex function on an ordinary
interval.  In the latter case, restrict to a compact subinterval and
extend the function to a finite convex function on \(\mathbf R\) by
affine supporting lines at its two endpoints.  Its heat regularizations
\(A_{1/n}\) are smooth and convex, decrease by
\eqref{eq:p071-infinite-heat-monotonicity}, and converge to the original
function.  Thus every section of \(\mathscr J_{-\infty}\) belongs locally
to \(\operatorname{Psh}_{S,1}\).  Hence
\[
 \mathscr J_{-\infty}\subseteq\operatorname{Psh}_{S,1}
 \subseteq\operatorname{Psh}_S.
\]
Combining this with
\eqref{eq:p071-infinite-generated-in-trace-class} proves
\eqref{eq:p071-infinite-first-stage-collapse} and all the stated
approximation assertions.  Finally, \(u(t)=t\), with endpoint value
\(-\infty\), satisfies the trace condition and is obtained explicitly
by the same construction.
\end{proof}

\endgroup

\begin{remark}
\label{rem:ch16-collapse-scope-firewall}
The four results have different exact scopes.  In
\Cref{prop:ch16-collapse-tropical-toric}, the space may have tropical
boundary, but it must be a tropical toric partial compactification equipped
with the standard Lagerberg structure to which the cited regularization
theorem applies.  In
\Cref{thm:ch16-collapse-finite-tropical-curves}, the inputs are finitely many
pairwise distinct rays in a finite-dimensional real vector space and the
stated identification of the stage-zero generator sheaf.  No lattice,
rationality, weight, balancing, calibration, harmony, or boundary
designation is used.  Thus ordinary ambient-induced one-valent endpoints,
unbalanced finite stars, and polar branches are included.

This last assertion is not a statement about every structure that the
literature calls a tropical boundary.  A finite star in the theorem has
sedentarity zero and the ambient-induced smooth sheaf.  The theorem does not
automatically apply to positive-sedentarity toric charts, which are covered
separately by \Cref{prop:ch16-collapse-tropical-toric}.  The two
constant-near-an-endpoint interval structures of
\cite[Example~6.7.2]{CLD25} are covered separately by
\Cref{thm:p071-cld-constant-collar-endpoint,%
thm:p071-cld-infinite-collar-endpoint}.  None of these results covers the
infinite-star structures of \cite[\S6.7.3]{CLD25}.

Finiteness is essential to the proof: it permits simultaneous shrinking of
all branches and smoothing through finitely many affine pieces.  In
dimensions at least two, the analogous Jensen reduction leaves a genuine
face-compatibility problem for polyhedral regularization.  The
regularity property studied by Botero--Mazzon--Pille-Schneider isolates the
piecewise-affine part of that problem and is established in dimension one,
but piecewise-affine approximation is not by itself approximation by the
smooth functions defining CLD stage zero
\cite[Theorem~6.4, Remark~6.6 and Proposition~6.13]{BMPS26}.  The
regularized-maximum construction above supplies that bridge for the finite
ambient-induced one-dimensional model, while the two endpoint theorems use
the order structure of a single interval and do not couple branches.
Arbitrary higher-dimensional settings, infinite-valence stars, and general
nonambient one-dimensional smooth structures therefore remain within
\Cref{prob:frontier-cld-hierarchy}.
\end{remark}

\subsection{Composition, maxima, and tropical convexity}
\label{sec:ch16-composition}

The next three statements are the working tools for producing psh
functions.  They are stated with the filtration index because that index is
what \Cref{sec:ch16-polar} needs.

\begin{proposition}
\label{prop:ch16-composition}
Let $u_1,\ldots,u_m$ be psh on $X$, let $\Omega\subseteq\mathbf R^m$ be open
and convex with $(u_1(x),\ldots,u_m(x))\in\Omega$ for all $x$, and let
$g:\Omega\to\mathbf R$ be smooth, convex, and nondecreasing in each
variable.  Then $g(u_1,\ldots,u_m)$ is psh, and it lies in
$\operatorname{Psh}_{X,\alpha}$ as soon as every $u_i$ does.
\end{proposition}

\begin{proof}
Transfinite induction on $\alpha$.  For $\alpha=0$ all $u_i$ are smooth,
hence so is $u=g(u_1,\ldots,u_m)$, and the chain rule gives
\begin{equation}
\label{eq:ch16-composition-hessian}
 \dprimeddprime u
 =\sum_{i,j}\frac{\partial^2g}{\partial y_i\partial y_j}(u)\,
    \dprime u_i\wedge\ddprime u_j
 +\sum_j\frac{\partial g}{\partial y_j}(u)\,\dprimeddprime u_j .
\end{equation}
The first sum is positive because the Hessian of $g$ is positive
semidefinite: writing that Hessian as $\sum_k\lambda_k w_k\otimes w_k$ with
$\lambda_k\geq0$, the sum becomes
$\sum_k\lambda_k\,\dprime\ell_k\wedge\ddprime\ell_k$ with
$\ell_k=\sum_iw_{k,i}u_i$, and each $\dprime\ell\wedge\ddprime\ell$ is
positive by \Cref{prop:ch09-positivity-dictionary}.  The second sum is a
nonnegative combination of the positive forms $\dprimeddprime u_j$.  Hence
$\dprimeddprime u\geq0$ and $u\in\operatorname{Psh}_{X,0}$.

For a limit ordinal, locality of the filtration reduces to a smaller
$\beta$.  For a successor $\alpha=\beta+1$, choose locally decreasing
sequences $v_{i,k}\downarrow u_i$ as $k\to\infty$, with
$v_{i,k}\in\operatorname{Psh}_{X,\beta}$.  Since $g$ is continuous and
nondecreasing in each variable, $g(v_{1,k},\ldots,v_{m,k})$ decreases to
$g(u_1,\ldots,u_m)$, and each term lies in $\operatorname{Psh}_{X,\beta}$ by
induction.  This is \cite[Proposition~8.2.1]{CLD25}.
\end{proof}

\begin{corollary}
\label{cor:ch16-max-psh}
Let $g:\Omega\to\mathbf R$ be convex and nondecreasing in each variable, and
let $u_1,\ldots,u_m$ be continuous psh with values in $\Omega$.  Then
$g(u_1,\ldots,u_m)$ is continuous psh, and it lies in
$\operatorname{CPsh}_{X,\sup(1,\alpha)}$ if all $u_i$ lie in
$\operatorname{CPsh}_{X,\alpha}$.  In particular
$\max\{u_1,\ldots,u_m\}$ is continuous psh, and lies in
$\operatorname{Psh}_{X,\sup(1,\alpha)}$ whenever all $u_i$ lie in
$\operatorname{Psh}_{X,\alpha}$.
\end{corollary}

\begin{proof}
A convex function on an open convex set is continuous, so
$g(u_1,\ldots,u_m)$ is continuous.  Fix $x$ and a relatively compact convex
neighbourhood $\Omega'$ of $(u_1(x),\ldots,u_m(x))$ in $\Omega$.  On
$\Omega'$, convolution with a smooth kernel writes $g$ as a uniform limit of
smooth functions $g_k$ that are again convex and nondecreasing in each
variable; that the two properties survive convolution is immediate, since
both are preserved by averaging translates.  By
\Cref{prop:ch16-composition} each $g_k(u_1,\ldots,u_m)$ lies in
$\operatorname{Psh}_{X,\alpha}$, and by
\Cref{prop:ch16-hierarchy-status}\,\cref{itm:ch16-hier:cont} the uniform
limit lies in $\operatorname{CPsh}_{X,\sup(1,\alpha)}$.  For the last
assertion use the smoothing
\begin{equation}
\label{eq:ch16-softmax}
 \max\{y_1,\ldots,y_m\}
 \;\leq\;
 a^{-1}\log\bigl(\rme^{ay_1}+\cdots+\rme^{ay_m}\bigr)
 \;\leq\;
 \max\{y_1,\ldots,y_m\}+a^{-1}\log m ,
\end{equation}
whose middle term is smooth, convex, and nondecreasing in each variable, and
converges uniformly to the maximum as $a\to+\infty$.  Convexity of the
log-sum-exp function is the computation recorded after
\Cref{cor:superform-convex-composition}: its Hessian is $a$ times the
covariance matrix of the coordinate directions under the probability weights
$\rme^{ay_k}/\sum_j\rme^{ay_j}$, hence positive semidefinite; monotonicity in
each variable is clear because each partial derivative is one of those
weights.  The two inequalities in \eqref{eq:ch16-softmax} are immediate from
$\rme^{a\max_ky_k}\leq\sum_k\rme^{ay_k}\leq m\,\rme^{a\max_ky_k}$.  The
resulting statement is \cite[Corollaries~8.2.2--8.2.4]{CLD25}.
\end{proof}

\begin{proposition}
\label{prop:ch16-tropically-convex}
\index{tropically convex function}
Let $f:X\to E$ be a tropical chart, let $P$ be a parlinear subspace of $E$
containing $f(X)$, and let $v$ be a continuous convex function on $P$.  Then
$f^*v$ is continuous psh.  A function that is locally of this form is called
\emph{tropically convex}; a local difference of two such functions is called
\emph{tropically $d$-convex}.
\end{proposition}

\begin{proof}
Consider the subsheaf of the sheaf of continuous functions on $P$ consisting
of those $v$ for which $f^*v$ is psh.  It contains every smooth $v$ with
$\dprimeddprime v\geq0$, because pullback along a morphism of $G$-tropical
spaces preserves smoothness and positivity of the Hessian
(\Cref{prop:cld-form-pullback}), and it is stable under uniform limits by
\Cref{prop:ch16-hierarchy-status}\,\cref{itm:ch16-hier:cont}.  Every
continuous convex function on $P$ is a local uniform limit of smooth convex
ones by \Cref{cor:ch09-monotone-smooth-approximation}, so the subsheaf
contains all of them.  This is \cite[Propositions~8.2.5
and~8.2.6]{CLD25}.
\end{proof}

\begin{proposition}
\label{prop:ch16-dpsh-smooth}
Every $G$-smooth function on $X$ is locally a difference of two tropically
convex functions; in particular it is locally a difference of two continuous
psh functions.  If it is parlinear, the two summands may be taken
parlinear.
\end{proposition}

\begin{proof}
The assertion is local, so fix $x$, a compact analytic domain $W$ containing
a neighbourhood of $x$, a tropical chart $f:W\to E$ tropicalizing $u$, and a
$G$-smooth $v$ near the compact polyhedron $f(W)$ with $u=f^*v$.  On a
neighbourhood of a compact polyhedron, a smooth function is a difference of
two smooth convex functions: choose $c>0$ so large that $v+c\|t\|^2$ has
positive semidefinite Hessian on a neighbourhood of the compact polyhedron,
which is possible because the Hessian of $v$ is bounded there, and write
$v=(v+c\|t\|^2)-c\|t\|^2$, both summands being smooth and convex.  Both summands pull back to tropically
convex functions by \Cref{prop:ch16-tropically-convex}.  The parlinear case
uses instead the decomposition of a parlinear function into a difference of
two convex parlinear functions on the convex hull of $f(W)$.  These are
\cite[Lemma~8.8.1 and Proposition~8.8.2]{CLD25}.
\end{proof}

\begin{definition}[dpsh functions]
\label{def:ch16-dpsh}
\index{dpsh function}
Let $\operatorname{CPsh}^\pm_X$ denote the subsheaf of $\mathbf R$-vector
spaces of the sheaf of continuous functions on $X$ generated by
$\operatorname{CPsh}_X$.  Its sections are called \emph{dpsh}.  By
\Cref{prop:ch16-dpsh-smooth} every $G$-smooth function is dpsh.
\end{definition}

\subsection{Four properties that must not be confused}
\label{sec:ch16-four-properties}

\begin{center}
\small
\begin{tabular}{p{.24\textwidth}|p{.36\textwidth}|p{.31\textwidth}}
\textbf{Property}&\textbf{Definition}&\textbf{Status in this chapter}\\
\hline
psh
 &section of $\operatorname{Psh}_X$
 &too weak for the recursion: $uT$ need not exist\\
continuous psh
 &section of $C^0_X\cap\operatorname{Psh}_X$
 &the operative class; all of
  \Crefrange{sec:ch16-one-factor}{sec:ch16-cln}\\
locally approachable
 &section of $\operatorname{Psh}_{X,1}$
 &needed only for polar negligibility,
  \Cref{thm:cld-polar-negligibility}\\
locally bounded
 &bounded on compact sets
 &\emph{not} used: no locally bounded theory is asserted\\
\end{tabular}
\end{center}

The fourth line deserves emphasis.  In complex pluripotential theory the
Bedford--Taylor product is constructed for locally bounded psh functions
\cite{BT76,BT82}, and continuity is a convenience, not a hypothesis.  The
reason is that a locally bounded complex psh function is a decreasing limit
of smooth psh functions, so the monotone convergence theorem for the product
applies.  In the CLD setting the analogous smoothing is unavailable
(\Cref{prop:ch16-hierarchy-status}\,\cref{itm:ch16-hier:approach}), and the
substitute---the sheaf-closure principle of the next section---requires
\emph{uniform} rather than merely monotone approximation.  Uniform limits
preserve continuity.  This is why the CLD product is a product of continuous
potentials, and why the reader must not import complex statements verbatim.

\begin{warning}{A locally bounded psh function need not be continuous, and
the recursion does not see it}
Nothing above forbids a bounded discontinuous psh function.  If $u$ is such
a function and $T$ a positive closed current, then $uT$ is not defined by
the order-zero pairing, because the pairing extends $T$ only to
\emph{continuous} compactly supported forms.  One may of course integrate
$u$ against the Radon measures that are the coefficients of $T$, but the
resulting object is not known to be weakly continuous along approximations,
and none of \Crefrange{thm:one-bt-factor}{thm:cld-cln} is asserted for it.
In particular the classical convergence theorem of \cite{BT82} for
decreasing sequences of locally bounded psh functions is \emph{not} being
transported.
\end{warning}

\section{Multiplying a positive current by a potential}
\label{sec:ch16-one-factor}

\subsection{The product \texorpdfstring{$uT$}{uT}}
\label{sec:ch16-uT}

Let $T$ be a positive current of bidegree $(p,p)$.  By
\Cref{prop:current-positivity-calculus}, $T$ has order zero on fixed compact
chart data: for every compact $K$ and every symmetric test form $\alpha$
supported in $K$ one has
\begin{equation}
\label{eq:ch16-order-zero}
 |\langle T,\alpha\rangle|
 \leq C_K\,\|\alpha\|_{C^0(K)} .
\end{equation}
Hence $T$ extends by continuity from smooth compactly supported test forms
to \emph{continuous} ones, and for $u$ continuous one may set
\begin{equation}
\label{eq:ch16-uT-definition}
 \langle uT,\alpha\rangle:=\langle T,u\alpha\rangle ,
\end{equation}
the right-hand side being evaluated through that order-zero extension, since
$u\alpha$ need not be smooth.  The resulting current $uT$ has the same
bidegree as $T$, and is positive whenever $u\geq0$.

There is nothing to prove in \eqref{eq:ch16-uT-definition}, but two remarks
prevent later confusion.  First, $uT$ depends on the values of $u$ only on
$\Supp(T)$.  Second, the assignment $u\mapsto uT$ is \emph{not} continuous
for weak convergence of $u$ against Lebesgue measure; it is continuous for
local uniform convergence, and that is exactly the topology in which the
whole chapter operates.

\begin{lemma}
\label{lem:continuous-times-positive-current}
Suppose $(T_a)$ is a weakly convergent net of positive currents with limit
$T$, and $(u_a)$ is a net of continuous functions converging locally
uniformly to $u$.  Then
\[
                              u_aT_a\longrightarrow uT.
\]
\end{lemma}
\begin{proof}
It suffices to work near the compact support of a test form $\alpha$.  Choose
a smooth $v$ with $\|u-v\|_\infty<\varepsilon$ and write
\[
 u_aT_a-uT=(u_a-v)T_a+v(T_a-T)+(v-u)T.
\]
Express $\alpha$ as a difference of two positive test forms.  Positivity
bounds the first and third terms by the uniform errors times the masses of
$T_a$ and $T$ on a fixed positive test form; these masses are eventually
bounded because $T_a\to T$.  The middle term tends to zero by weak
convergence.  First let $a$ tend along its directed set, then let
$\varepsilon\to0$.
\end{proof}

This is \cite[Lemma~8.5.1]{CLD25}.  We spell out the three steps, since the
same decomposition recurs in
\Crefrange{thm:cld-bt-continuity}{thm:cld-cln}.

\begin{proof}[Second proof, with the constants displayed]
Fix a smooth test form $\alpha$ with compact support $S$, and a compact
neighbourhood $K$ of $S$.  Write $\alpha=\alpha^+-\alpha^-$ with
$\alpha^\pm$ smooth, positive, supported in the interior of $K$
(\Cref{prop:ch09-positivity-dictionary} supplies such a decomposition
chartwise, and a partition of unity globalizes it), and put
$\beta=\alpha^++\alpha^-$.  For any continuous $w$ and any positive current
$S'$,
\begin{equation}
\label{eq:ch16-positivity-sandwich}
 |\langle wS',\alpha\rangle|
 \leq\|w\|_{C^0(K)}\,\langle S',\beta\rangle ,
\end{equation}
because $-\|w\|_{C^0(K)}\beta\leq w\alpha\leq\|w\|_{C^0(K)}\beta$ as
continuous forms and $S'$ is positive.  Now let $\varepsilon>0$ and pick a
smooth $v$ with $\|u-v\|_{C^0(K)}<\varepsilon$; such a $v$ exists by
\Cref{thm:cld-smooth-partition} and local uniform density of smooth
functions on chart data.  Then
\[
 |\langle u_aT_a-uT,\alpha\rangle|
 \leq\underbrace{\|u_a-v\|_{C^0(K)}\langle T_a,\beta\rangle}_{(\mathrm I)}
   +\underbrace{|\langle T_a-T,v\alpha\rangle|}_{(\mathrm{II})}
   +\underbrace{\|v-u\|_{C^0(K)}\langle T,\beta\rangle}_{(\mathrm{III})}.
\]
Term $(\mathrm{II})$ tends to $0$ because $v\alpha$ is a smooth test form.
The numbers $\langle T_a,\beta\rangle$ converge to $\langle T,\beta\rangle$,
so they are eventually bounded by $2\langle T,\beta\rangle+1=:M$.  Since
$\|u_a-v\|_{C^0(K)}\leq\|u_a-u\|_{C^0(K)}+\varepsilon\to\varepsilon$, we get
$\limsup_a(\mathrm I)\leq M\varepsilon$ and
$(\mathrm{III})\leq\varepsilon\langle T,\beta\rangle$.  Hence
\[
 \limsup_a|\langle u_aT_a-uT,\alpha\rangle|
 \leq\varepsilon\bigl(M+\langle T,\beta\rangle\bigr),
\]
and $\varepsilon$ was arbitrary.
\end{proof}

\subsection{The sheaf-closure principle}
\label{sec:ch16-sheaf-closure}

Almost every statement in
\Crefrange{sec:ch16-one-factor}{sec:ch16-continuity} is proved by verifying
a property first for smooth psh potentials and then observing that the set
of potentials for which it holds is a sheaf stable under local uniform
limits.  We isolate that observation once.

\begin{lemma}
\label{lem:ch16-cpsh-induction}
Let $\mathscr P$ be a subpresheaf of the sheaf $C_X$ of continuous
real-valued functions on $X$ such that
\begin{enumerate}
\item\label{itm:ch16-closure:sheaf} $\mathscr P$ is a sheaf: membership is a
 local condition;
\item\label{itm:ch16-closure:smooth} $\mathscr P(U)$ contains every smooth
 $u$ on $U$ with $\dprimeddprime u\geq0$;
\item\label{itm:ch16-closure:unif} $\mathscr P(U)$ is stable under local
 uniform limits.
\end{enumerate}
Then $\mathscr P$ contains $\operatorname{CPsh}_X$; that is, every
continuous psh function on every open subset of $X$ lies in $\mathscr P$.
\end{lemma}

\begin{proof}
By definition $\operatorname{CPsh}_X$ is the smallest subsheaf of $C_X$
satisfying \cref{itm:ch16-closure:smooth} and \cref{itm:ch16-closure:unif}.
Since $\mathscr P\cap\operatorname{CPsh}_X$ is again a subsheaf of $C_X$
satisfying both, minimality gives
$\mathscr P\cap\operatorname{CPsh}_X=\operatorname{CPsh}_X$.
\end{proof}

\begin{warning}{The principle is not a smoothing theorem}
\Cref{lem:ch16-cpsh-induction} says that a property propagates \emph{through}
the generation process.  It does \emph{not} say that a given continuous psh
function is a uniform limit of smooth psh functions on a fixed neighbourhood.
The distinction is the content of
\Cref{prop:ch16-hierarchy-status}\,\cref{itm:ch16-hier:approach}: only the
functions in $\operatorname{CPsh}_{X,1}$ are known to admit such an
approximation.  A proof that verified a property for smooth psh functions and
then said ``approximate an arbitrary continuous psh function by smooth psh
functions'' would be circular; every proof below instead verifies the three
hypotheses of the lemma.  Whenever a hypothesis fails to be stable under
uniform limits---for instance a quantitative statement whose constant is not
uniform---the principle is unavailable and a different argument is needed.
This is why \Cref{thm:cld-cln} is proved by an inductive integration by parts
rather than by closure.
\end{warning}

\subsection{Symmetric closed currents and smooth factors}
\label{sec:ch16-symmetric-leibniz}

The base case of every induction is the following Leibniz identity.  It is
where closedness of $T$ and symmetry of positive currents are consumed.

\begin{lemma}
\label{lem:ch16-leibniz-closed}
Let $S$ be a closed symmetric current of bidegree $(p,p)$, let
$q\leq p-1$, and let $\alpha$ be a smooth symmetric $(q,q)$-form.  Then
\[
              \dprimeddprime(\alpha\wedge S)
              =(\dprimeddprime\alpha)\wedge S .
\]
\end{lemma}

\begin{proof}
Let $\beta$ be a smooth test form of bidegree $(p-q-1,p-q-1)$.  Both
$(\dprimeddprime\alpha)\wedge S$ and $\alpha\wedge S$ are symmetric, so it
suffices to treat symmetric $\beta$.  By the definition of the current
differential and of the wedge with a smooth form,
\[
 \langle\dprimeddprime\alpha\wedge S,\beta\rangle
 -\langle\alpha\wedge S,\dprimeddprime\beta\rangle
 =\bigl\langle S,\;
   \dprimeddprime\alpha\wedge\beta-\alpha\wedge\dprimeddprime\beta
  \bigr\rangle .
\]
Now expand
\[
 \dprime\bigl(\alpha\wedge\ddprime\beta
              -\ddprime\alpha\wedge\beta\bigr)
 =\dprime\alpha\wedge\ddprime\beta
  +\alpha\wedge\dprimeddprime\beta
  -\dprimeddprime\alpha\wedge\beta
  +\ddprime\alpha\wedge\dprime\beta ,
\]
using $(\dprime)^2=0$ and the sign rules of \Cref{chap:superforms}.  The
$(p,p)$-form
$\dprime\alpha\wedge\ddprime\beta+\ddprime\alpha\wedge\dprime\beta$ is
\emph{antisymmetric} for the involution $J$: interchanging the two exterior
factors exchanges the two summands and introduces one sign, exactly as in
\Cref{lem:ch09-J-conjugation}.  A symmetric current annihilates
antisymmetric forms, so pairing the displayed identity with $S$ gives
\[
 \bigl\langle S,\;
  \dprime\bigl(\alpha\wedge\ddprime\beta-\ddprime\alpha\wedge\beta\bigr)
 \bigr\rangle
 =\bigl\langle S,\;
   \alpha\wedge\dprimeddprime\beta-\dprimeddprime\alpha\wedge\beta
  \bigr\rangle .
\]
The left-hand side vanishes because $\dprime S=0$.  Hence
$\langle\dprimeddprime\alpha\wedge S,\beta\rangle
 =\langle\alpha\wedge S,\dprimeddprime\beta\rangle
 =\langle\dprimeddprime(\alpha\wedge S),\beta\rangle$.
This is \cite[Lemma~8.5.3]{CLD25}.
\end{proof}

\begin{corollary}
\label{cor:ch16-smooth-case}
If $T$ is positive and closed and $u$ is smooth, then
\[
 \dprimeddprime(uT)=\dprimeddprime u\wedge T=T\wedge\dprimeddprime u .
\]
In particular, if $u$ is smooth psh then $\dprimeddprime(uT)$ is positive
and closed.
\end{corollary}

\begin{proof}
Apply \Cref{lem:ch16-leibniz-closed} with $q=0$ and $\alpha=u$, noting that
a positive current is symmetric by definition and that $\dprime T=0$ implies
$\ddprime T=0$ by symmetry (\Cref{prop:current-positivity-calculus}).
Positivity of $\dprimeddprime u\wedge T$ is
\Cref{prop:current-positivity-calculus}\,(1), since $\dprimeddprime u$ is
strongly positive in bidegree $(1,1)$ when $u$ is smooth psh
(\Cref{prop:ch09-positivity-inclusions}).  Closedness follows from
$(\dprime)^2=(\ddprime)^2=0$.
\end{proof}

\subsection{The one-factor theorem}
\label{sec:ch16-one-factor-theorem}

\begin{theorem}
\label{thm:one-bt-factor}
Let $T$ be positive and closed, and let $u$ be continuous CLD-psh.  Then
\[
                         \dprimeddprime(uT)
\]
is a positive closed current.  If $u$ is smooth, it equals
$\dprimeddprime u\wedge T$.
\end{theorem}
\begin{proof}
For smooth $u$, closedness of $T$ and the Leibniz rule give
$\dprimeddprime(uT)=\dprimeddprime u\wedge T$, which is positive.  The class
of continuous functions for which $\dprimeddprime(uT)$ is positive is local,
contains smooth psh functions, and is closed under local uniform limits by
\Cref{lem:continuous-times-positive-current} and continuity of the
current differential.  By the defining sheafwise uniform-closure property
of $\operatorname{CPsh}$, this class therefore contains every continuous
CLD-psh function.  This argument does not choose a direct smooth
approximation of a higher-stage function.  Closedness follows from
$(\dprime)^2=(\ddprime)^2=0$.
\end{proof}

\begin{proof}[Second proof, verifying
\Cref{lem:ch16-cpsh-induction} step by step]
For an open $U\subseteq X$ let
\[
 \mathscr P(U)=\bigl\{u\in C(U):\ \dprimeddprime(uT|_U)\geq0\bigr\}.
\]
\emph{Sheaf property.}  Positivity of a current is a local condition
(\Cref{prop:current-positivity-calculus}), and restriction of currents is
compatible with multiplication by a continuous function, so $u\in\mathscr
P(U)$ as soon as $u|_{U_i}\in\mathscr P(U_i)$ for an open cover.  Thus
$\mathscr P$ is a subsheaf of $C_X$.

\emph{Smooth psh functions belong.}  This is
\Cref{cor:ch16-smooth-case}.

\emph{Stability under local uniform limits.}  Let $u_k\in\mathscr P(U)$
converge locally uniformly to $u$.  By
\Cref{lem:continuous-times-positive-current}, applied with the constant net
$T_a=T|_U$, one has $u_kT|_U\to uT|_U$ weakly.  The operator
$\dprimeddprime$ is weakly continuous on currents
(\Cref{prop:current-differential-calculus}), so
$\dprimeddprime(u_kT|_U)\to\dprimeddprime(uT|_U)$.  A weak limit of positive
currents is positive, again by
\Cref{prop:current-positivity-calculus}.  Hence $u\in\mathscr P(U)$.

\Cref{lem:ch16-cpsh-induction} now gives
$\operatorname{CPsh}_X\subseteq\mathscr P$.  Closedness of
$\dprimeddprime(uT)$ is formal.
\end{proof}

This is \cite[Lemma~8.5.5]{CLD25}.  It uses only
continuous psh potentials; a product with an unbounded potential would
require a separate integrability theory.

\begin{remark}[Where each hypothesis is used]
\label{rem:ch16-hypothesis-audit}
Positivity of $T$ is used twice: to define $uT$ through
\eqref{eq:ch16-order-zero}, and to make $\dprimeddprime u\wedge T$ positive
in the smooth case.  Closedness of $T$ is used once, in
\Cref{lem:ch16-leibniz-closed}, and it is essential: if $\dprime T\neq0$
then $\dprimeddprime(uT)$ acquires the extra terms
$\dprime u\wedge\ddprime T$ and $\ddprime u\wedge\dprime T$, which have no
sign.  Continuity of $u$ is used once, in
\eqref{eq:ch16-uT-definition}.  Boundarylessness of $X$ is not used in this
subsection at all; it enters first in
\Cref{sec:ch16-cln}, through Green's formula.
\end{remark}

\begin{proposition}
\label{prop:ch16-smooth-factor}
Let $T$ be positive closed and let $\gamma$ be a smooth closed positive
form.  For every continuous psh $u$,
\[
 \dprimeddprime\bigl(u\,(T\wedge\gamma)\bigr)
 =\bigl(\dprimeddprime(uT)\bigr)\wedge\gamma .
\]
\end{proposition}

\begin{proof}
Both sides are weakly continuous in $u$ for local uniform convergence, by
\Cref{lem:continuous-times-positive-current} and the fact that wedging with
a fixed smooth form is weakly continuous.  Both are local in $u$.  For
smooth psh $u$ the identity reads
$\dprimeddprime u\wedge T\wedge\gamma
 =(\dprimeddprime u\wedge T)\wedge\gamma$, which is associativity of the
wedge product with smooth forms.  Apply
\Cref{lem:ch16-cpsh-induction} to the sheaf of continuous $u$ satisfying
the identity.
\end{proof}

\section{Bedford--Taylor currents and harmonic ambiguity}
\label{sec:ch16-bt-currents}

There is a subtle point.  If $\dprimeddprime u=\dprimeddprime v$, it is not
formal for a general positive current $T$ that
$\dprimeddprime(uT)=\dprimeddprime(vT)$.  CLD isolate the currents for which
this independence holds.

\subsection{The smooth case, and what it does not give}
\label{sec:ch16-ambiguity-smooth}

Suppose $u,v$ are psh and $h=u-v$ happens to be smooth.  Then, by
\Cref{cor:ch16-smooth-case},
\begin{equation}
\label{eq:ch16-smooth-difference}
 \dprimeddprime(uT)=\dprimeddprime(hT)+\dprimeddprime(vT)
 =\dprimeddprime h\wedge T+\dprimeddprime(vT),
\end{equation}
so that $\dprimeddprime u=\dprimeddprime v$, which for smooth $h$ means
$\dprimeddprime h=0$, forces
$\dprimeddprime(uT)=\dprimeddprime(vT)$.  Nothing in this computation
survives when $h$ is merely continuous: the first equality of
\eqref{eq:ch16-smooth-difference} is still valid, since $u\mapsto uT$ is
linear, but the second is exactly the identity
$\dprimeddprime(hT)=\dprimeddprime h\wedge T$ whose validity is in question.
This is \cite[\S8.6.1]{CLD25}, where the authors record that they do not
know whether the conclusion persists in general.

\begin{definition}
\label{def:ch16-bt-current}
\index{Bedford--Taylor current}
A positive closed current $T$ is a \emph{continuous Bedford--Taylor
current} if, on every open set and for all continuous psh functions $u,v$
satisfying
\[
                     \dprimeddprime[u]=\dprimeddprime[v],
\]
one has
\[
                     \dprimeddprime(uT)=\dprimeddprime(vT).
\]
For such a current it is unambiguous to write
$\dprimeddprime u\wedge T$ for $\dprimeddprime(uT)$.
\end{definition}

Throughout this chapter all potentials in the Bedford--Taylor predicate are
continuous, and we abbreviate ``continuous Bedford--Taylor current'' to
``Bedford--Taylor current.''  This convention is narrower than the literal
wording of \cite[Definition~8.6.2]{CLD25}, which may quantify over
discontinuous psh functions.  No result below silently identifies the two
predicates.

We shall use the notation $\dprimeddprime u\wedge T$ even when $T$ is not
known to be Bedford--Taylor, as CLD do; every current constructed in this
book is Bedford--Taylor by \Cref{prop:bt-current-stability}, so the abuse is
harmless, but a reader importing a current from elsewhere should check the
property before using the notation.

\subsection{Stability of the class}
\label{sec:ch16-bt-stability}

\begin{proposition}
\label{prop:bt-current-stability}
The Bedford--Taylor property is local.  The integration current $\delta_X$
has it.  It is preserved by wedging with a smooth positive closed form, by
weak limits of positive closed Bedford--Taylor currents, and by
\[
                     T\longmapsto\dprimeddprime(uT)
\]
for continuous psh $u$.
\end{proposition}
\begin{proof}
Locality is the sheaf property.  For $\delta_X$,
$\dprimeddprime(u\delta_X)=\dprimeddprime[u]$.  Smooth closed factors commute
with $\dprimeddprime$ by the Leibniz rule, and the resulting current stays
positive by \Cref{prop:current-positivity-calculus}.  For a weakly convergent net,
\Cref{lem:continuous-times-positive-current} permits passage to the
limit in the defining equality.  For the last assertion, the functions for
which the claimed current is Bedford--Taylor form a sheaf containing the
smooth psh functions and stable under uniform limits.  The defining closure
property of $\operatorname{CPsh}$ gives the result without a direct smooth
regularization.  These assertions are
\cite[Definition~8.6.2 and Propositions~8.6.3, 8.6.5--8.6.8]{CLD25}.
\end{proof}

\begin{proof}[Second proof, item by item]
\emph{Locality.}  Let $(U_i)$ cover $X$ with $T|_{U_i}$ Bedford--Taylor.  If
$u,v$ are continuous psh on an open $V$ with
$\dprimeddprime u=\dprimeddprime v$, then
$\dprimeddprime(uT)=\dprimeddprime(vT)$ holds on each $U_i\cap V$, hence on
$V$ because currents form a sheaf (\Cref{chap:cld-currents}).  Conversely
the restriction of a Bedford--Taylor current is Bedford--Taylor, since the
defining condition quantifies over all opens.  Thus the Bedford--Taylor
currents form a subsheaf of the sheaf of positive closed currents.

\emph{The current $\delta_X$.}  For a continuous $u$ one has
$u\delta_X=[u]$ by definition of the integration current, so
$\dprimeddprime(u\delta_X)=\dprimeddprime[u]$, which depends only on
$\dprimeddprime[u]$ tautologically.

\emph{Smooth positive closed factors.}  Let $\gamma$ be smooth, positive and
closed, and let $u,v$ be continuous psh with
$\dprimeddprime u=\dprimeddprime v$.  By
\Cref{prop:ch16-smooth-factor},
$\dprimeddprime(u(T\wedge\gamma))=\dprimeddprime(uT)\wedge\gamma$ and
likewise for $v$; the two right-hand sides agree because $T$ is
Bedford--Taylor.  Positivity of $T\wedge\gamma$ is
\Cref{prop:current-positivity-calculus}\,(2), and it is closed because
$\dprime(T\wedge\gamma)=\dprime T\wedge\gamma\pm T\wedge\dprime\gamma=0$.

\emph{Weak limits.}  Let $T_k\to T$ weakly with each $T_k$ Bedford--Taylor.
The limit is positive and closed by
\Cref{prop:current-positivity-calculus} and weak continuity of the
differential.  For continuous psh $u$,
\Cref{lem:continuous-times-positive-current} gives $uT_k\to uT$, hence
$\dprimeddprime(uT_k)\to\dprimeddprime(uT)$; the same for $v$.  Passing to
the limit in $\dprimeddprime(uT_k)=\dprimeddprime(vT_k)$ gives the required
equality.

\emph{Iteration.}  Fix a Bedford--Taylor current $T$ and let
\[
 \mathscr P(U)=\bigl\{f\in C(U):\ \dprimeddprime(fT|_U)
   \text{ is a Bedford--Taylor current}\bigr\}.
\]
This is a sheaf, by locality just proved.  If $f$ is smooth psh, then
$\dprimeddprime(fT)=\dprimeddprime f\wedge T$ by
\Cref{cor:ch16-smooth-case}, and $\dprimeddprime f$ is a smooth form that is
positive and closed, the latter because
$\dprime\dprimeddprime f=\ddprime\dprimeddprime f=0$; the preceding item
therefore applies and $\dprimeddprime f\wedge T$ is Bedford--Taylor.
Stability of $\mathscr P$ under local uniform limits is
the weak-limit item together with
\Cref{lem:continuous-times-positive-current}.  Now apply
\Cref{lem:ch16-cpsh-induction}.
\end{proof}

\begin{corollary}
\label{cor:ch16-bt-times-function}
Let $T$ be a Bedford--Taylor current and let $f$ be a nonnegative locally
constant function.  Then $fT$ is Bedford--Taylor.
\end{corollary}

\begin{proof}
Locality (\Cref{prop:bt-current-stability}) reduces to the case of a
constant $f=c\geq0$, where $cT$ is positive and closed and
$\dprimeddprime(u\,cT)=c\,\dprimeddprime(uT)$ by linearity.
\end{proof}

CLD state the stronger \cite[Corollary~8.6.7]{CLD25}, for an arbitrary
continuous nonnegative $f$ such that $fT$ is again positive and closed; the
extra input there is a local smooth approximation of $f$ combined with
\Cref{prop:bt-current-stability}.  Only the locally constant case is used
below, so we do not reproduce that argument.

\subsection{Dimension one}
\label{sec:ch16-bt-curves}

In analytic dimension one the ambiguity can be removed for a broad class
of spaces over an arbitrary complete nontrivially valued ground field.  On a
connected good separated strict rig-smooth boundaryless curve, every positive
closed current of bidegree $(0,0)$ is a constant multiple of the integration
current.  The same classification holds on a connected normal curve when
rig-smoothness fails only along a locally finite set of rigid points.  More
generally, on a possibly reducible or nonreduced pure curve for which every
normalized reduced irreducible component is rig-smooth away from such a set,
the constant may vary by irreducible component.  Nilpotent thickening leaves
ordinary CLD forms, currents, positivity, closedness, and continuous psh
sheaves unchanged, while the integration current continues to record the
generic component multiplicities.

The proof is arranged to keep the source curve visible throughout the
tropical calculation.  A first bridge constructs a full canonical
virtual-annular collar whose tropical ray has an exact source preimage; the
two examples following it show respectively the ordinary annular model and a
degree-two residual branch.  Proper tropical pushforward, full-star
balancing, and Riemann--Roch on the residual curve over its actual ground
field then give the smooth classification.  Finite normalization and a
rigid-point zero-atom estimate give the reduced theorem, and nilreduction
gives the nonreduced theorem and its multiplicity-weighted curvature
comparison.  Between these two steps, the vanishing of smooth top-degree test
forms at rigid points gives the removable-exception extension on a normal
curve.  A local zero-atom estimate under ordinary normality then prevents
curvature measures on different normalized components from cancelling at
their rigid conductor and intersection points.  This yields the full
componentwise statement just described.  The final nodal example displays
both componentwise classification and the effect of a double thickening.

For an integral algebraic curve $C$, we write
$\operatorname{Rat}(C)$ for its function field.
\index{function field!of a residual curve}
\index{residual curve!function field}

% LIVE-READY STAGING ONLY, 28 August 2026.
% Release gate: do not integrate this prose until the arbitrary-field
% nonreduced exact-live statement/proof bytes receive exact-live verification.
% The rendered prose below deliberately contains no gate-status language.

\paragraph{Why the source collar has to be identified.}
A tropical chart records the values of finitely many analytic functions; it
does not remember a source point by name.  Two different branches can
therefore have the same image ray, and a remote point can in principle map
into a short segment intended for a local branch calculation.  This matters
for currents.  Euclidean balancing determines the flux carried by a target
ray, whereas the classification argument needs the flux on one specified
source branch.  Equality of ray germs alone does not identify those two
numbers.

There is a second, more elementary, restriction.  A smooth partition
function on the source is not generally an analytic moment coordinate.  On
an annular skeleton the negative logarithm of the norm of an analytic unit
is piecewise integral affine.  A smooth plateau with a genuinely curved
transition cannot have that form.  If a cutoff is retained in a local
calculation, the legal replacement is the factorization in
\eqref{eq:p21p3-gf-source-cutoff-factorization}: one writes
\(\chi=\Phi\circ H_{\mathrm{trop}}\) for an analytic moment block \(H\)
and a smooth function \(\Phi\) on its tropical target.  One appends \(H\)
to the full analytic moment \(F\), pulls superforms back with the notation
\(F^*\beta\), and never declares \(\chi\) itself to be analytic.

\Cref{prop:p21p3-gf-virtual-annulus-bridge} resolves the source ambiguity
in three stages.  At an interior type-\(2\) point \(x\), every branch \(b\)
corresponds to a closed point \(p_b\) of the normal projective integral
residual curve \(C_x\).  A residue function with order one at \(p_b\) lifts
to an analytic unit whose branch-order coordinate is one.  This makes the
outgoing vector primitive over the original field, even when \(p_b\) has
residue degree greater than one.  Weak approximation at finitely many closed
points then supplies a finite analytic block separating every selected ray
from every other active branch.

The source neighborhood is chosen before the final proper localization.  All
virtual-disc components are retained whole and are constant for the first
moment block; only finitely many virtual-annular components can be active.
After those finitely many annular arms have been shortened, the localization
lemma is applied again inside the resulting one-vertex star.  For a selected
branch it produces the full canonical collar specified in
\eqref{eq:p21p3-gf-canonical-full-source-collar} and proves the exact
preimage identity
\eqref{eq:p21p3-gf-source-separated-ray-preimage}.  Consequently a test
supported on the corresponding relative-open ray segment receives source
contribution from precisely that collar.

One may split a virtual annulus after a finite separable extension in order
to compute its slopes, but no current is pulled back to that extension.  If
the split annulus has expansion index \(e_b\), the ordinary current flux and
the calibrated integration-current flux both acquire the same factor
\(e_b\).  Their ratio is therefore defined downstairs over \(K\) and is
independent of the chosen moment presentation.  This is why the canonical
collar, rather than a chosen geometric branch after scalar extension, is the
right source object.


% Canonical-collar repair of the source-legal arbitrary-field
% virtual-annulus/tropical-ray comparison.
% Date: 28 August 2026.
% Standalone production-verification candidate; not integrated into the book.
% The frozen 20260827 files, the rejected c5a643/292bb repaired artifact, and
% the narrowly rejected 0315d3/c903cf finite-annulus artifact are not modified.

\begin{lemma}
\label{lem:p21p3-gf-legal-source-localization}
Let \(K\) be a complete nontrivially valued non-Archimedean field, let
\(Y\) be a good separated strict rig-smooth \(K\)-analytic curve, let
\(y\in Y\), and let \(W\) be an open neighborhood of \(y\).  Let \(\alpha\)
be an ordinary CLD-smooth form near \(y\), and let \(G\) be a finite analytic
moment block near \(y\).  If a smooth source cutoff \(\chi\) is to be retained,
also regard \(\chi\) as part of the data.

After shrinking around \(y\), there are a compact analytic neighborhood
\(V\subset W\) of \(y\), analytic moments \(F_{\mathrm{bd}}\),
\(F_{\mathrm{form}}\), and, when \(\chi\) is present, \(H\), and a smooth
superform \(\eta\) such that, for
\[
 F=(F_{\mathrm{bd}},F_{\mathrm{form}},G,H):V\longrightarrow
       \mathbf G_{m,K}^{N},
\]
with the last block omitted when there is no \(\chi\), one has
\begin{align}
 \alpha|_V&=F^*(\pi_{\mathrm{form}}^*\eta),
 \label{eq:p21p3-gf-source-local-form}\\
 F_{\mathrm{trop}}(y)&\notin
       F_{\mathrm{trop}}(\partial(V/W)),
 \label{eq:p21p3-gf-source-local-boundary}
\end{align}
and, when \(\chi\) is present, there is a smooth function \(\Phi\) on a
neighborhood of \(H_{\mathrm{trop}}(V)\) such that
\begin{equation}
 \chi|_V=\Phi\circ H_{\mathrm{trop}}.
 \label{eq:p21p3-gf-source-cutoff-factorization}
\end{equation}
Here \(\pi_{\mathrm{form}}\) denotes projection to the
\(F_{\mathrm{form}}\)-coordinates.  In particular,
\begin{equation}
 F_{\mathrm{trop}}(z)=F_{\mathrm{trop}}(z')
       \quad\Longrightarrow\quad \chi(z)=\chi(z')
       \qquad(z,z'\in V).
 \label{eq:p21p3-gf-source-cutoff-fibre-constancy}
\end{equation}

Put
\begin{equation}
 \Omega=\mathbf R^N\setminus
       F_{\mathrm{trop}}(\partial(V/W)),
 \qquad U=F_{\mathrm{trop}}^{-1}(\Omega)\subset V.
 \label{eq:p21p3-gf-source-local-omega}
\end{equation}
Then \(U\) is an open neighborhood of \(y\) in \(W\), and
\(F_{\mathrm{trop}}|_U:U\to\Omega\) is proper.  If \(\lambda\) is a smooth
function compactly supported in \(\Omega\), with support in the domain of
\(\pi_{\mathrm{form}}^*\eta\), then
\[
 F^*(\lambda\,\pi_{\mathrm{form}}^*\eta)
\]
has compact support in \(U\) and extends by zero to an ordinary smooth form
on \(W\).  It agrees with \(\alpha\) wherever \(\lambda=1\) on the tropical
image.

Consequently the construction can be performed inside each member of any
prescribed finite source cover of a compact test support.  In particular, if
the members are chosen to be single virtual-annular collars and
single-vertex source neighborhoods, every resulting local current comparison
takes place in one such source neighborhood.  Smooth partitions of unity may
be used to pass between the local comparisons, but a partition function is
never asserted to be an analytic tropical coordinate: if it must occur in a
local moment presentation, \eqref{eq:p21p3-gf-source-cutoff-factorization}
is used and the analytic block \(H\), not \(\chi\), is appended.
\end{lemma}

\begin{proposition}
\label{prop:p21p3-gf-virtual-annulus-bridge}
Let \(K\) be a complete nontrivially valued non-Archimedean field, let
\(Y\) be a good separated strict rig-smooth \(K\)-analytic curve, and let \(x\) be an
interior point of type \(2\).  Write \(C_x\) for the normal projective
integral residual curve at \(x\).  If \(b\) is a branch at \(x\), write
\(p_b\in C_x\) for its corresponding closed point.

For \(g\in\mathcal O_{Y,x}^{\times}\), let
\(\partial_b(g)\in\mathbf Z\) be its exponent in the rank-one quotient of
the refined branch value group.  More explicitly, choose
\(s_b\in\operatorname{Rat}(C_x)^{\times}\) with
\(\operatorname{ord}_{p_b}(s_b)=1\), lift \(s_b\) to a norm-one analytic
unit \(f_b\in\mathcal O_{Y,x}^{\times}\), and require
\begin{equation}
\label{eq:p21p3-gf-branch-order-definition}
 \overline{|g(x)|_b}
 =\overline{|f_b(x)|_b}^{\,\partial_b(g)}
 \quad\text{in}\quad
 |\mathcal H(x)^{\times}|_b/|\mathcal H(x)^{\times}|.
\end{equation}
Then \(\partial_b(g)\) is independent of \(s_b\) and its lift subject to
\(\operatorname{ord}_{p_b}(s_b)=1\), and
\begin{equation}
\label{eq:p21p3-gf-branch-order-residue}
 \partial_b(g)=\operatorname{ord}_{p_b}(\widetilde g)
\end{equation}
whenever \(g\) has norm one at \(x\) and \(\widetilde g\) denotes its
degree-one residue.

The following assertions hold.
\begin{enumerate}
\item\label{itm:p21p3-gf-bridge:ray}
Let \(g_1,\ldots,g_m\in\mathcal O_{Y,x}^{\times}\), and assume that one
of them, say \(g_m\), satisfies \(\partial_b(g_m)=1\).  Put
\[
 G_{\mathrm{trop}}
   =(-\log|g_1|,\ldots,-\log|g_m|),
 \qquad
 v_b=(\partial_b(g_1),\ldots,\partial_b(g_m)).
\]
There are a cofinal virtual-annular section \(Z_b\) of \(b\) and a number
\(\varepsilon_b>0\) such that every \(|g_j|\) factors through the canonical
retraction \(\rho_b:Z_b\to S^{\mathrm{an}}(Z_b)\), the restriction of
\(G_{\mathrm{trop}}\) to \(S^{\mathrm{an}}(Z_b)\) is a homeomorphism, and
\begin{equation}
\label{eq:p21p3-gf-exact-branch-ray}
 G_{\mathrm{trop}}(Z_b)
 =G_{\mathrm{trop}}(x)+(0,\varepsilon_b)v_b.
\end{equation}
In particular, \(v_b\in\mathbf Z^m\) is the primitive outgoing tangent of
this Euclidean ray.

\item\label{itm:p21p3-gf-bridge:separation}
Let \(b_1,\ldots,b_s\) be distinct branches at \(x\), and let \(F_0\) be
any moment defined near \(x\).  There is a finite family of analytic units
\(f_1,\ldots,f_N\in\mathcal O_{Y,x}^{\times}\) such that, for the
concatenated moment \(F=(F_0,f_1,\ldots,f_N)\), the ray germs associated
with \(b_1,\ldots,b_s\) are pairwise distinct primitive rays and no branch
outside \(\{b_1,\ldots,b_s\}\) has the same outgoing ray germ as any
\(b_i\).

\item\label{itm:p21p3-gf-bridge:localization}
Let \(T\) be a positive \(\mathrm d'\)- and \(\mathrm d''\)-closed CLD
current of bidimension \((1,1)\), equivalently of bidegree \((0,0)\), near
\(x\), and let \(\alpha\) be an ordinary CLD-smooth \((1,1)\)-test summand
in the local current calculation.  Fix distinct selected branches
\(b_1,\ldots,b_s\) at \(x\).  Choose a one-vertex source neighborhood
\(W\) of \(x\); it may have
infinitely many branches, but it contains no other retained vertex.  First use
\cref{lem:p21p3-gf-legal-source-localization} to obtain, on a germ inside
\(W\), a moment \(F_0\) containing the form and every retained-cutoff block.
Append the branch-separating block from
\cref{itm:p21p3-gf-bridge:separation}, and put
\(G_0=(F_0,F_{\mathrm{br}})\).  On a common open domain \(O\subset W\) of
its coordinates, choose a connected open elementary neighborhood
\(V_*\subset O\) of \(x\), with compact-tree closure, for which
\[
 V_*\setminus\{x\}
   =\coprod_{\lambda\in\Lambda}B_\lambda
      \amalg\coprod_{\nu=1}^{r}A_\nu .
\]
Here every \(B_\lambda\) is a virtual disc and the \(A_\nu\) are the
finitely many gently virtual-annular components.  Choose a skeleton cutpoint
\(q_\nu\in A_\nu\) in each of the latter, and let \(W_1\) be the component
containing \(x\) of
\(V_*\setminus\{q_1,\ldots,q_r\}\).  The cutpoints can be chosen so close to
\(x\) that \(W_1\) is an open one-vertex source neighborhood with
\[
 W_1\setminus\{x\}
   =\coprod_{\lambda\in\Lambda}B_\lambda
      \amalg\coprod_{\nu=1}^{r}A_\nu^0,
\]
where \(A_\nu^0\) is the retained virtual subannulus.  The moment
\(G_{0,\mathrm{trop}}\) is identically \(G_{0,\mathrm{trop}}(x)\) on every
entire \(B_\lambda\); on each \(A_\nu^0\) it is either constant or has as
image a short Euclidean ray segment prescribed by its branch-order vector.
Each selected ray is distinct from the ray of every other active branch.
If \(A_{\nu(i)}^0\) represents \(b_i\), there are \(\delta_i>0\) and a
branch-order vector \(w_i\) such that
\begin{equation}
 G_{0,\mathrm{trop}}^{-1}\!\left(
   G_{0,\mathrm{trop}}(x)+(0,\delta_i)w_i\right)\cap W_1
   =A_{\nu(i)}^0.
 \label{eq:p21p3-gf-preliminary-exact-source-locator}
\end{equation}
Reapply
\cref{lem:p21p3-gf-legal-source-localization} inside \(W_1\), taking \(G_0\)
as the prescribed block \(G\); call its output \(V,F,\Omega,U\).  Thus
\(V\subset W_1\), and \(\Omega\) and \(U\) are given by
\eqref{eq:p21p3-gf-source-local-omega} for the resulting full moment \(F\).

For each selected branch \(b_i\), let
\(t_i:A_{\nu(i)}^0\to(0,\delta_i)\) be the radial parameter determined by
the order-one separator coordinate.  There is
\(0<\varepsilon_i<\delta_i\) such that the full canonical virtual subannulus
\[
 Z_{b_i}:=t_i^{-1}((0,\varepsilon_i))\subset U
\]
satisfies, with \(v_i=(\partial_{b_i}(f))_f\) denoting the full branch-order
vector of the coordinates of \(F\),
\begin{align}
 F_{\mathrm{trop}}(Z_{b_i})
   &=F_{\mathrm{trop}}(x)+(0,\varepsilon_i)v_i,
 \label{eq:p21p3-gf-source-separated-ray-image}\\
 F_{\mathrm{trop}}^{-1}\!\left(
   F_{\mathrm{trop}}(x)+(0,\varepsilon_i)v_i\right)\cap U
   &=Z_{b_i}.
 \label{eq:p21p3-gf-source-separated-ray-preimage}
\end{align}
Thus the one-cell of the localized one-dimensional tropical image issuing
from \(F_{\mathrm{trop}}(x)\) in the \(b_i\)-direction is exactly the image
of that source collar, and its primitive tangent is \(v_i\).

After a residue function \(g\in\operatorname{Rat}(C_x)^\times\) and a
norm-one analytic lift have also been appended to the moment, write again
\(v_i\) for the enlarged primitive vector.  Its lift coordinate is
\(\operatorname{ord}_{p_{b_i}}(g)\).  Hence, if
the edge flux is \(a_i\), its term in that scalar component of current
balancing is
\(a_i\operatorname{ord}_{p_{b_i}}(g)\).  The exact preimage identity
\eqref{eq:p21p3-gf-source-separated-ray-preimage}, rather than a purported
tropicalization of a smooth plateau, is the source-separation input.  This
comparison takes place entirely over \(K\); no pullback or descent of the
current is involved.
\end{enumerate}
\end{proposition}

\begin{remark}
\label{rem:p21p3-gf-no-plateau-coordinate}
On a strict annulus with skeleton coordinate \(t=-\log|T|\), the restriction
of \(-\log|f|\) for an analytic unit \(f\) is piecewise affine with integral
slopes.  A nonconstant \(C^\infty\) plateau with a transition interval on
which its second derivative is nonzero is therefore not an analytic moment
coordinate.  The legal expression is \(\chi=\Phi\circ H_{\mathrm{trop}}\),
as in \eqref{eq:p21p3-gf-source-cutoff-factorization}; one appends \(H\), not
\(\chi\).
\end{remark}

\begin{remark}
\label{rem:p21p3-gf-virtual-annulus-index}
After a finite separable extension splitting \(Z_b\), a source annulus
coordinate can map to the primitive target coordinate with a positive
integer expansion index \(e_b\).  The proof shows that all coordinate
exponents are multiplied by the same \(e_b\), while the saturated target
tangent remains \(v_b\).  This is precisely why ordinary-current flux and
canonical calibrated flux acquire the same factor under tropical
pushforward.
\end{remark}

\paragraph{Scope.}
The source-localization assertion concerns ordinary CLD-smooth forms on good
separated strict rig-smooth curves over a nontrivially valued field.  It does
not itself produce the finite source cover used by an application.  Nor does
it assert the corresponding statement for weakly smooth forms presented only
by harmonic tropicalizations, trivially valued or non-strict spaces, boundary
points, or singular curves before normalization.

% Proof of the canonical-collar repair of the source-legal arbitrary-field
% virtual-annulus comparison.
% Date: 28 August 2026.
% Standalone production-verification candidate; not integrated into the book.
% The rejected c5a643/292bb and 0315d3/c903cf artifacts are preserved
% unchanged in separate files.

\begin{proof}[Proof of
\cref{lem:p21p3-gf-legal-source-localization}]
Use the ordinary local moment description in
\cite[Proposition~15.1.1(b)]{CLD25}.  After shrinking \(W\) about \(y\),
there are an analytic moment \(F_{\mathrm{form}}\) and a smooth superform
\(\eta\) such that \(\alpha=F_{\mathrm{form}}^*\eta\).  Shrink once more so that
the prescribed finite block \(G\) is defined.  When \(\chi\) is present,
\cite[Corollary~15.1.2]{CLD25} gives, after one further shrinking, an analytic
moment \(H\) and a smooth function \(\Phi\) such that
\(\chi=\Phi\circ H_{\mathrm{trop}}\).  Relabel this common neighborhood as
\(W\).

Apply \cite[Lemma~12.3.6]{CLD25} to the analytic space \(W\) at \(y\).  It
gives a compact analytic neighborhood \(V\) of \(y\) in \(W\) and a moment
\(F_{\mathrm{bd}}\) such that
\[
 F_{\mathrm{bd,trop}}(y)\notin
 F_{\mathrm{bd,trop}}(\partial(V/W)).
\]
Concatenate \(F_{\mathrm{bd}}\), \(F_{\mathrm{form}}\), \(G\), and \(H\).
Projection to the first block proves
\eqref{eq:p21p3-gf-source-local-boundary}; projection to the form and cutoff
blocks proves \eqref{eq:p21p3-gf-source-local-form} and
\eqref{eq:p21p3-gf-source-cutoff-factorization}.  Equality of two full
tropical values implies equality of their \(H\)-projections, so applying
\(\Phi\) proves \eqref{eq:p21p3-gf-source-cutoff-fibre-constancy}.

The remainder is the localization used in the proof of
\cite[Theorem~16.8.5]{CLD25}.  The set \(\Omega\) in
\eqref{eq:p21p3-gf-source-local-omega} is open because the relative boundary
is compact.  The set \(U\) misses \(\partial(V/W)\), hence is contained in the
relative interior of \(V\) and is open in \(W\).  It contains \(y\) by
\eqref{eq:p21p3-gf-source-local-boundary}.  Moreover
\(F_{\mathrm{trop}}|_U\) is obtained from the proper map
\(F_{\mathrm{trop}}:V\to\mathbf R^N\) by base change along
\(\Omega\hookrightarrow\mathbf R^N\), and is therefore proper.  Thus the
inverse image of the compact support of \(\lambda\) is compact in \(U\).
The displayed pullback form consequently extends by zero across
\(W\setminus U\), and the equality with \(\alpha\) where \(\lambda=1\)
follows from \eqref{eq:p21p3-gf-source-local-form}.

For the last assertion, apply the construction at each point of the compact
test support, starting inside the prescribed member of the source cover.
The resulting open sets \(U\) cover that support, so finitely many suffice.
Smooth cutoffs subordinate to this finite cover, with sum one near the
compact support, exist by \cite[Corollary~15.2.9]{CLD25}.  Current identities
are local and therefore glue after this finite decomposition.  If one of
those cutoffs is retained in a
single local presentation, the preceding application of
\cite[Corollary~15.1.2]{CLD25} supplies \(H\); at no point was the cutoff
itself used as an analytic coordinate.
\end{proof}

\begin{proof}[Proof of
\cref{prop:p21p3-gf-virtual-annulus-bridge}]
The residual curve and the branch--valuation correspondence over an
arbitrary complete field are
\cite[\S\S3.3.5.2, 4.2.6, 4.2.9, and 4.2.11.1]
{Ducros24CurveStructure}.  Since \(x\)
is interior, every closed point of \(C_x\) occurs as a branch.  For the branch
\(b\) corresponding to \(p_b\), the refined value group is the lexicographic
extension of \(|\mathcal H(x)^{\times}|\) by the discrete valuation
\(\operatorname{ord}_{p_b}\).  If \(s_b\) is a uniformizer of the discrete
valuation ring \(\mathcal O_{C_x,p_b}\), then
\cite[\S4.2.11.1]{Ducros24CurveStructure} says precisely that the class of
its refined value generates
\[
 |\mathcal H(x)^{\times}|_b/|\mathcal H(x)^{\times}|\simeq\mathbf Z.
\]
Since \(Y\) is normal and \(x\) is nonrigid,
\cite[\S4.2.18]{Ducros24CurveStructure} identifies
\(\mathcal O_{Y,x}\) with \(\kappa(x)\), and
\cite[\S4.2.6]{Ducros24CurveStructure} says that \(\kappa(x)\) is dense in
\(\mathcal H(x)\).  Hence \(s_b\), viewed as a homogeneous residue element
of degree one, has a norm-one lift \(f_b\in\mathcal O_{Y,x}\).  This is also
the finite residue-lifting step made explicit in the proof of
\cite[Lemma~6.14]{GJR26}.  Its value in \(\mathcal H(x)\) is nonzero, so
after shrinking it is an analytic unit.  This proves that
\eqref{eq:p21p3-gf-branch-order-definition} is meaningful.  Replacing a
uniformizer by another one multiplies it by a unit in
\(\mathcal O_{C_x,p_b}\), whose class in the displayed quotient is trivial.
This proves independence and also gives
\eqref{eq:p21p3-gf-branch-order-residue}.

We next prove the exact ray assertion.  By
\cite[Proposition~4.3.9(2) and \S\S4.3.10.2--4.3.10.3]
{Ducros24CurveStructure}, \(b\) has a cofinal family of virtual-annular
sections.  Choose one and call it \(Z\).  Its analytic skeleton and canonical
retraction are compatible with every complete scalar extension by
\cite[\S\S3.6.35.5--3.6.35.10]{Ducros24CurveStructure}.  After a finite
separable extension and then passage to a completed algebraic closure, a
component above \(Z\) is an ordinary annulus whose skeleton maps
homeomorphically to \(S^{\mathrm{an}}(Z)\).

Let \(\beta\) be the branch of this split annulus above \(b\), and choose an
annulus coordinate \(\theta\) so that \(-\log|\theta|\) increases away from
the endpoint over \(x\).
The annular calculation in
\cite[\S\S3.6.22--3.6.24 and \S4.3.6.4]
{Ducros24CurveStructure} shows that, after shortening the annulus once for
the finite family \(g_1,\ldots,g_m\),
\begin{equation}
\label{eq:p21p3-gf-split-annulus-monomials}
                  |g_j|=c_j|\theta|^{n_j}
                  \qquad(1\leq j\leq m)
\end{equation}
for constants \(c_j>0\) and integers \(n_j\).  Naturality of the refined
branch valuation is \cite[\S4.2.9]{Ducros24CurveStructure}.  Consequently,
restriction from \(b\) to \(\beta\) sends the generator defined by \(g_m\) to
a positive integral multiple of the primitive annular generator.  Hence
there is one integer \(e_b>0\), independent of \(j\), such that
\begin{equation}
\label{eq:p21p3-gf-common-expansion-index}
                       n_j=e_b\partial_b(g_j)
                       \qquad(1\leq j\leq m).
\end{equation}
Indeed, this is just the homomorphism between the two rank-one quotient
value groups; \(n_m=e_b\) because \(\partial_b(g_m)=1\).  The same integer
is the tangent-lattice expansion index of this split annular presentation
by \cite[\S3.6.24.2]{Ducros24CurveStructure}.  In the finite-morphism
setting of \cite[Theorem~4.3.13(i)]{Ducros24CurveStructure}, that annular
index is identified with the corresponding residual ramification index.

Put
\[
 t=-\log|g_m|+\log|g_m(x)|.
\]
The orientation and
\cite[\S4.2.11.1]{Ducros24CurveStructure} make \(t\) strictly positive
away from \(x\) on a sufficiently short section.  Equations
\eqref{eq:p21p3-gf-split-annulus-monomials} and
\eqref{eq:p21p3-gf-common-expansion-index}, followed by passage to the
limit at \(x\), give on the skeleton
\begin{equation}
\label{eq:p21p3-gf-coordinate-ray-formula}
 -\log|g_j|+\log|g_j(x)|
       =\partial_b(g_j)t
       \qquad(1\leq j\leq m).
\end{equation}
Every invertible analytic function on a virtual annulus has the same norm
as at its canonical retraction by
\cite[\S3.6.35.8]{Ducros24CurveStructure}.  Thus
\eqref{eq:p21p3-gf-coordinate-ray-formula} holds on the whole shortened
section \(Z_b\), and \(t\) identifies its skeleton with an interval
\((0,\varepsilon_b)\).  This proves
\eqref{eq:p21p3-gf-exact-branch-ray}.  Since the last coordinate of \(v_b\)
is one, \(v_b\) is primitive in the saturated ambient lattice.  The
one-dimensional skeleton attached to the moment, restricted to any compact
subannular band, has the unique integral affine structure for which the
moment maps its faces integral-affinely onto their images by
\cite[\S8.13]{GJR26}; hence this is also the primitive tangent of the
corresponding source one-cell.  This proves
\cref{itm:p21p3-gf-bridge:ray}.

We prove the finite separation assertion.  Let \(p_i=p_{b_i}\).  Weak
approximation on the normal projective curve \(C_x\) gives rational functions
\(s_1,\ldots,s_s\) with
\begin{equation}
\label{eq:p21p3-gf-kronecker-orders}
                \operatorname{ord}_{p_i}(s_j)=\delta_{ij}.
\end{equation}
For completeness, choose a closed point \(q\) away from the \(p_i\).  For
\(n\gg0\), Serre vanishing gives
\[
 H^1\!\left(C_x,
   \mathcal O_{C_x}\!\left(nq-2\sum_i p_i\right)\right)=0.
\]
The divisor exact sequence makes the simultaneous value-and-first-jet map
from \(H^0(C_x,\mathcal O_{C_x}(nq))\) onto
\(\bigoplus_i\mathcal O_{C_x,p_i}/\mathfrak m_{p_i}^2\) surjective.
Prescribing a nonzero first jet and value zero at \(p_j\), and value one at
the other \(p_i\), yields \(s_j\) satisfying
\eqref{eq:p21p3-gf-kronecker-orders}.  This argument is valid for closed
points of arbitrary residue degree.

Lift the \(s_j\) to norm-one analytic units by the density argument above
(equivalently, the lifting step in the proof of
\cite[Lemma~6.14]{GJR26}), and first append those lifts to \(F_0\).  By
\cite[Proposition~4.2.4(ii)]{Ducros24CurveStructure}, every
analytic unit has nonzero branch order in only finitely many branches.
The preliminary outgoing vector in branch \(b_i\) therefore has, in the new
\(s_1,\ldots,s_s\) block, the standard basis vector \(e_i\); in particular it
is primitive, and the selected vectors are pairwise nonparallel.

There can still be an unselected branch whose preliminary vector is a
positive multiple of the vector for some \(b_i\).  There are only finitely
many such branches: any one of them has positive \(s_i\)-order and hence
belongs to the finite exceptional set supplied by
\cite[Proposition~4.2.4(ii)]{Ducros24CurveStructure}.  Call their union
\(Q\).  For every \(c\in Q\), weak approximation as above gives a rational
function \(h_c\) satisfying
\begin{equation}
\label{eq:p21p3-gf-collision-killers}
 \operatorname{ord}_{p_{b_i}}(h_c)=0\quad(1\leq i\leq s),
 \qquad
 \operatorname{ord}_{p_c}(h_c)=1.
\end{equation}
Lift all \(h_c\) and append them as well.  The new coordinates vanish in
branch order on every selected branch, whereas the \(h_c\)-coordinate has
order one on \(c\).  Thus each old collision is destroyed.  A newly active
branch cannot create a new collision: projection to the preliminary
coordinates would already have made it an element of \(Q\).  This proves
\cref{itm:p21p3-gf-bridge:separation}.

We finally prove \cref{itm:p21p3-gf-bridge:localization}.  First use
\cref{lem:p21p3-gf-legal-source-localization} inside the prescribed source
\(W\), with no branch block yet, and retain the resulting moment germ
\(F_0\).  Apply
\cref{itm:p21p3-gf-bridge:separation} to \(F_0\).  On a common neighborhood
of \(x\), let \(F_{\mathrm{br}}\) be the resulting separating block and put
\(G_0=(F_0,F_{\mathrm{br}})\).  Let \(O\subset W\) be a common open
neighborhood of \(x\) on which every coordinate of \(G_0\) is an analytic
unit.  We choose the source star before adding any further moment coordinates.
By \cite[Theorem~4.5.4]{Ducros24CurveStructure}, after shrinking inside
\(O\), there is a connected open elementary neighborhood \(V_*\) of \(x\)
whose closure is a compact tree and for which
\begin{equation}
 V_*\setminus\{x\}
   =\coprod_{\lambda\in\Lambda}B_\lambda
      \amalg\coprod_{\nu=1}^{r}A_\nu .
 \label{eq:p21p3-gf-elementary-star-decomposition}
\end{equation}
The \(B_\lambda\) are virtual discs, whereas
\(A_1,\ldots,A_r\) are the finitely many gently virtual-annular components.
Every component in \eqref{eq:p21p3-gf-elementary-star-decomposition}
represents one branch at \(x\), and the compact-tree condition prevents an
annular component from returning to \(x\) at its other end.

An invertible analytic function has constant norm on a virtual disc by
\cite[\S3.6.32.3]{Ducros24CurveStructure}.  Continuity at \(x\) therefore
gives
\begin{equation}
 G_{0,\mathrm{trop}}|_{B_\lambda}
       =G_{0,\mathrm{trop}}(x)
       \qquad(\lambda\in\Lambda).
 \label{eq:p21p3-gf-disc-components-at-vertex}
\end{equation}
This retains all, possibly infinitely many, virtual-disc components whole.

Only the finitely many annular components need to be shortened.  Let
\(b_\nu\) be the branch represented by \(A_\nu\), and put
\(w_\nu=(\partial_{b_\nu}(g))_g\), where \(g\) runs through the coordinates
of \(G_0\).  For this calculation only, choose a norm-one analytic unit
\(u_\nu\in\mathcal O_{Y,x}^{\times}\) with
\(\partial_{b_\nu}(u_\nu)=1\); it is not appended to the moment.  Since there
are only finitely many \(A_\nu\), all the \(u_\nu\) are units on a common
sufficiently short collar.  After a finite separable extension splitting
\(A_\nu\), choose an annulus coordinate \(\theta_\nu\) oriented away from the
end over \(x\).  Apply \eqref{eq:p21p3-gf-split-annulus-monomials}
simultaneously to the coordinates of \(G_0\) and to \(u_\nu\).  Naturality of
the refined branch valuation, \cite[\S4.2.9]{Ducros24CurveStructure}, and the
rank-one quotient argument of
\eqref{eq:p21p3-gf-common-expansion-index} give one integer \(e_\nu>0\) such
that the split-annulus exponents are
\[
 n_{\nu,g}=e_\nu\partial_{b_\nu}(g),
 \qquad n_{\nu,u}=e_\nu.
\]
Consequently, after choosing a skeleton point \(q_\nu\) sufficiently close to
the end accumulating at \(x\),
\begin{equation}
 G_{0,\mathrm{trop}}(z)
   =G_{0,\mathrm{trop}}(x)+t_\nu(z)w_\nu,
 \qquad 0<t_\nu(z)<\delta_\nu,
 \label{eq:p21p3-gf-finite-annulus-ray-formula}
\end{equation}
on the retained collar, where
\[
 t_\nu(z)=-\log|u_\nu(z)|+\log|u_\nu(x)|.
\]
The orientation makes \(t_\nu\) positive and strictly increasing along the
retained skeleton.  By \cite[\S3.6.35.8]{Ducros24CurveStructure} it factors
through the canonical retraction.  If \(w_\nu=0\), the displayed formula says
that the whole retained collar maps to the vertex.
Sections~3.6.35.9--3.6.35.10 of \cite{Ducros24CurveStructure} identify the
inverse image \(A_\nu^0\) of the retained skeleton interval with an analytic
virtual subannulus, compatibly with scalar extension.

Let \(W_1\) be the component containing \(x\) of
\(V_*\setminus\{q_1,\ldots,q_r\}\).  Analytic curves are locally connected
graphs by \cite[Theorem~3.5.1]{Ducros24CurveStructure}; hence \(W_1\) is an
open connected analytic neighborhood of \(x\).  The tree description and
the choice of the \(q_\nu\) give
\begin{equation}
 W_1\setminus\{x\}
   =\coprod_{\lambda\in\Lambda}B_\lambda
      \amalg\coprod_{\nu=1}^{r}A_\nu^0.
 \label{eq:p21p3-gf-trimmed-star-decomposition}
\end{equation}
Thus no independent truncation has been made in infinitely many directions.

Equation \eqref{eq:p21p3-gf-disc-components-at-vertex} and the order-one
separator coordinate show that a selected branch cannot be represented by a
virtual-disc component.  Hence it is represented by a unique
\(A_{\nu(i)}^0\).  Put \(w_i=w_{\nu(i)}\) and
\(\delta_i=\delta_{\nu(i)}\).  The residue-unit block contains a coordinate
of order one on \(b_i\), so \(w_i\ne0\), and
\cref{itm:p21p3-gf-bridge:separation} says that
\(\mathbf R_{>0}w_i\) is different from the ray of every other active
component.  After moving the finitely many \(q_\nu\) closer to \(x\), if
necessary, equations \eqref{eq:p21p3-gf-disc-components-at-vertex} and
\eqref{eq:p21p3-gf-finite-annulus-ray-formula} therefore give the exact
preliminary locator
\begin{equation}
 G_{0,\mathrm{trop}}^{-1}\!\left(
 G_{0,\mathrm{trop}}(x)+(0,\delta_i)w_i\right)\cap W_1
 = A_{\nu(i)}^0.
\end{equation}
Indeed, every virtual disc and every inactive annulus maps to the vertex,
every other active annulus maps into a distinct open ray, and on the selected
annulus the order-one coordinate makes \(t_\nu\) a genuine radial parameter.
This excludes remote source collisions without an infinite simultaneous
choice of branch sections.

Reapply \cref{lem:p21p3-gf-legal-source-localization} inside \(W_1\), with
\(G=G_0\), and call its compact domain and full moment \(V\) and \(F\).  The
lemma gives boundary avoidance for this full moment.  Hence the open set
\(U\) in \eqref{eq:p21p3-gf-source-local-omega} contains \(x\), lies in
\(W_1\), and \(F_{\mathrm{trop}}|_U\) is proper.  Let
\(P:\mathbf R^N\to\mathbf R^{N_0}\) be projection to the \(G_0\)-block, so
that \(P\circ F_{\mathrm{trop}}=G_{0,\mathrm{trop}}\).

For a selected branch \(b_i\), the full moment \(F\) still contains the
order-one separator coordinate.  On \(A_{\nu(i)}^0\), the formula above for
that coordinate identifies its normalized negative logarithm with the radial
parameter \(t_i=t_{\nu(i)}\).  Since \(U\) is an open neighborhood of \(x\),
\cite[\S4.3.10.3]{Ducros24CurveStructure} and
\cite[\S\S3.6.35.9--3.6.35.10]{Ducros24CurveStructure} allow us to choose
\(0<\varepsilon_i<\delta_i\) so that
\begin{equation}
 Z_{b_i}:=t_i^{-1}((0,\varepsilon_i))
       \subset A_{\nu(i)}^0\cap U
 \label{eq:p21p3-gf-canonical-full-source-collar}
\end{equation}
is the full canonical virtual subannulus over that skeleton interval.  After
decreasing \(\varepsilon_i\) once more, the calculation in
\cref{itm:p21p3-gf-bridge:ray}, applied simultaneously to all coordinates of
\(F\), gives
\[
 F_{\mathrm{trop}}(z)=F_{\mathrm{trop}}(x)+t_i(z)v_i
 \qquad(z\in Z_{b_i}),
\]
where \(v_i\) is the full primitive branch-order vector.  This proves
\eqref{eq:p21p3-gf-source-separated-ray-image}; it also gives
\(P(v_i)=w_i\).  If \(z\in U\) and
\(F_{\mathrm{trop}}(z)=F_{\mathrm{trop}}(x)+t v_i\) for
\(0<t<\varepsilon_i\), projection and
\eqref{eq:p21p3-gf-preliminary-exact-source-locator} put \(z\) in
\(A_{\nu(i)}^0\).  The order-one separator coordinate then gives
\(t_i(z)=t\).  By the definition
\eqref{eq:p21p3-gf-canonical-full-source-collar}, \(z\in Z_{b_i}\).  The
reverse inclusion follows from the displayed ray formula, and hence
\eqref{eq:p21p3-gf-source-separated-ray-preimage} holds.

Consequently the Euclidean current
\[
 \langle T_{\Omega},\beta\rangle
       =\langle T|_U,F^*\beta\rangle
\]
used in the proof of \cite[Theorem~16.8.5]{CLD25}, when tested on a
sufficiently short relative-open subsegment of the selected ray, receives
source contribution only from \(Z_{b_i}\).  If a norm-one lift of
\(g\in\operatorname{Rat}(C_x)^\times\) is appended, projection to the
previous coordinates preserves the exact separation, while
\eqref{eq:p21p3-gf-branch-order-residue} identifies its scalar slope with
\(\operatorname{ord}_{p_{b_i}}(g)\).  Multiplication by the edge flux \(a_i\)
gives the asserted balancing term.  No smooth plateau was used as a moment
coordinate, and no scalar extension of \(T\) occurred.

Equation \eqref{eq:p21p3-gf-common-expansion-index} also proves the claim in
\cref{rem:p21p3-gf-virtual-annulus-index}.  Under pushforward, the rank-one
matrix current has one factor \(e_b\): its tangent tensor contributes
\(e_b^2\) and one-dimensional change of measure contributes \(e_b^{-1}\).
The canonical calibration has the same expansion factor by
\cite[\S14.3.2, Lemma~14.3.4, and Proposition~14.4.2]{CLD25}.  No arbitrary
current was extended to a splitting field at any point of the argument.
\end{proof}


% LIVE-READY STAGING ONLY, 28 August 2026.
% Release gate: do not integrate these examples until the arbitrary-field
% nonreduced exact-live statement/proof bytes receive exact-live verification.
% The rendered examples below deliberately contain no gate-status language.

\begin{example}
\label{ex:p21p3-ordinary-annulus-collar}
Let \(A\) be a strict open-annular branch collar issuing from an interior
type-\(2\) point \(x\) of an ambient rig-smooth analytic curve.  Choose an
annular coordinate \(T\) represented by an analytic unit near \(x\), and
orient the collar so that
\[
                  t=-\log|T|+\log|T(x)|
\]
increases away from \(x\).  Thus \(T\) and all values at \(x\) below are
taken from analytic germs on an ambient neighborhood of \(x\), not from a
fictitious value at a missing endpoint of an abstract open annulus.  Let
\(f_1,\ldots,f_m\) be analytic units on such a neighborhood and suppose
that \(f_m\) has branch order one.  After shortening \(A\), the annular norm
calculation gives
\[
 -\log|f_j|+\log|f_j(x)|=n_jt,
 \qquad n_j\in\mathbf Z,
\]
and \(n_m=1\).  Hence \(v=(n_1,\ldots,n_m)\) is primitive and, on the full
collar \(Z=t^{-1}((0,\varepsilon))\),
\[
 F_{\mathrm{trop}}(Z)=F_{\mathrm{trop}}(x)+(0,\varepsilon)v,
 \qquad F=(f_1,\ldots,f_m).
\]
On the one-branch subset \(A\cup\{x\}\), the \(f_m\)-coordinate gives the
exact preimage equality as well.  In a genuine open star around \(x\), the
bridge appends branch-separating coordinates before asserting the same
equality.  The virtual-annular statement differs only by a common positive
expansion index after a splitting extension; that index cancels in the ratio
of ordinary and calibrated fluxes.

Choose a radial smooth cutoff \(\chi=\Phi(t)\) supported in this collar.  It
is a smooth function of the tropical coordinate of the analytic moment
\(T\), exactly as in
\eqref{eq:p21p3-gf-source-cutoff-factorization}.  Unless \(\Phi\) is
locally integral affine, however, \(\chi\) itself is not the negative
logarithm of the norm of an analytic unit.  This elementary annulus is the
local model for the distinction between a legal smooth function of tropical
coordinates and an illegal ``smooth moment coordinate.''
\end{example}

\begin{example}
\label{ex:p21p3-higher-degree-residual-branch}
Take \(K=\mathbf Q_3\), let \(Y=(\mathbf P^1_K)^{\mathrm{an}}\), and let
\(x\) be the Gauss point of the closed unit disc in the affine coordinate
\(z\).  The residual curve at \(x\) is
\(C_x=\mathbf P^1_{\mathbf F_3}\).  Its closed points, rather than merely
its \(\mathbf F_3\)-rational points, parametrize the branches at \(x\).

The polynomial
\[
                              q(Z)=Z^2+1
\]
is irreducible over \(\mathbf F_3\).  It therefore defines a closed point
\(p_q\in C_x\) of degree two.  On \(C_x\),
\begin{equation}
                    \operatorname{div}(q)=[p_q]-2[\infty].
 \label{eq:p21p3-quadratic-divisor}
\end{equation}
The lift \(g(z)=z^2+1\) is a norm-one analytic unit in the stalk at \(x\).
Its branch orders are one on the \(p_q\)-branch, minus two on the
\(\infty\)-branch, and zero on every other branch.  In particular, the
\(p_q\)-ray has a coordinate of slope one and is primitive over \(K\); one
does not split \(p_q\) into geometric points in order to define it.

Select the \(p_q\)- and \(\infty\)-branches in
\cref{prop:p21p3-gf-virtual-annulus-bridge}.  Its separating block and exact
preimage identity produce two named full source collars.  Let
\(m_q,m_\infty\) be their calibrated fluxes, and let
\(a_q,a_\infty\) be the ratios of the fluxes of a positive closed current
\(T\) to those calibrated fluxes.  Applying scalar balancing to the
coordinate \(g\), first for the integration current and then for \(T\),
gives
\begin{align}
                   m_q-2m_\infty&=0,
 \label{eq:p21p3-quadratic-calibrated-balance}\\
             a_qm_q-2a_\infty m_\infty&=0.
 \label{eq:p21p3-quadratic-current-balance}
\end{align}
Consequently \(m_q=2m_\infty\) and \(a_q=a_\infty\).  The residue degree
has not disappeared: it is exactly the factor two in
\eqref{eq:p21p3-quadratic-divisor} and in the calibrated flux.  The
Riemann--Roch argument in the proof of
\cref{prop:p21p3-gfc-smooth-rigidity} repeats this comparison for arbitrary
closed points and proves \(m_p=c_Y\deg(p)\) and
\(a_pm_p=c_T\deg(p)\).  This is the mechanism that replaces the unavailable
assumption that every residual branch be rational.
\end{example}

% LIVE-READY STAGING ONLY, 28 August 2026.
% Release gate: do not integrate this prose until the arbitrary-field
% nonreduced exact-live statement/proof bytes receive exact-live verification.
% The rendered prose below deliberately contains no gate-status language.

\paragraph{From local flux to a global current.}
Proper tropical pushforward sends a localized positive closed current to a
positive closed Euclidean current supported on a one-dimensional polyhedral
set.  On the interior of an edge, positivity and the two closedness equations
force a constant nonnegative multiple of the integration current.  At a
vertex the endpoint terms satisfy the full balancing equation; no ray of the
active target star may be discarded.

Divide the edge flux of the given current by the calibrated flux of the
integration current.  The expansion-index comparison makes this ratio
intrinsic.  At a type-\(3\) point, the two opposite rays have the same ratio.
At a type-\(2\) point, write \(\operatorname{Rat}(C_x)\) for the function
field of the residual curve \(C_x\).  Appending a lift of
\(g\in\operatorname{Rat}(C_x)^\times\) turns the scalar component of
balancing into
\begin{equation}
        \sum_{p\in C_x} a_p m_p\operatorname{ord}_p(g)=0,
 \label{eq:p21p3-principal-balance}
\end{equation}
where \(m_p>0\) is the calibrated branch flux and \(a_p\) is the intrinsic
ratio.  Equivalently, define additive functionals on the divisor group by
\[
                        L_T([p])=a_pm_p,
                        \qquad L_Y([p])=m_p.
\]
Both are nonnegative on effective divisors, and
\eqref{eq:p21p3-principal-balance} says that \(L_T\) annihilates principal
divisors; the same balancing calculation for the integration current gives
this property for \(L_Y\).  Riemann--Roch over the actual residue field
forces each functional to be a nonnegative multiple of degree.  Comparing
the two multiples shows that all \(a_p\) are equal, without a rational point
on \(C_x\) and without algebraic closedness of \(K\).

The ratios propagate through all type-\(2\) and type-\(3\) points of a
connected rig-smooth curve.  The remaining type-\(1\) and type-\(4\) set has
tropical dimension zero.  By the tropical-dimension-zero test-form theorem,
the germ of every ordinary smooth \((1,1)\)-test form vanishes at each of
its points.  A current remainder already zero off that set therefore pairs
to zero with every test form.  This endpoint argument is distinct from the
rigid-point zero-atom assertion: the latter is used only at the locally
finite rigid conductor points when one descends through normalization.

This proves the connected classification in
\cref{prop:p21p3-gfc-smooth-rigidity}.  Passing through the finite
normalizations of the irreducible components gives
\eqref{eq:p21p3-curve-classification}.  Finally,
\cref{lem:p21p3-gfc-rigid-zero-atom} prevents the positive curvature
measures on different normalization sheets from cancelling at the
conductor.  That last step is why
\cref{prop:p21p3-curve-rigidity} proves Bedford--Taylor independence for
continuous psh pairs on the reduced curve itself, rather than only on its
normalization.  The nilreduction result
\cref{prop:p21p3-curve-rigidity-nonreduced} then shows that the ordinary
CLD calculus and the continuous-pair independence property are unchanged by
a nilpotent thickening, while the fundamental integration current retains
the generic component multiplicities.


% Live-book statement adaptation of the production-verified arbitrary-field
% reduced curve packet.  Date: 28 August 2026.

\begin{lemma}
\label{lem:p21p3-gfc-rigid-zero-atom}
Let $K$ be a complete nontrivially valued non-Archimedean field, let
$Y$ be a good separated strict rig-smooth $K$-analytic curve without
boundary, let $y\in Y$ be a rigid point, and let $u$ be a continuous
CLD-plurisubharmonic function on $Y$.  Then the positive curvature
measure
\[
                         \mu_u:=\mathrm d'\mathrm d''[u]
\]
has no atom at $y$.
\end{lemma}

\begin{proposition}
\label{prop:p21p3-gfc-smooth-rigidity}
Let $K$ be a
complete nontrivially valued non-Archimedean field and let $Y$ be a
connected good separated strict rig-smooth $K$-analytic curve without
boundary.  Every positive $\mathrm d'$- and $\mathrm d''$-closed CLD
current $T$ of bidegree $(0,0)$ on $Y$ has the form
\begin{equation}
\label{eq:p21p3-gfc-smooth-rigidity}
                              T=c\delta_Y
\end{equation}
for a unique $c\in\mathbf R_{\geq0}$.
\end{proposition}

\begin{proposition}
\label{prop:p21p3-curve-rigidity}
Let $K$ be a
complete nontrivially valued non-Archimedean field and let $X$ be a good
separated strict reduced $K$-analytic space, pure of dimension one and
without boundary.  Assume that the normalization of every irreducible
component of $X$ is rig-smooth.  For each irreducible component $X_i$,
write
\[
        \nu_i:\widetilde X_i\longrightarrow X_i\hookrightarrow X,
        \qquad
        \delta_{X_i}:=(\nu_i)_*\delta_{\widetilde X_i}.
\]
Then the following assertions hold.

\begin{enumerate}
\item\label{itm:p21p3-curve:classification}
Every positive $\mathrm d'$- and $\mathrm d''$-closed CLD current $T$
of bidegree $(0,0)$ on $X$ has a unique expression
\begin{equation}
\label{eq:p21p3-curve-classification}
                         T=\sum_i c_i\delta_{X_i},
                         \qquad c_i\in\mathbf R_{\geq0}.
\end{equation}
The sum is locally finite.

\item\label{itm:p21p3-curve:bedford-taylor}
Let $W\subseteq X$ be open and let $U,V$ be continuous
CLD-plurisubharmonic functions on $W$.  If
\begin{equation}
\label{eq:p21p3-curve-equal-curvature}
           \mathrm d'\mathrm d''[U]
           =\mathrm d'\mathrm d''[V],
\end{equation}
then, for every $T$ as in
\cref{itm:p21p3-curve:classification},
\begin{equation}
\label{eq:p21p3-curve-bt-conclusion}
             \mathrm d'\mathrm d''\bigl((U-V)T|_W\bigr)=0.
\end{equation}
Thus every such $T$ has the Bedford--Taylor independence property for
continuous psh pairs in the CLD sense.  This is the continuous-psh version
of the condition in \cite[Definition~8.6.2]{CLD25}.
\end{enumerate}
\end{proposition}

\begin{remark}
\label{rem:p21p3-curve-scope}
No algebraic-closedness hypothesis is made.  The residual curve at a
type-$2$ point is allowed to be merely normal, projective, and integral,
and it need not have a closed point of degree one.  The result does not
treat a trivially valued field, a non-strict or non-good curve, boundary
points, a nonreduced curve, a reduced curve with a non-rig-smooth component
normalization, or analytic dimension at least two.  The zero-atom assertion
is only needed and only stated at rigid points; over a general ground field
not every type-$1$ point is rigid.  Since CLD Definition~8.6.2 literally
quantifies over all psh functions, including possibly discontinuous ones,
the proposition does not by itself assert that stronger property.
\end{remark}

% Live-book proof adaptation of the production-verified arbitrary-field
% reduced curve packet.  Date: 28 August 2026.

\paragraph{Imported facts and the local bridge.}
Currents, positivity, multiplication of an order-zero current by a
continuous function, and differentials of currents are used as in
\cite[\S\S7.1--7.2, Example~7.2.12, and \S7.7]{CLD25}.  In particular,
\cite[Lemma~7.7.5]{CLD25} gives
\begin{equation}
\label{eq:p21p3-gfc-transpose-dd}
 \bigl\langle \mathrm d'\mathrm d''S,\varphi\bigr\rangle
 =
 \bigl\langle S,\mathrm d'\mathrm d''\varphi\bigr\rangle .
\end{equation}
Positive currents of bidimension $(0,0)$ are Radon measures by
\cite[Proposition~7.2.8]{CLD25}, and positive currents in every bidegree
have order zero by \cite[Proposition~7.2.9]{CLD25}.  We use the local
moment description of smooth forms and functions from
\cite[Proposition~15.1.1(b) and Corollary~15.1.2]{CLD25}, smooth
partitions of unity and subordinate cutoffs from
\cite[Corollary~15.2.4, Proposition~15.2.8, and
Corollary~15.2.9]{CLD25}, boundary-avoiding moments from
\cite[Lemma~12.3.6]{CLD25}, and the proper tropical pushforward in the
proof of \cite[Theorem~16.8.5]{CLD25}.  We also use
\cite[Proposition~3.3.6, Lemma~12.3.6, Theorem~12.2.4,
Corollary~12.2.5, \S14.3.2, Lemma~14.3.4,
Proposition~14.4.2, \S16.6, and Proposition~16.6.3]{CLD25}.

For an arbitrary complete ground field, the curve structure used below is
the branch--valuation correspondence, the virtual-annulus description,
and triangulation theory in
\cite[\S\S3.3.5.2, 4.2.6, 4.2.9, 4.2.10--4.2.11, and
4.3.9--4.3.10]{Ducros24CurveStructure}.  Thus an interior
type-$2$ point $x$ has a normal projective integral residual curve
$C_x$, and its branches correspond to all closed points of $C_x$.
The branch corresponding to $p$ carries the discrete valuation
$\operatorname{ord}_p$.  An interior type-$3$ point has two branches
and a virtual-annular neighborhood.  For the arbitrary-field endpoint
statement used later, \cite[Theorem~4.2.5]{Ducros24CurveStructure}
shows that an interior nonrigid point with two branches must have type
$2$ or $3$, while \cite[Theorem~4.5.4, especially \S4.5.4.1, and
\S4.5.19.1]{Ducros24CurveStructure} supplies the virtual-disc local
model and confirms that type-$1$ and type-$4$ points of a rig-smooth
curve are unibranch.  The finite-fibre and integral-affine
properties of tropical skeleta are
\cite[Corollary~3.18 and Remark~3.20]{GJR26}; finite families of
homogeneous residue functions are lifted to analytic germs by the
construction in the proof of \cite[Lemma~6.14]{GJR26}.

The proof uses \cref{lem:p21p3-gf-legal-source-localization} and exactly the
three assertions of
\cref{prop:p21p3-gf-virtual-annulus-bridge}.  In particular, it uses the
exact downstairs branch ray, finite separation of prescribed type-$2$
branches, the canonical-collar elementary star in
\cref{itm:p21p3-gf-bridge:localization}, its
global-in-\(U\) exact-preimage identity, and compatibility with the
source-separated localization of \cite[Theorem~16.8.5]{CLD25}.  It never
pulls an arbitrary current to a field extension.

\begin{proof}[Proof of \cref{lem:p21p3-gfc-rigid-zero-atom}]
The regular one-dimensional local ring $\mathcal O_{Y,y}$ has a
uniformizer represented, after shrinking, by an analytic function $f$
with $f(y)=0$.  Regard $f$ as a morphism to
$\mathbf A^{1,\mathrm{an}}_K$.  Its fibre germ at $y$ is
zero-dimensional.  Since source and target are boundaryless, the
boundary-composition identity \cite[(11.4.1.3)]{CLD25} gives
$y\in\operatorname{Int}(Y/\mathbf A^{1,\mathrm{an}}_K)$, and
\cite[Lemma~11.6.5(a)]{CLD25} makes $f$ finite at $y$.  The induced
finite homomorphism from the discrete valuation ring at $0$ to
$\mathcal O_{Y,y}$ is torsion-free and hence flat.  After shrinking as
in \cite[\S11.6.3]{CLD25}, choose a compact analytic domain $Q_{R_0}$
containing $y$, an admissible radius $R_0>0$, and a finite flat map of
constant degree $d>0$
\begin{equation}
\label{eq:p21p3-gfc-finite-rigid-tube}
             f:Q_{R_0}\longrightarrow D(0,R_0)
\end{equation}
such that $f^{-1}(0)=\{y\}$.  A finite map has empty relative boundary;
the same boundary identity puts the relative boundary of $Q_{R_0}$
over $\{|T|=R_0\}$.

Choose an admissible radius $0<R<R_0$ and fix
\[
 0<\varepsilon_0<\tfrac12\log(R_0/R).
\]
For an admissible $0<r<R$ with
$\log(R/r)>4\varepsilon_0$, put $s_R=-\log R$ and
$s_r=-\log r$.  Choose a $C^\infty$ function
$\rho_r:\mathbf R\to[0,1]$ which is zero on
\[
 (-\infty,s_R-\varepsilon_0]\cup[s_r+\varepsilon_0,\infty),
\]
is one on $[s_R,s_r]$, is nondecreasing on
$[s_R-\varepsilon_0,s_R]$, and is nonincreasing on
$[s_r,s_r+\varepsilon_0]$.  Set
\[
 q_r(t):=\int_{-\infty}^{t}\rho_r(\tau)\,\mathrm d\tau,
 \qquad
 H_r:=\int_{\mathbf R}\rho_r(\tau)\,\mathrm d\tau .
\]
Then
\begin{equation}
\label{eq:p21p3-gfc-radial-data}
 H_r\geq\log(R/r),
 \qquad
 \int_{\mathbf R}|q_r''(t)|\,\mathrm dt=2.
\end{equation}
On the punctured tube put
$\chi_r=q_r(-\log|f|)$.  Extend it by zero across the outer region and
by the constant $H_r$ near $y$.  The extensions occur where $q_r$
is constant, so $\chi_r$ is a compactly supported CLD-smooth function,
with support in the fixed compact domain $Q_{R_0}$.

On each of the two transition tubes, the map $s=-\log|f|$ projects onto
the corresponding target interval and
\begin{equation}
\label{eq:p21p3-gfc-radial-dd}
 \mathrm d'\mathrm d''\chi_r
   =q_r''(s)\,\mathrm d's\wedge\mathrm d''s .
\end{equation}
The sign of $q_r''$ is fixed on each transition band.  Restrict proper
calibrated pushforward under
\eqref{eq:p21p3-gfc-finite-rigid-tube} to either transition tube.  By
\cite[Proposition~14.4.2]{CLD25} it is the corresponding target annulus
with total local degree $d$; in particular all sheets and residue-degree
factors are included.  If $\mu_f$ and $\mu_D$ denote the source and radial
calibrated measures, then $f_*\mu_f=d\mu_D$ on each transition band.
Since the two bands have disjoint supports and $q_r''$ has one sign on
each of them, \eqref{eq:p21p3-gfc-radial-data} gives
\begin{equation}
\label{eq:p21p3-gfc-radial-tv}
 \left\|\mathrm d'\mathrm d''\chi_r\wedge\delta_Y\right\|_{\mathrm{TV}}
 =d\int_{\mathbf R}|q_r''(t)|\,\mathrm dt=2d.
\end{equation}

The measure $\mu_u=\mathrm d'\mathrm d''[u]$ is positive by
\cite[Lemma~8.5.5]{CLD25}.  Positivity,
\eqref{eq:p21p3-gfc-transpose-dd}, and
\eqref{eq:p21p3-gfc-radial-tv} yield
\begin{align}
 \log(R/r)\,\mu_u(\{y\})
 &\leq \int_Y\chi_r\,\mathrm d\mu_u                                      \notag\\
 &=\bigl\langle[u],\mathrm d'\mathrm d''\chi_r\bigr\rangle                \notag\\
 &\leq 2d\sup_{Q_{R_0}}|u|.
\label{eq:p21p3-gfc-radial-bound}
\end{align}
The right side is independent of $r$.  Nontriviality of the valuation
supplies admissible radii tending to zero, so
\eqref{eq:p21p3-gfc-radial-bound} proves $\mu_u(\{y\})=0$.
\end{proof}

\begin{proof}[Proof of \cref{prop:p21p3-gfc-smooth-rigidity}]
We first record three ingredients and then globalize them.

\paragraph{The Euclidean calculation.}
Let $\Gamma$ be a locally finite one-dimensional polyhedral set in an
open subset of $\mathbf R^N$, and let $S$ be a positive current of
bidimension $(1,1)$, supported on $\Gamma$, with
$\mathrm d'S=\mathrm d''S=0$.  Positivity and
\cite[Proposition~7.2.9]{CLD25} write $S$ as a symmetric positive
semidefinite matrix $M=(M_{jk})$ of Radon measures.  Closedness is the
distributional system
\begin{equation}
\label{eq:p21p3-gfc-matrix-divergence}
                   \sum_j\partial_jM_{jk}=0
                   \quad(1\leq k\leq N).
\end{equation}
On the relative interior of an edge, choose affine coordinates
$(t,z_2,\ldots,z_N)$ in which the edge is
$\{z_2=\cdots=z_N=0\}$.  Testing
\eqref{eq:p21p3-gfc-matrix-divergence} against
$\psi(t)z_\ell$, for $\ell>1$, kills every transverse row and column;
testing against $\psi(t)$ then makes the remaining measure a constant
multiple of Lebesgue measure.  Thus
\begin{equation}
\label{eq:p21p3-gfc-edge-density}
                     S|_e=b_e[\,e\,],
                     \qquad b_e\geq0,
\end{equation}
with $b_e$ constant.  Endpoint terms at a vertex give
\begin{equation}
\label{eq:p21p3-gfc-balancing}
                       \sum_{e\ni v}b_eu_{e,v}=0,
\end{equation}
where $u_{e,v}$ is the primitive outgoing tangent.  A positive remainder
at $v$ would have the form $P\delta_v$ with $P$ positive
semidefinite.  Once the endpoint terms balance, tests vanishing at $v$
with arbitrary first jet isolate the derivatives of $P\delta_v$ in
\eqref{eq:p21p3-gfc-matrix-divergence}; hence $P=0$.  In the vertex
applications below,
\cref{itm:p21p3-gf-bridge:localization} produces a one-vertex source
neighborhood on which the full nonconstant target star of the fixed finite
moment is finite.  We apply \eqref{eq:p21p3-gfc-balancing} to that entire
finite target star.  We never infer a balancing identity by deleting rays
and passing to an arbitrary finite substar.

\paragraph{Legitimate source-separated localization.}
We make explicit why no smooth plateau is being treated as an analytic
moment coordinate.  Let $\alpha$ be a compactly supported smooth
$(1,1)$-test form.  Start by localizing in the source.  The local moment
description \cite[Proposition~15.1.1(b)]{CLD25}, a smooth partition of
unity, compactness of $\operatorname{Supp}(\alpha)$, the finite-fibre
assertion \cite[Corollary~3.18]{GJR26}, and subdivision give finitely many
summands, each contained either in one relatively compact
virtual-annular collar or in one star neighborhood of a retained vertex.
The two collars at a type-$3$ point may be retained in its single star.
An edge neighborhood meets no other retained source edge, and a star
neighborhood contains no other retained vertex.  This source-first
operation is made before any Euclidean pushforward; distinct edge pieces
are never put into one compact source domain merely in order to separate
them by a cutoff.

Fix one summand in a source neighborhood $W$.  On an edge collar, apply
\cref{lem:p21p3-gf-legal-source-localization} directly.  At a type-$3$
center, first include in its prescribed block a $K$-analytic unit with
nonzero annular slope and shorten the two collars; the existence and
two-sided monomial norm formula are part of
\cite[\S\S4.2.11.2 and 4.3.9--4.3.10]{Ducros24CurveStructure}.  Apply the
same lemma after this shortening.

Suppose now that $W$ is a one-vertex star about a type-$2$ point $x$.
For the comparison at hand, let $A$ be the finite family of branches that
must be isolated; it contains the branches on which the form block is
nonconstant and, in the divisor calculation below, the finite support of
the principal divisor being tested.  We invoke
\cref{itm:p21p3-gf-bridge:localization} with $T$, this exact test summand
$\alpha$, this exact $W$, and the selected family $A$, and spell out its
two-stage output.  Its first application of
\cref{lem:p21p3-gf-legal-source-localization} produces a legal form and
retained-cutoff moment $F_0$.  After the branch block is appended, write
$G_0=(F_0,F_{\mathrm{br}})$.  The bridge then shrinks to its canonical-collar
one-vertex source neighborhood $W_1\subset W$.  Every virtual-disc
component is constant for $G_{0,\mathrm{trop}}$ on its entirety; only
finitely many retained virtual-annular components can be active, and every
selected active ray is separated from every other active branch.

Crucially, the bridge next reapplies
\cref{lem:p21p3-gf-legal-source-localization} \emph{inside $W_1$}, with
$G_0$ as its prescribed block.  We use its re-localized output
\begin{equation}
\label{eq:p21p3-gfc-relocalized-output}
 V\subset W_1,\qquad F,\qquad
 \Omega=\mathbf R^N\setminus
 F_{\mathrm{trop}}\bigl(\partial(V/W_1)\bigr),\qquad
 U=F_{\mathrm{trop}}^{-1}(\Omega).
\end{equation}
For every selected branch $b_i\in A$, it supplies the full canonical collar
$Z_{b_i}\subset U$ and the exact identity
\begin{equation}
\label{eq:p21p3-gfc-imported-exact-preimage}
 F_{\mathrm{trop}}^{-1}\!\left(
 F_{\mathrm{trop}}(x)+(0,\varepsilon_i)v_i\right)\cap U
 =Z_{b_i},
\end{equation}
which is precisely
\eqref{eq:p21p3-gf-source-separated-ray-preimage}.  Thus the isolated
target ray receives exactly its named source collar, not merely a branch
with the same ray germ.  This is the source-separation fact used below.

The full moment $F$ has finitely many analytic-unit coordinates.  By
\cite[Proposition~4.2.4(ii)]{Ducros24CurveStructure}, only finitely many
branches at $x$ are nonconstant for at least one of them.  Hence the
\emph{full} nonconstant target star in this localized moment is finite.
All of its rays, including rays activated by the second localization, are
retained when we apply Euclidean balancing.  We do not discard the
remaining geometric branches by a finite-fibre argument: they are constant
for $F$ on cofinal collars and map to the vertex.

In every case, the legal source-localization lemma gives a genuine analytic
block $H$ and a smooth $\Phi$ with
$\chi=\Phi\circ H_{\mathrm{trop}}$ whenever a source cutoff $\chi$ is
retained.  We append $H$, never $\chi$.  It also supplies the compact
domain $V$, the full moment $F$, and the proper map
$F_{\mathrm{trop}}|_U:U\to\Omega$; in the type-$2$ case these are exactly
the re-localized objects in
\eqref{eq:p21p3-gfc-relocalized-output}.  Thus no nonlinear cutoff is
declared to be a moment coordinate, and the current is never multiplied by
a cutoff.  The proper Euclidean pushforward
\begin{equation}
\label{eq:p21p3-gfc-euclidean-pushforward}
 \langle T_\Omega,\beta\rangle
       :=\langle T|_U,F^*\beta\rangle,
 \qquad \beta\in A_c^{1,1}(\Omega),
\end{equation}
defines a positive closed Euclidean current supported on a locally finite
one-dimensional polyhedral set.  Restriction to the open set $U$, not
multiplication by a cutoff, preserves closedness.  In a type-$2$ star,
\eqref{eq:p21p3-gfc-imported-exact-preimage} ensures that an isolated
selected ray receives the contribution of exactly its named source collar.

On such a collar, let $b_e\geq0$ be the flux of $T_\Omega$ from
\eqref{eq:p21p3-gfc-edge-density}, and let $m_e>0$ be the corresponding
flux for the canonical integration current.  Put
\begin{equation}
\label{eq:p21p3-gfc-branch-ratio}
                              a_e:=b_e/m_e.
\end{equation}
On a common source-separated concatenation of two moments, tropical
pushforward multiplies $b_e$ by the tangent-lattice expansion index:
the tangent tensor supplies its square and one-dimensional measure change
supplies its inverse.  The calibrated flux $m_e$ is multiplied by the
same index by \cite[\S14.3.2, Lemma~14.3.4, and
Proposition~14.4.2]{CLD25}.  Therefore $a_e$ is independent of the
moment, subdivision, and localized presentation, and is constant along
the virtual-annular collar.  This comparison is entirely downstairs over
$K$.

\paragraph{Comparison at type-$3$ and type-$2$ points.}
At an interior type-$3$ point the two branch rays in the refined moment
are opposite rays of one line.  Applying
\eqref{eq:p21p3-gfc-balancing} first to $T$ and then to $\delta_Y$
shows that the ratios \eqref{eq:p21p3-gfc-branch-ratio} on the two sides
are equal.  Possible virtual-annulus expansion factors occur equally in
the ordinary and calibrated fluxes and have already cancelled in the
ratio.

Let $x$ now be an interior type-$2$ point, and let $C_x$ be its
normal projective integral residual curve.  For a closed point
$p\in C_x$, use a rational function $s_p$ with
$\operatorname{ord}_p(s_p)=1$, a norm-one analytic lift, and
\cref{itm:p21p3-gf-bridge:ray}.  Apply
\cref{itm:p21p3-gf-bridge:localization} with the $p$-branch among the
selected branches.  Its re-localized exact-preimage identity
\eqref{eq:p21p3-gfc-imported-exact-preimage} makes the resulting ray the
image of precisely that source collar.  The outgoing downstairs vector
has a coordinate equal to one and is therefore primitive.  Denote the
calibrated flux on this isolated branch ray by $m_p>0$, and its ratio by
$a_p$.  These quantities are independent of the normalized moment: in
a common concatenation the source tangent and both projected tangents
retain a coordinate equal to one, so both projection indices are one.

Define additive maps on the divisor group of $C_x$ by
\begin{equation}
\label{eq:p21p3-gfc-divisor-functionals}
 L_T([p])=a_pm_p,
 \qquad
 L_Y([p])=m_p.
\end{equation}
They are nonnegative on effective divisors.  If
$g\in\operatorname{Rat}(C_x)^\times$, append a norm-one analytic lift
of $g$, and apply the source-separated localization in
\cref{itm:p21p3-gf-bridge:localization} simultaneously to the finite set
$\operatorname{Supp}(\operatorname{div}(g))$.  In the scalar coordinate
of this lift, the primitive slope on the $p$-branch is
$\operatorname{ord}_p(g)$; every branch outside that finite support has
slope zero.  The scalar component of
\eqref{eq:p21p3-gfc-balancing} gives
\begin{equation}
\label{eq:p21p3-gfc-T-principal}
       0=\sum_p a_pm_p\operatorname{ord}_p(g)
        =L_T(\operatorname{div}(g)).
\end{equation}
The integration current is closed by \cite[\S7.7.4]{CLD25}; applying the
same calculation to $\delta_Y$ gives
\begin{equation}
\label{eq:p21p3-gfc-Y-principal}
       0=\sum_p m_p\operatorname{ord}_p(g)
        =L_Y(\operatorname{div}(g)).
\end{equation}
No infinite sum and no cutoff on an infinite-valence star is used.

We use the following elementary consequence of Riemann--Roch over the
actual residue field.  Let $C$ be a normal projective integral curve over
a field $k$, and let
$L:\operatorname{Div}(C)\to\mathbf R$ be additive, nonnegative on
effective divisors, and zero on principal divisors.  If $D$ has degree
zero, choose an effective divisor $A$ such that
$\deg A+\chi(C,\mathcal O_C)>0$.
For every $n\in\mathbf Z$, Riemann--Roch gives
\[
 h^0(C,\mathcal O_C(A+nD))
 \geq \deg A+\chi(C,\mathcal O_C)>0.
\]
Thus $A+nD$ is linearly equivalent to an effective divisor and
\[
                         0\leq L(A)+nL(D)
                         \qquad(n\in\mathbf Z).
\]
Letting $n$ tend to both signs of infinity gives $L(D)=0$.  Choose
any closed point $p_0$ of $C$.  Since
$\deg(p_0)D-\deg(D)[p_0]$
has degree zero,
\begin{equation}
\label{eq:p21p3-gfc-degree-functional}
                 L(D)=c_L\deg(D),
 \qquad
                 c_L:=\frac{L([p_0])}{\deg(p_0)}\geq0.
\end{equation}
This argument neither assumes geometric integrality nor uses a degree-one
closed point.

Apply \eqref{eq:p21p3-gfc-degree-functional} to $L_Y$ and $L_T$.
There are $c_Y>0$ and $c_T\geq0$ such that
\[
                  m_p=c_Y\deg(p),
 \qquad
                  a_pm_p=c_T\deg(p)
\]
for every closed point $p$.  Hence
\begin{equation}
\label{eq:p21p3-gfc-equal-star-ratios}
                              a_p=c_T/c_Y
\end{equation}
on every branch at $x$.

\paragraph{Globalization on the smooth curve.}
Let $E\subset Y$ be the set of type-$1$ and type-$4$ points.  We justify
the endpoint assertion over the original field.  At a nonrigid point,
\cite[Theorem~4.2.5]{Ducros24CurveStructure} says that the existence of
two branches forces type $2$ or $3$; purity and the local graph structure
give at least one branch.  At a rigid point the virtual-disc neighborhood
of \cite[Theorem~4.5.4 and \S4.5.4.1]{Ducros24CurveStructure} gives the
same conclusion.  Thus every point of $E$ is unibranch, as is also
recorded for nonrigid type-$1$/type-$4$ points in
\cite[\S4.5.19.1]{Ducros24CurveStructure}.

Connected analytic spaces are arcwise connected by
\cite[\S12.1.12]{CLD25}.  If two points lie outside $E$, choose an arc
joining them.  By \cite[\S1.3.12]{Ducros24CurveStructure}, a point lying
in the open interval of a graph has at least the two branches supplied by
the two sides.  Hence no unibranch point of $E$ can occur in the interior
of the chosen arc, and $Y\setminus E$ is connected.  The ratio is
constant on each annular collar, propagates through type-$3$ points by
the two-ray comparison, and propagates through type-$2$ points by
\eqref{eq:p21p3-gfc-equal-star-ratios}.  It is therefore one constant
$c\geq0$ on $Y\setminus E$.

The source-separated local moments tropicalize every local test form.
The edge-density calculation and the absence of a vertex remainder then
show $T=c\delta_Y$ on a neighborhood of each point of $Y\setminus E$.
Every point of $E$ has tropical dimension zero, so
\cite[Proposition~3.3.6]{CLD25} says that the germ of every smooth
$(1,1)$-test form vanishes there.  Thus the support of each such test
form is disjoint from $E$, and the current $T-c\delta_Y$, already zero
off $E$, pairs to zero with every test form.  This proves
\eqref{eq:p21p3-gfc-smooth-rigidity}.  Uniqueness follows by testing on a
nonempty virtual-annular band.  All arguments were made on compact test
supports, so $Y$ need not be compact.
\end{proof}

\begin{proof}[Proof of \cref{prop:p21p3-curve-rigidity}]
We first pass from the smooth normalizations to the reduced curve.  On a
good affinoid domain write
\[
 \mathcal Q:=\nu_*\mathcal O_{\widetilde X}/\mathcal O_X,
 \qquad
 \mathfrak c:=\operatorname{Ann}_{\mathcal O_X}(\mathcal Q),
\]
and let $Z$ be the union of
\[
 \operatorname{Supp}(\mathcal Q)
   =V(\mathfrak c)
   =\operatorname{Supp}(\mathcal O_X/\mathfrak c)
\]
with the pairwise intersections of distinct irreducible components.
Analytic local rings on a good space are excellent by
\cite[\S11.1.10]{CLD25}; normalization is finite and its components
correspond to the irreducible components by
\cite[\S11.3.1]{CLD25}.  Hence $\mathcal Q$ and the conductor are
coherent.  Their support misses every generic point of the reduced pure
curve and is therefore zero-dimensional.  The same dimension argument
applies to proper intersections of distinct components.  Irreducible
components are locally finite by \cite[\S11.2.5]{CLD25}; consequently
$Z$ is closed and locally finite.

Strictness is used here.  A zero-dimensional strict affinoid algebra is
Artinian, and the strict affinoid Nullstellensatz makes the completed
residue field of each of its points finite over $K$.  Thus every point
of $Z$ is rigid.  A finite morphism induces a finite extension of
completed residue fields, so every point of
\begin{equation}
\label{eq:p21p3-gfc-conductor-preimage}
                         Z_i:=\nu_i^{-1}(Z),
\end{equation}
is rigid as well, though it need not be $K$-rational.  Moreover, a finite
algebra over a strict affinoid algebra is strict.  Since normalization is
finite, every $\widetilde X_i$ is therefore strict.

The finite normalization map is topologically proper, has empty relative
boundary, and is an isomorphism away from $Z_i$.  Applying the
boundary-composition identity \cite[(11.4.1.3)]{CLD25} to
$X_i\to X\to\mathcal M(K)$ and then to
$\widetilde X_i\to X_i\to\mathcal M(K)$ shows that
$\widetilde X_i$ is boundaryless.  It is connected because it is the
normalization of one irreducible component.

The fundamental-cycle identity is
\begin{equation}
\label{eq:p21p3-gfc-fundamental-cycle}
             \delta_X=\sum_i(\nu_i)_*\delta_{\widetilde X_i}.
\end{equation}
Indeed, a reduced pure curve has generic multiplicity one on every
irreducible component.  Integration agrees with the associated cycle by
\cite[Proposition~16.6.3]{CLD25}, and normalization has calibrated degree
one over the normal locus.  The complement is the locally finite rigid
set $Z$, where every $(1,1)$-test-form germ vanishes by
\cite[Proposition~3.3.6]{CLD25}; this proves
\eqref{eq:p21p3-gfc-fundamental-cycle} everywhere.

The complement $\widetilde X_i\setminus Z_i$ is open, good, separated,
strict, rig-smooth, and boundaryless.  Deleting $Z_i$ does not disconnect
$\widetilde X_i$.  Indeed, two
remaining points can be joined by an arc, and
\cite[\S1.3.12]{Ducros24CurveStructure} shows that no rigid unibranch
point can lie in the interior of that arc.  On $X\setminus Z$, the component
images are disjoint and normalization is an analytic isomorphism.
Applying \cref{prop:p21p3-gfc-smooth-rigidity} to each connected
$\widetilde X_i\setminus Z_i$ gives
\[
                         T=\sum_i c_i\delta_{X_i},
                         \qquad c_i\geq0,
\]
there.  Set $S=T-\sum_i c_i\delta_{X_i}$.  It has zero germ off $Z$,
while every $(1,1)$-test form has support disjoint from $Z$ by
\cite[Proposition~3.3.6]{CLD25}.  Hence $S=0$, proving
\eqref{eq:p21p3-curve-classification}.  The sum is locally finite,
and uniqueness follows by restricting to an open normal locus of $X_i$
meeting no other component.

It remains to prove the Bedford--Taylor assertion.  Put
$\widetilde W_i=\nu_i^{-1}(W)$.  Multiplying the order-zero identity
\eqref{eq:p21p3-gfc-fundamental-cycle} by a continuous function $F$ on
$W$, using \cite[Example~7.2.12]{CLD25}, gives
\begin{equation}
\label{eq:p21p3-gfc-function-cycle-decomposition}
 [F]=F\delta_W
     =\sum_i(\nu_i)_*
        \bigl((\nu_i^*F)\delta_{\widetilde W_i}\bigr).
\end{equation}
The projection identity used here follows directly by testing against a
smooth form, and the sum is finite on each test support.

Pullback preserves continuous plurisubharmonic functions by
\cite[Proposition~8.1.5]{CLD25}.  On $\widetilde W_i$, set
\[
 \mu_i^U=\mathrm d'\mathrm d''[\nu_i^*U],
 \qquad
 \mu_i^V=\mathrm d'\mathrm d''[\nu_i^*V].
\]
These are positive Radon measures by \cite[Lemma~8.5.5]{CLD25}.
Proper pushforward commutes with the differential on currents by
transposition; equivalently, use \cite[Proposition~8.5.6]{CLD25}.
Applying \eqref{eq:p21p3-gfc-function-cycle-decomposition} to $U$ and
$V$, the curvature equality
\eqref{eq:p21p3-curve-equal-curvature} becomes
\begin{equation}
\label{eq:p21p3-gfc-component-measure-sum}
                    \sum_i(\nu_i)_*(\mu_i^U-\mu_i^V)=0.
\end{equation}

Away from $Z$, the component images are disjoint and $\nu_i$ is an
isomorphism.  Hence each signed measure $\mu_i^U-\mu_i^V$ is supported
on the locally finite rigid set $Z_i\cap\widetilde W_i$.  At any point of
this set, apply \cref{lem:p21p3-gfc-rigid-zero-atom} on a connected open
neighborhood contained in $\widetilde W_i$.  Neither positive measure has
an atom there.  A signed Radon measure supported on a locally finite set
is the locally finite sum of its point masses.  Therefore
\begin{equation}
\label{eq:p21p3-gfc-component-measure-equality}
                              \mu_i^U=\mu_i^V
                              \quad\text{for every }i.
\end{equation}

Finally put $h=U-V$.  By
\eqref{eq:p21p3-curve-classification},
\[
 hT|_W=\sum_i c_i(\nu_i)_*
       \bigl((\nu_i^*h)\delta_{\widetilde W_i}\bigr).
\]
Proper pushforward commutes with $\mathrm d'\mathrm d''$, and
\eqref{eq:p21p3-gfc-component-measure-equality} gives
\begin{align*}
 \mathrm d'\mathrm d''(hT|_W)
 &=\sum_i c_i(\nu_i)_*
       \mathrm d'\mathrm d''[\nu_i^*U-\nu_i^*V] \\
 &=\sum_i c_i(\nu_i)_*(\mu_i^U-\mu_i^V)=0.
\end{align*}
This is \eqref{eq:p21p3-curve-bt-conclusion}.  It proves the
independence condition of \cite[Definition~8.6.2]{CLD25} for the
continuous psh pairs stated here; no claim about discontinuous psh
functions is made.
\end{proof}


The next statement isolates a different one-dimensional enlargement.  It
does not pass through normalization: normality is assumed from the outset,
and a locally finite rigid exceptional set is removed by the support theorem
for smooth $(1,1)$-forms.  The following statement and proof are the exact
bytes accepted by production verification.

% Exact staging-only statement for production Rethlas verification.
% Date: 29 August 2026.
% This file is not a live-book input and makes no novelty or priority claim.

\begin{proposition}
\label{prop:p21p3-rigid-exception-extension}
Let $K$ be a complete nontrivially valued non-Archimedean field, and let
$X$ be a nonempty connected good separated strict normal $K$-analytic curve
without boundary.  Let $E\subset X$ be a locally finite set of rigid points, put
$Y:=X\setminus E$, and assume that $Y$ is rig-smooth.  Then the following
assertions hold.

\begin{enumerate}
\item\label{itm:p21p3-rigid-exception:connected}
The open analytic curve $Y$ is nonempty and connected.

\item\label{itm:p21p3-rigid-exception:classification}
Every positive $\mathrm d'$- and $\mathrm d''$-closed CLD current $T$ of
bidegree $(0,0)$ on $X$ has the form
\begin{equation}
\label{eq:p21p3-rigid-exception-classification}
                              T=c\delta_X
\end{equation}
for a unique $c\in\mathbf R_{\geq0}$.

\item\label{itm:p21p3-rigid-exception:bedford-taylor}
Let $W\subseteq X$ be open, and let $u,v$ be continuous
CLD-plurisubharmonic functions on $W$.  If
\begin{equation}
\label{eq:p21p3-rigid-exception-equal-curvature}
                  \mathrm d'\mathrm d''[u]
                  =\mathrm d'\mathrm d''[v],
\end{equation}
then every current $T$ as in
\cref{itm:p21p3-rigid-exception:classification} satisfies
\begin{equation}
\label{eq:p21p3-rigid-exception-bt-conclusion}
              \mathrm d'\mathrm d''\bigl((u-v)T|_W\bigr)=0.
\end{equation}
Thus these currents have the Bedford--Taylor independence property for
continuous psh pairs in the CLD sense.
\end{enumerate}
\end{proposition}

\begin{remark}
\label{rem:p21p3-rigid-exception-scope}
The statement does not remove nontrivial valuation, goodness,
separatedness, strictness, normality, or boundarylessness.  The exceptional
set consists only of rigid points and is locally finite, and its complement
is assumed rig-smooth.  The Bedford--Taylor conclusion concerns continuous
psh pairs; it does not assert the literal possibly discontinuous predicate
of \cite[Definition~8.6.2]{CLD25}.  No assertion is made in analytic
dimension at least two.
\end{remark}

% Exact staging-only proof for production Rethlas verification.
% Date: 29 August 2026.
% This file is not a live-book input and makes no novelty or priority claim.

\begin{proof}[Proof of
\cref{prop:p21p3-rigid-exception-extension}]
Because $X$ is separated, its underlying topological space is Hausdorff.
A locally finite union of closed subsets of a Hausdorff space is closed;
applied to the singletons in $E$, this shows that $E$ is closed.  Hence
$Y=X\setminus E$ is an open analytic subspace of $X$.  Goodness,
separatedness, and strictness pass to open analytic subspaces.  An open
immersion has empty relative boundary, and the boundary-composition formula
\cite[(11.4.1.3)]{CLD25} therefore shows that $Y$ is boundaryless.
We shall also use the corresponding separation property for currents.
Every good separated analytic space is locally holomorphically separated by
\cite[\S15.2.2]{CLD25}, and is therefore locally differentially separated by
\cite[Corollary~15.2.4]{CLD25}.  This applies to $X$ and to its open analytic
subspaces, in particular to $Y$ and to the open sets considered below.

We next prove the assertion about connectedness.  By
\cite[Theorem~3.5.1]{Ducros24CurveStructure}, the underlying topological
space of every analytic curve is a graph.  A connected graph is connected by
injective arcs by \cite[\S1.3.5]{Ducros24CurveStructure}.  On the other hand,
\cite[Proposition~3.5.2]{Ducros24CurveStructure} states that the number of
branches at a rigid point of an analytic curve equals the cardinality of its
inverse image in the normalization.  Since $X$ is normal, every point of $E$
is consequently unibranch.  By
\cite[\S1.3.12]{Ducros24CurveStructure}, every point lying in the interior of
an injective arc is pluribranch.

The set $Y$ is nonempty.  Indeed, if $E$ is empty this is immediate; if
$x\in E$, the fact that $x$ is unibranch gives a nonempty branch issuing
from $x$.  Choose a connected component of $U\setminus\{x\}$ representing
that branch for a sufficiently small neighborhood $U$ of $x$, and choose a
point $z$ in that component.  After shrinking $U$, the graph structure
supplies an injective arc from $x$ to $z$ with nonempty interior.  Any point
in that interior is pluribranch, hence does not belong to $E$.  Now let
$y_0,y_1\in Y$.  If $y_0=y_1$, there is nothing to prove.
Otherwise choose an injective arc $I$ in $X$ joining them.  If $I$ met $E$,
any point of $I\cap E$ would be an interior point of $I$, because its
endpoints lie in $Y$.  Such a point would be both pluribranch and unibranch,
a contradiction.  Thus $I\subset Y$.
This proves \cref{itm:p21p3-rigid-exception:connected}, in fact by showing
that $Y$ is connected by injective arcs.

Let $T$ be as in
\cref{itm:p21p3-rigid-exception:classification}.  Restriction to the open
subspace $Y$ is defined by extension by zero of compactly supported test
forms, as in \cite[\S7.1.6]{CLD25}.  It preserves positivity; by the locality
and differential identities of \cite[\S\S7.7.1--7.7.3]{CLD25}, it also
preserves $\mathrm d'$- and $\mathrm d''$-closedness.  The curve $Y$ has all
the hypotheses of the already established rig-smooth classification
\cref{prop:p21p3-gfc-smooth-rigidity}.  Hence there is a unique
$c\in\mathbf R_{\geq0}$ such that
\begin{equation}
\label{eq:p21p3-rigid-exception-on-complement}
                            T|_Y=c\delta_Y.
\end{equation}
By the definition of the integration current in
\cite[\S7.3.1]{CLD25} and the restriction operation of
\cite[\S7.1.6]{CLD25}, one has $\delta_X|_Y=\delta_Y$.  Put
$S:=T-c\delta_X$.  Then $S|_Y=0$ by
\eqref{eq:p21p3-rigid-exception-on-complement}.  Let $\alpha$ be an
arbitrary compactly supported smooth $(1,1)$-test form on $X$.  A rigid
point $x$ satisfies $d_K(x)=0$: indeed, rigid means that
$\mathscr H(x)/K$ is finite by \cite[\S11.1.10]{CLD25}, and both the residue
transcendence-degree term and the rational-rank term in the dimension formula
\cite[(11.2.3.1)]{CLD25} then vanish.  This remains true for a finite
inseparable extension.  The analytic form-support statement
\cite[Proposition~15.2.1(a)]{CLD25}, which is the analytic consequence of
the tropical support vanishing in \cite[Proposition~3.3.6]{CLD25}, gives
\begin{equation}
\label{eq:p21p3-rigid-exception-test-support}
 \operatorname{Supp}(\alpha)
 \subseteq\{x\in X:d_K(x)\geq1\}
 \subseteq X\setminus E=Y.
\end{equation}
Thus $\alpha|_Y$ is a compactly supported test form on $Y$, and therefore
\[
                   \langle S,\alpha\rangle
                   =\langle S|_Y,\alpha|_Y\rangle=0.
\]
Since $\alpha$ was arbitrary, $S=0$.  This proves
\eqref{eq:p21p3-rigid-exception-classification}.  Uniqueness of $c$ follows
after restriction to $Y$ from uniqueness in
\cref{prop:p21p3-gfc-smooth-rigidity}.

Finally, let $W,u,v$ satisfy
\eqref{eq:p21p3-rigid-exception-equal-curvature}, and put $h=u-v$.
The local differential separation recorded in the first paragraph puts us
within the hypotheses of the order-zero result.  Positive currents have
order zero by
\cite[Proposition~7.2.9]{CLD25}, so multiplication by the continuous
function $h$ is defined as in \cite[Example~7.2.12]{CLD25}.  Corollary~7.2.10
of \cite{CLD25} defines $[h]$ for the locally integrable function $h$.  For
every compactly supported smooth $(1,1)$-test form $\alpha$ on $W$, the
definitions in \cite[Corollary~7.2.10, Example~7.2.12, and
\S7.3.1]{CLD25} give
\[
 \bigl\langle h(\delta_X|_W),\alpha\bigr\rangle
 =\int_W h\alpha
 =\langle[h],\alpha\rangle.
\]
Consequently $h(\delta_X|_W)=[h]$.  Using
\eqref{eq:p21p3-rigid-exception-classification} and linearity gives
\begin{align*}
 \mathrm d'\mathrm d''\bigl(hT|_W\bigr)
 &=c\,\mathrm d'\mathrm d''[h]\\
 &=c\bigl(\mathrm d'\mathrm d''[u]
          -\mathrm d'\mathrm d''[v]\bigr)=0.
\end{align*}
This is \eqref{eq:p21p3-rigid-exception-bt-conclusion} and proves
\cref{itm:p21p3-rigid-exception:bedford-taylor}.
\end{proof}


% Live-book statement adaptation of the arbitrary-ground-field nonreduced
% addendum.  Date: 28 August 2026.

\begin{proposition}
\label{prop:p21p3-curve-rigidity-nonreduced}
Let \(K\) be a complete
nontrivially valued non-Archimedean field and let \(X\) be a good separated
strict \(K\)-analytic space, pure of dimension one and without boundary.
The space \(X\) is not assumed reduced.  Write

\[
             \iota:X_{\mathrm{red}}\hookrightarrow X
\]

for the canonical nilpotent closed immersion.  Let \((X_i)_i\) be the
irreducible components of \(X_{\mathrm{red}}\), assume that their
normalizations are rig-smooth, and write

\[
 \nu_i:\widetilde X_i\longrightarrow X_i
       \hookrightarrow X_{\mathrm{red}}\xrightarrow{\iota}X,
 \qquad
 \delta_i:=(\nu_i)_*\delta_{\widetilde X_i}.
\]

Let \(m_i\in\mathbf Z_{>0}\) be the coefficient of \([X_i]\) in the
fundamental cycle \([X]\).  Then the following assertions hold.

\begin{enumerate}
\item\label{itm:p21p3-nr:calculus}
For every \((p,q)\), pullback by \(\iota\) identifies the sheaves of
ordinary CLD-smooth forms on \(X\) and \(X_{\mathrm{red}}\), compatibly
with wedge product, \(J\), \(\mathrm d'\), \(\mathrm d''\), support, and all
three positivity cones.  On compactly supported test forms this is an
isomorphism of topological vector spaces.  Its transpose identifies the
CLD current sheaves, their weak topologies, positivity, and
\(\mathrm d'\)- and \(\mathrm d''\)-closedness.  It identifies multiplication
by smooth functions for all currents.  If \(S\) is positive and \(f\) is a
locally bounded Borel function, it also identifies the CLD product \(fS\);
by linearity the same holds for currents locally expressible as differences
of positive currents.  No multiplication by a Borel function is asserted
for an arbitrary order-zero current.  Under the common underlying
topological space, the CLD-psh and continuous CLD-psh sheaves on \(X\) and
\(X_{\mathrm{red}}\) coincide.  The analogous pullback identification
holds for \(G\)-forms.

\item\label{itm:p21p3-nr:integration}
The integration current is not, in general, invariant under reduction.
Precisely,

\begin{equation}
\label{eq:p21p3-nr-fundamental-cycle}
        \delta_X=\sum_i m_i\delta_i,
        \qquad
        \iota_*\delta_{X_{\mathrm{red}}}=\sum_i\delta_i.
\end{equation}

Both sums are locally finite.  In particular,
\(\delta_X=\iota_*\delta_{X_{\mathrm{red}}}\) if and only if \(X\) is
generically reduced along every \(X_i\).

\item\label{itm:p21p3-nr:classification}
Every positive \(\mathrm d'\)- and \(\mathrm d''\)-closed CLD current \(T\)
of bidegree \((0,0)\) on \(X\) has a unique expression

\begin{equation}
\label{eq:p21p3-nr-classification}
                     T=\sum_i c_i\delta_i,
                     \qquad c_i\in\mathbf R_{\geq0}.
\end{equation}

The sum is locally finite.  If \(X_{\mathrm{red}}\) is connected and
rig-smooth, there is one generic multiplicity \(m\), and every such current
has the form \(T=c\delta_X\) for a unique \(c\geq0\).

\item\label{itm:p21p3-nr:bedford-taylor}
Let \(W\subseteq X\) be open and let \(U,V\) be continuous CLD-psh
functions on \(W\).  Identify \(W\) with its inverse image in
\(X_{\mathrm{red}}\), and, for every continuous function \(F\) on \(W\),
put

\[
 [F]_X:=F(\delta_X|_W),\qquad
 [F]_{X_{\mathrm{red}}}:=F(\delta_{X_{\mathrm{red}}}|_W).
\]

With the identifications in \cref{itm:p21p3-nr:calculus}, one has the
exact equivalence

\begin{equation}
\label{eq:p21p3-nr-curvature-equivalence}
 \mathrm d'\mathrm d''[U]_X=\mathrm d'\mathrm d''[V]_X
 \quad\Longleftrightarrow\quad
 \mathrm d'\mathrm d''[U]_{X_{\mathrm{red}}}
   =\mathrm d'\mathrm d''[V]_{X_{\mathrm{red}}}.
\end{equation}

If these equivalent conditions hold, then every \(T\) in
\cref{itm:p21p3-nr:classification} satisfies

\begin{equation}
\label{eq:p21p3-nr-bt-conclusion}
       \mathrm d'\mathrm d''\bigl((U-V)T|_W\bigr)=0.
\end{equation}

Consequently every positive closed CLD current of bidegree \((0,0)\) on
\(X\) has the Bedford--Taylor independence property for continuous psh
pairs.  Under the current identification in
\cref{itm:p21p3-nr:calculus}, this continuous-pair property is invariant
under the nilpotent thickening
\(X_{\mathrm{red}}\hookrightarrow X\).
\end{enumerate}
\end{proposition}

\begin{remark}
\label{rem:p21p3-nr-scope}
No algebraic-closedness hypothesis is made.  Strictness is used
to ensure that the zero-dimensional conductor locus and its inverse image
in every normalization consist of locally finite rigid points; the rigid
zero-atom lemma then replaces the algebraically closed/type-\(1\) input.
The addendum does not treat a trivially valued field, a non-strict or
non-good curve, boundary points, a component with non-rig-smooth
normalization, or analytic dimension at least two.  The last item concerns
only continuous psh pairs.  Since CLD Definition~8.6.2 literally
quantifies over psh functions that may be discontinuous, no assertion of
that stronger predicate is made here.  Nilpotents enter only through the
generic multiplicities \(m_i\); no nilpotence exponent or transverse
vanishing-order estimate is used.
\end{remark}

\begin{remark}
\label{rem:p21p3-nr-double-line}
For

\[
 X=\bigl(\operatorname{Spec}K[T,\epsilon]/(\epsilon^2)\bigr)^{\mathrm{an}},
 \qquad X_{\mathrm{red}}=\mathbf A_K^{1,\mathrm{an}},
\]

the unique generic multiplicity is \(2\).  Hence all smooth forms and all
positive closed currents are identified with those on the analytic line,
whereas
\(\delta_X=2\iota_*\delta_{\mathbf A_K^{1,\mathrm{an}}}\).
Thus generic multiplicity, and not a new nilpotent differential form or
current, is the precise obstruction to reduction-invariance of the
integration current.  The proof uses only the normalizations of the
components of \(X_{\mathrm{red}}\); no normalization of the nonreduced
space is needed.
\end{remark}

% Live-book proof adaptation of the arbitrary-ground-field nonreduced
% addendum.  Date: 28 August 2026.

\begin{proof}
The reduction has the same underlying topological space and the same
completed residue fields as \(X\).  Affinoid charts pass to quotients by
nilpotent ideals, so goodness, separatedness, strictness, and pure
dimension one pass to \(X_{\mathrm{red}}\).  The immersion \(\iota\) is
finite and hence has empty relative boundary by
\cite[\S11.4.4]{CLD25}.  The boundary-composition identity
\cite[(11.4.1.3)]{CLD25} therefore shows that
\(X_{\mathrm{red}}\) is boundaryless.  Thus \(X_{\mathrm{red}}\)
satisfies the ambient hypotheses of the reduced theorem
\cref{prop:p21p3-curve-rigidity}.

\smallskip
\emph{Forms, currents, and psh functions.}
By \cref{itm:ch13-pb:red} of
\cref{prop:ch13-pullback-properties}, proved from
\cite[Proposition~15.1.3(b) and the following form variant]{CLD25},
pullback gives an isomorphism

\begin{equation}
\label{eq:p21p3-nr-form-isomorphism}
       \iota^*:A_X^{p,q}\xrightarrow{\ \sim\ }
                       A_{X_{\mathrm{red}}}^{p,q}.
\end{equation}

The point needed for currents is topological, not merely algebraic.  A
datum for the compact-test topology consists of a compact support and a
finite family of moment charts.  Restriction sends such a datum on \(X\)
to the same datum on \(X_{\mathrm{red}}\).  Conversely, the cited
almost-good result applies because \(X\) is good.  At every point of the
compact support, a finite tuple of invertible functions defining a reduced
moment chart lifts through the nilpotent quotient to a tuple of unit germs
on \(X\).  Each lift has the same absolute value as its reduction at every
analytic point.  After shrinking the ordinary open neighbourhoods and
using compactness, one obtains a finite lifted refinement; product moment
maps give a common refinement on overlaps.

Restriction and these finite lifted refinements are cofinal systems of
compact chart data for the final locally convex topologies.  On a lifted
chart, the tropical images, Lagerberg representatives, polyhedral
refinements, and every \(C^r\)-seminorm are identical.  Consequently

\begin{equation}
\label{eq:p21p3-nr-test-homeomorphism}
 \iota^*:A^{p,q}_{X,c}\xrightarrow{\ \sim\ }
                         A^{p,q}_{X_{\mathrm{red}},c}
\end{equation}

is a homeomorphism.  The same chart description proves compatibility with
wedge product, \(J\), the differentials, supports, and all positivity cones
of forms.  Taking continuous duals gives

\begin{equation}
\label{eq:p21p3-nr-current-isomorphism}
 \iota_*:D^{a,b}(X_{\mathrm{red}})
              \xrightarrow{\ \sim\ }D^{a,b}(X),
 \qquad
 \langle\iota_*S,\alpha\rangle
       =\langle S,\iota^*\alpha\rangle .
\end{equation}

This isomorphism identifies the weak topologies.  It preserves positivity
because the test-form cones agree, and it preserves closedness because
\(\iota^*\) commutes with \(\mathrm d'\) and \(\mathrm d''\).
Multiplication by a smooth function commutes with \(\iota_*\) directly by
transposition.  If \(S\) is positive, the Borel product of
\cite[Example~7.2.12]{CLD25} is defined through the measures
\(\mu_{S,\omega}\) attached to positive test forms \(\omega\).  Under
\(\iota^*\), the positive test forms and these attached measures agree.
Thus \(\iota_*(fS)=f\iota_*S\) for every locally bounded Borel function
\(f\).

If \(S=S_1-S_2\) locally with \(S_1,S_2\) positive, define
\(fS=fS_1-fS_2\).  This is independent of the decomposition.  Indeed, from
\(S_1-S_2=S'_1-S'_2\) one gets
\(S_1+S'_2=S'_1+S_2\); equality of these positive currents gives equality
of all their attached measures, and hence
\(fS_1+fS'_2=fS'_1+fS_2\).  The local definitions glue, and the same
calculation proves compatibility with \(\iota_*\).  Nothing in this
argument defines \(fS\) for an arbitrary order-zero current.

For the psh assertion, stage zero of the CLD filtration is the same on the
two spaces: smoothness and positivity of
\(\mathrm d'\mathrm d''u\) agree by
\eqref{eq:p21p3-nr-form-isomorphism}.  Successor stages take decreasing
limits, the continuous filtration takes local uniform limits, and limit
ordinal stages take local unions.  The underlying open sets and functions
are identical.  Transfinite induction in
\cref{def:ch16-psh-filtration} therefore identifies the psh and continuous
psh sheaves at every stage and after stabilization.  The \(G\)-form
argument is the same, using the \(G\)-form variant immediately after
\cite[Proposition~15.1.3]{CLD25}.  This proves
\cref{itm:p21p3-nr:calculus}.

\smallskip
\emph{The integration current and generic multiplicity.}
By \cite[Definition~13.2.5 and \S13.5.2]{CLD25}, the fundamental cycles
are

\[
                 [X]=\sum_i m_i[X_i],
                 \qquad
                 [X_{\mathrm{red}}]=\sum_i[X_i].
\]

Let
\(\bar\nu_i:\widetilde X_i\to X_i\hookrightarrow X_{\mathrm{red}}\),
so that \(\nu_i=\iota\circ\bar\nu_i\).  The normalization is finite and
generically an isomorphism.  The compatibility of integration with
analytic cycles in \cref{thm:current-fundamental-cycle}, namely
\cite[Proposition~16.6.3]{CLD25}, identifies the current of \([X_i]\) with
\((\bar\nu_i)_*\delta_{\widetilde X_i}\) on \(X_{\mathrm{red}}\), and
with \((\nu_i)_*\delta_{\widetilde X_i}\) after pushforward to \(X\).
It follows that

\[
 \delta_X=\sum_i m_i\delta_i,
 \qquad
 \iota_*\delta_{X_{\mathrm{red}}}=\sum_i\delta_i.
\]

The irreducible components are locally finite by
\cite[\S11.2.5]{CLD25}, so both sums are locally finite.  Their equality
forces \(m_i=1\) for every \(i\): restrict to a nonempty normal open subset
of \(X_i\) meeting no other component and use the nonzero integration
current there.  Conversely, \(m_i=1\) for every \(i\) visibly gives
equality.  Since the length of the generic local ring along \(X_i\) is
\(m_i\), this is equivalent to generic reducedness along every component.
This proves \cref{itm:p21p3-nr:integration}.

\smallskip
\emph{Classification.}
Let \(T\) be positive and closed on \(X\).  By
\eqref{eq:p21p3-nr-current-isomorphism}, there is a unique positive
closed current \(T_{\mathrm{red}}\) on \(X_{\mathrm{red}}\) with
\(T=\iota_*T_{\mathrm{red}}\).  The reduced classification
\cref{itm:p21p3-curve:classification} of
\cref{prop:p21p3-curve-rigidity} gives

\[
                  T_{\mathrm{red}}=\sum_i c_i
                    (\bar\nu_i)_*\delta_{\widetilde X_i},
                  \qquad c_i\geq0.
\]

Pushing forward by \(\iota\) proves
\eqref{eq:p21p3-nr-classification}, including local finiteness and
uniqueness.  If \(X_{\mathrm{red}}\) is connected and rig-smooth, its
regular local rings are domains, so its irreducible components are
disjoint open-and-closed subsets.  It therefore has one component, with
one generic multiplicity \(m\).  Then
\(\delta_X=m\delta_1\), and
\(c_1\delta_1=(c_1/m)\delta_X\).  This proves
\cref{itm:p21p3-nr:classification}.

\smallskip
\emph{The strict-rigid exceptional locus.}
Let
\(\bar\nu:\coprod_i\widetilde X_i\to X_{\mathrm{red}}\) be the
normalization map.  On a good affinoid domain put

\[
 \mathcal Q:=\bar\nu_*\mathcal O_{\coprod_i\widetilde X_i}
                 /\mathcal O_{X_{\mathrm{red}}},
 \qquad
 \mathfrak c:=\operatorname{Ann}_{\mathcal O_{X_{\mathrm{red}}}}
                         (\mathcal Q).
\]

Let \(Z\) be the union of the conductor locus

\[
 \operatorname{Supp}(\mathcal Q)
   =V(\mathfrak c)
   =\operatorname{Supp}(\mathcal O_{X_{\mathrm{red}}}/\mathfrak c)
\]

and the pairwise intersections of distinct irreducible components.
Analytic local rings on a good space are excellent by
\cite[\S11.1.10]{CLD25}; normalization is finite and its components
correspond to the irreducible components by
\cite[\S11.3.1]{CLD25}.  Hence \(\mathcal Q\) and the conductor are
coherent.  Their support misses every generic point of the reduced pure
curve and is zero-dimensional.  The same dimension argument applies to
the proper intersections of distinct components.  Local finiteness of the
components, again \cite[\S11.2.5]{CLD25}, shows that \(Z\) is closed and
locally finite.

Here strictness replaces algebraic closedness.  Locally, a
zero-dimensional closed analytic subspace of a strict affinoid domain is
the spectrum of a zero-dimensional strict affinoid algebra \(B\).  By
\cite[Theorem~4.5]{Tate71Rigid}, the algebra \(B\) is noetherian; having
Krull dimension zero, it is therefore Artinian.
For a point \(z\) of its Berkovich spectrum, the kernel of the associated
seminorm is maximal; the strict affinoid Nullstellensatz
\cite[Theorem~4.5]{Tate71Rigid} makes
\(B/\ker(z)\) finite over \(K\).  This finite-dimensional \(K\)-algebra is
already complete, so
\(\mathcal H(z)=B/\ker(z)\) is finite over \(K\).  Thus every point of
\(Z\) is rigid, although it need not be \(K\)-rational.

Put

\begin{equation}
\label{eq:p21p3-nr-conductor-preimage}
                         Z_i:=\bar\nu_i^{-1}(Z).
\end{equation}

A finite morphism induces finite extensions of completed residue fields,
so every point of \(Z_i\) is rigid.  A finite algebra over a strict
affinoid algebra is strict by \cite[Theorem~4.4]{Tate71Rigid};
consequently each \(\widetilde X_i\) is strict.  The finite maps
\(X_i\to X_{\mathrm{red}}\) and
\(\widetilde X_i\to X_i\) have empty relative boundary by
\cite[\S11.4.4]{CLD25}.  Applying
\cite[(11.4.1.3)]{CLD25} twice shows that every
\(\widetilde X_i\) is boundaryless.  It is good and separated, and it is
rig-smooth by hypothesis.  Away from \(Z\), the component images are
disjoint and each \(\bar\nu_i\) is an analytic isomorphism.  These are
exactly the hypotheses needed below for the rigid zero-atom input; no
assertion about arbitrary type-\(1\) points is used.

\smallskip
\emph{The multiplicity-weighted curvature identity.}
Fix \(W,U,V\) as in \cref{itm:p21p3-nr:bedford-taylor}, identify \(W\)
with the corresponding open subset of \(X_{\mathrm{red}}\), and put
\(\widetilde W_i=\bar\nu_i^{-1}(W)\).  Denote the base-changed maps by

\[
 \nu_{i,W}^{X}:\widetilde W_i\longrightarrow W\subset X,
 \qquad
 \nu_{i,W}^{\mathrm{red}}:\widetilde W_i\longrightarrow
                         W\subset X_{\mathrm{red}}.
\]

Pullback preserves continuous psh functions by
\cite[Proposition~8.1.5]{CLD25}.  On \(\widetilde W_i\), define positive
Radon measures

\[
 \mu_i^U:=\mathrm d'\mathrm d''[\bar\nu_i^*U],
 \qquad
 \mu_i^V:=\mathrm d'\mathrm d''[\bar\nu_i^*V],
 \qquad
 \sigma_i:=\mu_i^U-\mu_i^V.
\]

Their positivity is \cite[Lemma~8.5.5]{CLD25}.  For a continuous function
\(F\) on \(W\), multiply the two locally finite identities in
\eqref{eq:p21p3-nr-fundamental-cycle} by \(F\), using
\cite[Example~7.2.12]{CLD25}.  Evaluation on a compactly supported test
form and the defining transpose pairing for finite pushforward give

\begin{align}
 [F]_X
   &=\sum_i m_i(\nu_{i,W}^{X})_*
       \bigl(((\nu_{i,W}^{X})^*F)\delta_{\widetilde W_i}\bigr),
 \label{eq:p21p3-nr-function-current-X}\\
 [F]_{X_{\mathrm{red}}}
   &=\sum_i(\nu_{i,W}^{\mathrm{red}})_*
       \bigl(((\nu_{i,W}^{\mathrm{red}})^*F)
                    \delta_{\widetilde W_i}\bigr).
 \label{eq:p21p3-nr-function-current-red}
\end{align}

For \(F=U,V\), differentiate these identities and use the continuous-psh
projection formula \cite[Proposition~8.5.6]{CLD25}; equivalently, proper
pushforward commutes with the current differentials by transposition.
Thus equality of the curvature currents computed with \(\delta_X\) is
equivalent to

\begin{equation}
\label{eq:p21p3-nr-weighted-measure-sum}
                       \sum_i m_i(\nu_{i,W}^{X})_*\sigma_i=0,
\end{equation}

whereas equality computed with \(\delta_{X_{\mathrm{red}}}\) is
equivalent to

\begin{equation}
\label{eq:p21p3-nr-unweighted-measure-sum}
                       \sum_i
                         (\nu_{i,W}^{\mathrm{red}})_*\sigma_i=0.
\end{equation}

Write \(Z_W=Z\cap W\) and
\(Z_{i,W}=Z_i\cap\widetilde W_i=\bar\nu_i^{-1}(Z_W)\).  Away from \(Z_W\),
the component images are disjoint and the relevant normalization map is
an isomorphism.  Restricting either
\eqref{eq:p21p3-nr-weighted-measure-sum} or
\eqref{eq:p21p3-nr-unweighted-measure-sum} to such an open set, and
using \(m_i>0\) in the weighted case, yields

\begin{equation}
\label{eq:p21p3-nr-component-support}
                       \operatorname{Supp}(\sigma_i)
                       \subseteq Z_{i,W}
                       \quad\text{for every }i.
\end{equation}

Every point of \(Z_{i,W}\) is rigid.  Around such a point, choose an open
neighbourhood contained in \(\widetilde W_i\).  It is a good separated
strict rig-smooth boundaryless curve, so
\cref{lem:p21p3-gfc-rigid-zero-atom} applies to the restrictions of
\(\bar\nu_i^*U\) and \(\bar\nu_i^*V\).  Neither positive measure
\(\mu_i^U\) nor \(\mu_i^V\) has an atom there.  A signed Radon measure
supported on a locally finite set is the locally finite sum of its point
masses.  Hence \eqref{eq:p21p3-nr-component-support} implies

\begin{equation}
\label{eq:p21p3-nr-component-measure-equality}
                         \sigma_i=0\qquad\text{for every }i.
\end{equation}

This argument works starting from either the weighted or the unweighted
identity.  Conversely,
\eqref{eq:p21p3-nr-component-measure-equality} implies both identities.
This proves the exact curvature equivalence
\eqref{eq:p21p3-nr-curvature-equivalence}; in particular, generic
multiplicities cannot conceal cancellation at a component intersection.

Finally put \(h=U-V\) and write \(T=\sum_i c_i\delta_i\).  Multiplication
by \(h\) is defined for the positive current \(T\) by
\cite[Example~7.2.12]{CLD25}; continuity makes \(h\) locally bounded.
Proper pushforward, linearity of the current differentials, and
\eqref{eq:p21p3-nr-component-measure-equality} give

\begin{align*}
 \mathrm d'\mathrm d''(hT|_W)
 &=\sum_i c_i(\nu_{i,W}^{X})_*
     \mathrm d'\mathrm d''
       [\bar\nu_i^*U-\bar\nu_i^*V]\\
 &=\sum_i c_i(\nu_{i,W}^{X})_*(\mu_i^U-\mu_i^V)=0.
\end{align*}

This is \eqref{eq:p21p3-nr-bt-conclusion}.  It proves precisely the
Bedford--Taylor independence identity for the continuous psh pairs stated
in the proposition.  The proof does not invoke, and the proposition does
not assert, the additional discontinuous-psh quantifier in
\cite[Definition~8.6.2]{CLD25}.  This completes the proof.
\end{proof}


To pass from a single normal curve to several normalized components, one
must rule out cancellation at points above their conductor and intersection
loci.  The next local statement does exactly this.  It assumes normality
only at the rigid point under consideration and works inside the analytic
interior.  The statement and proof below are the exact bytes accepted by
production verification.

% Quarantine-only candidate statement.
% Date: 29 August 2026.
% This file is not a live-book input and makes no novelty or priority claim.

\begin{lemma}
\label{lem:p21p3-componentwise-normal-rigid-zero-atom}
Let $K$ be a complete nontrivially valued non-Archimedean field, let $X$ be
a good separated strict $K$-analytic space, and let $y\in X$ be a rigid
point.  Assume that
\[
                         y\in\operatorname{Int}(X/K),
\]
and that the ordinary analytic local ring $\mathcal O_{X,y}$ is normal.
Let $\Omega\subseteq X$ be an open neighborhood of $y$ on which $X$ is
purely one-dimensional, and let $u$ be a continuous
CLD-plurisubharmonic function on $\Omega$.  Put
\[
 \Omega^\circ:=\Omega\cap\operatorname{Int}(X/K),
 \qquad
 \mu_u:=\mathrm d'\mathrm d''[u|_{\Omega^\circ}].
\]
Then $\mu_u$ is a positive Radon measure on $\Omega^\circ$ and has no atom
at $y$:
\begin{equation}
\label{eq:p213-normal-rigid-zero-atom}
                         \mu_u(\{y\})=0.
\end{equation}
\end{lemma}

\begin{remark}
\label{rem:p213-normal-rigid-zero-atom-scope}
Normality is required only for the local ring at $y$; neither $X$ nor
$\Omega$ is assumed normal away from $y$.  The point is required only to
lie in the analytic interior; the ambient space need not be boundaryless
away from a neighborhood of $y$.  The normality is ordinary normality, not
geometric normality.  Thus the statement does not restore rig-smoothness
over an imperfect ground field.  No assertion is made for a boundary point,
a nonrigid point, a nonnormal local ring, a trivially valued ground field,
or a discontinuous psh function.
\end{remark}

% Quarantine-only source-aware proof.
% Date: 29 August 2026.
% This file is not a live-book input and makes no novelty or priority claim.

\begin{proof}[Proof of
\cref{lem:p21p3-componentwise-normal-rigid-zero-atom}]
We first reduce normality to the precise local coordinate needed in the
estimate.  Since $X$ is good, $\mathcal O_{X,y}$ is a noetherian henselian
excellent local ring by \cite[\S11.1.10]{CLD25}.  A rigid point is Zariski
closed in a sufficiently small strict affinoid neighborhood.  Indeed, the
fraction field of the quotient by the evaluation kernel embeds densely in
the finite-dimensional $K$-vector space $\mathscr H(y)$.  It is therefore
finite-dimensional over $K$, as is its $K$-subalgebra given by the affinoid
quotient.  That quotient is a finite-dimensional domain over $K$, hence a
field, so the evaluation kernel is maximal.  The central dimension of the
germ is consequently zero, and the local dimension formula
\cite[(11.5.5.1)]{CLD25}, together with the assumed pure local dimension,
gives
\[
                  \dim_{\mathrm{Krull}}\mathcal O_{X,y}=1.
\]
The normal local ring $\mathcal O_{X,y}$ is a domain.  The equivalence
between a one-dimensional noetherian normal local domain and a discrete
valuation ring is
\cite[\href{https://stacks.math.columbia.edu/tag/00PD}{Tag~00PD}]
{stacks-project}.  We may therefore choose a uniformizer
$\pi\in\mathcal O_{X,y}$.

Put $X^\circ:=\operatorname{Int}(X/K)$ and
$\Omega^\circ:=\Omega\cap X^\circ$.  These sets are open, and the
boundary-composition identity \cite[(11.4.1.3)]{CLD25} shows that every open
analytic subspace of $X^\circ$ is boundaryless.  After shrinking around
$y$, we may consequently work in a boundaryless good separated strict open
set
\[
                         U\subseteq\Omega^\circ
\]
on which $X$ is purely one-dimensional and normal and on which $\pi$ is
represented by an analytic function $f$ satisfying $f(y)=0$.  Such a normal
shrinking exists because the ordinary normal locus is open by
\cite[\S11.3.1]{CLD25}.  The germ of the fibre of
$f:U\to\mathbf A_K^{1,\mathrm{an}}$ over $0$ is zero-dimensional at $y$,
because
\[
                  \mathcal O_{X,y}/(f)
                  =\mathcal O_{X,y}/(\pi)
\]
is a field.  Formula (11.4.1.3) also gives
$y\in\operatorname{Int}(U/\mathbf A_K^{1,\mathrm{an}})$.  Thus
\cite[Lemma~11.6.5(a)]{CLD25} makes $f$ finite at $y$.

We spell out the flatness point, since ordinary module-flatness must not be
silently substituted for analytic flatness.  Let
\[
       A:=\mathcal O_{\mathbf A_K^{1,\mathrm{an}},0},
       \qquad B:=\mathcal O_{X,y}.
\]
The ring $A$ is a discrete valuation ring with uniformizer $T$, and the
local homomorphism $A\to B$ sends $T$ to $\pi$.  It is injective: if
$0\ne a\in A$, write $a=T^m v$ with $v$ a unit; then
$a(f)=\pi^m v(f)\ne0$ in the domain $B$.  Since $f$ is finite at $y$, the
$A$-module $B$ is finite by \cite[\S11.4.4]{CLD25}.  It is torsion-free,
hence finite free over the discrete valuation ring $A$.  Thus $B$ is flat
over $A$.  Since $y\in\operatorname{Int}(U/\mathbf A_K^{1,\mathrm{an}})$,
the local-ring criterion of \cite[Remark~11.4.7]{CLD25} upgrades this
module-flatness to analytic flatness of $f$ at $y$.  The local finite-flat
construction of \cite[\S11.6.3]{CLD25} now gives compact analytic
neighborhoods on which $f$ is finite and flat and on which $y$ is the only
point over $0$.  Restricting the target to a sufficiently small closed disc
of admissible radius $R_0>0$, and the source to its inverse image, therefore
produces a compact analytic neighborhood $Q_{R_0}\subset U$ of $y$ and a
finite flat map of constant degree $d>0$
\begin{equation}
\label{eq:p213-normal-rigid-finite-tube}
          f:Q_{R_0}\longrightarrow D(0,R_0),
          \qquad f^{-1}(0)=\{y\}.
\end{equation}
A finite morphism has empty relative boundary.  By
\cite[(11.4.1.3) and Example~11.4.2]{CLD25}, the boundary of $Q_{R_0}$ is
therefore contained in the inverse image of the outer Gauss point
$\{|T|=R_0\}$.

Choose an admissible radius $0<R<R_0$ and a number
\[
             0<\varepsilon_0<\tfrac12\log(R_0/R).
\]
For every admissible $0<r<R$ such that
$\log(R/r)>4\varepsilon_0$, set
\[
                    s_R=-\log R,
                    \qquad s_r=-\log r.
\]
Choose a function $\rho_r\in C^\infty(\mathbf R,[0,1])$ which vanishes on
\[
 (-\infty,s_R-\varepsilon_0]\cup
 [s_r+\varepsilon_0,\infty),
\]
equals $1$ on $[s_R,s_r]$, is nondecreasing on
$[s_R-\varepsilon_0,s_R]$, and is nonincreasing on
$[s_r,s_r+\varepsilon_0]$.  Define
\[
 q_r(t):=\int_{-\infty}^{t}\rho_r(\tau)\,\mathrm d\tau,
 \qquad
 H_r:=\int_{\mathbf R}\rho_r(\tau)\,\mathrm d\tau.
\]
Then
\begin{equation}
\label{eq:p213-normal-rigid-radial-data}
 H_r\geq\log(R/r),
 \qquad
 \int_{\mathbf R}|q_r''(t)|\,\mathrm dt=2.
\end{equation}

On $Q_{R_0}\setminus\{y\}$ put
\[
                         \chi_r=q_r(-\log|f|).
\]
The function $q_r$ is zero on an outer collar of $Q_{R_0}$ and is constant
equal to $H_r$ for sufficiently large arguments.  We may therefore extend
$\chi_r$ by $H_r$ across $y$ and by zero to $U\setminus Q_{R_0}$.  The
extensions are made where the function is constant; on the remaining locus
$f$ is invertible.  The moment description of smooth functions in
\cite[Corollary~15.1.2]{CLD25} therefore shows that $\chi_r$ is a
nonnegative compactly supported CLD-smooth function on $U$.  Moreover,
\begin{equation}
\label{eq:p213-normal-rigid-value-at-y}
                         \chi_r(y)=H_r.
\end{equation}

Let $s=-\log|f|$.  On the two transition tubes one has
\begin{equation}
\label{eq:p213-normal-rigid-radial-dd}
       \mathrm d'\mathrm d''\chi_r
       =q_r''(s)\,\mathrm d's\wedge\mathrm d''s.
\end{equation}
The two tubes are disjoint, and $q_r''$ has a fixed sign on each one.
Apply the calibrated pushforward formula
\cite[Proposition~14.4.2]{CLD25} to the restriction of
\eqref{eq:p213-normal-rigid-finite-tube} over each corresponding target
annulus.  The identification of calibrated degree with finite-flat degree
in \cite[Lemma~13.3.8]{CLD25} shows that the pushforward factor is $d$;
this includes all local and residue-degree multiplicities.  Consequently
\begin{align}
 \left\|\mathrm d'\mathrm d''\chi_r\wedge\delta_U
 \right\|_{\mathrm{TV}}
 &=d\int_{\mathbf R}|q_r''(t)|\,\mathrm dt \notag\\
 &=2d.
\label{eq:p213-normal-rigid-radial-tv}
\end{align}
Neither \cite[Proposition~14.4.2]{CLD25} nor
\cite[Lemma~13.3.8]{CLD25} requires the source to be rig-smooth.

Because $\Omega^\circ$ is boundaryless, its integration current is closed
by \cite[\S7.7.4]{CLD25}.  Applying \cite[Lemma~8.5.5]{CLD25} with
$T=\delta_{\Omega^\circ}$ shows that the current
\[
                         \mu_u
       =\mathrm d'\mathrm d''[u|_{\Omega^\circ}]
\]
is positive on $\Omega^\circ$; as a positive current of bidimension
$(0,0)$ it is a Radon measure by \cite[Proposition~7.2.8]{CLD25}.  Its
restriction to the boundaryless open set $U$ is
\[
                 \mu_u|_U=\mathrm d'\mathrm d''[u|_U].
\]
Equivalently, one may apply the same two results with $T=\delta_U$
directly.  Positivity,
\eqref{eq:p213-normal-rigid-radial-data},
\eqref{eq:p213-normal-rigid-value-at-y}, the transposition identity
\cite[Lemma~7.7.5]{CLD25}, and
\eqref{eq:p213-normal-rigid-radial-tv} yield
\begin{align}
 \log(R/r)\,\mu_u(\{y\})
 &\leq H_r\mu_u(\{y\})                                      \notag\\
 &\leq\int_U\chi_r\,\mathrm d\mu_u                          \notag\\
 &=\bigl\langle[u|_U],\mathrm d'\mathrm d''\chi_r\bigr\rangle \notag\\
 &\leq 2d\sup_{Q_{R_0}}|u|.
\label{eq:p213-normal-rigid-final-bound}
\end{align}
The last line follows by first taking the absolute value of the pairing and
then using its total-variation bound.  Its right-hand side is independent of
$r$.  Since the valuation of $K$ is nontrivial, admissible positive radii
tend to zero.  Letting $r\to0$ in
\eqref{eq:p213-normal-rigid-final-bound} proves
\eqref{eq:p213-normal-rigid-zero-atom}.
\end{proof}


Combining this local estimate with the normal-curve classification,
nilreduction, and the fundamental-cycle multiplicities gives the
componentwise result for reducible and nonreduced curves.  In particular,
the curvature equality defined using the multiplicity-weighted fundamental
current is equivalent to the corresponding equality on the reduction.  The
following statement and proof are again the exact production-verified bytes.

% Quarantine statement for the componentwise rigid-exception endpoint.
% Date: 29 August 2026.
% This file is not a live-book input and has not been submitted to production
% Rethlas.  Its proof uses the separately frozen exact local zero-atom lemma.

\begin{proposition}
\label{prop:p21p3-componentwise-rigid-exceptions}
Let $K$ be a complete nontrivially valued non-Archimedean field, and let $X$
be a good separated strict boundaryless $K$-analytic space, pure of dimension
one.
The space $X$ is not assumed reduced.  Write
\[
                 \iota:X_{\mathrm{red}}\hookrightarrow X
\]
for the canonical nilpotent closed immersion.  Let $(X_i)_i$ be the
irreducible components of $X_{\mathrm{red}}$, and write
\[
 \bar\nu_i:\widetilde X_i\longrightarrow X_i
       \hookrightarrow X_{\mathrm{red}},
 \qquad
 \nu_i:=\iota\circ\bar\nu_i,
 \qquad
 \delta_i:=(\nu_i)_*\delta_{\widetilde X_i}.
\]
For every $i$, assume that there is a locally finite set
$E_i\subset\widetilde X_i$ of rigid points such that
$\widetilde X_i\setminus E_i$ is rig-smooth.  Let
$m_i\in\mathbf Z_{>0}$ be the coefficient of $[X_i]$ in the fundamental
cycle $[X]$.  Then the following assertions hold.

\begin{enumerate}
\item\label{itm:p21p3-componentwise:calculus}
For every $(p,q)$, pullback by $\iota$ identifies the sheaves of ordinary
CLD-smooth forms on $X$ and $X_{\mathrm{red}}$, compatibly with wedge
product, $J$, $\mathrm d'$, $\mathrm d''$, support, and the three positivity
cones.  On compactly supported test forms this is an isomorphism of
topological vector spaces.  Its transpose identifies the CLD current
sheaves, their weak topologies, positivity, and
$\mathrm d'$- and $\mathrm d''$-closedness.  It identifies multiplication
by smooth functions for all currents.  For positive currents it also
identifies multiplication by locally bounded Borel functions, and by
linearity the same is true for currents locally expressible as differences
of positive currents.  No Borel product is asserted for an arbitrary
order-zero current.  Under the common underlying topological space, the
CLD-psh and continuous CLD-psh sheaves coincide.  The same pullback
identification holds for $G$-forms.

\item\label{itm:p21p3-componentwise:integration}
The integration currents retain the generic multiplicities:
\begin{equation}
 \label{eq:p21p3-componentwise-fundamental-cycles}
       \delta_X=\sum_i m_i\delta_i,
       \qquad
       \iota_*\delta_{X_{\mathrm{red}}}=\sum_i\delta_i.
\end{equation}
Both sums are locally finite.  Consequently
$\delta_X=\iota_*\delta_{X_{\mathrm{red}}}$ if and only if $X$ is
generically reduced along every $X_i$.

\item\label{itm:p21p3-componentwise:classification}
Every positive $\mathrm d'$- and $\mathrm d''$-closed CLD current $T$ of
bidegree $(0,0)$ on $X$ has a unique expression
\begin{equation}
 \label{eq:p21p3-componentwise-classification}
                         T=\sum_i c_i\delta_i,
                         \qquad c_i\in\mathbf R_{\geq0}.
\end{equation}
The sum is locally finite.

\item\label{itm:p21p3-componentwise:bedford-taylor}
Let $W\subseteq X$ be open, and let $U,V$ be continuous CLD-psh functions
on $W$.  Identify $W$ with its inverse image in $X_{\mathrm{red}}$, and,
for a continuous function $F$ on $W$, put
\[
 [F]_X:=F(\delta_X|_W),
 \qquad
 [F]_{X_{\mathrm{red}}}:=
      F(\delta_{X_{\mathrm{red}}}|_W).
\]
Then
\begin{equation}
 \label{eq:p21p3-componentwise-curvature-equivalence}
 \mathrm d'\mathrm d''[U]_X=\mathrm d'\mathrm d''[V]_X
 \quad\Longleftrightarrow\quad
 \mathrm d'\mathrm d''[U]_{X_{\mathrm{red}}}
  =\mathrm d'\mathrm d''[V]_{X_{\mathrm{red}}}.
\end{equation}
If these equivalent conditions hold, then every current in
\cref{itm:p21p3-componentwise:classification} satisfies
\begin{equation}
 \label{eq:p21p3-componentwise-bt-conclusion}
             \mathrm d'\mathrm d''\bigl((U-V)T|_W\bigr)=0.
\end{equation}
Thus every such current has Bedford--Taylor independence for continuous
psh pairs in the CLD sense.
\end{enumerate}
\end{proposition}

\begin{remark}
\label{rem:p21p3-componentwise-residual}
The local input
\cref{lem:p21p3-componentwise-normal-rigid-zero-atom} is used only to prevent
component curvature measures from cancelling at the rigid conductor and
intersection locus; the classification does not consume it.  Consequently,
reducibility, nilpotent thickenings, generic multiplicities, and locally
finite rigid failures of rig-smoothness on the component normalizations are
no longer residual one-dimensional cases of P21.3.  Within the ambient
hypotheses of that problem, the remaining domain is analytic dimension at
least two, a trivially valued ground field, and one-dimensional component
normalizations whose non-rig-smooth locus is not contained in a locally
finite set of rigid points.  The proposition does not remove goodness,
separatedness, strictness, or boundarylessness, and it does not assert the
possibly discontinuous psh predicate of
\cite[Definition~8.6.2]{CLD25}.
\end{remark}

% Quarantine source-aware proof for the componentwise rigid-exception
% endpoint.  Date: 29 August 2026.
% This file uses the separately frozen exact local zero-atom lemma.

\begin{proof}[Proof of
\cref{prop:p21p3-componentwise-rigid-exceptions}]
We separate the reductions because only the last, measure-theoretic one uses
\cref{lem:p21p3-componentwise-normal-rigid-zero-atom}.

\paragraph{Nilreduction and the ambient geometry.}
An analytic seminorm kills every nilpotent.  Hence $X$ and
$X_{\mathrm{red}}$ have the same underlying topological space, the same
analytic points, and the same completed residue fields.  Affinoid charts of
the reduction are quotients of affinoid charts of $X$ by their nilradicals.
It follows directly that goodness, separatedness, strictness, and pure
dimension one pass to $X_{\mathrm{red}}$.  The immersion $\iota$ is finite
and has empty relative boundary by \cite[\S11.4.4]{CLD25}.  The
boundary-composition identity \cite[(11.4.1.3)]{CLD25} therefore shows that
$X_{\mathrm{red}}$ is boundaryless.

Both $X$ and $X_{\mathrm{red}}$ are good and topologically separated, hence
locally holomorphically separated by \cite[\S15.2.2]{CLD25} and locally
differentially separated by \cite[Corollary~15.2.4]{CLD25}.  Thus the Radon
measure and Borel-multiplication statements below apply by
\cite[Lemma~7.2.4, Proposition~7.2.8, and Example~7.2.12]{CLD25}.

We next justify the full calculus assertion, including its topology.  If
$A\to A/I$ is an affinoid quotient with $I$ nilpotent, every unit of $A/I$
lifts to a unit of $A$: if $ab=1+n$ with $n$ nilpotent, the finite geometric
series for $(1+n)^{-1}$ corrects $b$ to an inverse.  Every analytic seminorm
on $A$ kills $I$, so a lifted unit and its reduction have exactly the same
absolute value at every analytic point.  Consequently a moment map on a
reduced chart lifts, after shrinking, to a moment map on $X$ with the same
tropicalization.  Conversely every moment on $X$ restricts to one on the
reduction.  This is the nilpotent-immersion case of
\cite[Proposition~15.1.3(b) and the following form variant]{CLD25}.

The construction of ordinary CLD forms from moment charts now gives
\begin{equation}
 \label{eq:p21p3-componentwise-form-isomorphism}
 \iota^*:A_X^{p,q}\xrightarrow{\ \sim\ }
                    A_{X_{\mathrm{red}}}^{p,q}.
\end{equation}
The equality of tropicalizations makes
\eqref{eq:p21p3-componentwise-form-isomorphism} compatible with wedge
product, $J$, $\mathrm d'$, $\mathrm d''$, support, and each positivity
cone.  For the topological assertion, fix a compact test support.  A
compact-test datum is a finite collection of moment charts and Lagerberg
representatives.  Restriction and the lifted charts just constructed are
cofinal collections of such data on the two spaces.  On corresponding
charts the tropical images, polyhedral subdivisions, representatives, and
every $C^r$-seminorm are identical.  Thus
\begin{equation}
 \label{eq:p21p3-componentwise-test-homeomorphism}
 \iota^*:A^{p,q}_{X,c}\xrightarrow{\ \sim\ }
                    A^{p,q}_{X_{\mathrm{red}},c}
\end{equation}
is a homeomorphism of locally convex spaces.

Taking continuous duals in
\eqref{eq:p21p3-componentwise-test-homeomorphism} gives
\begin{equation}
 \label{eq:p21p3-componentwise-current-isomorphism}
 \iota_*:D^{a,b}(X_{\mathrm{red}})
       \xrightarrow{\ \sim\ }D^{a,b}(X),
 \qquad
 \langle\iota_*S,\alpha\rangle
       =\langle S,\iota^*\alpha\rangle.
\end{equation}
The test-form positivity cones agree, so this isomorphism preserves
positivity.  It preserves $\mathrm d'$- and $\mathrm d''$-closedness because
$\iota^*$ commutes with the two differentials.  It also identifies the weak
topologies and multiplication by smooth functions.

If $S$ is positive, its product with a locally bounded Borel function is
defined through the Radon measures attached to positive test forms, as in
\cite[Example~7.2.12]{CLD25}.  Those test forms and attached measures agree
under \eqref{eq:p21p3-componentwise-form-isomorphism}, so
$\iota_*(fS)=f\iota_*S$.  If locally $S=S_1-S_2$ with $S_1,S_2$ positive,
the same formula defines $fS$ by linearity.  It is independent of the
decomposition: an equality
$S_1+S'_2=S'_1+S_2$ of positive currents is an equality of all attached
measures, which remains true after multiplication by $f$.  This argument
does not define a Borel product for an arbitrary order-zero current.

Finally, the first stage in the CLD psh filtration is the same on $X$ and
$X_{\mathrm{red}}$, because smoothness and positivity of
$\mathrm d'\mathrm d''u$ agree.  Successor stages take decreasing limits,
the continuous filtration takes local uniform limits, and limit stages take
local unions.  The open sets and their real-valued functions are the same.
Transfinite induction therefore identifies the psh and continuous psh
sheaves.  The $G$-form assertion follows identically from the $G$-form
variant after \cite[Proposition~15.1.3]{CLD25}.  This proves
\cref{itm:p21p3-componentwise:calculus}.

\paragraph{Normalization and the locally finite rigid locus.}
Let
\[
 \bar\nu:\coprod_i\widetilde X_i\longrightarrow X_{\mathrm{red}}
\]
be the componentwise normalization.  On a good affinoid domain define
\[
 \mathcal Q:=\bar\nu_*\mathcal O_{\coprod_i\widetilde X_i}
                  /\mathcal O_{X_{\mathrm{red}}},
 \qquad
 \mathfrak c:=\operatorname{Ann}_{\mathcal O_{X_{\mathrm{red}}}}
                  (\mathcal Q).
\]
Let $Z_0$ be the conductor locus
\[
 \operatorname{Supp}(\mathcal Q)
   =V(\mathfrak c)
   =\operatorname{Supp}(\mathcal O_{X_{\mathrm{red}}}/\mathfrak c),
\]
and let $Z$ be the union of $Z_0$ with all pairwise intersections of
distinct irreducible components.  Analytic local rings on a good analytic
space are excellent by \cite[\S11.1.10]{CLD25}; normalization is finite and
its components correspond to the irreducible components by
\cite[\S11.3.1]{CLD25}.  Thus $\mathcal Q$ and the conductor are coherent.
Their support misses the generic point of every reduced component and is
zero-dimensional.  The same dimension argument applies to a proper
intersection of two components.  Irreducible components are locally finite
by \cite[\S11.2.5]{CLD25}; on a relatively compact affinoid neighborhood
only finitely many occur.  Hence $Z$ is closed and locally finite.

Strictness turns this zero-dimensional set into a rigid set.  Locally a
zero-dimensional closed analytic subspace of a strict affinoid is the
spectrum of a zero-dimensional strict affinoid algebra $B$.  The algebra is
noetherian by \cite[Theorem~4.5]{Tate71Rigid}, hence Artinian.  The strict
affinoid Nullstellensatz in the same theorem makes every residue field
$B/\mathfrak m$ finite over $K$; it is already complete.  Therefore every
point of $Z$ is rigid, not necessarily $K$-rational.  Put
\begin{equation}
 \label{eq:p21p3-componentwise-conductor-preimage}
                         Z_i:=\bar\nu_i^{-1}(Z).
\end{equation}
A finite morphism induces a finite extension of completed residue fields,
so every point of $Z_i$ is rigid.  Properness of a finite map and local
finiteness of $Z$ show that $Z_i$ is locally finite.

A finite algebra over a strict affinoid algebra is strict by
\cite[Theorem~4.4]{Tate71Rigid}; hence $\widetilde X_i$ is strict.  The maps
$X_i\to X_{\mathrm{red}}$ and
$\widetilde X_i\to X_i$ are finite and have empty relative boundary.
Applying \cite[(11.4.1.3)]{CLD25} twice shows that
$\widetilde X_i$ is boundaryless.  It is good and separated, normal by
construction, and connected because it normalizes one irreducible
component.

We shall apply the already verified normal rigid-exception proposition on
\begin{equation}
 \label{eq:p21p3-componentwise-normal-open}
                         U_i:=\widetilde X_i\setminus Z_i.
\end{equation}
For completeness, we check all of its topological hypotheses.  At a rigid
point of a normal analytic curve the normalization fibre has cardinality
one, so the point is unibranch by
\cite[Proposition~3.5.2]{Ducros24CurveStructure}.  Every connected analytic
curve is a graph and is connected by injective arcs by
\cite[Theorem~3.5.1 and \S1.3.5]{Ducros24CurveStructure}, whereas an
interior point of an injective arc is pluribranch by
\cite[\S1.3.12]{Ducros24CurveStructure}.  It follows that an injective arc
whose endpoints avoid $Z_i$ cannot meet $Z_i$; hence $U_i$ is connected.
It is nonempty as well: if one starts at a point of $Z_i$, its unique branch
contains an injective subarc with a pluribranch interior point, and that
point is not rigid and hence not in $Z_i$.

The space $U_i$ is therefore nonempty, connected, good, separated, strict,
normal, and boundaryless.  The set
\[
                         E_i':=E_i\setminus Z_i\subset U_i
\]
is locally finite and rigid, while
\[
 U_i\setminus E_i'
   =\widetilde X_i\setminus(E_i\cup Z_i)
\]
is an open subspace of the rig-smooth curve
$\widetilde X_i\setminus E_i$ and is therefore rig-smooth.  Thus every
hypothesis of \cref{prop:p21p3-rigid-exception-extension} is verified for
the pair $(U_i,E_i')$.

\paragraph{Generic multiplicities.}
By \cite[Definition~13.2.5 and \S13.5.2]{CLD25},
\[
                   [X]=\sum_i m_i[X_i],
 \qquad
                   [X_{\mathrm{red}}]=\sum_i[X_i].
\]
Integration agrees with the current attached to the fundamental cycle by
\cite[Proposition~16.6.3]{CLD25}.  The normalization is finite and has
calibrated degree one over the normal locus.  Its exceptional image lies in
$Z$, and every smooth $(1,1)$-test-form germ vanishes at a rigid point by
\cite[Proposition~3.3.6 and Proposition~15.2.1(a)]{CLD25}.  Consequently
the cycle identities hold as current identities on all of $X$:
\[
       \delta_X=\sum_i m_i(\nu_i)_*\delta_{\widetilde X_i},
 \qquad
       \iota_*\delta_{X_{\mathrm{red}}}
          =\sum_i(\nu_i)_*\delta_{\widetilde X_i}.
\]
This is \eqref{eq:p21p3-componentwise-fundamental-cycles}.  Local finiteness
of the components makes both sums locally finite.  If all $m_i$ equal one,
the two currents are equal.  Conversely, restrict an asserted equality to a
nonempty rig-smooth open subset of $X_i$ meeting neither $Z$ nor another
component.  The nonzero integration current there gives $m_i=1$.  Since
$m_i$ is the length of the generic local ring along $X_i$, this is
equivalent to generic reducedness along every component.  This proves
\cref{itm:p21p3-componentwise:integration}.

\paragraph{Current classification and uniqueness.}
Let $T$ be a positive $\mathrm d'$- and $\mathrm d''$-closed current of
bidegree $(0,0)$ on $X$.  By
\eqref{eq:p21p3-componentwise-current-isomorphism}, there is a unique
positive closed current $T_{\mathrm{red}}$ on $X_{\mathrm{red}}$ with
$T=\iota_*T_{\mathrm{red}}$.  The open set
$X_{\mathrm{red}}\setminus Z$ is the disjoint union of the normal component
loci $X_i\setminus Z$, and
\[
 \bar\nu_i|_{U_i}:U_i\xrightarrow{\ \sim\ }X_i\setminus Z.
\]
Transport the restriction of $T_{\mathrm{red}}$ on $X_i\setminus Z$ to a
current $T_i$ on $U_i$.  Restriction to an open subspace preserves
positivity and closedness by \cite[\S7.1.6 and \S\S7.7.1--7.7.3]{CLD25}.
The verified classification
\cref{itm:p21p3-rigid-exception:classification} applied to $(U_i,E_i')$
therefore gives a unique $c_i\in\mathbf R_{\geq0}$ such that
\begin{equation}
 \label{eq:p21p3-componentwise-classification-off-Z}
                            T_i=c_i\delta_{U_i}.
\end{equation}

Write
$\bar\delta_i:=(\bar\nu_i)_*\delta_{\widetilde X_i}$ on
$X_{\mathrm{red}}$, and set
\[
                         S:=T_{\mathrm{red}}-\sum_i c_i\bar\delta_i.
\]
The sum is locally finite, and
\eqref{eq:p21p3-componentwise-classification-off-Z} says that $S$ vanishes
on $X_{\mathrm{red}}\setminus Z$.  Let $\alpha$ be a compactly supported
smooth $(1,1)$-test form.  A rigid point $z$ has $d_K(z)=0$, since
$\mathscr H(z)/K$ is finite and both terms in the dimension formula
\cite[(11.2.3.1)]{CLD25} vanish.  The analytic support theorem
\cite[Proposition~15.2.1(a)]{CLD25} gives
\[
 \operatorname{Supp}(\alpha)
    \subseteq\{x:d_K(x)\geq1\}
    \subseteq X_{\mathrm{red}}\setminus Z.
\]
Thus $\langle S,\alpha\rangle=0$, and $S=0$.  Pushing forward by $\iota$
proves \eqref{eq:p21p3-componentwise-classification}.

For uniqueness, fix $i$ and restrict an equality of two such sums to a
nonempty rig-smooth open subset of $X_i\setminus Z$ that meets no other
component.  Such a subset exists because
$U_i\setminus E_i'$ is nonempty by
\cref{itm:p21p3-rigid-exception:connected}.  The integration current on it
is nonzero, so the two coefficients of $\delta_i$ agree.  This proves both
uniqueness and \cref{itm:p21p3-componentwise:classification}.  Notice that
the zero-atom input has not been used.

\paragraph{Componentwise splitting of the curvature measures.}
Let $W,U,V$ be as in
\cref{itm:p21p3-componentwise:bedford-taylor}, and regard $W$ also as an
open subset of $X_{\mathrm{red}}$.  Put
\[
             \widetilde W_i:=\bar\nu_i^{-1}(W),
\]
and denote the restrictions of $\nu_i$ and $\bar\nu_i$ by
$\nu_{i,W}^X$ and $\nu_{i,W}^{\mathrm{red}}$, respectively.  The good
separated spaces $\widetilde X_i$ are locally differentially
separated by \cite[\S15.2.2 and Corollary~15.2.4]{CLD25}.  Pullback preserves
continuous psh functions by
\cite[Proposition~8.1.5]{CLD25}.  Hence
\begin{equation}
 \label{eq:p21p3-componentwise-measures}
 \mu_i^U:=\mathrm d'\mathrm d''[\bar\nu_i^*U],
 \qquad
 \mu_i^V:=\mathrm d'\mathrm d''[\bar\nu_i^*V]
\end{equation}
are positive Radon measures by \cite[Lemma~8.5.5]{CLD25}.  Put
$\sigma_i:=\mu_i^U-\mu_i^V$.

For every continuous function $F$ on $W$, multiply the two current
identities in \eqref{eq:p21p3-componentwise-fundamental-cycles} by $F$.
The product is defined by \cite[Example~7.2.12]{CLD25}, and direct testing
against a compactly supported form gives
\begin{align}
 [F]_X
 &=\sum_i m_i(\nu_{i,W}^X)_*
       \bigl(((\nu_{i,W}^X)^*F)\delta_{\widetilde W_i}\bigr),
 \label{eq:p21p3-componentwise-function-current-X}\\
 [F]_{X_{\mathrm{red}}}
 &=\sum_i(\nu_{i,W}^{\mathrm{red}})_*
       \bigl(((\nu_{i,W}^{\mathrm{red}})^*F)
                              \delta_{\widetilde W_i}\bigr).
 \label{eq:p21p3-componentwise-function-current-red}
\end{align}
Finite proper pushforward commutes with the two current differentials by
transposition; for continuous psh functions this is also
\cite[Proposition~8.5.6]{CLD25}.  Differentiating
\eqref{eq:p21p3-componentwise-function-current-X} for $F=U,V$ shows that
the left equality in
\eqref{eq:p21p3-componentwise-curvature-equivalence} is equivalent to
\begin{equation}
 \label{eq:p21p3-componentwise-weighted-sum}
                 \sum_i m_i(\nu_{i,W}^X)_*\sigma_i=0.
\end{equation}
Since $\nu_{i,W}^X=\iota\circ\nu_{i,W}^{\mathrm{red}}$ and $\iota_*$ is
injective by \eqref{eq:p21p3-componentwise-current-isomorphism}, this is
equivalent to
\begin{equation}
 \label{eq:p21p3-componentwise-weighted-sum-red}
       \sum_i m_i(\nu_{i,W}^{\mathrm{red}})_*\sigma_i=0
       \qquad\text{on }X_{\mathrm{red}}.
\end{equation}
Likewise the right equality is equivalent to
\begin{equation}
 \label{eq:p21p3-componentwise-unweighted-sum}
                 \sum_i(\nu_{i,W}^{\mathrm{red}})_*\sigma_i=0.
\end{equation}

Set $Z_W:=Z\cap W$ and
$Z_{i,W}:=Z_i\cap\widetilde W_i$.  Away from $Z_W$, the component images in
$X_{\mathrm{red}}$ are disjoint and every $\nu_{i,W}^{\mathrm{red}}$ is an
analytic isomorphism onto its component image.  Restricting
\eqref{eq:p21p3-componentwise-weighted-sum-red} to the image of one component,
and using $m_i>0$, gives
\begin{equation}
 \label{eq:p21p3-componentwise-support}
                    \operatorname{Supp}(\sigma_i)\subseteq Z_{i,W}.
\end{equation}
Exactly the same conclusion follows from
\eqref{eq:p21p3-componentwise-unweighted-sum}.

Every point of $Z_{i,W}$ is rigid.  The space $\widetilde X_i$ is
boundaryless, pure one-dimensional, and normal, so every such point lies in
its interior and has a normal analytic local ring.  Apply
\cref{lem:p21p3-componentwise-normal-rigid-zero-atom} on the open curve
$\widetilde W_i$ to each of the two continuous psh functions in
\eqref{eq:p21p3-componentwise-measures}.  It gives
\[
                         \mu_i^U(\{z\})
                         =\mu_i^V(\{z\})=0
                         \qquad(z\in Z_{i,W}).
\]
We spell out why this kills the entire signed measure.  Since $Z_{i,W}$ is
locally finite, every $z\in Z_{i,W}$ has a neighborhood meeting that set
only at $z$.  On such a neighborhood,
\eqref{eq:p21p3-componentwise-support} makes $\sigma_i$ a multiple of
$\delta_z$; its coefficient is
$\sigma_i(\{z\})=0$.  The measure also vanishes off $Z_{i,W}$.  Hence
\begin{equation}
 \label{eq:p21p3-componentwise-component-equality}
                            \sigma_i=0
                            \qquad\text{for every }i.
\end{equation}
Starting from either the weighted or the unweighted sum therefore gives
\eqref{eq:p21p3-componentwise-component-equality}; conversely that equality
implies both sums.  This proves the exact equivalence
\eqref{eq:p21p3-componentwise-curvature-equivalence}.  In particular,
neither different generic multiplicities nor component intersections can
hide cancellation.

Finally put $h=U-V$ and use
\eqref{eq:p21p3-componentwise-classification}.  Multiplication by $h$ is
defined because $T$ is positive and $h$ is locally bounded.  On compact
test supports the component sum is finite.  Proper pushforward,
\eqref{eq:p21p3-componentwise-component-equality}, and linearity give
\begin{align*}
 \mathrm d'\mathrm d''(hT|_W)
 &=\sum_i c_i(\nu_{i,W}^X)_*
       \mathrm d'\mathrm d''[\bar\nu_i^*U-\bar\nu_i^*V]\\
 &=\sum_i c_i(\nu_{i,W}^X)_*(\mu_i^U-\mu_i^V)=0.
\end{align*}
This is \eqref{eq:p21p3-componentwise-bt-conclusion}.  The proof uses only
continuous psh pairs and makes no claim about the stronger discontinuous
quantifier in \cite[Definition~8.6.2]{CLD25}.
\end{proof}


% LIVE-READY STAGING ONLY, 28 August 2026.
% Release gate: do not integrate this example until the arbitrary-field
% nonreduced exact-live statement/proof bytes receive exact-live verification.
% The rendered example below deliberately contains no gate-status language.

\begin{example}
\label{ex:p21p3-reducible-node-bt}
Let \(K\) be a complete nontrivially valued non-Archimedean field and let
\[
                 X=\{XY=0\}\subset\mathbf P^2_K.
\]
This is the reduced union of the two projective lines
\(L_1=\{X=0\}\) and \(L_2=\{Y=0\}\), meeting transversely at the
\(K\)-rational node \(o=[0:0:1]\).  Its analytification is good, separated,
strict, pure of dimension one, and boundaryless.  The normalization is
\[
 \nu:L_1^{\mathrm{an}}\amalg L_2^{\mathrm{an}}
          \longrightarrow X^{\mathrm{an}},
\]
write \(\nu_i\) for its restriction to \(L_i^{\mathrm{an}}\).  The two
points \(o_i\in L_i^{\mathrm{an}}\) above \(o\) are rigid.

By \cref{prop:p21p3-curve-rigidity}, every positive closed current of
bidegree \((0,0)\) on \(X^{\mathrm{an}}\) is
\begin{equation}
                         T=c_1\delta_{L_1}+c_2\delta_{L_2},
                         \qquad c_1,c_2\geq0.
 \label{eq:p21p3-node-current}
\end{equation}
The coefficients need not agree: the node does not create a balancing law
between distinct irreducible components.

Let \(W\subseteq X^{\mathrm{an}}\) be open, let \(u,v\) be continuous psh
functions on \(W\), and assume
\(
 \mathrm d'\mathrm d''[u]=\mathrm d'\mathrm d''[v].
\)
Write \(W_i=\nu_i^{-1}(W)\) and
\[
 \mu_i^u=\mathrm d'\mathrm d''[\nu_i^*u],
 \qquad
 \mu_i^v=\mathrm d'\mathrm d''[\nu_i^*v].
\]
If \(o\notin W\), the normalization-sheet images are already disjoint over
\(W\), and the curvature equality is sheetwise.  Suppose that \(o\in W\).
Away from \(o\), the images of the two normalization sheets are disjoint.
The curvature equality therefore shows that the signed measure
\(\mu_i^u-\mu_i^v\) is supported at \(o_i\) for each \(i\).  By
\cref{lem:p21p3-gfc-rigid-zero-atom}, neither positive measure
\(\mu_i^u\) nor \(\mu_i^v\) has an atom at the rigid point \(o_i\).  Hence
\begin{equation}
                              \mu_i^u=\mu_i^v
                              \qquad(i=1,2).
 \label{eq:p21p3-node-sheetwise-equality}
\end{equation}
Using \eqref{eq:p21p3-node-current}, proper pushforward from the two
normalization sheets, and \eqref{eq:p21p3-node-sheetwise-equality}, one
obtains
\[
 \mathrm d'\mathrm d''\bigl((u-v)T|_W\bigr)
   =\sum_{i=1}^2c_i(\nu_i)_*(\mu_i^u-\mu_i^v)=0.
\]
Thus equality of the two curvature measures cannot arise from cancellation
of opposite atoms on the two branches of the node.

There is also a concrete nonreduced variant.  Put
\[
 \mathfrak X^{(2)}
   :=\operatorname{Proj}K[X,Y,Z]/\bigl((XY)^2\bigr).
\]
Its reduction is \(X\), and its generic multiplicity along each \(L_i\) is
two.  By \cref{prop:p21p3-curve-rigidity-nonreduced}, its analytification
has, under the common current identification, the same ordinary CLD forms,
currents, positive cone, and continuous-pair Bedford--Taylor predicate as
\(X^{\mathrm{an}}\), whereas its fundamental integration current is
\begin{equation}
        \delta_{(\mathfrak X^{(2)})^{\mathrm{an}}}
          =2\delta_{L_1}+2\delta_{L_2}.
 \label{eq:p21p3-double-node-fundamental-current}
\end{equation}
Thus nilpotents change the generic cycle multiplicities, not the normalized
component currents that occur in \eqref{eq:p21p3-node-current}.

The example also shows why continuity and the rigid-point zero-atom input
are used in the descent: the argument establishes the continuous-pair
Bedford--Taylor property, not the literal possibly discontinuous scope of
\cite[Definition~8.6.2]{CLD25}.
\end{example}


\subsection{Reading the definition}
\label{sec:ch16-reading-bt}

The definition shows exactly what ``harmonic ambiguity'' means here.  For a
smooth $h$ with $\dprimeddprime h=0$, one always has
$\dprimeddprime(hT)=0$.  For a merely continuous difference $h=u-v$, the
same conclusion is guaranteed when $T$ is Bedford--Taylor.  This should not
be confused with the curve-test definition of pluriharmonicity from
\Cref{chap:cld-poincare-lelong}.

Three separate conditions on a continuous function $h$ occur in this
chapter, and they are logically distinct:
\begin{center}
\small
\begin{tabular}{p{.27\textwidth}|p{.36\textwidth}|p{.28\textwidth}}
\textbf{Condition}&\textbf{Meaning}&\textbf{Implications proved here}\\
\hline
$\dprimeddprime[h]=0$
 &vanishing of one current on $X$
 &affinity on cell interiors
  (\Cref{prop:ch16-cell-rigidity})\\
$h$ harmonic in the CLD sense
 &locally an $\mathbf R$-combination of parlinear $h_i$ with
  $\dprimeddprime[h_i]=0$
 &Stokes has no boundary defect
  (\Cref{thm:bt-harmonic-compatibility})\\
$h$ classically pluriharmonic
 &$h$ and $-h$ pass every analytic curve test
 &stability under pullback and scalar extension
  (\Cref{chap:cld-poincare-lelong})\\
\end{tabular}
\end{center}
None of the three implications is reversed anywhere in this book without a
theorem, and \Cref{sec:ch16-harmonic-ambiguity} exhibits the exact failure
that separates the first from the second.

\section{The mixed product}
\label{sec:ch16-iteration}

\subsection{The recursion and its first properties}
\label{sec:ch16-recursion}

Let $T$ be a positive closed current and let $u_1,\ldots,u_q$ be continuous
psh functions, with the resulting bidegree not exceeding $(n,n)$.  Define
recursively
\[
 \dprimeddprime u_1\wedge\cdots\wedge\dprimeddprime u_q\wedge T
 :=\dprimeddprime\!\left(
 u_1\,\dprimeddprime u_2\wedge\cdots\wedge
 \dprimeddprime u_q\wedge T\right).
\]
\Cref{thm:one-bt-factor} shows inductively that this is positive and
closed.  When $T$ is Bedford--Taylor, it depends only on the current factors,
not the chosen potentials.

\begin{proposition}
\label{prop:ch16-iterate-basic}
Let $T$ be positive closed of bidegree $(r,r)$ and let $u_1,\ldots,u_q$ be
continuous psh with $q+r\leq n$.  Then:
\begin{enumerate}
\item\label{itm:ch16-iter:posclosed}
 $\dprimeddprime u_1\wedge\cdots\wedge\dprimeddprime u_q\wedge T$ is a
 positive closed current of bidegree $(q+r,q+r)$;
\item\label{itm:ch16-iter:exact}
 for $q\geq1$ it is $\dprimeddprime$-exact;
\item\label{itm:ch16-iter:smooth}
 if all $u_i$ are smooth, it coincides with the ordinary wedge product
 $\dprimeddprime u_1\wedge\cdots\wedge\dprimeddprime u_q\wedge T$ of smooth
 forms with $T$;
\item\label{itm:ch16-iter:bt}
 if $T$ is Bedford--Taylor, so is the iterate, and it depends only on the
 currents $\dprimeddprime u_1,\ldots,\dprimeddprime u_q$ and on $T$;
\item\label{itm:ch16-iter:delta}
 for $T=\delta_X$ one has
 $\dprimeddprime u_1\wedge\cdots\wedge\dprimeddprime u_q\wedge\delta_X
  =\dprimeddprime u_1\wedge\cdots\wedge\dprimeddprime u_q$.
\end{enumerate}
\end{proposition}

\begin{proof}
\cref{itm:ch16-iter:posclosed} is $q$ applications of
\Cref{thm:one-bt-factor}, each legitimate because the previous stage is
positive and closed.  \cref{itm:ch16-iter:exact} is immediate from the
outermost $\dprimeddprime$ in the definition.  \cref{itm:ch16-iter:smooth}
follows by $q$ applications of \Cref{cor:ch16-smooth-case}, unwinding the
recursion from the inside out.  \cref{itm:ch16-iter:bt} is the last assertion
of \Cref{prop:bt-current-stability} applied $q$ times, together with the
definition of a Bedford--Taylor current for the independence statement.
\cref{itm:ch16-iter:delta} is \Cref{prop:bt-current-stability} for
$\delta_X$, since $u\delta_X=[u]$.  These are the assertions collected in
\cite[\S8.7.1]{CLD25}.
\end{proof}

\subsection{Symmetry and locality}
\label{sec:ch16-symmetry}

\begin{theorem}
\label{thm:bt-symmetry-locality}
The mixed current is symmetric in $u_1,\ldots,u_q$ and is local in all its
arguments.  If all $u_i$ are smooth, it is the ordinary wedge product of
their Hessian forms with $T$.
\end{theorem}
\begin{proof}
Locality follows from restriction of currents and of multiplication by a
continuous function.  It suffices to interchange two adjacent factors.  Fix
$v$ and consider the sheaf of continuous psh $u$ for which
\[
 \dprimeddprime\bigl(u\,\dprimeddprime(vT)\bigr)
 =\dprimeddprime\bigl(v\,\dprimeddprime(uT)\bigr).
\]
For smooth $u$ this follows from the differential calculus and closedness.
The class is a sheaf stable under local uniform limits and contains the
smooth psh functions; the defining closure property of
$\operatorname{CPsh}$ therefore puts every continuous psh $u$ in it.
Adjacent transpositions generate the symmetric group.  The smooth formula
follows at every recursive step from the Leibniz rule.  This is
\cite[Proposition~8.7.4]{CLD25}.
\end{proof}

\begin{proof}[Second proof of the transposition step]
Fix a continuous psh $v$ and set
\[
 \mathscr P(U)=\Bigl\{u\in C(U):\
  \dprimeddprime\bigl(u\,\dprimeddprime(vT|_U)\bigr)
  =\dprimeddprime\bigl(v\,\dprimeddprime(uT|_U)\bigr)\Bigr\}.
\]
\emph{Sheaf.}  Equality of two currents is local.

\emph{Smooth psh functions.}  Let $u$ be smooth psh.  The current
$S:=\dprimeddprime(vT)$ is positive and closed by
\Cref{thm:one-bt-factor}, hence symmetric, so
\Cref{lem:ch16-leibniz-closed} applies and gives
\[
 \dprimeddprime\bigl(u\,S\bigr)=\dprimeddprime u\wedge S
 =\dprimeddprime u\wedge\dprimeddprime(vT).
\]
On the other hand $\dprimeddprime u$ is a smooth closed positive
$(1,1)$-form, so \Cref{prop:ch16-smooth-factor} gives
\[
 \dprimeddprime\bigl(v\,\dprimeddprime(uT)\bigr)
 =\dprimeddprime\bigl(v\,(T\wedge\dprimeddprime u)\bigr)
 =\dprimeddprime(vT)\wedge\dprimeddprime u .
\]
The two right-hand sides agree because the wedge of a current with two
smooth forms is associative and commutative in even bidegrees.  Hence
$u\in\mathscr P(U)$.

\emph{Uniform limits.}  If $u_k\to u$ locally uniformly with
$u_k\in\mathscr P(U)$, then $u_kS\to uS$ and
$u_kT\to uT$ by \Cref{lem:continuous-times-positive-current}, hence
$\dprimeddprime(u_kT)\to\dprimeddprime(uT)$; a second application of that
lemma, with the convergent net of positive closed currents
$\dprimeddprime(u_kT)$ and the constant net $v$, gives
$v\,\dprimeddprime(u_kT)\to v\,\dprimeddprime(uT)$.  Applying
$\dprimeddprime$ once more and using its weak continuity, both sides of the
defining equality pass to the limit.

By \Cref{lem:ch16-cpsh-induction}, $\mathscr P\supseteq\operatorname{CPsh}$.
Since the recursion inserts the potentials one at a time, exchanging
$u_i$ and $u_{i+1}$ in the $q$-fold product is exactly the case just proved
applied to the positive closed current
$\dprimeddprime u_{i+2}\wedge\cdots\wedge\dprimeddprime u_q\wedge T$;
transpositions of adjacent letters generate $\mathfrak S_q$.
\end{proof}

\begin{proposition}
\label{prop:ch16-vanishing-above-dimension}
Let $f:Y\to X$ be a morphism of tropical spaces, let
$p>\dim X$, and let $u_1,\ldots,u_p$ be continuous psh on $X$.  Then
\[
 \dprimeddprime(f^*u_1)\wedge\cdots\wedge\dprimeddprime(f^*u_p)=0 .
\]
\end{proposition}

\begin{proof}
Pullbacks of continuous psh functions are continuous psh
(\Cref{prop:cld-psh-calculus}).  If all $u_i$ are smooth, the product is the
current associated with the smooth form
$f^*(\dprimeddprime u_1\wedge\cdots\wedge\dprimeddprime u_p)$; the form
inside the pullback is a $(p,p)$-form on $X$, whose dimension is $<p$, hence
is zero, and so therefore is its pullback.  The set
of $p$-tuples of continuous psh functions for which the conclusion holds is
a subsheaf of $(\operatorname{CPsh}_X)^p$ containing the smooth tuples and
stable under local uniform limits, by
\Cref{thm:cld-bt-continuity} below; applying
\Cref{lem:ch16-cpsh-induction} in each variable in turn gives the general
case.  This is \cite[Remark~8.8.9]{CLD25}.
\end{proof}

\subsection{Extension to differences of psh functions}
\label{sec:ch16-dpsh-extension}

The recursion produces positive currents from psh potentials.  For the
comparison theorems of \Cref{chap:cld-metrized-bundles} one needs to insert
$G$-smooth, in particular parlinear, potentials that are not psh.  The
extension is by multilinearity, and it is unique.

\begin{theorem}
\label{thm:ch16-dpsh-product}
Let $T$ be a positive closed current of bidegree $(r,r)$ and let $p$ satisfy
$p+r\leq n$.  There is a unique $(p+1)$-linear map
\[
\begin{gathered}
 C^0_X\times\bigl(\operatorname{CPsh}^\pm_X\bigr)^p
 \longrightarrow
 D^{p+r,p+r}(X),\\
 (u_0,u_1,\ldots,u_p)\longmapsto
   u_0\,\dprimeddprime u_1\wedge\cdots\wedge\dprimeddprime u_p\wedge T,
\end{gathered}
\]
which agrees with the recursive product when $u_1,\ldots,u_p$ are psh and
$u_0$ is continuous.  It is symmetric in $u_1,\ldots,u_p$, takes values in
differences of positive currents, and satisfies
\begin{equation}
\label{eq:ch16-dpsh-outer}
 \dprimeddprime\bigl(u_0\,\dprimeddprime u_1\wedge\cdots
   \wedge\dprimeddprime u_p\wedge T\bigr)
 =\dprimeddprime u_0\wedge\dprimeddprime u_1\wedge\cdots
   \wedge\dprimeddprime u_p\wedge T
\end{equation}
when $u_0$ is dpsh.  If $T$ is Bedford--Taylor, the value depends only on
$u_0$ and on the currents $\dprimeddprime u_1,\ldots,\dprimeddprime u_p$.
Finally, if $h^1,\ldots,h^p$ are smooth with $\dprimeddprime h^i=0$, then
\begin{equation}
\label{eq:ch16-harmonic-quotient}
 u_0\,\dprimeddprime(u_1+h^1)\wedge\cdots
    \wedge\dprimeddprime(u_p+h^p)\wedge T
 =u_0\,\dprimeddprime u_1\wedge\cdots
    \wedge\dprimeddprime u_p\wedge T ,
\end{equation}
whether or not $T$ is Bedford--Taylor.
\end{theorem}

\begin{proof}
\emph{Uniqueness} is forced: $\operatorname{CPsh}^\pm_X$ is generated as a
sheaf of $\mathbf R$-vector spaces by $\operatorname{CPsh}_X$, so a
multilinear map is determined by its values on psh arguments, and
$C^0_X$ is spanned by its nonnegative sections.

\emph{Existence} requires that the recursive product be additive in each psh
argument, which is what makes the multilinear extension well defined.
Additivity in $u_1$ is immediate from linearity of $u\mapsto uT$ and of
$\dprimeddprime$.  Additivity in $u_i$ for $i\geq2$ follows from that case
and the symmetry of \Cref{thm:bt-symmetry-locality}.  Compatibility of two
presentations $u=u'-u''=w'-w''$ of the same dpsh function follows from
additivity, since $u'+w''=w'+u''$ is an identity between psh functions.

\emph{Positivity structure.}  Expanding each $u_i$ as a difference of psh
functions writes the value as an alternating sum of $2^p$ positive currents,
each multiplied by $u_0^+$ or $u_0^-$; hence it is a difference of two
positive currents.

\emph{Formula \eqref{eq:ch16-dpsh-outer}.}  Both sides are multilinear, so it
suffices to treat $u_0$ psh, where it is the definition of the recursion
with one more factor inserted at the outside, together with
\Cref{thm:bt-symmetry-locality}.

\emph{Formula \eqref{eq:ch16-harmonic-quotient}.}  By multilinearity it
suffices to insert one $h=h^1$, that is, to prove
$u_0\,\dprimeddprime h\wedge S=0$ for the positive closed current
$S=\dprimeddprime u_2\wedge\cdots\wedge\dprimeddprime u_p\wedge T$.  Since
$h$ is smooth, \Cref{lem:ch16-leibniz-closed} gives
$\dprimeddprime(hS)=\dprimeddprime h\wedge S=0$, and multiplying by the
continuous function $u_0$ preserves the vanishing.  Note that no
Bedford--Taylor hypothesis was used: the point is that $h$ is \emph{smooth},
not merely $\dprimeddprime$-harmonic.
\end{proof}

\begin{corollary}
\label{cor:ch16-harmonic-quotient}
Let $\operatorname{Harm}_X$ denote the sheaf of smooth functions $h$ with
$\dprimeddprime h=0$.  The map of \Cref{thm:ch16-dpsh-product} factors
through
\[
 C^0_X\times\bigl(\operatorname{CPsh}^\pm_X/\operatorname{Harm}_X\bigr)^p .
\]
\end{corollary}

\begin{proof}
This is \eqref{eq:ch16-harmonic-quotient}, which is
\cite[\S8.8.5]{CLD25}.
\end{proof}

\begin{warning}{The quotient is by \emph{smooth} harmonic functions}
\Cref{cor:ch16-harmonic-quotient} quotients by $\operatorname{Harm}_X$, the
sheaf of \emph{smooth} functions killed by $\dprimeddprime$.  It does not
quotient by the larger sheaf of continuous functions $h$ with
$\dprimeddprime[h]=0$; for that, one needs $T$ to be a Bedford--Taylor
current, or the approximation hypothesis of
\Cref{thm:ch16-pluriharmonic-criterion}.  The whole of
\Cref{sec:ch16-harmonic-ambiguity} is about the gap between these two
quotients, and \Cref{ex:ch16-cross-fan} shows that the gap is not empty.
\end{warning}

\subsection{Local and relative Monge--Amp\`ere measures}
\label{sec:ch16-ma-measures}

\begin{definition}[Relative Monge--Amp\`ere operator]
\label{def:ch16-relative-ma}
If $T$ is positive closed and $u$ is continuous psh, define
\[
                         \MA_T(u)=(\dprimeddprime u)^q\wedge T
\]
whenever the top degree is reached.  For $T=\delta_X$ and $q=n$, write
\[
                         \MA(u)=(\dprimeddprime u)^n.
\]
This is a positive Radon measure.
\end{definition}

No global mass normalization is built into this local definition.  On a
proper polarized variety the total mass is later fixed by an intersection
number; on an arbitrary open analytic space it depends on support and
boundary behavior.

\begin{proposition}
\label{prop:ch16-ma-elementary}
Let $u,v$ be continuous psh on $X$ and $T$ positive closed of complementary
bidegree.
\begin{enumerate}
\item\label{itm:ch16-ma:radon}
 $\MA_T(u)$ is a positive Radon measure, and for $u$ dpsh it is a signed
 Radon measure of locally finite total variation.
\item\label{itm:ch16-ma:homog}
 $\MA_T(\lambda u)=\lambda^q\MA_T(u)$ for $\lambda\geq0$, and
 $\MA_T(u+c)=\MA_T(u)$ for a constant $c$.
\item\label{itm:ch16-ma:local}
 $\MA_T(u)$ depends only on the restriction of $u$ to a neighbourhood of any
 given point, and $\MA_T(u)=\MA_T(v)$ on any open set where $u-v$ is smooth
 and $\dprimeddprime$-harmonic.
\item\label{itm:ch16-ma:mixed}
 The polarization
 $(u_1,\ldots,u_q)\mapsto\dprimeddprime u_1\wedge\cdots
  \wedge\dprimeddprime u_q\wedge T$ is the unique symmetric multilinear form
 on $\operatorname{CPsh}^\pm_X$ whose diagonal restriction is $\MA_T$.
\end{enumerate}
\end{proposition}

\begin{proof}
\cref{itm:ch16-ma:radon}: a positive current of bidimension $(0,0)$ is a
positive Radon measure (\Cref{chap:cld-currents}); the dpsh case follows from
\Cref{thm:ch16-dpsh-product}, which writes the value as a difference of two
positive currents.  \cref{itm:ch16-ma:homog} is multilinearity together with
$\dprimeddprime c=0$ for a constant.  \cref{itm:ch16-ma:local} is
\Cref{thm:bt-symmetry-locality} and \eqref{eq:ch16-harmonic-quotient}.
\cref{itm:ch16-ma:mixed} is the usual polarization identity, valid because
the form is multilinear over $\mathbf R$ and the characteristic is zero.
\end{proof}

\subsection{Support on a skeleton}
\label{sec:ch16-support-skeleton}

The following statement is the reason why every Monge--Amp\`ere measure in
this book is supported on a polyhedral set of the right dimension.

\begin{proposition}
\label{prop:ch16-skeletal-support}
Let $u_1,\ldots,u_p$ be tropically $d$-convex on $X$ and let $\alpha$ be a
smooth $(n-p,n-p)$-form.  Then
\[
 u_0\longmapsto
 \bigl\langle
   u_0\,\dprimeddprime u_1\wedge\cdots\wedge\dprimeddprime u_p,\alpha
 \bigr\rangle
\]
is a signed Radon measure on $X$ whose support is locally contained in a
skeletal subset.  More precisely, if $f:U\to E$ is a tropical chart with
$u_j=f^*a_j$ for $d$-convex $a_j$ on a parlinear $P\supseteq f(U)$ and
$\alpha=f^*\omega$, then the support of that measure meets $U$ inside
$\Sigma_f^{(n)}$.
\end{proposition}

\begin{proof}
The measure statement is
\Cref{prop:ch16-ma-elementary}\,\cref{itm:ch16-ma:radon}.  For the support,
assume first that the $a_j$ are smooth and convex.  Then the restriction of
the measure to $U$ is the measure attached to the smooth positive form
$f^*(\dprimeddprime a_1\wedge\cdots\wedge\dprimeddprime a_p\wedge\omega)$,
whose support is contained in $\Sigma_f^{(n)}$ by the definition of the
local skeleton (\Cref{def:trop-local-skeleton}) and the fact that an
$f$-tropical form vanishes off $\Sigma_f^{(n)}$
(\Cref{thm:trop-sigma-characterization}).  Containment of a support in a
fixed closed set is preserved under weak limits, so
\Cref{thm:cld-bt-continuity} extends the conclusion to all continuous convex
$a_j$, using \Cref{cor:ch09-monotone-smooth-approximation} to approximate
them uniformly on compact subsets by smooth convex functions.
Multilinearity extends it to $d$-convex $a_j$.  This is
\cite[Remark~8.8.6]{CLD25}.
\end{proof}

\begin{warning}{The $G$-form and the current do not agree}
If $u_0,\ldots,u_p$ happen to be $G$-$C^2$, the expression
$u_0\,\dprimeddprime u_1\wedge\cdots\wedge\dprimeddprime u_p$ can also be
read as a $G$-form of type $(p,p)$.  The current defined by
\Cref{thm:ch16-dpsh-product} is in general \emph{not} the current associated
with that $G$-form: the difference is a boundary term supported on
codimension-one faces, and it is exactly the term computed in
\Cref{prop:ch16-skeletal-defect}.  The two objects coincide precisely in the
harmonic situation of \Cref{thm:bt-harmonic-compatibility}.  This is
\cite[Remark~8.8.7]{CLD25}, and it is the single most common source of
error in hand computations.
\end{warning}

\section{Continuity: the exact theorem}
\label{sec:ch16-continuity}

\begin{theorem}
\label{thm:cld-bt-continuity}
Let $(T_a)$ be a weakly convergent net of positive closed currents with limit
$T$.  For $i=0,\ldots,q$, let $(u_i^a)$ be a net of continuous functions
converging locally uniformly to $u_i$, and assume that $u_i^a$ is psh for
$i\geq1$.  Then
\[
 u_0^a\,\dprimeddprime u_1^a\wedge\cdots\wedge
 \dprimeddprime u_q^a\wedge T_a
 \longrightarrow
 u_0\,\dprimeddprime u_1\wedge\cdots\wedge
 \dprimeddprime u_q\wedge T.
\]
\end{theorem}
\begin{proof}
Induct on $q$.  The case $q=0$ is
\Cref{lem:continuous-times-positive-current}.  For the induction step,
the inner $(q-1)$-fold products converge and remain positive closed.  Multiply
them by $u_1^a$, apply the lemma, differentiate twice, and then multiply by
$u_0^a$.  Each operation is weakly continuous under the stated hypotheses.
\end{proof}

This is \cite[Proposition~8.7.2]{CLD25}.  It is worth writing the induction
out once, because the order in which the four operations are applied is the
whole content.

\begin{proof}[Second proof]
Put
\[
\begin{gathered}
 S_a=\dprimeddprime\bigl(u_2^a\,\dprimeddprime(\cdots
   \dprimeddprime(u_q^aT_a)\cdots)\bigr),\\
 S=\dprimeddprime\bigl(u_2\,\dprimeddprime(\cdots
   \dprimeddprime(u_qT)\cdots)\bigr).
\end{gathered}
\]
By the inductive hypothesis applied with $q-1$ factors and with the constant
outer function $u_0\equiv1$, the nets $S_a$ and $S$ are positive closed and
$S_a\to S$ weakly.  Four steps now finish the induction:
\begin{enumerate}
\item $u_1^aS_a\to u_1S$, by
 \Cref{lem:continuous-times-positive-current} applied to the positive net
 $(S_a)$ and the uniformly convergent net $(u_1^a)$;
\item $\dprimeddprime(u_1^aS_a)\to\dprimeddprime(u_1S)$, by weak continuity
 of $\dprimeddprime$ (\Cref{prop:current-differential-calculus});
\item each $\dprimeddprime(u_1^aS_a)$ is positive and closed by
 \Cref{thm:one-bt-factor}, hence so is the limit
 $\dprimeddprime(u_1S)$;
\item $u_0^a\,\dprimeddprime(u_1^aS_a)\to u_0\,\dprimeddprime(u_1S)$, by a
 second application of \Cref{lem:continuous-times-positive-current}, this
 time to the positive net produced in step (3).
\end{enumerate}
Step (3) is not decoration: without it, step (4) would multiply a net of
currents that is not known to be positive, and
\Cref{lem:continuous-times-positive-current} would not apply.
\end{proof}

Two useful consequences are worth stating precisely.

\begin{corollary}
\label{cor:ch16-decreasing-nets}
If nets $u_i^a\downarrow u_i$ share a directed set and every $u_i^a,u_i$ is
continuous psh, then the mixed products converge locally weakly.
\end{corollary}
\begin{proof}
On each compact set, the net form of Dini's theorem,
\Cref{prop:dini-nets}, makes the decreasing convergence uniform; apply
\Cref{thm:cld-bt-continuity}.
\end{proof}

\begin{corollary}
\label{cor:ch16-monotone-ma}
Let $u_k$ be continuous psh with $u_k\downarrow u$ pointwise, $u$
continuous.  Then $\MA_T(u_k)\to\MA_T(u)$ weakly, and for every open
relatively compact $U$,
\[
 \MA_T(u)(U)\ \leq\ \liminf_k\MA_T(u_k)(U),
 \qquad
 \MA_T(u)(\overline U)\ \geq\ \limsup_k\MA_T(u_k)(\overline U).
\]
\end{corollary}

\begin{proof}
Weak convergence is \Cref{cor:ch16-decreasing-nets} with a sequence.  The
two inequalities are the portmanteau inequalities for weakly convergent
positive Radon measures, valid for open and for compact sets respectively.
\end{proof}

\begin{remark}[What is not claimed]
\label{rem:ch16-not-claimed}
If the decreasing limit is discontinuous or merely locally bounded, Dini's
argument fails and \Cref{thm:cld-bt-continuity} does not apply.  Thus
the classical complex Bedford--Taylor convergence theorem for arbitrary
locally bounded psh functions is not being asserted in the CLD setting.
Likewise, the full comparison principle
\[
 \int_{\{u<v\}}(\dprimeddprime v)^n
 \leq\int_{\{u<v\}}(\dprimeddprime u)^n
\]
is not a formal consequence of the results quoted here and is not stated in
this generality in the CLD source.  Global comparison principles used later
come from the BFJ/BJ energy theories under their own hypotheses.
\end{remark}

There are nevertheless rigorous comparison statements already available:
the product is unchanged when a potential is replaced by another with the
same Hessian against a Bedford--Taylor current; it is compatible with proper
pushforward; and in harmonic or toric charts it agrees with the corresponding
real convex product.  These are the interfaces proved below.
\Cref{sec:ch16-comparison} proves the comparison principle in the case where
it is a theorem---on a tropical chart, where convexity is available---and
records exactly what that case does and does not give.

\subsection{A variant adapted to harmonic charts}
\label{sec:ch16-harmonic-continuity}

The next statement will be needed in \Cref{sec:ch16-harmonic-charts}, where
the potentials are pullbacks of convex functions under a harmonic chart and
are therefore \emph{not} psh in general: only their products with a positive
closed current are positive.

\begin{proposition}
\label{prop:ch16-harmonic-continuity}
Let $(T_a)$ be a weakly convergent net of positive closed currents with
limit $T$, let $f:X\to\mathbf R^m$ be a harmonic chart, and let
$v_i^a\to v_i$ locally uniformly with all $v_i^a,v_i$ continuous convex on
$\mathbf R^m$.  Put $u_i^a=f^*v_i^a$ and $u_i=f^*v_i$.  Then
\[
 u_0^a\,\dprimeddprime u_1^a\wedge\cdots\wedge\dprimeddprime u_q^a\wedge T_a
 \longrightarrow
 u_0\,\dprimeddprime u_1\wedge\cdots\wedge\dprimeddprime u_q\wedge T .
\]
\end{proposition}

\begin{proof}
The proof of \Cref{thm:cld-bt-continuity} is repeated verbatim, with a
single change at step (3): positivity and closedness of
$\dprimeddprime(u_1^aS_a)$ are no longer supplied by
\Cref{thm:one-bt-factor}, whose hypothesis is that $u_1^a$ be psh, but by
\Cref{cor:ch16-harmonic-convex-positive}, whose hypothesis is that $u_1^a$
be the pullback of a convex function by a harmonic chart.  All other steps
use only continuity of the potentials and positivity of the currents.  This
is \cite[Proposition~8.7.3]{CLD25}; the reader should note the forward
reference, which is not circular, since
\Cref{cor:ch16-harmonic-convex-positive} is proved from
\Cref{thm:bt-harmonic-compatibility} and
\Cref{thm:cld-bt-continuity}, neither of which uses the present statement.
\end{proof}

\section{Chern--Levine--Nirenberg estimates}
\label{sec:ch16-cln}

The Chern--Levine--Nirenberg inequality \cite{CLN69} bounds the local mass of
a Monge--Amp\`ere current by a product of sup-norms of its potentials.  It is
the only quantitative statement in the chapter, and the only one whose proof
does not go through the sheaf-closure principle: the constant produced by the
induction is uniform in the potentials, and it is that uniformity, not a
closure property, that lets the estimate pass to limits.

\subsection{The basic inequality}
\label{sec:ch16-cln-basic}

\begin{theorem}
\label{thm:cld-cln}
Let $T$ be positive closed, let $q\geq0$, and let $\omega$ be a smooth test
form of the complementary bidegree.  For every compact neighborhood $K$ of
$\Supp(\omega)$ there is a constant $C_{K,\omega}$ such that, for smooth
functions $u_0,\ldots,u_q$ with $u_1,\ldots,u_q$ psh,
\[
 \left|\left\langle
 u_0\,\dprimeddprime u_1\wedge\cdots\wedge
 \dprimeddprime u_q\wedge T,\omega\right\rangle\right|
 \leq C_{K,\omega}\prod_{i=0}^q\|u_i\|_K.
\]
The same estimate holds for continuous psh potentials in the CLD product.
\end{theorem}
\begin{proof}
Proceed by induction on $q$.  Write $\omega=\alpha_0-\beta_0$ with
$\alpha_0,\beta_0$ positive and compactly supported in the interior of $K$,
and set $\omega_0=\alpha_0+\beta_0$.  Positivity gives
\[
 \left|\left\langle
 u_0\,\dprimeddprime u_1\wedge\cdots\wedge
 \dprimeddprime u_q\wedge T,\omega\right\rangle\right|
 \leq \|u_0\|_K
 \left\langle\dprimeddprime u_1\wedge\cdots\wedge
 \dprimeddprime u_q\wedge T,\omega_0\right\rangle.
\]
Closedness and Green's formula identify the last pairing with
\[
 \left\langle
 u_1\,\dprimeddprime u_2\wedge\cdots\wedge
 \dprimeddprime u_q\wedge T,\dprimeddprime\omega_0\right\rangle.
\]
The induction hypothesis, applied to the fixed test form
$\dprimeddprime\omega_0$, proves the estimate.  This is the induction in
\cite[Proposition~8.4.1]{CLD25}.  For the final
extension, the estimate defines a sheafwise condition stable at each uniform
limit stage by \Cref{thm:cld-bt-continuity}; transfinite induction
through the continuous psh hierarchy gives the stated version.  No direct
smooth approximation of a higher-stage potential is used.
\end{proof}

The proof produces the constant, and it is worth extracting it, because in
every later application the dependence on $T$ must be visible.

\begin{proposition}
\label{prop:ch16-cln-constant}
Fix a compact $K$ and a smooth $\omega$ supported in the interior of $K$.
Define smooth compactly supported forms $\Omega_0,\ldots,\Omega_q$ by
\[
 \Omega_0=\alpha_0+\beta_0,
 \qquad
 \Omega_{j+1}=(\dprimeddprime\Omega_j)^++(\dprimeddprime\Omega_j)^-,
\]
where $\gamma=\gamma^+-\gamma^-$ denotes a fixed decomposition of a smooth
form into positive smooth forms supported in the interior of $K$, and
$\omega=\alpha_0-\beta_0$ is such a decomposition of $\omega$.  Then one may
take
\[
 C_{K,\omega}=\langle T,\Omega_q\rangle .
\]
In particular $C_{K,\omega}$ is linear in $T$, is finite because $T$ has
order zero, and does not depend on $u_0,\ldots,u_q$.
\end{proposition}

\begin{proof}
Induction on $q$.  For $q=0$ the estimate is
\eqref{eq:ch16-positivity-sandwich} with $S'=T$ and $w=u_0$, whose right-hand
side is $\|u_0\|_K\langle T,\Omega_0\rangle$.  Assume the statement for
$q-1$ and all test forms.  With $R:=\dprimeddprime u_1\wedge\cdots
\wedge\dprimeddprime u_q\wedge T$, positivity of $R$
(\Cref{prop:ch16-iterate-basic}\,\cref{itm:ch16-iter:posclosed}) and
\eqref{eq:ch16-positivity-sandwich} give
\[
 |\langle u_0R,\omega\rangle|\leq\|u_0\|_K\langle R,\Omega_0\rangle .
\]
Since $R=\dprimeddprime\bigl(u_1\,\dprimeddprime u_2\wedge\cdots
\wedge\dprimeddprime u_q\wedge T\bigr)$ and $X$ is boundaryless, the pairing
formula of \Cref{prop:current-differential-calculus} gives
\[
 \langle R,\Omega_0\rangle
 =\bigl\langle u_1\,\dprimeddprime u_2\wedge\cdots
   \wedge\dprimeddprime u_q\wedge T,\ \dprimeddprime\Omega_0\bigr\rangle .
\]
Apply the induction hypothesis with $q-1$ factors and the test form
$\dprimeddprime\Omega_0$, whose associated first form is
$(\dprimeddprime\Omega_0)^++(\dprimeddprime\Omega_0)^-=\Omega_1$.  The
constant is $\langle T,\Omega_q\rangle$ and the sup-norm factors multiply.
\end{proof}

\begin{corollary}
\label{cor:ch16-cln-oscillation}
With the notation of \Cref{prop:ch16-cln-constant} and $q\geq1$,
\[
 \left|\left\langle
   u_0\,\dprimeddprime u_1\wedge\cdots\wedge\dprimeddprime u_q\wedge T,
   \omega\right\rangle\right|
 \leq 2^{-q}\,\langle T,\Omega_q\rangle\,\|u_0\|_K
   \prod_{i=1}^q\operatorname{osc}_K(u_i),
\]
where $\operatorname{osc}_K(u)=\sup_Ku-\inf_Ku$.  In particular the estimate
is invariant under adding constants to the psh potentials, as it must be.
\end{corollary}

\begin{proof}
For $i\geq1$ the current
$\dprimeddprime u_1\wedge\cdots\wedge\dprimeddprime u_q\wedge T$ is unchanged
if $u_i$ is replaced by $u_i+c_i$, by
\eqref{eq:ch16-harmonic-quotient} applied to the smooth harmonic function
$c_i$.  Choosing
$c_i=-\tfrac12(\sup_Ku_i+\inf_Ku_i)$ makes
$\|u_i+c_i\|_K=\tfrac12\operatorname{osc}_K(u_i)$, and
\Cref{thm:cld-cln} applies to the shifted potentials.
\end{proof}

\begin{corollary}
\label{cor:ch16-cln-mass}
Let $K\Subset L$ be compact subsets of $X$ with $K$ contained in the
interior of $L$, and fix $q$.  There is a constant $C=C(K,L,q,T)$ such that
for all continuous psh $u_1,\ldots,u_q$,
\[
 \bigl\|\dprimeddprime u_1\wedge\cdots\wedge
   \dprimeddprime u_q\wedge T\bigr\|_K
 \leq C\prod_{i=1}^q\operatorname{osc}_L(u_i).
\]
Consequently, if $(u_i^{(k)})_k$ are continuous psh with uniformly bounded
oscillations on $L$, the corresponding products form a family of positive
currents with uniformly bounded mass on $K$, and every weak cluster point is
positive.
\end{corollary}

\begin{proof}
Choose a smooth cutoff $\chi$ with $0\leq\chi\leq1$, $\chi=1$ on $K$ and
$\Supp\chi$ in the interior of $L$
(\Cref{thm:cld-smooth-partition}), and a smooth strongly positive form
$\Theta$ of complementary bidegree on a chart neighbourhood.  Apply
\Cref{cor:ch16-cln-oscillation} with $u_0=1$ and $\omega=\chi\Theta$.  The
mass of a positive current on $K$ is bounded by a fixed multiple of its
pairing with $\chi\Theta$, by
\Cref{prop:positive-current-compactness}.  The last assertion is
\Cref{prop:positive-current-compactness} again.
\end{proof}

\subsection{A barrier, and weighted forms}
\label{sec:ch16-barrier}

The polar theorem of \Cref{sec:ch16-polar} requires a version of the
estimate in which the outer function is an \emph{unbounded} psh function.
The device is to replace the cutoff $\chi$ by a family of weights built from
a psh exhaustion.  We first record the existence of such an exhaustion; this
is the second of the three imported inputs of \Cref{sec:ch16-imported}.

\begin{theorem}
\label{thm:ch16-barrier}
Let $X$ be a good, boundaryless $n$-dimensional $K$-analytic space.  Every
point of $X$ has a basis of open neighbourhoods $U$ on which there is a
smooth psh function $\rho:U\to\mathbf R$ with $\rho<0$ and
$\{\rho\leq0\}$ compact.
\end{theorem}

\begin{proof}
This is \cite[Proposition~16.9.1]{CLD25}.  We recall its shape, since the
hypotheses matter.  Goodness supplies an affinoid neighbourhood $V$ of the
given point $x$, and inside the relative interior of $V$ a rational
subdomain $W\ni x$ cut out by $|f_i|\leq\lambda_i|g|$ for analytic
$f_1,\ldots,f_m,g$ on $V$ with no common zero and $\lambda_i>0$.  On the
locus $U\subseteq V$ where $g$ is invertible the functions
$u_i=\log(|f_i/g|/\lambda_i)$ are psh and parlinear off their polar sets,
and $u=\max_iu_i$ is psh with $\{u\leq t\}$ a compact rational domain for
every $t\geq0$, because the $f_i$ and $g$ have no common zero.  Smoothing
$u$ by the regularized maxima of \eqref{eq:ch16-softmax} produces smooth psh
$u_\varepsilon$ with $u\leq u_\varepsilon\leq u+\varepsilon$; taking
$\varepsilon=1$ and $\rho=u_1-u_1(x)-1$ gives a smooth psh function with
$\rho(x)=-1<0$ and $\{\rho\leq0\}\subseteq\{u\leq u_1(x)+1\}$ compact.
\end{proof}

\begin{definition}[The weight forms]
\label{def:ch16-weight-forms}
Let $\rho$ be smooth psh on $X$ with $U=\{\rho<0\}$ relatively compact, put
$\rho_-=\max\{-\rho,0\}$, and fix a real number $a>1$.  For
$0\leq k\leq n$ define an $(n-k,n-k)$-form on $U$ by
\begin{equation}
\label{eq:ch16-beta-k}
 \beta_k=\rho_-^{k+a}\,(\dprimeddprime\rho)^{n-k}
   +(k+a)\,\rho_-^{k+a-1}\,
     \dprime\rho\wedge\ddprime\rho\wedge(\dprimeddprime\rho)^{n-k-1},
\end{equation}
the second term being omitted when $k=n$.
\end{definition}

\begin{lemma}
\label{lem:ch16-beta-properties}
With the notation of \Cref{def:ch16-weight-forms}:
\begin{enumerate}
\item\label{itm:ch16-beta:reg}
 $\beta_k$ has coefficients of class $C^{\lfloor a\rfloor-1}$, is positive,
 and has compact support contained in $\overline U$;
\item\label{itm:ch16-beta:dd}
 $\ddprime\beta_k
   =-2(k+a)\rho_-^{k+a-1}\ddprime\rho\wedge(\dprimeddprime\rho)^{n-k}$ and
 \[
 \begin{aligned}
  \dprimeddprime\beta_k
  =2(k+a)\Bigl[&-\rho_-^{k+a-1}(\dprimeddprime\rho)^{n-k+1}\\
   &+(k+a-1)\rho_-^{k+a-2}
     \dprime\rho\wedge\ddprime\rho\wedge(\dprimeddprime\rho)^{n-k}\Bigr];
 \end{aligned}
 \]
\item\label{itm:ch16-beta:sandwich}
 $-2(k+a)\beta_{k-1}\leq\dprimeddprime\beta_k\leq2(k+a)\beta_{k-1}$ for
 $1\leq k\leq n$.
\end{enumerate}
\end{lemma}

\begin{proof}
\cref{itm:ch16-beta:reg}  For $a>1$ the function
$\theta_a(t)=\max\{-t,0\}^a$ is of class $C^{\lfloor a\rfloor}$ with
$\theta_a'=-a\theta_{a-1}$, whence
$\dprime\rho_-^{\,b}=-b\rho_-^{\,b-1}\dprime\rho$ and similarly for
$\ddprime$, valid for $b>1$; since $k+a-1>0$, both terms of
\eqref{eq:ch16-beta-k} are of class $C^{\lfloor a\rfloor-1}$ and vanish
where $\rho\geq0$.  Positivity: $\dprimeddprime\rho$ is positive because
$\rho$ is smooth psh, $\dprime\rho\wedge\ddprime\rho$ is positive by
\Cref{prop:ch09-positivity-dictionary}, and products of positive forms with
nonnegative coefficients are positive
(\Cref{prop:ch09-positivity-products}).

\cref{itm:ch16-beta:dd}  Differentiate \eqref{eq:ch16-beta-k}.  In the first
term, $\ddprime\rho_-^{k+a}=-(k+a)\rho_-^{k+a-1}\ddprime\rho$.  In the
second, the coefficient contributes
$-(k+a)(k+a-1)\rho_-^{k+a-2}\ddprime\rho\wedge\dprime\rho\wedge
\ddprime\rho\wedge(\cdots)$, which vanishes because
$\ddprime\rho\wedge\ddprime\rho=0$; and
$\ddprime(\dprime\rho\wedge\ddprime\rho)
 =-\dprimeddprime\rho\wedge\ddprime\rho$ by $\ddprime\dprime=-\dprimeddprime$
and $(\ddprime)^2=0$.  Adding the two contributions gives the stated
$\ddprime\beta_k$.  Applying $\dprime$ once more and using
$\dprime\rho_-^{k+a-1}=-(k+a-1)\rho_-^{k+a-2}\dprime\rho$ gives the formula
for $\dprimeddprime\beta_k$.

\cref{itm:ch16-beta:sandwich}  Write
$\beta_{k-1}=\rho_-^{k-1+a}(\dprimeddprime\rho)^{n-k+1}
 +(k-1+a)\rho_-^{k+a-2}\dprime\rho\wedge\ddprime\rho\wedge
  (\dprimeddprime\rho)^{n-k}$; both summands are positive.  The upper bound
follows by discarding the negative first term of $\dprimeddprime\beta_k$ and
comparing the second with $2(k+a)$ times the second summand of
$\beta_{k-1}$, the coefficients being $(k+a-1)$ and $(k-1+a)$, which are
equal.  The lower bound follows by discarding the positive second term and
comparing the first with $-2(k+a)$ times the first summand of
$\beta_{k-1}$.  This is the computation of \cite[\S8.4.2]{CLD25}.
\end{proof}

\subsection{The weighted inequality with its constant}
\label{sec:ch16-cln-weighted}

\begin{theorem}
\label{thm:ch16-weighted-cln}
Assume $a\geq3$ and put
\begin{equation}
\label{eq:ch16-cln-constant}
 C(k,a)=(4\rme)^k\,a^{-1}\,(k+a)\,(2+a)(3+a)\cdots(k+1+a).
\end{equation}
Let $V$ be a smooth negative psh function on $U=\{\rho<0\}$ and let
$u_1,\ldots,u_k$ be smooth psh functions on $X$.  Then
\begin{equation}
\label{eq:ch16-weighted-cln-main}
 \int_U\beta_k\wedge(-V)\,\dprimeddprime u_1\wedge\cdots\wedge
   \dprimeddprime u_k
 \;\leq\;
 C(k,a)\,\|u_1\|_U\cdots\|u_k\|_U\int_U(-V)\beta_0 .
\end{equation}
\end{theorem}

\begin{proof}
Write $J_k$ for the left-hand side and
$\psi_k=\dprimeddprime u_1\wedge\cdots\wedge\dprimeddprime u_k$, so that
$J_k=\int_U(-V)\beta_k\wedge\psi_k$ and $J_0=\int_U(-V)\beta_0$.  We prove
\begin{equation}
\label{eq:ch16-cln-recursion}
 J_k\ \leq\ 4\rme\,\frac{(k+a)(k+a+1)}{k+a-1}\,\|u_k\|_U\,J_{k-1},
 \qquad k\geq1,
\end{equation}
from which \eqref{eq:ch16-weighted-cln-main} follows by iteration: the
telescoping product
$\prod_{j=1}^k\frac{(j+a)(j+a+1)}{j+a-1}$ equals
$a^{-1}(k+a)(2+a)(3+a)\cdots(k+1+a)$, because
$\prod_{j=1}^k(j+a)\big/\prod_{j=1}^k(j+a-1)=(k+a)/a$ and
$\prod_{j=1}^k(j+a+1)=(2+a)\cdots(k+1+a)$; multiplying by $(4\rme)^k$ gives
\eqref{eq:ch16-cln-constant}.  For $k=0$ the inequality
\eqref{eq:ch16-weighted-cln-main} is an equality, and $C(0,a)=1$.

\emph{Step 1: two integrations by parts.}  Since $\beta_k$ vanishes off $U$
and $X$ is boundaryless, Green's formula
(\Cref{cor:analytic-stokes-green}) gives
\[
 J_k=\int_X(-V\beta_k)\wedge\dprimeddprime(u_k\psi_{k-1})
    =\int_X\dprimeddprime(-V\beta_k)\wedge u_k\psi_{k-1}.
\]
Expanding $\dprimeddprime(V\beta_k)$ by the Leibniz rule and using symmetry
of $V$, $\beta_k$ and $\psi_{k-1}$ to cancel the two mixed first-order terms
against each other, then integrating by parts once more in the surviving
mixed term $2\,\dprime V\wedge\ddprime\beta_k$, one obtains
\begin{equation}
\label{eq:ch16-cln-after-parts}
 J_k=\int_Uu_k\bigl(V\dprimeddprime\beta_k
     -\dprimeddprime V\wedge\beta_k\bigr)\wedge\psi_{k-1}
   +2\int_UV\,\dprime u_k\wedge\ddprime\beta_k\wedge\psi_{k-1}.
\end{equation}

\emph{Step 2: the Cauchy--Schwarz trick.}  For $C^1$ functions $u,v$ and
$h>0$ the positivity of $(\dprime u-h\dprime v)\wedge(\ddprime u-h\ddprime
v)$ gives
\begin{equation}
\label{eq:ch16-cln-cauchy}
 \dprime u\wedge\ddprime v+\dprime v\wedge\ddprime u
 \leq h^{-1}\dprime u\wedge\ddprime u+h\,\dprime v\wedge\ddprime v .
\end{equation}
Assume for the moment that $u_k>0$.  Using
\Cref{lem:ch16-beta-properties}\,\cref{itm:ch16-beta:dd} to expand
$\ddprime\beta_k$, applying \eqref{eq:ch16-cln-cauchy} with $h=u_k$ and
$v=\rho$, and using $\dprimeddprime V\geq0$, the two displayed terms of
\eqref{eq:ch16-cln-after-parts} are dominated by
\[
\begin{aligned}
 &\int_Uu_k(-V)\Bigl[-\dprimeddprime\beta_k
  +4(k+a)^2\rho_-^{k+a-2}\dprime\rho\wedge\ddprime\rho\wedge
   (\dprimeddprime\rho)^{n-k}\Bigr]\wedge\psi_{k-1}\\
 &\qquad+\int_U(-V)\beta_k\wedge\psi_{k-1}\wedge u_k^{-1}
   \dprime u_k\wedge\ddprime u_k .
\end{aligned}
\]
By \Cref{lem:ch16-beta-properties}\,\cref{itm:ch16-beta:dd} the bracket
equals
\[
 2(k+a)\rho_-^{k+a-1}(\dprimeddprime\rho)^{n-k+1}
 +(k+a+1)\,2(k+a)\,\rho_-^{k+a-2}
   \dprime\rho\wedge\ddprime\rho\wedge(\dprimeddprime\rho)^{n-k},
\]
which is at most $c\,\beta_{k-1}$ with
$c=2(k+a)(k+a+1)/(k+a-1)$.  Hence
\begin{equation}
\label{eq:ch16-cln-positive-case}
 J_k\ \leq\ c\,\|u_k\|_U\,J_{k-1}
   +\int_U(-V)\beta_k\wedge\psi_{k-1}\wedge
     u_k^{-1}\dprime u_k\wedge\ddprime u_k ,
\end{equation}
valid whenever $u_k$ is strictly positive.

\emph{Step 3: removing the last term by an exponential substitution.}  Let
$M=\|u_k\|_U$, let $B>0$, and apply \eqref{eq:ch16-cln-positive-case} to
$v_k=\exp(Bu_k)$, which is smooth, psh by
\Cref{prop:ch16-composition}, and strictly positive.  From
$\dprime v_k=Bv_k\dprime u_k$ one gets
\[
 \dprimeddprime v_k
 =v_k^{-1}\dprime v_k\wedge\ddprime v_k+Bv_k\,\dprimeddprime u_k
 \ \geq\ B\rme^{-BM}\dprimeddprime u_k
   +v_k^{-1}\dprime v_k\wedge\ddprime v_k ,
\]
using $u_k\geq-M$ and positivity of $\dprimeddprime u_k$.  Therefore
\[
\begin{aligned}
 B\rme^{-BM}J_k
 &\leq\int_U\beta_k\wedge(-V)\psi_{k-1}\wedge\dprimeddprime v_k
   -\int_U\beta_k\wedge(-V)\psi_{k-1}\wedge
      v_k^{-1}\dprime v_k\wedge\ddprime v_k\\
 &\leq c\,\|v_k\|_UJ_{k-1}\ \leq\ c\,\rme^{BM}J_{k-1},
\end{aligned}
\]
the middle inequality being \eqref{eq:ch16-cln-positive-case} applied to
$v_k$.  Thus $J_k\leq cB^{-1}\rme^{2BM}J_{k-1}$, and the choice
$B=1/(2M)$ gives $J_k\leq2\rme\,c\,M\,J_{k-1}$, which is
\eqref{eq:ch16-cln-recursion}.  This is
\cite[Proposition~8.4.3]{CLD25}.
\end{proof}

\begin{remark}[Size of the constant]
\label{rem:ch16-cln-constant-size}
For the standard choice $a=3$ used in \Cref{sec:ch16-polar} one has
\[
 C(k,3)=\frac{(4\rme)^k}{3}\,(k+3)\cdot5\cdot6\cdots(k+4)
 =\frac{(4\rme)^k(k+3)(k+4)!}{72}.
\]
Thus $C(1,3)=4\rme\cdot4\cdot120/72=\tfrac{80}{3}\rme\approx72.5$.  The
factorial growth in $k$ is an artifact of the exponential substitution in
Step 3 and is irrelevant for the applications, in which $k\leq n$ is fixed
and only finiteness matters.  What \emph{is} essential is that $C(k,a)$
depends on neither $V$ nor the $u_i$; this is what allows both limits in
\Cref{thm:cld-polar-negligibility} to be taken.
\end{remark}

The CLN bound is the local mass estimate behind weak compactness and polar
negligibility.  It does not control energy tails of unbounded potentials;
that belongs to the global finite-energy theory.

\section{Comparison principles: what holds, and where}
\label{sec:ch16-comparison}

\Cref{rem:ch16-not-claimed} recorded that the comparison principle is not a
formal consequence of the CLD construction.  That is a statement about
\emph{generality}, not a statement that no comparison principle is
available.  On a tropical chart, where psh means convex
(\Cref{prop:cld-psh-positive-current} and
\Cref{thm:ch09-distributional-convexity}), the comparison principle is a
theorem of convex analysis with a short proof, and it is in that form that
the book uses it: in \Cref{chap:cld-examples} for toric metrics, and in
\Cref{part:bfj} as the model for the model-metric maximum principle
\eqref{eq:model-maximum-principle}.  This section proves the chart case,
the metric-graph case, and one genuinely global statement---the domination
of masses---that does hold in the CLD generality.

\subsection{The subgradient lemma}
\label{sec:ch16-subgradient}

Throughout this subsection $\Omega\subseteq\mathbf R^n$ is a bounded open
set.  For $u$ convex on $\Omega$ and $E\subseteq\Omega$ write
$\partial u(E)=\bigcup_{x\in E}\partial u(x)$ for the union of the
subdifferentials, and recall from \Cref{app:convexity} that
$E\mapsto|\partial u(E)|$, the Lebesgue measure of the subgradient image, is
the Alexandrov Monge--Amp\`ere measure $\MA_{\mathbf R}(u)$, and that
\begin{equation}
\label{eq:ch16-alexandrov-normalization}
 (\dprimeddprime u)^n=n!\,\MA_{\mathbf R}(u)
\end{equation}
by \Cref{prop:ch09-ma-coordinates} and
\Cref{thm:appD-alexandrov-measure}.

\begin{lemma}
\label{lem:ch16-subgradient-inclusion}
Let $u,v:\overline\Omega\to\mathbf R$ be continuous, with $u\leq v$ on
$\Omega$ and $u=v$ on $\partial\Omega$.  Then
$\partial v(\Omega)\subseteq\partial u(\Omega)$.
\end{lemma}

\begin{proof}
Let $p\in\partial v(x_0)$ with $x_0\in\Omega$, so that
$v(x)\geq\ell(x):=v(x_0)+\langle p,x-x_0\rangle$ for all
$x\in\overline\Omega$.  The continuous function $\ell-u$ attains its maximum
$m$ on the compact set $\overline\Omega$, say at $x_1$.  Then
$\ell-m\leq u$ on $\overline\Omega$ with equality at $x_1$; hence
$p\in\partial u(x_1)$ provided $x_1\in\Omega$.

Suppose instead that $x_1\in\partial\Omega$.  There
$u(x_1)=v(x_1)\geq\ell(x_1)$, so $m=\ell(x_1)-u(x_1)\leq0$.  On the other
hand $m\geq\ell(x_0)-u(x_0)=v(x_0)-u(x_0)\geq0$.  Hence $m=0$ and the
maximum is also attained at $x_0\in\Omega$, so $p\in\partial u(x_0)$.  In
both cases $p\in\partial u(\Omega)$.
\end{proof}

\begin{theorem}
\label{thm:ch16-comparison-chart}
Let $u,v$ be convex and continuous on $\overline\Omega$ with $u\geq v$ on
$\partial\Omega$, and put $G=\{x\in\Omega:u(x)<v(x)\}$.  Then
\[
 \int_G(\dprimeddprime v)^n\ \leq\ \int_G(\dprimeddprime u)^n .
\]
\end{theorem}

\begin{proof}
If $G=\emptyset$ there is nothing to prove.  Otherwise $G$ is open and
bounded, and $\partial G\subseteq\{u=v\}$: a boundary point of $G$ lying in
$\Omega$ satisfies $u=v$ by continuity, and a boundary point lying on
$\partial\Omega$ satisfies $u\geq v$ by hypothesis and $u\leq v$ as a limit
from $G$.  Apply \Cref{lem:ch16-subgradient-inclusion} on $G$ to get
$\partial v(G)\subseteq\partial u(G)$, hence
$\MA_{\mathbf R}(v)(G)\leq\MA_{\mathbf R}(u)(G)$; multiply by $n!$ and use
\eqref{eq:ch16-alexandrov-normalization}.
\end{proof}

\begin{corollary}
\label{cor:ch16-domination-chart}
Let $u,v$ be convex and continuous on $\overline\Omega$ with $u=v$ on
$\partial\Omega$.  If $(\dprimeddprime u)^n\geq(\dprimeddprime v)^n$ as
measures on $\Omega$, then $u\leq v$ on $\Omega$.
\end{corollary}

\begin{proof}
Suppose not, and let $d=\operatorname{diam}\Omega$.  Choose $x_0\in\Omega$
with $u(x_0)>v(x_0)$ and $\varepsilon>0$ so small that
$\varepsilon d^2<u(x_0)-v(x_0)$.  Put
$u_\varepsilon(x)=u(x)+\varepsilon(|x-x_0|^2-d^2)$, a convex continuous
function on $\overline\Omega$ with $u_\varepsilon\leq u$ and
$u_\varepsilon(x_0)>v(x_0)$.  On $\partial\Omega$,
$v=u\geq u_\varepsilon$, so
$G=\{x\in\Omega:v(x)<u_\varepsilon(x)\}$ is a nonempty open set with
$u_\varepsilon=v$ on $\partial G$, and $G$ has positive Lebesgue measure.
\Cref{lem:ch16-subgradient-inclusion}, applied on $G$ with the smaller
function $v$ and the larger $u_\varepsilon$, gives
$\MA_{\mathbf R}(u_\varepsilon)(G)\leq\MA_{\mathbf R}(v)(G)$.  On the other
hand the Minkowski determinant inequality
$\det(A+B)\geq\det A+\det B$ for positive semidefinite $A,B$, applied to the
Hessians and extended to general convex functions by
\Cref{cor:ch09-monotone-smooth-approximation} and
\Cref{cor:ch16-monotone-ma}, gives
\[
 \MA_{\mathbf R}(u_\varepsilon)\ \geq\ \MA_{\mathbf R}(u)
   +(2\varepsilon)^n\lambda
 \ \geq\ \MA_{\mathbf R}(v)+(2\varepsilon)^n\lambda ,
\]
$\lambda$ denoting Lebesgue measure.  Evaluating on $G$ yields
$\MA_{\mathbf R}(v)(G)\geq\MA_{\mathbf R}(v)(G)+(2\varepsilon)^n|G|$, a
contradiction.
\end{proof}

\begin{warning}{Two things this does not prove}
\Cref{thm:ch16-comparison-chart} is a statement about a bounded open subset
of $\mathbf R^n$ with a genuine boundary, and it is proved by convex
analysis, not by the Bedford--Taylor recursion.  Two extensions are
\emph{not} claimed.  First, it does not follow that the same inequality
holds on an arbitrary boundaryless tropical or analytic $X$ for arbitrary
continuous psh $u,v$: the hypothesis $u\geq v$ on $\partial\Omega$ has no
substitute, and a chartwise argument would have to control the transition
between charts, where the analogue of convexity is not available.  Second,
it does not extend to unbounded potentials without an integrability
hypothesis; the corresponding statements in the global theories carry their
own finite-energy assumptions and are proved in
\Cref{chap:bfj-ma,chap:bj-finite-energy}.
\end{warning}

\subsection{Comparison on a metric graph}
\label{sec:ch16-comparison-graph}

In dimension one the comparison principle is the maximum principle, and it
holds on any compact metric graph with boundary.

\begin{proposition}
\label{prop:ch16-comparison-graph}
Let $\Gamma$ be a compact metric graph, let $\partial\subseteq\Gamma$ be a
nonempty finite set meeting every connected component, and let $f,g$ be
continuous, piecewise affine on $\Gamma$.  If $\Delta_\Gamma f\geq
\Delta_\Gamma g$ on $\Gamma\setminus\partial$ and $f\leq g$ on $\partial$,
then $f\leq g$ on $\Gamma$.
\end{proposition}

\begin{proof}
Put $h=f-g$, a continuous piecewise affine function with
$\Delta_\Gamma h\geq0$ on $\Gamma\setminus\partial$ and $h\leq0$ on
$\partial$.  Suppose $m:=\max_\Gamma h>0$, and let $F=\{h=m\}$, a nonempty
closed set disjoint from $\partial$.  Choose a connected component $C$ of
$\Gamma$ meeting $F$.  Since $\partial\cap C\neq\emptyset$ and $h\leq0$
there, $F\cap C$ is a nonempty proper closed subset of the connected space
$C$, so it has a topological boundary point $p$ in $C$.  Every outgoing
slope of $h$ at $p$ is $\leq0$, because $h(p)=m$ is the global maximum.
Moreover some tangent direction $\nu$ at $p$ contains points not in $F$
arbitrarily close to $p$; since $h$ is affine on the germ of that direction,
$h(p+t\nu)=m+t\,d_\nu h(p)$ for small $t$, and the existence of points with
$h<m$ forces $d_\nu h(p)<0$.  Hence
$\bigl(\Delta_\Gamma h\bigr)(\{p\})<0$ by \eqref{eq:graph-laplacian},
contradicting $\Delta_\Gamma h\geq0$ off $\partial$.
\end{proof}

\begin{corollary}
\label{cor:ch16-graph-mass-comparison}
With $\Gamma$, $\partial$, $f$, $g$ as above, if $f=g$ on $\partial$ and
$f\leq g$ on $\Gamma$, then
$\Delta_\Gamma g(\Gamma\setminus\partial)
 \leq\Delta_\Gamma f(\Gamma\setminus\partial)$.
\end{corollary}

\begin{proof}
Total mass vanishes on a compact graph
(\Cref{thm:ch11-graph-green}\,\cref{itm:ch11-graph:mass}), so the assertion
is equivalent to
$\Delta_\Gamma f(\partial)\leq\Delta_\Gamma g(\partial)$, that is, to
$\Delta_\Gamma(f-g)(\partial)\leq0$.  Since $h=f-g\leq0$ with $h=0$ on
$\partial$, every outgoing slope of $h$ at a point of $\partial$ is
$\leq0$, so the sum of outgoing slopes at each such point is $\leq0$.
\end{proof}

\subsection{What does hold in full CLD generality}
\label{sec:ch16-comparison-general}

Three statements survive without any convexity or boundary hypothesis.  All
three are consequences of results already proved, and they are exactly what
is used in \Cref{chap:cld-metrized-bundles,chap:cld-examples}.

\begin{proposition}
\label{prop:ch16-mass-monotone}
Let $T$ be positive closed of bidegree $(r,r)$, let $q+r\leq n$, and fix a
smooth test form $\omega$ of complementary bidegree with support in the
interior of a compact $K$.  There is a constant $C=C(K,\omega,q,T)$ such
that for all continuous psh $u,v$ on $X$,
\[
 \Bigl|\bigl\langle
   (\dprimeddprime u)^q\wedge T-(\dprimeddprime v)^q\wedge T,\ \omega
 \bigr\rangle\Bigr|
 \ \leq\ C\,\|u-v\|_{C^0(K)}\,
   \max\bigl\{\operatorname{osc}_K(u),\operatorname{osc}_K(v)\bigr\}^{q-1}.
\]
\end{proposition}

\begin{proof}
Write the difference as the telescoping sum
\[
 \sum_{j=0}^{q-1}(\dprimeddprime u)^j\wedge\dprimeddprime(u-v)\wedge
   (\dprimeddprime v)^{q-1-j}\wedge T ,
\]
each summand being a dpsh product in the sense of
\Cref{thm:ch16-dpsh-product}.  Since $X$ is boundaryless and every factor
other than $\dprimeddprime(u-v)$ is closed, the pairing formula of
\Cref{prop:current-differential-calculus} moves that one differential onto
the test form:
\[
 \begin{aligned}
 \bigl\langle(\dprimeddprime u)^j\wedge\dprimeddprime(u-v)\wedge
   (\dprimeddprime v)^{q-1-j}\wedge T,\ \omega\bigr\rangle
 &=\bigl\langle(u-v)\,(\dprimeddprime u)^j\wedge\\
 &\qquad
   (\dprimeddprime v)^{q-1-j}\wedge T,\ \dprimeddprime\omega\bigr\rangle .
 \end{aligned}
\]
Now apply \Cref{cor:ch16-cln-oscillation} with outer function $u-v$, the
fixed test form $\dprimeddprime\omega$, and the $q-1$ psh potentials
remaining; summing the $q$ terms gives the estimate.
\end{proof}

\begin{proposition}
\label{prop:ch16-bt-comparison}
Let $T$ be positive closed on an open $U$, and let $u,v$ be continuous psh
on $U$ such that $h=v-u$ is locally a uniform limit of smooth functions
$h_j$ with $\dprimeddprime h_j\geq0$.  Then
$\dprimeddprime(uT)\leq\dprimeddprime(vT)$ on $U$.  In particular
$\dprimeddprime[u]\leq\dprimeddprime[v]$, obtained by taking $T=\delta_U$.
\end{proposition}

\begin{proof}
By hypothesis $h_j\to h$ locally uniformly, so $h_jT\to hT$ by
\Cref{lem:continuous-times-positive-current}; and
$\dprimeddprime(h_jT)=\dprimeddprime h_j\wedge T\geq0$ by
\Cref{cor:ch16-smooth-case} together with
\Cref{prop:current-positivity-calculus}.  Weak continuity of the
differential and weak closedness of the positive cone give
$\dprimeddprime(hT)\geq0$.  Linearity of $w\mapsto wT$ gives
$\dprimeddprime(vT)-\dprimeddprime(uT)=\dprimeddprime(hT)\geq0$.  No
Bedford--Taylor hypothesis is needed.
\end{proof}

\begin{warning}{The approximation hypothesis is the whole content}
\Cref{prop:ch16-bt-comparison} converts an approximation hypothesis on
$v-u$ into a comparison of products.  The converse implication is not
available: knowing only $\dprimeddprime[u]\leq\dprimeddprime[v]$ tells one
that the continuous function $h=v-u$ has $\dprimeddprime[h]\geq0$ as a
current, and it is not known that such an $h$ is CLD-psh, still less that it
admits smooth psh approximants.  That gap is the ``classically psh implies
CLD-psh'' question recorded in \Cref{chap:cld-poincare-lelong}.  The same
caveat governs \Cref{thm:ch16-pluriharmonic-criterion} below, with the
inequality replaced by an equality.
\end{warning}

\section{Worked computations}
\label{sec:ch16-worked}

Every construction of this chapter is now tested on an explicit object.  The
computations are carried to numbers, and each number is checked against an
independent normalization---an intersection number, a lattice volume, or a
mass identity.

\subsection{A metric graph}
\label{sec:ch16-worked-graph}

\begin{example}[A metric graph]
\label{ex:ch16-metric-graph}
For $n=1$ and a continuous piecewise affine subharmonic function $u$ on a
metric graph,
\[
                   \MA(u)=\dprimeddprime u
                    =\sum_x\left(\sum_{\nu\text{ at }x}
                         \partial_\nu u(x)\right)\delta_x.
\]
Positivity is precisely nonnegativity of the outgoing-slope sums.  If $h$ is
harmonic, adding $h$ does not change this measure.  On a compact connected
graph, total mass is zero for a globally defined function, illustrating why
local positive curvature and global boundary conditions must be kept
separate.
\end{example}

\begin{example}[The tripod, with numbers]
\label{ex:ch16-tripod-numbers}
Let $\Gamma$ be the tripod of \Cref{ex:ch11-graph-tripod-check}: a central
vertex $p$ joined by three edges $e_1,e_2,e_3$ of length $1$ to leaves
$q_1,q_2,q_3$.  Let $u$ be the piecewise affine
function with $u(p)=0$ and $u(q_i)=a_i$, affine on each edge.  The outgoing
slope of $u$ at $p$ along $e_i$ is $a_i$, and at $q_i$ along $e_i$ it is
$-a_i$.  Hence, by \eqref{eq:graph-laplacian},
\[
 \MA(u)=\dprimeddprime u
 =\Bigl(\sum_ia_i\Bigr)\delta_p-\sum_ia_i\delta_{q_i}.
\]
For $(a_1,a_2,a_3)=(1,2,3)$ this is
$6\delta_p-\delta_{q_1}-2\delta_{q_2}-3\delta_{q_3}$, of total mass $0$, in
accordance with
\Cref{thm:ch11-graph-green}\,\cref{itm:ch11-graph:mass}.  The function $u$ is
psh on the open tripod $\Gamma\setminus\{q_1,q_2,q_3\}$, since the only
interior slope sum is $6\geq0$; it is not psh at the leaves, and indeed the
Laplacian is negative there.

Now take $T=\delta_\Gamma$ and compute $\MA_T$ against a second potential.
Since $\dim\Gamma=1$, the mixed product has only one factor and
$\MA_{\delta_\Gamma}(u)=\dprimeddprime u$; the recursion terminates
immediately.  The relative operator becomes interesting only after
multiplying by a positive closed current of bidegree $(0,0)$, that is, by a
locally constant nonnegative function.  If $c$ is such a function, then
$\dprimeddprime(u\,c\delta_\Gamma)=c\,\dprimeddprime u$ on each component,
which is the trivial case of \Cref{cor:ch16-bt-times-function}.  The point
of the example is negative: in dimension one the Bedford--Taylor recursion
has no content beyond the Laplacian, and every phenomenon of this chapter
that is genuinely nonlinear must be sought in dimension at least two.
\end{example}

\subsection{A two-dimensional fan}
\label{sec:ch16-worked-fan}

\begin{example}[The fan of $\mathbf P^2$]
\label{ex:ch16-fan-p2}
Regard $\mathbf R^2$ with the lattice $\mathbf Z^2$ as a tropical space of
level $2$, and set
\[
 g(x,y)=\max\{0,x,y\}.
\]
By \Cref{cor:ch16-max-psh}, $g$ is continuous psh, and it lies in
$\operatorname{CPsh}_{\mathbf R^2,1}$ by
\Cref{rem:ch16-collapse-euclidean}.  Its subdifferential at the origin is
the convex hull of the gradients of the three affine pieces,
\[
 \partial g(0)=\operatorname{conv}\{(0,0),(1,0),(0,1)\},
\]
of lattice area $\tfrac12$, and $\partial g(x)$ has measure zero for
$x\neq0$ because $g$ is affine near such an $x$.  Hence
$\MA_{\mathbf R}(g)=\tfrac12\delta_0$ and, by
\eqref{eq:ch16-alexandrov-normalization},
\begin{equation}
\label{eq:ch16-fan-p2-mass}
 \MA(g)=(\dprimeddprime g)^2=2!\cdot\tfrac12\,\delta_0=\delta_0,
 \qquad\text{total mass }1.
\end{equation}
The check: $g$ is the support function of the polytope
$\operatorname{conv}\{(0,0),(1,0),(0,1)\}$ of $\cO(1)$ on $\mathbf P^2$, and
$(\cO(1)^2)=1$; \eqref{eq:ch16-fan-p2-mass} agrees with
\Cref{thm:toric-ma-comparison}.

The same computation may be performed by the recursion, which is
instructive.  Put $u=\max\{0,x\}$ and $v=\max\{0,y\}$.  Then
$\dprimeddprime u=\delta_{\{x=0\}}$ and $\dprimeddprime v=\delta_{\{y=0\}}$
are the integration currents of the two coordinate lines with weight one, in
the sense of \Cref{def:ch11-integration-current}; each is closed because a
line with constant weight is balanced
(\Cref{thm:ch11-current-closed}).  The Bedford--Taylor product is
\[
 \dprimeddprime u\wedge\dprimeddprime v
 =\dprimeddprime\bigl(u\,\delta_{\{y=0\}}\bigr)
 =\dprimeddprime\bigl(\max\{0,x\}\big|_{\{y=0\}}\bigr)
 =\delta_{(0,0)},
\]
the middle equality because $u\delta_{\{y=0\}}$ is the current of the
function $x\mapsto\max\{0,x\}$ on the line $\{y=0\}$, whose Laplacian is the
unit atom at the origin.  Total mass $1$, which is the transverse tropical
intersection number of the two lines (\Cref{thm:tropical-stable-intersection}).

Polarization now recovers \eqref{eq:ch16-fan-p2-mass}: since
$\partial u(0)=[0,1]\times\{0\}$ and $\partial v(0)=\{0\}\times[0,1]$ have
zero area, $(\dprimeddprime u)^2=(\dprimeddprime v)^2=0$, while
$\partial(u+v)(0)=[0,1]^2$ has area $1$, so
$(\dprimeddprime(u+v))^2=2\delta_0$.  Expanding the square gives
$2\,\dprimeddprime u\wedge\dprimeddprime v=2\delta_0$, in agreement with the
recursion.  On the toric side $u+v$ is the support function of the unit
square, the polytope of $\cO(1,1)$ on $\mathbf P^1\times\mathbf P^1$, and
$(\cO(1,1)^2)=2$.
\end{example}

\begin{example}[A two-dimensional fan that is not a vector space]
\label{ex:ch16-fan-product}
Let $L\subset\mathbf R^2$ be the standard tropical line: the union of the
rays $\mathbf R_{\geq0}v_1$, $\mathbf R_{\geq0}v_2$,
$\mathbf R_{\geq0}v_3$ with $v_1=(1,0)$, $v_2=(0,1)$, $v_3=(-1,-1)$, all
weights $1$; it is balanced because $v_1+v_2+v_3=0$
(\Cref{ex:ch10-standard-line}).  Put
\[
 X=L\times\mathbf R\subset\mathbf R^3,
\]
a balanced weighted fan of pure dimension $2$, hence a tropical space of
level $2$.  Define, for $(z,t)\in L\times\mathbf R$,
\[
 u(z,t)=\max\{0,x(z),y(z)\},
 \qquad
 v(z,t)=|t| ,
\]
the restrictions to $X$ of ambient convex functions; both are continuous psh
by \Cref{prop:ch16-tropically-convex}.

\emph{First factor.}  On the ray $\mathbf R_{\geq0}v_1$ one has
$u(sv_1)=\max\{0,s,0\}=s$; on $\mathbf R_{\geq0}v_2$, $u(sv_2)=s$; on
$\mathbf R_{\geq0}v_3$, $u(sv_3)=\max\{0,-s,-s\}=0$.  The outgoing slopes at
the vertex are therefore $1,1,0$, and $u$ is affine on each ray.  Hence, on
$L$,
\[
 \dprimeddprime\bigl(u|_L\bigr)=(1+1+0)\,\delta_0=2\delta_0 ,
\]
so on $X$
\[
 \dprimeddprime u=2\,\delta_{\{0\}\times\mathbf R},
\]
the integration current of the vertical line with weight $2$.

\emph{Second factor.}  The function $|t|$ has outgoing slopes $1$ and $1$ at
$t=0$, so $\dprimeddprime v=2\,\delta_{L\times\{0\}}$.

\emph{Product.}  Since $X$ is a product and $u$, $v$ depend on separate
factors, the recursion gives
\[
 \dprimeddprime u\wedge\dprimeddprime v
 =\dprimeddprime\bigl(u\cdot2\delta_{L\times\{0\}}\bigr)
 =2\,\dprimeddprime\bigl(u|_{L\times\{0\}}\bigr)
 =2\cdot2\,\delta_{(0,0)}
 =4\,\delta_{(0,0)},
\]
of total mass $4$.  The check is the product formula for Laplacians: if
$a$ is psh on $L$ and $b$ is psh on $\mathbf R$, then on $L\times\mathbf R$
\[
 \dprimeddprime(a\oplus b)^{\wedge2}
 =2\,\bigl(\dprimeddprime a\bigr)\otimes\bigl(\dprimeddprime b\bigr),
\]
which for $a=u|_L$, $b=v$ gives $2\cdot(2\delta_0)\otimes(2\delta_0)
=8\delta_{(0,0)}$, that is $\MA(u+v)=8\delta_0$; and
$\MA(u)=\MA(v)=0$ because each of $u,v$ is affine along the line supporting
the corresponding current, so expanding $\MA(u+v)$ gives again
$\dprimeddprime u\wedge\dprimeddprime v=4\delta_0$.

The example is worth keeping because $X$ is not a Euclidean chart: the
vertex of $L$ is a point at which no local mollification is available, and
\Cref{rem:ch16-collapse-euclidean} does not apply to $X$.  Nevertheless both
potentials are restrictions of ambient convex functions, hence lie in
$\operatorname{CPsh}_{X,1}$ by
\Cref{prop:ch16-tropically-convex}, and the whole computation stays inside
the first stage of the filtration.
\end{example}

\begin{example}[Toric normalization]
\label{ex:ch16-toric-normalization}
On a torus chart let $u=F\circ\operatorname{trop}$ with $F$ finite convex.
For smooth strictly convex $F$,
\[
          (\dprimeddprime u)^n
          =n!\det(D^2F)\,\rd u
\]
on the calibrated skeleton.  Uniform approximation by smooth convex
functions and \Cref{thm:cld-bt-continuity} give
\[
          (\dprimeddprime u)^n
          =n!\operatorname{MA}_{\mathbf R}(F)
\]
for continuous convex $F$, with lattice-normalized Lebesgue measure.  This
agrees with the convention and computations of
\cref{chap:model-calculations}.  The factor $n!$
comes from expanding the labelled wedge, not from the definition of the
Alexandrov measure.
\end{example}

\subsection{A Chern--Levine--Nirenberg inequality verified}
\label{sec:ch16-worked-cln}

\begin{example}[The estimate, its constant, and its sharpness in dimension
one]
\label{ex:ch16-cln-verified}
Take $X=\mathbf R$ with its standard lattice, $T=\delta_X$, and $q=1$.  A
test form of complementary bidegree is a smooth compactly supported
function $\chi$, and for convex $u_1$ and continuous $u_0$,
\[
 \bigl\langle u_0\,\dprimeddprime u_1,\chi\bigr\rangle
 =\int_{\mathbf R}u_0\,\chi\,\rd\mu,
 \qquad \mu=\MA(u_1)=u_1''\,\rd t
\]
in the distributional sense.  Suppose $\chi\geq0$ and $\Supp\chi\subseteq K$.
The recipe of \Cref{prop:ch16-cln-constant} gives $\Omega_0=\chi$ and
$\Omega_1=(\chi'')^++(\chi'')^-=|\chi''|$, hence the constant
$C_{K,\chi}=\int_{\mathbf R}|\chi''|$.  The estimate asserts
\[
 \Bigl|\int u_0\chi\,\rd\mu\Bigr|
 \leq\|u_0\|_K\|u_1\|_K\int|\chi''| ,
\]
and this is verified directly: positivity of $\mu$ and of $\chi$ give
$|\int u_0\chi\,\rd\mu|\leq\|u_0\|_K\int\chi\,\rd\mu$; two integrations by
parts give $\int\chi\,\rd\mu=\int u_1\chi''$; and the last integral is at
most $\|u_1\|_K\int|\chi''|$.  Each step of the general proof is visible.

For a bound not involving a test function, take
$K=[-1,1]\subset L=[-2,2]$ and apply
\Cref{cor:ch16-cln-mass}.  Here the constant can be computed by hand:
for $u$ convex on $L$,
\[
 \MA(u)\bigl([-1,1]\bigr)=u'(1)-u'(-1)
 \leq\bigl(u(2)-u(1)\bigr)-\bigl(u(-1)-u(-2)\bigr)
 \leq2\operatorname{osc}_L(u),
\]
the first inequality by monotonicity of $u'$ and the mean value theorem on
$[1,2]$ and $[-2,-1]$.  Thus $C(K,L,1,\delta_X)=2$ is admissible.  It is not
optimal: for $u(t)=\tfrac12|t|$ one has $\operatorname{osc}_L(u)=1$ and
$\MA(u)([-1,1])=1$, so the best constant lies in $[1,2]$.  The example is a
reminder that the CLN constants of \Cref{sec:ch16-cln} are structural, not
sharp; no argument in this book uses their size.
\end{example}

\subsection{A product that fails to extend across a polar set}
\label{sec:ch16-worked-polar}

\begin{example}[A Bedford--Taylor product with two extensions]
\label{ex:ch16-polar-failure}
Let $K$ be algebraically closed and nontrivially valued, let
$X=\mathbf A^{2,\An}_K$ with coordinates $z_1,z_2$, and put
\[
 u=\log\max\{|z_1|,|z_2|\},
 \qquad
 Y=\{z_1=z_2=0\} .
\]
The pair $(z_1,z_2)$ is a regular sequence, $u$ is psh, and
$u=-\infty$ exactly on $Y$.  On $X\setminus Y$ the function $u$ is
real-valued and continuous, so the Bedford--Taylor product
$(\dprimeddprime u)^2$ is defined there by
\Cref{prop:ch16-iterate-basic}.  We claim that
\begin{equation}
\label{eq:ch16-polar-vanishes}
 (\dprimeddprime u)^2=0\qquad\text{on }X\setminus Y .
\end{equation}
Indeed, near a point with $z_1\neq0$ set $w=z_2/z_1$ and
$t_1=\log|z_1|$, $t_2=\log|w|$; these are the coordinates of a tropical
chart $f$, and
\[
 u=t_1+\max\{0,t_2\}=f^*\bigl(t_1\bigr)+p^*\bigl(\max\{0,\cdot\}\bigr),
\]
where $p$ is the moment map $-\log|w|$ with one-dimensional target.  The
first summand is parlinear, so it contributes nothing to
$\dprimeddprime u$; the second is the pullback of a convex function of one
variable along a map to $\mathbf R$, so its square vanishes by
\Cref{prop:ch16-vanishing-above-dimension}.  The region $z_2\neq0$ is
symmetric.  This proves \eqref{eq:ch16-polar-vanishes}.

The trivial extension of \eqref{eq:ch16-polar-vanishes} across $Y$ is the
zero current.  It is \emph{not} the canonical extension.  Indeed, by the
Bochner--Martinelli formula \Cref{thm:cld-local-bochner-martinelli} applied
with $d=2$, the canonical extension $\Lambda$ of $u\,\dprimeddprime u$
satisfies
\[
 \dprimeddprime\Lambda=\delta_Y ,
\]
a positive measure of mass $1$ on any neighbourhood of the origin.  The two
extensions differ by $\delta_Y$.

The mechanism is transparent in the truncations.  For $r>0$ put
$u_r=\log\max\{r,|z_1|,|z_2|\}$, a continuous PL psh function with
$u_r\downarrow u$.  As recalled in the proof of
\Cref{thm:cld-local-bochner-martinelli}, $(\dprimeddprime u_r)^2$ is the
integration current over the fibre above the Gauss point of radius $r$, of
mass $1$ for every $r$; letting $r\downarrow0$ gives $\delta_Y$.  Thus a
decreasing sequence of continuous psh functions has products of constant
mass $1$, while the limit function is not real-valued and the products of
the limit, computed off the polar set, have mass $0$.  There is no
contradiction with \Cref{thm:cld-bt-continuity}, whose hypothesis---local
uniform convergence to a \emph{continuous} function---fails here; and there
is no contradiction with \Cref{thm:cld-polar-negligibility}, which concerns
integrability of a psh function against a product of \emph{continuous} psh
potentials, not the definition of a product of unbounded ones.

Two morals.  First, ``the product of $\dprimeddprime u$ with itself'' is not
determined by the restriction of the product to the complement of the polar
locus, even when that complement has full measure; the canonical extension
is extra data supplied by a theorem, as
\Cref{chap:cld-poincare-lelong} already warned.  Second, the mass that
disappears is exactly the mass carried by $Y$, which is where the
nonpluripolar products of \Cref{part:bfj,part:bj} choose the other
convention: they discard it by construction.  The comparison is made precise
in \Cref{sec:ch16-local-vs-global}.
\end{example}

\section{Proper pushforward}
\label{sec:ch16-pushforward}

\begin{proposition}
\label{prop:bt-projection-formula}
Let $j:Y\to X$ be topologically proper, let $S$ be positive closed on $Y$,
and let $u_1,\ldots,u_q$ be continuous psh on $X$.  Then
\[
\begin{split}
 &\dprimeddprime u_1\wedge\cdots\wedge\dprimeddprime u_q\wedge j_*S\\
 &\qquad=j_*\bigl(
 \dprimeddprime(j^*u_1)\wedge\cdots\wedge
 \dprimeddprime(j^*u_q)\wedge S\bigr).
\end{split}
\]
\end{proposition}
\begin{proof}
For one factor, evaluate both sides on a test form and use the definition of
pushforward together with
$j^*(u\alpha)=j^*u\,j^*\alpha$.  The current differential commutes with
pushforward.  Induction gives the formula.  This is
\cite[Proposition~8.7.5]{CLD25}.
\end{proof}

\begin{proof}[Second proof of the one-factor case]
Fix a continuous psh $u$ on $X$ and consider
\[
 \mathscr P(U)=\bigl\{f\in C(U):\
   \dprimeddprime\bigl(f\,j_*S|_U\bigr)
   =j_*\dprimeddprime\bigl((j^*f)\,S|_{j^{-1}U}\bigr)\bigr\}.
\]
This is a sheaf of $\mathbf R$-vector spaces, since pushforward, restriction
and the differential are all local on the base and compatible with sums.
It is stable under local uniform limits: if $f_k\to f$ uniformly on
compacta then $j^*f_k\to j^*f$ uniformly on compacta of $Y$ meeting
$\Supp S$, because $j$ is topologically proper; now apply
\Cref{lem:continuous-times-positive-current} on $Y$ and weak continuity of
$j_*$ (\Cref{prop:current-pushforward-formal}).  Finally, let $f$ be smooth
psh and let $\omega$ be a smooth compactly supported test form on an open
$U$.  Then
\[
\begin{aligned}
 \langle\dprimeddprime(fj_*S),\omega\rangle
 &=\langle fj_*S,\dprimeddprime\omega\rangle
  =\langle j_*S,f\,\dprimeddprime\omega\rangle
  =\langle S,\ j^*f\ j^*\dprimeddprime\omega\rangle\\
 &=\langle (j^*f)S,\ \dprimeddprime j^*\omega\rangle
  =\langle\dprimeddprime\bigl((j^*f)S\bigr),j^*\omega\rangle\\
 &=\bigl\langle j_*\dprimeddprime\bigl((j^*f)S\bigr),\omega\bigr\rangle,
\end{aligned}
\]
using $j^*\dprimeddprime=\dprimeddprime j^*$ for smooth forms
(\Cref{prop:cld-form-pullback}).  Thus $\mathscr P$ satisfies the three
hypotheses of \Cref{lem:ch16-cpsh-induction}, whence
$\mathscr P\supseteq\operatorname{CPsh}_X$.  Iterating $q$ times, and using
that $j^*$ carries continuous psh functions to continuous psh functions
(\Cref{prop:cld-psh-calculus}), gives the stated formula.
\end{proof}

This formal formula and \cref{thm:cld-fibre-projection} are complementary.
The former pushes an
already defined current; the latter identifies the pushforward of a smooth
vertical form with actual integration over analytic fibres.

\begin{proposition}
\label{prop:ch16-degree-formula}
Let $X$ be a boundaryless good $n$-dimensional $K$-analytic space, let
$f:Y\to X$ be a topologically compact morphism from a boundaryless
$n$-dimensional space which is finite over a dense Zariski-open subset, and
let $\delta:X\to\mathbf R$ extend the local degree function of $f$.  For
continuous psh $u_1,\ldots,u_m$ that are locally approachable,
\[
 f_*\bigl(\dprimeddprime(f^*u_1)\wedge\cdots\wedge
    \dprimeddprime(f^*u_m)\bigr)
 =\delta\cdot\dprimeddprime u_1\wedge\cdots\wedge\dprimeddprime u_m .
\]
\end{proposition}

\begin{proof}
The left-hand side is a positive current, being a proper pushforward of one.
The right-hand side is the product of the positive current
$\dprimeddprime u_1\wedge\cdots\wedge\dprimeddprime u_m$ by the bounded
Borel function $\delta$; that this product is again a current requires that
the measure attached to the product does not charge the locus where
$\delta$ fails to be locally constant, which is a Zariski-closed set of
dimension $<n$.  That is exactly
\Cref{thm:cld-polar-negligibility}.  Off that locus $f$ is finite flat of
degree $\delta$, and the identity is the projection formula
\Cref{prop:bt-projection-formula} together with the multiplicity count of
\Cref{thm:trop-finite-pushforward}.  This is
\cite[Corollary~16.9.5]{CLD25}; the locally approachable hypothesis enters
only through the negligibility theorem.
\end{proof}

\section{Skeletal currents and their products}
\label{sec:ch16-skeletal}

A current is \emph{skeletal} if, locally in each chart, it is represented by
integration over a calibrated polyhedral subset.  Smooth CLD forms can be
multiplied with it directly, whereas nonsmooth potentials require the
Bedford--Taylor recursion.  This section computes the discrepancy exactly.
The outcome---\Cref{prop:ch16-skeletal-defect} and
\Cref{thm:ch16-skeletal-product}---is what makes the Monge--Amp\`ere measures
of \Cref{chap:cld-examples} computable by hand.

\subsection{Notation for a skeletal current}
\label{sec:ch16-skeletal-notation}

Let $T$ be a closed skeletal current of bidimension $(p,p)$ on $X$.  For a
tropical chart $f$ on an open $U$ we write $P^T_f\subseteq U$ for the
$p$-calibrated skeletal subset representing $T$ against $f$-tropical test
forms, in the sense of \Cref{chap:analytic-tropicalization}; thus
\[
 \langle T,\rho\wedge\omega\rangle=\int_{P^T_f}\rho\wedge\omega
\]
for $f$-tropical $\omega$ and smooth $\rho$ with $\rho\omega$ compactly
supported.  The symbol $\int^{\partial}$ denotes the boundary integral of
\Cref{chap:tropical-integration}, so that Stokes reads
$\int_C\dprime\eta=\int^\partial_C\eta$.

\begin{proposition}
\label{prop:ch16-skeletal-first-derivative}
Let $T$ be closed skeletal of bidimension $(p,p)$ and let $u$ be $G$-smooth
on $X$.  Then, for a tropical chart $f$ on $U$, an integer $r\leq p$, an
$f$-tropical smooth $\omega$ of type $(p-r,p)$ and a compactly supported
smooth $\rho$ of type $(r,0)$,
\[
 \bigl\langle\dprime(uT),\rho\wedge\omega\bigr\rangle
 =\int_{P^T_f}\dprime u\wedge\rho\wedge\omega .
\]
In particular $\dprime(uT)=T\wedge\dprime u=\dprime u\wedge T$.
\end{proposition}

\begin{proof}
By the boundary description of the differential of a skeletal current,
\[
 \bigl\langle\dprime(uT),\rho\wedge\omega\bigr\rangle
 =\int_{P^T_f}\dprime u\wedge\rho\wedge\omega
   -\int^{\partial}_{P^T_f}u\,\rho\wedge\omega .
\]
Closedness of $T$ says that
$\int^\partial_{P^T_f}v\,\rho\wedge\omega=\langle\dprime T,
 v\rho\wedge\omega\rangle=0$ for every \emph{smooth} $v$, and smooth
functions are dense in continuous ones for the topology in which the
boundary integral is continuous; hence the boundary term vanishes for $u$ as
well.  The last assertion is the case $r=0$.  This is
\cite[Proposition~8.10.1]{CLD25}.
\end{proof}

\subsection{The boundary defect of the second derivative}
\label{sec:ch16-skeletal-defect}

\begin{proposition}
\label{prop:ch16-skeletal-defect}
Let $T$ be closed skeletal of type $(p,p)$, let $f:X\to\mathbf R^m$ be a
tropical chart, let $v$ be $G$-smooth on $\mathbf R^m$ and $u=f^*v$.  For
every symmetric $f$-tropical smooth $\omega$ of type $(p-1,p-1)$ and every
smooth $\rho$ with $\rho\omega$ compactly supported,
\begin{equation}
\label{eq:ch16-skeletal-defect}
 \bigl\langle\dprimeddprime(uT),\rho\omega\bigr\rangle
 =\int_{P^T_f}\dprimeddprime u\wedge\rho\omega
   -\int^{\partial}_{P^T_f}\ddprime u\wedge\rho\omega .
\end{equation}
\end{proposition}

\begin{proof}
Choose a compact domain $W$ containing a neighbourhood of $\Supp(\rho\omega)$
and a $p$-calibrated skeletal $Q\supseteq P^T_f\cap W$ against which $T$ may
be integrated for all forms that are tropical for a common refinement of $f$
and of a chart tropicalizing $\rho$.  Then
\[
 \bigl\langle\dprimeddprime(uT),\rho\omega\bigr\rangle
 =\langle uT,\dprimeddprime(\rho\omega)\rangle
 =\int_Qu\,\dprimeddprime(\rho\omega).
\]
Green's formula on $Q$ (\Cref{thm:balanced-stokes}), legitimate because $u$
and $\rho\omega$ are symmetric, gives
\[
 \int_Qu\,\dprimeddprime(\rho\omega)
 =\int_Q\dprimeddprime u\wedge\rho\omega
  +\int^\partial_Q\bigl(u\,\ddprime(\rho\omega)-\ddprime u\wedge\rho\omega
  \bigr).
\]
The first term equals $\int_{P^T_f}\dprimeddprime u\wedge\rho\omega$, since
$\dprimeddprime u\wedge\omega$ is $f$-tropical of type $(p,p)$.  In the
boundary terms, $\int^\partial_Qu\,\ddprime(\rho\omega)$ vanishes: for
smooth $u_k\to u$ locally uniformly one has
$\int^\partial_Qu_k\ddprime(\rho\omega)
 =\langle\dprime T,u_k\ddprime(\rho\omega)\rangle=0$ by closedness, and one
passes to the limit.  The remaining boundary term may be computed on
$P^T_f$ instead of $Q$, by applying Stokes to
$\dprime(\ddprime u\wedge\rho\omega)$ on both and comparing the interior
integrals, all of which are pairings of $T$ with $f$-tropical forms of the
appropriate types.  This is \cite[Proposition~8.10.2]{CLD25}.
\end{proof}

The two terms of \eqref{eq:ch16-skeletal-defect} are the two halves of the
warning after \Cref{prop:ch16-skeletal-support}: the first is the naive
product of $T$ with the $G$-form $\dprimeddprime u$, and the second is the
defect.  We now compute the defect for a parlinear potential, where the
first term vanishes.

\begin{theorem}
\label{thm:ch16-skeletal-product}
Let $T$ be a closed skeletal current of bidimension $(p,p)$ and let $u$ be
parlinear on $X$.  Then $\dprimeddprime(uT)$ is skeletal, of bidimension
$(p-1,p-1)$.  Explicitly, choose a cellular decomposition $\mathscr C$ of
$P^T_f$ into $f$-cells on which $u$ is affine.  For a cell
$D\in\mathscr C_{p-1}$, a $(p-1)$-vector $\nu_D$ and an orientation
$\varepsilon_D$ of $f(D)$, and for each $p$-cell $C\supseteq D$ the outward
data $\lambda_C\in\mathbf R$ and $a_C=\partial_{e_C}(u|_C)$ determined by
$\mu_C=\lambda_C[\,e_C\wedge\nu_D,\varepsilon_C]$ as in
\Cref{chap:tropical-integration}, the current $\dprimeddprime(uT)$ is
integration over $\bigcup_{D}D$ with the $(p-1)$-calibration
\begin{equation}
\label{eq:ch16-skeletal-weights}
 \mu_D=\Bigl(\sum_{C\supseteq D}a_C\lambda_C\Bigr)
        f^*[\,\nu_D,\varepsilon_D\,].
\end{equation}
\end{theorem}

\begin{proof}
The assertion is local, so assume $X$ paracompact and $u$ $f$-tropical.
Since $u$ is parlinear, $\dprimeddprime u=0$ as a $G$-form, and
\eqref{eq:ch16-skeletal-defect} reduces to
\[
 \bigl\langle\dprimeddprime(uT),\rho\omega\bigr\rangle
 =-\int^\partial_{P^T_f}\ddprime u\wedge\rho\omega
 =-\sum_{C\in\mathscr C_p}\int^\partial_C\ddprime u\wedge\rho\omega .
\]
Each boundary integral is a sum of integrals over $(p-1)$-cells, so the
functional $\rho\mapsto\langle\dprimeddprime(uT),\rho\omega\rangle$ is a
measure supported on $\bigcup_{D\in\mathscr C_{p-1}}D$ and absolutely
continuous with respect to Lebesgue measure on each such cell.  It therefore
suffices to compute it after restricting to the relative interior of one
$D$, and we may assume $\Supp\rho$ meets no other $(p-1)$-cell.

Subtracting from $u$ an affine function agreeing with it on $D$ changes
nothing: an affine $f^*v_D$ is smooth with $\dprimeddprime(f^*v_D)=0$, so
$\dprimeddprime(f^*v_D\,T)=\dprimeddprime(f^*v_D)\wedge T=0$ by
\Cref{lem:ch16-leibniz-closed}.  Hence assume $u|_D=0$.

Work in coordinates adapted to a $p$-cell $C\supseteq D$: write
$D=\{0\}\times\mathbf R^{p-1}$ and $C=\mathbf R_{\leq0}\times\mathbf
R^{p-1}$, with $\nu_D=e_2\wedge\cdots\wedge e_p$,
$\nu_C=e_1\wedge\cdots\wedge e_p$, and $u(x)=a_Cx_1$ on $C$.  Decompose
\[
 \omega|_C=\omega_0+\dprime x_1\wedge\omega_1
   +\ddprime x_1\wedge\omega_2
   +\dprime x_1\wedge\ddprime x_1\wedge\omega_3
\]
with $\omega_i$ involving only $x_2,\ldots,x_p$.  Then
$\ddprime u|_C=a_C\ddprime x_1$, and contracting with
$\mu_C=\lambda_C[e_1\wedge\nu_D,\varepsilon_C]$ kills every term of
$\omega|_C$ except $\omega_0$, because $\ddprime x_1\wedge\ddprime x_1=0$
and because a term containing $\dprime x_1$ pairs to zero against a
$(p-1)$-vector on $D$.  Since $\omega|_D=\omega_0|_D$, one obtains
\[
 \int^\partial_C\langle\ddprime u\wedge\rho\omega,\mu_C\rangle
 =a_C\lambda_C\int_D\langle\rho\omega|_D,\nu_D\rangle .
\]
Summing over the $p$-cells $C\supseteq D$ and reinstating the global sign
gives \eqref{eq:ch16-skeletal-weights}.  Letting $D$ vary produces a
$(p-1)$-calibration on the union of the $(p-1)$-cells, which is exactly the
statement that $\dprimeddprime(uT)$ is skeletal.  This is
\cite[Theorem~8.10.3]{CLD25}.
\end{proof}

\begin{corollary}
\label{cor:ch16-skeletal-iterated}
Let $T$ be closed skeletal and let $u_1,\ldots,u_m$ be parlinear.  Then
$\dprimeddprime u_1\wedge\cdots\wedge\dprimeddprime u_m\wedge T$ is
parlinear, in particular skeletal.
\end{corollary}

\begin{proof}
Induction on $m$ using \Cref{thm:ch16-skeletal-product} and the definition
of the iterated product.  This is \cite[Corollary~8.10.5]{CLD25}.
\end{proof}

\begin{remark}[The $G$-smooth case]
\label{rem:ch16-skeletal-gsmooth}
For $u$ merely $G$-smooth the same argument applies with $a_C$ replaced by
the partial derivative $\partial u/\partial x_1$, which is a smooth function
on each cell rather than a constant.  The conclusion is that
$\dprimeddprime(uT)$ is the sum of the parlinear current $T\wedge
\dprimeddprime u$ and a current supported on the codimension-one skeleton
carrying a \emph{$G$-smooth calibration}.  This is
\cite[Remark~8.10.4]{CLD25}, and it is the precise form of the discrepancy
warned about after \Cref{prop:ch16-skeletal-support}.
\end{remark}

\begin{example}[Reading \eqref{eq:ch16-skeletal-weights} on a graph]
\label{ex:ch16-skeletal-weights-graph}
Take $p=1$, $X$ a metric graph with unit weights, $T=\delta_X$, and $u$
piecewise affine.  A $0$-cell $D$ is a point $x$; the $1$-cells $C\supseteq
D$ are the edges at $x$; $\lambda_C=1$ for the canonical calibration; and
$a_C=\partial_{e_C}(u|_C)$ is the outgoing slope of $u$ along that edge.
Formula \eqref{eq:ch16-skeletal-weights} becomes
$\mu_x=\sum_{C\ni x}\partial_{e_C}u$, which is
\eqref{eq:graph-laplacian}.  The sign bookkeeping of
\Cref{lem:ch11-graph-trace} is exactly the passage from
$\int^\partial$ to outgoing slopes.
\end{example}

\section{Harmonic charts and toric products}
\label{sec:ch16-harmonic-charts}

\subsection{Harmonic functions and harmonic tropicalizations}
\label{sec:ch16-harmonic-definitions}

\begin{definition}[Harmonic parlinear function]
\label{def:ch16-harmonic-function}
\index{harmonic function!CLD}
A parlinear function $u$ on $X$ is \emph{harmonic} if $\dprimeddprime[u]=0$
as a current.  A function is harmonic if it is locally a real linear
combination of harmonic parlinear functions.  A $G$-tropical chart
$u=(u_1,\ldots,u_d):X\to\mathbf R^d$ is a \emph{harmonic chart} if every
$u_i$ is harmonic.  A $G$-form $\alpha$ is \emph{locally harmonically
tropicalizable} if locally $\alpha=u^*\beta$ for a harmonic parlinear
$u:X\to\mathbf R^d$ and a smooth $\beta$ near a polyhedron containing
$u(X)$.
\end{definition}

Harmonic functions form a sheaf of $\mathbf R$-vector spaces containing the
affine functions.  The definition is not the naive one: as CLD stress in
\cite[Remark~7.8.3]{CLD25}, there is no reason in general for a continuous
$u$ with $\dprimeddprime[u]=0$ to be harmonic in the above sense, and it is
not even known whether local integrability plus $\dprimeddprime[u]=0$ forces
continuity.  \Cref{sec:ch16-harmonic-ambiguity} makes the gap explicit.

The point of harmonicity is a Stokes statement: for a harmonically
tropicalizable $G$-form $\alpha$, the boundary current $[\alpha]^\partial$
vanishes, so that $\dprime[\alpha]=[\dprime\alpha]$ and
$\ddprime[\alpha]=[\ddprime\alpha]$.  This is
\cite[Theorem~7.8.5 and Corollary~7.8.8]{CLD25}, the third of the imported
inputs listed in \Cref{sec:ch16-imported}, and it is proved in the
delta-form chapter \Cref{sec:harmonic-trop} in its Gubler--Jell--Rabinoff
form.

\subsection{The compatibility theorem}
\label{sec:ch16-harmonic-compatibility}

\begin{theorem}
\label{thm:bt-harmonic-compatibility}
Let $T$ be a skeletal Bedford--Taylor current and let $\alpha$ be a $G$-form
locally tropicalized by harmonic coordinates.  Then
\[
 \dprime(T\wedge\alpha)=T\wedge\dprime\alpha,
 \qquad
 \ddprime(T\wedge\alpha)=T\wedge\ddprime\alpha.
\]
Consequently, if $T$ is locally a difference of skeletal Bedford--Taylor
currents and $u_1,\ldots,u_q$ are $G$-smooth functions locally tropicalized
by harmonic coordinates, their recursive Bedford--Taylor product equals the
direct product with the $G$-form
$\bigwedge_i\dprimeddprime u_i$.
\end{theorem}
\begin{proof}
Work in a common harmonic tropical chart and subdivide so that the skeletal
support of $T$ is a calibrated weighted complex $C$.  By the definition of
the differential of a current, the difference
$\dprime(T\wedge\alpha)-T\wedge\dprime\alpha$, evaluated on a test form, is
the sum of the boundary terms obtained by applying Stokes on the facets of
$C$.  At a codimension-one face $\tau$, this term is contraction with
\[
             \sum_{\sigma\supset\tau}
                    m_\sigma u_{\sigma/\tau}.
\]
The Bedford--Taylor property of the skeletal current supplies the weighted
balancing relation, while harmonicity says that every harmonic coordinate
has zero total normal slope.  Hence the displayed class vanishes after
pairing with every coefficient of $\alpha$.  There is no residual boundary
current, which proves the $\dprime$ identity.  Applying the involution $J$
gives the $\ddprime$ identity.

For the consequence, start with $T$ and define the recursive product one
factor at a time.  The two identities just proved allow each differential
to pass through $T$ and through the previously inserted closed factors.
Thus inserting $u_i$ by the Bedford--Taylor recursion has the same effect as
wedging with the harmonic-chart form $\dprimeddprime u_i$.  Induction on $q$
gives the stated direct product.  Differences are handled by linearity.
This calibrated Stokes computation is
\cite[Proposition~8.11.1 and Corollary~8.11.2]{CLD25}.
\end{proof}

The heart of that proof is a single identity, which we isolate because
\Cref{sec:ch16-harmonic-ambiguity} shows that it can fail when the
hypotheses are weakened.

\begin{lemma}
\label{lem:ch16-normal-slope-vanishes}
Let $T$ be a skeletal Bedford--Taylor current of type $(p,p)$, let $v$ be a
harmonic parlinear function on $X$, and fix a cellular decomposition
$\mathscr C$ of a compact skeletal neighbourhood $\Sigma\subseteq P^T_f$ on
which $v$ is affine on each cell.  Let $D\in\mathscr C_{p-1}$, and for each
$p$-cell $C\supseteq D$ let $e_C$ be the primitive vector pointing from $D$
into $C$ and normalized so that the canonical calibration of $C$ is
$[\,e_C\wedge\mu_D,\varepsilon_C]$.  Then
\begin{equation}
\label{eq:ch16-normal-slope-sum}
 \sum_{C\supseteq D}\partial_{e_C}\bigl(v|_C\bigr)=0 .
\end{equation}
\end{lemma}

\begin{proof}
Harmonicity of $v$ and the Bedford--Taylor property give
$\dprimeddprime(vT)=\dprimeddprime[v]\wedge T=0$ on the interior of the
chart domain.  Since $v$ is parlinear, $\dprimeddprime v=0$ as a $G$-form,
so \eqref{eq:ch16-skeletal-defect} reduces to
\[
 0=\bigl\langle\dprimeddprime(vT),\rho\omega\bigr\rangle
 =-\int^\partial_\Sigma\ddprime v\wedge\rho\omega
\]
for every $f$-tropical symmetric $\omega$ of type $(p-1,p-1)$ and every
smooth $\rho$ with $\rho\omega$ compactly supported in the chart interior.
Fix $x$ in the relative interior of $D$ and choose $\omega=f^*q^*\gamma$,
where $q$ is a torus morphism for which $(q\circ f)_{\mathrm{trop}}$ is a
piecewise immersion of a neighbourhood of $x$ in $D$ into $\mathbf R^{p-1}$,
and take $\rho$ supported so as to meet no $p$-cell other than those
containing $D$.  The computation of
\Cref{thm:ch16-skeletal-product}, performed with $a_C=\partial_{e_C}(v|_C)$,
turns the displayed vanishing into
\[
 \int_D\Bigl(\sum_{C\supseteq D}\partial_{e_C}(v|_C)\Bigr)
   \rho\,\langle f^*q^*\gamma,\mu_D\rangle=0 .
\]
As $\rho$ and $\gamma$ are arbitrary and $q_*f_*\mu_D\neq0$, the bracket
vanishes, which is \eqref{eq:ch16-normal-slope-sum}.
\end{proof}

\begin{proof}[Completion of the proof of
\Cref{thm:bt-harmonic-compatibility}]
With $\alpha=u^*\beta$ for a harmonic parlinear $u:W\to\mathbf R^d$, Stokes
on the cells of $\mathscr C$ gives
\[
 \int_{P^T_f}\dprime(\rho\alpha)
 =\sum_{D\in\mathscr C_{p-1}}\int_D
   \Bigl(\sum_{C\in\mathscr C_D}\langle(\rho u^*\beta)_C,\mu_C\rangle\Bigr),
\]
where $\mathscr C_D$ is the set of $p$-cells $C\supseteq D$ on which
$u|_C$ is injective; cells on which $u|_C$ is not injective contribute zero
because $u^*\beta$ vanishes there.  Contracting and pushing forward through
$u$ turns the inner sum into
$\rho\circ u_D^{-1}\bigl\langle\beta,\sum_{C\in\mathscr C_D}u_*\mu_C
\bigr\rangle$, so it suffices to prove
$\sum_{C\in\mathscr C_D}u_*\mu_C=0$.  Since $\mu_C=e_C\wedge\mu_D$, one has
$u_*\mu_C=u_*(e_C)\wedge u_*(\mu_D)$ with $u_*(\mu_D)$ independent of $C$;
and the coordinates of $u_*(e_C)$ are the numbers
$\partial_{e_C}(v|_C)$ as $v$ runs through the components of $u$.  By
\Cref{lem:ch16-normal-slope-vanishes} each coordinate of
$\sum_{C\supseteq D}u_*(e_C)$ vanishes, whence
$\sum_{C\in\mathscr C_D}u_*\mu_C=0$.
\end{proof}

\begin{corollary}
\label{cor:ch16-harmonic-convex-positive}
Let $f:U\to\mathbf R^r$ be a harmonic chart and let $F$ be continuous convex.
Then $u=F\circ f$ is locally a difference of psh functions and
$\dprimeddprime u\wedge T$ is positive for every positive closed $T$ for
which the product is defined.
\end{corollary}
\begin{proof}
On every relatively compact convex subset of the tropical target, choose
smooth convex functions $F_j$ converging uniformly to $F$.  The Hessian of
$F_j$ is positive semidefinite, so
$\dprimeddprime(F_j\circ f)$ is a positive harmonic-chart form.  By the
first part of \Cref{thm:bt-harmonic-compatibility}, its wedge with the
positive closed current $T$ is positive.  Bedford--Taylor continuity passes
to the uniform limit and gives positivity of
$\dprimeddprime(F\circ f)\wedge T$.

The same local regularization writes $F\circ f$ as a locally approachable
psh function; in particular it is a difference of psh functions, trivially
by subtracting the psh constant zero.  If one works with a signed harmonic
chart presentation, decompose it into the difference of the two convex
presentations supplied by the chart.  This is
\cite[Corollary~8.11.3]{CLD25}.  The conclusion concerns convex functions in
harmonic charts and does not identify arbitrary classically psh functions
with the CLD approximation sheaf.
\end{proof}

\section{Harmonic ambiguity}
\label{sec:ch16-harmonic-ambiguity}

\Cref{def:ch16-bt-current} isolated the property that makes
$\dprimeddprime u\wedge T$ depend only on $\dprimeddprime u$.  This section
analyses that property.  It contains one positive result---a sufficient
criterion in terms of smooth pluriharmonic approximation---one regularity
statement on cell interiors, and one counterexample showing that the two do
not combine into a general theorem.  All three are recorded in the audit
ledger of \Cref{chap:verified-extensions}, which is where the book's
provenance for them is kept; the proofs are given here because this is where
they are used.

\subsection{A sufficient criterion}
\label{sec:ch16-pluriharmonic-criterion}

\begin{theorem}
\label{thm:ch16-pluriharmonic-criterion}
Let $U$ be open in $X$, let $T$ be a positive closed current on $U$, and let
$u,v$ be continuous psh on $U$ with
\begin{equation}
\label{eq:ch16-equal-hessians}
 \dprimeddprime[u]=\dprimeddprime[v].
\end{equation}
Put $h=u-v$ and assume that there are smooth functions $h_j$ on $U$ with
\begin{equation}
\label{eq:ch16-pluriharmonic-approximation}
 h_j\longrightarrow h\ \text{locally uniformly},
 \qquad
 \dprimeddprime h_j=0 .
\end{equation}
Then $\dprimeddprime(hT)=0$; equivalently
$\dprimeddprime(uT)=\dprimeddprime(vT)$.
\end{theorem}

\begin{proof}
Let $\alpha$ be a smooth compactly supported test form of complementary
bidegree, with support in a compact $K\Subset U$.  Positivity gives
$T$ finite mass on $K$ in the sense of \eqref{eq:ch16-positivity-sandwich},
so
\[
 \bigl|\langle(h_j-h)T,\alpha\rangle\bigr|
 \leq\|h_j-h\|_{C^0(K)}\,\langle T,\beta\rangle
\]
with $\beta$ as in \eqref{eq:ch16-positivity-sandwich}.  The right-hand side
tends to zero, so $h_jT\to hT$ as currents.  Differentiation is weakly
continuous (\Cref{prop:current-differential-calculus}), and for each $j$,
$h_j$ being smooth and $T$ closed,
\Cref{lem:ch16-leibniz-closed} gives
$\dprimeddprime(h_jT)=\dprimeddprime h_j\wedge T=0$.  Passing to the limit,
$\dprimeddprime(hT)=0$.  Linearity of $w\mapsto wT$ turns this into
$\dprimeddprime(uT)=\dprimeddprime(vT)$.  This is
\Cref{thm:verified-cld-criterion}.
\end{proof}

\begin{corollary}
\label{cor:ch16-smooth-h}
If $h=u-v$ is itself smooth and $\dprimeddprime h=0$, then
$\dprimeddprime(hT)=0$.
\end{corollary}

\begin{proof}
Take $h_j=h$.  This is \Cref{cor:verified-cld-smooth}, and it recovers the
computation of \Cref{sec:ch16-ambiguity-smooth}.
\end{proof}

\begin{warning}{The criterion is conditional}
\Cref{thm:ch16-pluriharmonic-criterion} does not produce the approximants
$h_j$; \eqref{eq:ch16-equal-hessians} alone does not supply them.  Whether
every continuous $h$ with $\dprimeddprime[h]=0$ admits smooth
pluriharmonic approximants is exactly the question left open in
\Cref{prob:frontier-cld-hierarchy}'s neighbourhood and recorded in
\Cref{rem:verified-cld-sufficient}.  The criterion is therefore a tool for
verifying the Bedford--Taylor property in situations where the approximants
are constructed by hand---harmonic charts, formal models, curve
skeleta---and not a proof that the property always holds.
\end{warning}

\subsection{Cell-interior regularity}
\label{sec:ch16-cell-interior}

On the interior of a cell there is no ambiguity at all: harmonicity in the
weakest sense already forces affinity.

\begin{proposition}
\label{prop:ch16-cell-rigidity}
Let $\sigma$ be a tropical cell of dimension $q$, let $\sigma^\circ$ be its
relative interior with affine coordinates $x_1,\ldots,x_q$, and let
$w\in L^1_{\loc}(\sigma^\circ)$ satisfy $\dprimeddprime[w]=0$ as a current
on $\sigma^\circ$.  Then $w$ agrees almost everywhere with an affine
function, and that affine function is its unique continuous representative.
\end{proposition}

\begin{proof}
In the affine coordinates the coefficients of $\dprimeddprime[w]$ are the
distributional second derivatives $\partial_i\partial_jw$, so the hypothesis
reads $\partial_i\partial_jw=0$ for all $i,j$.  Hence each first
distributional derivative $\partial_iw$ has vanishing distributional
gradient on the connected open set $\sigma^\circ$ and is therefore a
constant $a_i$.  The distributional gradient of $w-\sum_ia_ix_i$ vanishes,
so this difference is almost everywhere equal to a constant.  The resulting
affine function is continuous, and two continuous functions agreeing almost
everywhere on $\sigma^\circ$ agree everywhere.  This is
\Cref{prop:verified-cell-rigidity}.
\end{proof}

Thus, on a cell interior, the three conditions of the table in
\Cref{sec:ch16-reading-bt} collapse: $w$ is affine, hence parlinear, hence
smooth, and \Cref{cor:ch16-smooth-h} applies.  The failure must therefore be
sought at a face where several cells meet, and this is precisely what the
next example does.

\subsection{The cross-fan obstruction}
\label{sec:ch16-cross-fan}

\begin{example}[The cross-fan: cellwise harmonicity is not enough]
\label{ex:ch16-cross-fan}
Work on the four-ray fan
\begin{equation}
\label{eq:ch16-cross-fan}
 \Sigma=\mathbf R_{\geq0}e_1\cup\mathbf R_{\geq0}(-e_1)
   \cup\mathbf R_{\geq0}e_2\cup\mathbf R_{\geq0}(-e_2)
   \ \subset\ \mathbf R^2,
\end{equation}
regarded as an abstract calibrated tropical fan of pure dimension one, and
put
\begin{equation}
\label{eq:ch16-cross-functions}
 u=2|x_1|,\qquad \gamma=|x_1|+|x_2|,\qquad h=u-\gamma .
\end{equation}

\emph{Step 1: both potentials are psh for unit weights.}  Give every ray
weight $1$; this fan is balanced, since opposite rays cancel.  Set
$u_j(x)=2\sqrt{x_1^2+\varepsilon_j}$ and
$\gamma_j(x)=\sqrt{x_1^2+\varepsilon_j}+\sqrt{x_2^2+\varepsilon_j}$ for
$\varepsilon_j\downarrow0$.  These are smooth convex on $\mathbf R^2$, so
their restrictions to $\Sigma$ are smooth psh, and they decrease locally
uniformly to $u$ and $\gamma$.  Hence $u,\gamma$ are continuous psh, indeed
locally approachable.

\emph{Step 2: the two Laplacians agree.}  The outgoing slopes at the origin
are
\begin{equation}
\label{eq:ch16-cross-slopes}
\begin{array}{c|cccc|c}
 &+e_1&-e_1&+e_2&-e_2&\text{sum}\\ \hline
u&2&2&0&0&4\\
\gamma&1&1&1&1&4\\
h&1&1&-1&-1&0
\end{array}
\end{equation}
and all three functions are affine on the interior of each ray.  With unit
weights, \eqref{eq:graph-laplacian} gives
$\dprimeddprime[u]=\dprimeddprime[\gamma]=4\delta_0$ and
$\dprimeddprime[h]=0$.  In particular $h$ is harmonic on each cell interior
in the sense of \Cref{prop:ch16-cell-rigidity}---it is affine there---and its
unweighted balance at the origin vanishes.

\emph{Step 3: a different weighting.}  Now give each horizontal ray weight
$2$ and each vertical ray weight $1$.  Opposite rays still carry equal
weights, so the weighted fan is balanced; its weights are positive, so the
associated skeletal current $T$ is positive and closed.  For the current
$hT$, the outgoing defect at the origin is the \emph{weighted} sum of the
slopes of $h$:
\begin{equation}
\label{eq:ch16-cross-weighted-sum}
 2\cdot1+2\cdot1+1\cdot(-1)+1\cdot(-1)=2 .
\end{equation}
By \Cref{thm:ch16-skeletal-product}, with $p=1$ and the weights of $T$
entering through $\lambda_C$, this is the total mass of
$\dprimeddprime(hT)$, and there is no other contribution because $h$ is
affine on each ray interior.  Hence
\begin{equation}
\label{eq:ch16-cross-defect}
 \dprimeddprime(hT)=2\delta_0\neq0 ,
\end{equation}
even though $\dprimeddprime[u]=\dprimeddprime[\gamma]$.  Consequently $T$ is
\emph{not} a Bedford--Taylor current, and the recursive products
$\dprimeddprime(uT)$ and $\dprimeddprime(\gamma T)$ differ by $2\delta_0$.
This is \Cref{prop:verified-cross-fan}; the failed implication is analysed in
\Cref{rem:verified-cross-failed}.
\end{example}

\begin{remark}[Exactly which implication fails]
\label{rem:ch16-cross-reading}
Nothing in \Cref{ex:ch16-cross-fan} contradicts any theorem above.
\Cref{prop:ch16-cell-rigidity} remains true on each ray, and it is used.
What fails is the passage from
\begin{center}
``$h$ is affine on every cell and the \emph{unweighted} balance vanishes''
\end{center}
to
\begin{center}
``the balance weighted by the calibration of $T$ vanishes''.
\end{center}
The two balances agree when all weights are equal, and diverge as soon as
the weighting of $T$ differs from the fundamental calibration.
\Cref{lem:ch16-normal-slope-vanishes} shows that the second balance is
exactly what the Bedford--Taylor property supplies, and
\Cref{thm:ch16-pluriharmonic-criterion} shows that smooth pluriharmonic
approximants supply it too, because a smooth $h_j$ has
$\dprimeddprime h_j=0$ as a \emph{form}, an identity that is insensitive to
weights.  In the example no such approximants exist: a smooth function on
$\Sigma$ is the restriction of a smooth function on $\mathbf R^2$, and its
slope vector at the origin along $\pm e_i$ is $\pm\langle\rd F(0),e_i\rangle$,
so its slope sums are automatically zero for \emph{every} weighting, whereas
$h$ has weighted slope sum $2$.  Thus $h$ is a limit of smooth functions in
no way compatible with $\dprimeddprime h_j=0$, which is precisely the
hypothesis \eqref{eq:ch16-pluriharmonic-approximation} that fails.
\end{remark}

\begin{warning}{The analytic realization is only after pushforward}
\Cref{ex:ch16-cross-fan} is itself a counterexample on an abstract
calibrated tropical fan.  The weighted fan and its positive closed current
do admit the noncompact analytic realization of
\Cref{prop:verified-cross-analytic-realization}: a proper moment
tropicalization has the full cross as image, is exhausted by compact
harmonic-tropicalization charts, and carries an analytic cycle with weights
\((2,2,1,1)\).  The full cross cannot itself be the image of one harmonic
tropicalization in the definition used here, because that definition has a
compact source; see \Cref{prop:verified-cross-compact-obstruction}.

This still does \emph{not} exhibit a positive closed analytic current that
fails the Bedford--Taylor property.  Upstairs, the two pulled-back
potentials have the distinct curvature measures displayed in
\eqref{eq:verified-cross-upstairs-curvatures}; they become equal only after
proper tropical pushforward identifies two different Gauss points.
Accordingly CLD's remark that every current they construct is
Bedford--Taylor is not contradicted.  The example and its realization show
that cellwise harmonicity and balancing alone do not prove the property;
one needs the further compatibility isolated in
\Cref{thm:bt-harmonic-compatibility}.
\end{warning}

\subsection{The three levels of ambiguity}
\label{sec:ch16-ambiguity-levels}

\begin{center}
\small
\begin{tabular}{p{.24\textwidth}|p{.34\textwidth}|p{.33\textwidth}}
\textbf{Level}&\textbf{Ambiguity is resolved by}&\textbf{Reference}\\
\hline
CLD Bedford--Taylor currents
 &the defining independence from a potential with fixed
  $\dprimeddprime$
 &\Cref{def:ch16-bt-current},
  \Cref{prop:bt-current-stability}\\
harmonic tropicalizations
 &cancellation of normal slopes across calibrated faces
 &\Cref{lem:ch16-normal-slope-vanishes},
  \Cref{thm:bt-harmonic-compatibility}\\
global energy theories
 &a normalization removing constants, and sometimes
  model-harmonic functions
 &\Cref{chap:bfj-energy,chap:bj-energy}\\
\end{tabular}
\end{center}

The comparison chapters state the hypotheses under which these quotients
agree; \Cref{chap:master-dictionary} collects them.

\section{Polar sets}
\label{sec:ch16-polar}

Call a psh function \emph{locally approachable} if locally it is a decreasing
limit of smooth psh functions; for a continuous such function the convergence
is locally uniform.  This is a first-stage approximation hypothesis, stronger
than merely belonging abstractly to the transfinite CLD sheaf.

In the notation of \Cref{def:ch16-psh-filtration}, locally approachable means
being a section of $\operatorname{Psh}_{X,1}$, which is
\cite[Definition~8.1.13]{CLD25}.  The hypothesis is used twice in the proof
below, and in both places it is used for what it literally says: to produce
a decreasing sequence of \emph{smooth} psh functions to which
\Cref{thm:ch16-weighted-cln} applies.  Since the collapse of the filtration
at stage one is open (\Cref{prop:ch16-hierarchy-status}), the hypothesis
cannot at present be removed.

\begin{theorem}
\label{thm:cld-polar-negligibility}
Let $X$ be a good, pure $n$-dimensional boundaryless analytic space.  Let
$u_1,\ldots,u_m$ be continuous psh functions that are locally approachable,
let $V$ be locally approachable psh and locally integrable, and let $\alpha$
be a smooth $(n-m,n-m)$-form.  Then $V$ is locally integrable for the signed
measure
\[
          \dprimeddprime u_1\wedge\cdots\wedge
          \dprimeddprime u_m\wedge\alpha.
\]
If $\alpha$ is positive, this measure does not charge the polar locus
$\{V=-\infty\}$.  In particular it charges no analytic subset of dimension
$<n$.
\end{theorem}
\begin{proof}
The assertion is local.  Choose a rational affinoid neighborhood $W$ of a
point and a smaller neighborhood $W'\Subset W$.  Finitely many functions
defining $W'$ inside $W$, followed by a smooth regularized maximum of their
logarithmic absolute values, give a smooth psh barrier $\rho$ whose relevant
sublevel set is relatively compact in $W$.  This is the barrier construction
of \cite[Proposition~16.9.1]{CLD25}.

Choose decreasing smooth psh approximants $V_k\downarrow V$ and uniformly
convergent smooth psh approximants $u_{i,k}\to u_i$ on $W$.  Insert cutoff
functions obtained from $\rho$.  Repeated integration by parts and positivity
give the weighted Chern--Levine--Nirenberg estimate
\begin{equation}\label{eq:cld-weighted-cln}
 \int_{W'}(-V_k)
   \dprimeddprime u_{1,k}\wedge\cdots\wedge
   \dprimeddprime u_{m,k}\wedge\alpha
 \le C,
\end{equation}
after adding a harmless constant so that $V_k\le0$.  The constant depends on
the compact inclusion, the barrier, the local sup norms of the $u_i$, and
$\alpha$, but not on $k$.  To see the independence, move one
$\dprimeddprime$ at a time from $u_{i,k}$ to the fixed cutoff; the boundary
terms vanish because the cutoff has compact support, and each remaining term
is bounded by positivity and the fixed oscillation of $u_{i,k}$.  At the last
step the same operation moves $\dprimeddprime$ from $V_k$ to the barrier.

Uniform Bedford--Taylor continuity sends the mixed measures in
\eqref{eq:cld-weighted-cln} to
$\dprimeddprime u_1\wedge\cdots\wedge\dprimeddprime u_m\wedge\alpha$.
Fatou's lemma and then monotone convergence give local integrability of
$-V$, hence of $V$, for the resulting signed measure; decompose a general
smooth $\alpha$ into a difference of positive forms on the fixed compact
chart.  If $\alpha$ is positive, the measure is positive.  Since an
integrable extended-real function is finite almost everywhere, it assigns
zero mass to $\{V=-\infty\}$.

Finally, let $Z$ be an analytic subset of dimension less than $n$.  Locally
choose an analytic function $f$ vanishing on $Z$ but on no top-dimensional
component.  The functions
$\max\{\log|f|,-k\}$ admit decreasing smooth psh regularizations, so
$\log|f|$ is locally approachable, locally integrable, and equals
$-\infty$ on $Z$.  The preceding paragraph shows that the mixed measure does
not charge $Z$.  The estimate and limiting argument are
\cite[Theorem~16.9.2 and Corollaries~16.9.3--16.9.4]{CLD25}.
\end{proof}

The barrier is \Cref{thm:ch16-barrier}, and the weighted estimate is
\Cref{thm:ch16-weighted-cln}.  We record the two limits explicitly, since
the order in which they are taken is the only delicate point.

\begin{proof}[Second proof of the integrability statement, for
$m=n$ and $\alpha=1$]
\leavevmode\par\noindent
By \Cref{thm:ch16-barrier} we may assume that there is a smooth psh $\rho$
on $X$ with $U=\{\rho<0\}$ a relatively compact neighbourhood of the given
point; fix $a=3$ and let $\beta_0,\ldots,\beta_n$ be as in
\Cref{def:ch16-weight-forms}.  Local approachability gives smooth psh
$u_{i,\ell}\to u_i$ locally uniformly and smooth psh $V_p\downarrow V$; after
subtracting constants, $V_p\leq0$.  \Cref{thm:ch16-weighted-cln} gives, for
all $\ell,p$,
\[
 \int_X\beta_n\wedge(-V_p)\,\dprimeddprime u_{1,\ell}\wedge\cdots
   \wedge\dprimeddprime u_{n,\ell}
 \ \leq\ C(n,3)\,\|u_{1,\ell}\|_U\cdots\|u_{n,\ell}\|_U
   \int_U(-V_p)\beta_0 ,
\]
the extension of the integral from $U$ to $X$ being harmless because
$\beta_n$ vanishes off $U$.  Let $\ell\to\infty$ with $p$ fixed: the
left-hand side converges by \Cref{thm:cld-bt-continuity}, the outer factor
$\beta_n\wedge(-V_p)$ being a fixed continuous compactly supported weight,
and the sup-norms on the right converge to $\|u_i\|_U$.  Hence
\[
 \int_X\beta_n\wedge(-V_p)\,\dprimeddprime u_1\wedge\cdots
   \wedge\dprimeddprime u_n
 \ \leq\ C(n,3)\,\|u_1\|_U\cdots\|u_n\|_U\int_U(-V_p)\beta_0 .
\]
Now let $p\to\infty$.  On the left, $-V_p\uparrow-V$ and monotone
convergence applies to the positive measure
$\beta_n\wedge\dprimeddprime u_1\wedge\cdots\wedge\dprimeddprime u_n$; on
the right, monotone convergence applies to the positive measure
$\beta_0$, and the limit $\int_U(-V)\beta_0$ is finite because $V$ is
locally integrable and $\beta_0$ has bounded density.  Since $\beta_n$ is
continuous and strictly positive on $U$, we conclude that $-V$ is locally
integrable on $U$ for $\dprimeddprime u_1\wedge\cdots\wedge
\dprimeddprime u_n$.  The general $m$ and general smooth $\alpha$ follow by
comparing $\alpha$ with $\pm c(\dprimeddprime\|t\|^2)^{n-m}$ in a fixed
affine chart, as in \cite[Corollary~8.12.2]{CLD25}.
\end{proof}

The goodness and local approachability assumptions are retained here.  The
theorem is not a blanket assertion that every pluripolar set is negligible
for every possible non-Archimedean Monge--Amp\`ere construction.

\begin{warning}{Negligibility of polar sets is not extension across them}
\Cref{thm:cld-polar-negligibility} says that a \emph{given} product of
continuous psh potentials assigns no mass to the polar locus of an auxiliary
psh function.  It says nothing about defining a product of potentials that
are themselves unbounded, and it does not assert that a product defined off
a polar set extends canonically across it.
\Cref{ex:ch16-polar-failure} shows that the second assertion is false as
stated: there, the product computed off $Y$ is zero, while the canonical
extension supplied by the Bochner--Martinelli theorem has mass one on $Y$.
The two statements are compatible because the potentials in the example are
not continuous on $X$.
\end{warning}

\section{Local theory versus global finite energy}
\label{sec:ch16-local-vs-global}

\Cref{part:bfj} and \cref{part:bj} construct Monge--Amp\`ere measures for
unbounded potentials.  Those constructions and the one of this chapter share
notation and a name, and they are not the same operation.  Because several
later theorems consist precisely in identifying them on a restricted locus,
this section states the delimitation as a table rather than as a slogan.

\subsection{Three levels}
\label{sec:ch16-three-levels}

It is useful to end by separating three levels.
\begin{enumerate}
\item The present CLD product is local and differential.  It treats
continuous psh potentials and gives local measures, CLN estimates, and the
polar theorem above.
\item Model and PL metrics on a proper variety have global intersection
masses.  Delta-forms make their polyhedral curvature products more explicit;
\Cref{chap:delta-forms} develops that comparison.
\item The BFJ and BJ classes $\mathcal E^1$ contain unbounded global
potentials.  Their Monge--Amp\`ere measures are defined by energy-controlled
approximation and non-pluripolar extension, not by declaring the local CLD
continuous product to apply verbatim.
\end{enumerate}
Boundary is equally decisive.  Stokes on an open domain produces boundary
terms; on a proper boundaryless space those terms vanish and total masses can
be identified with intersection numbers.  Consequently, a local formula
cannot by itself prove a global mass or comparison principle.

\subsection{The delimitation table}
\label{sec:ch16-delimitation-table}

\begin{center}
\small
\begin{tabular}{p{.14\textwidth}|p{.19\textwidth}|p{.19\textwidth}|p{.16\textwidth}|p{.18\textwidth}}
\textbf{Feature}&\textbf{CLD product}
 &\textbf{BFJ $\cE^1$ product}&\textbf{BJ $\cE^1$ product}
 &\textbf{Where proved}\\ \hline
ambient space
 &tropical or analytic, boundaryless, any dimension
 &proper $X$ over a nontrivially valued $K$, ample $L$
 &proper $X$ over a trivially valued $K$
 &\Cref{sec:ch16-hypotheses}; \Cref{chap:bfj-ma}; \Cref{chap:bj-energy}\\
potentials
 &continuous psh, locally defined
 &global $\theta$-psh with finite energy, possibly unbounded
 &global psh with finite energy, homogeneous
 &\Cref{sec:ch16-four-properties}; \Cref{def:e1}; \Cref{def:bj-e1}\\
definition mechanism
 &recursion \eqref{eq:ch16-recursion-schema} on potentials
 &monotone extension from bounded potentials, then
  non-pluripolar truncation
 &monotone extension and energy limits
 &\Cref{sec:ch16-recursion}; \Cref{thm:bounded-ma-extension},
  \Cref{def:nonpluripolar-ma}; \Cref{thm:bj-mixed-ma}\\
mass control
 &local, by CLN
 &global, total mass $(L^n)$
 &global, total mass fixed by the class
 &\Cref{thm:cld-cln}; \Cref{chap:bfj-energy}; \Cref{chap:bj-energy}\\
polar sets
 &not charged, under local approachability
 &not charged, by construction
 &not charged, by construction
 &\Cref{thm:cld-polar-negligibility};
  \Cref{prop:bounded-ma-no-polar-mass}; \Cref{thm:bj-quasicontinuity}\\
continuity
 &local uniform convergence of potentials
 &monotone and strong-topology convergence
 &monotone and strong-topology convergence
 &\Cref{thm:cld-bt-continuity}; \Cref{prop:e1-monotone-continuity};
  \Cref{chap:bj-strong-topologies}\\
comparison principle
 &chartwise only
 &available, with energy hypotheses
 &available, with energy hypotheses
 &\Cref{thm:ch16-comparison-chart};
  \eqref{eq:model-maximum-principle}; \eqref{eq:bj-comparison-principle}\\
\end{tabular}
\end{center}

\subsection{What transfers, and what does not}
\label{sec:ch16-transfer}

Three transfers are proved in this book, and each is stated with its
hypotheses.
\begin{enumerate}
\item\label{itm:ch16-transfer:model}
 For semipositive model metrics, the CLD product of the associated formal
 Chern currents, the Chambert--Loir measure, and the BFJ intersection measure
 all agree; and the agreement persists for continuous semipositive metrics
 obtained by uniform approximation.  This is
 \Cref{thm:bfj-cld-ma-common-locus}.
\item\label{itm:ch16-transfer:toric}
 In the toric case both operators equal $n!$ times the real Alexandrov
 measure of the associated convex function, by
 \Cref{thm:toric-ma-comparison} and
 \Cref{ex:ch16-toric-normalization}.
\item\label{itm:ch16-transfer:curve}
 Let $K$ be algebraically closed, complete, and nontrivially valued, and
 let $C$ be a smooth proper algebraic curve over $K$.  On $C^{\An}$,
 Wanner identifies positivity of the CLD current of a continuous potential
 with Thuillier subharmonicity; for model or skeletal potentials the
 curvature is the metric-graph Laplacian, and every continuous current-psh
 potential lies locally in the first ordinary-smooth decreasing-limit
 stage.  This is
 \Cref{thm:wanner-curve-comparison,thm:p063-current-psh-ordinary-regularization}
 together with
 \Cref{prop:skeletal-curvature-curve}.
\end{enumerate}
Two failures of transfer must be kept in view.  First, the CLD product is
local: a bounded global BFJ-psh function need not be locally CLD-psh, and no
theorem in this book asserts it; this is the local-versus-global metric
question \Cref{prob:frontier-cld-global-metric}.  Second, the two
constructions differ on the mass carried by pluripolar sets:
\Cref{ex:ch16-polar-failure} exhibits a canonical extension of a CLD product
with mass one on a polar set, whereas the non-pluripolar products of
\Cref{chap:bfj-energy,chap:bj-finite-energy} discard that mass by
definition.  Conflating them would produce a wrong total mass in
\Cref{chap:compare-ma}.

Finally, harmonic ambiguity is resolved differently at the three levels.
For CLD Bedford--Taylor currents it is encoded by the defining independence
from a potential with fixed $\dprimeddprime$.  For harmonic tropicalizations
it is encoded by cancellation of slopes across calibrated faces.  For the
global energy theories constants, and sometimes model-harmonic functions,
are removed by a normalization.  The comparison chapters will state the
hypotheses under which these quotients agree.

\section{Chapter dictionary}
\label{sec:ch16-dictionary}

\subsection{From the objects of this chapter to their use}
\label{sec:ch16-dict-forward}

\begin{center}
\small
\begin{tabular}{p{.28\textwidth}|p{.34\textwidth}|p{.30\textwidth}}
\textbf{Object}&\textbf{Established property}&\textbf{First use}\\ \hline
the product $uT$
 &defined for $T$ positive, $u$ continuous
  (\eqref{eq:ch16-uT-definition})
 &\Cref{thm:one-bt-factor}\\
sheaf-closure principle
 &propagates a property from smooth psh to all continuous psh
  (\Cref{lem:ch16-cpsh-induction})
 &\Cref{thm:one-bt-factor,thm:bt-symmetry-locality}\\
one-factor theorem
 &$\dprimeddprime(uT)$ positive closed
  (\Cref{thm:one-bt-factor})
 &the recursion, \Cref{sec:ch16-recursion}\\
Bedford--Taylor current
 &independence from the potential
  (\Cref{def:ch16-bt-current})
 &\Cref{thm:bt-harmonic-compatibility}\\
mixed product
 &symmetric, local, positive closed
  (\Cref{thm:bt-symmetry-locality})
 &\Cref{chap:cld-metrized-bundles}\\
dpsh extension
 &unique multilinear extension
  (\Cref{thm:ch16-dpsh-product})
 &\Cref{thm:cld-fibre-projection}\\
continuity theorem
 &joint weak continuity
  (\Cref{thm:cld-bt-continuity})
 &\Cref{thm:toric-ma-comparison},
  \Cref{thm:cld-polar-negligibility}\\
CLN inequality
 &local mass $\leq$ product of sup-norms
  (\Cref{thm:cld-cln})
 &\Cref{prop:positive-current-compactness}\\
weighted CLN inequality
 &constant \eqref{eq:ch16-cln-constant}
 &\Cref{thm:cld-polar-negligibility}\\
skeletal product formula
 &weights \eqref{eq:ch16-skeletal-weights}
 &\Cref{chap:cld-examples}, \Cref{chap:delta-forms}\\
harmonic compatibility
 &recursion equals the naive $G$-form product
  (\Cref{thm:bt-harmonic-compatibility})
 &\Cref{sec:harmonic-trop},
  \Cref{thm:abelian-harmonic-tropicalization}\\
polar negligibility
 &no mass on $\{V=-\infty\}$
  (\Cref{thm:cld-polar-negligibility})
 &\Cref{prop:ch16-degree-formula},
  \Cref{chap:compare-ma}\\
\end{tabular}
\end{center}

\subsection{From the later theory back to this chapter}
\label{sec:ch16-dict-backward}

\begin{center}
\small
\begin{tabular}{p{.30\textwidth}|p{.32\textwidth}|p{.30\textwidth}}
\textbf{Later object}&\textbf{Smallest instance}&\textbf{Computed in}\\
\hline
Monge--Amp\`ere measure of a toric metric
 &$\max\{0,x,y\}$ on the fan of $\mathbf P^2$, mass $1$
 &\Cref{ex:ch16-fan-p2}\\
mixed Monge--Amp\`ere measure
 &$\dprimeddprime\max\{0,x\}\wedge\dprimeddprime\max\{0,y\}
   =\delta_0$
 &\Cref{ex:ch16-fan-p2}\\
product on a nonlinear tropical space
 &$L\times\mathbf R$, mass $4$
 &\Cref{ex:ch16-fan-product}\\
Chambert--Loir measure on a curve
 &tripod, atoms $6,-1,-2,-3$
 &\Cref{ex:ch16-tripod-numbers}\\
failure of extension across a polar set
 &$\log\max\{|z_1|,|z_2|\}$ on
  $\mathbf A^{2,\An}$
 &\Cref{ex:ch16-polar-failure}\\
harmonic ambiguity
 &cross-fan, defect $2\delta_0$
 &\Cref{ex:ch16-cross-fan}\\
CLN constant in practice
 &$\MA(u)([-1,1])\leq2\operatorname{osc}_{[-2,2]}(u)$
 &\Cref{ex:ch16-cln-verified}\\
\end{tabular}
\end{center}

\subsection{Hypothesis ledger}
\label{sec:ch16-ledger}

\begin{center}
\small
\begin{tabular}{p{.30\textwidth}|p{.62\textwidth}}
\textbf{Hypothesis}&\textbf{Where it is used, and what fails without it}\\
\hline
$T$ positive
 &defines $uT$ and gives positivity of the product; without it the
  recursion has no sign and \Cref{lem:continuous-times-positive-current}
  fails\\
$T$ closed
 &\Cref{lem:ch16-leibniz-closed}; otherwise
  $\dprimeddprime(uT)$ acquires $\dprime u\wedge\ddprime T$ and
  $\ddprime u\wedge\dprime T$\\
$u$ continuous
 &\eqref{eq:ch16-uT-definition}; a locally bounded discontinuous psh
  function is not admitted, see the warning in
  \Cref{sec:ch16-four-properties}\\
$X$ boundaryless
 &Green's formula in \Cref{thm:cld-cln} and
  \Cref{thm:ch16-weighted-cln}; with boundary, a flux term survives\\
$T$ Bedford--Taylor
 &independence of the potential; \Cref{ex:ch16-cross-fan} shows the
  conclusion genuinely fails otherwise\\
harmonic chart
 &\Cref{thm:bt-harmonic-compatibility}; without harmonicity the boundary
  current $[\alpha]^\partial$ need not vanish\\
locally approachable
 &\Cref{thm:cld-polar-negligibility}; the proof needs decreasing
  \emph{smooth} approximants and stage-one collapse is open in general\\
$X$ good
 &\Cref{thm:ch16-barrier}; no smooth psh exhaustion is available in
  general\\
\end{tabular}
\end{center}

\section*{Source notes}

\paragraph{The construction.}
The recursive product, the notion of Bedford--Taylor current, symmetry,
local uniform continuity, the projection formula, the Chern--Levine--%
Nirenberg estimates, the skeletal product formulae, harmonic-chart
compatibility, and the polar theorem all follow
\cite[\S\S8.1--8.12 and \S16.9]{CLD25}.  Our numbering of the imported
statements records the 2025 revision: the local results are
\cite[\S\S8.4--8.12]{CLD25} and their analytic transcriptions are
\cite[\S16.9]{CLD25}, which is why \Cref{thm:cld-polar-negligibility} is
cited from the latter while \Cref{thm:ch16-weighted-cln} is cited from the
former.  The reader comparing with the source will find the same theorems in
a different order: we have isolated the sheaf-closure principle
(\Cref{lem:ch16-cpsh-induction}) as a lemma, whereas the source repeats the
argument at each application; we have made the constant of
\Cref{thm:cld-cln} explicit as a mass of $T$
(\Cref{prop:ch16-cln-constant}), which the source does not do; and we have
placed the harmonic-ambiguity material in one section rather than
distributing it between \S8.6 and \S8.11.

\paragraph{The complex prototype.}
The recursion \eqref{eq:ch16-recursion-schema} is Bedford and Taylor's
\cite{BT76,BT82}, where it is carried out for locally bounded psh functions
on a domain in $\mathbf C^n$; the convergence theorem for decreasing
sequences and the capacity that makes it sharp are in \cite{BT82,BT87}.
Chern, Levine and Nirenberg's inequality \cite{CLN69} predates the product
and was proved for smooth potentials; its role as the mass estimate behind
compactness is due to Bedford and Taylor.  Standard modern accounts of the
complex theory are \cite[Chapter~3]{Kl91}, \cite[Chapter~III]{DemBook} and
\cite[Chapters~3 and~4]{GZ17}, and the reader is warned that all of them
prove statements---continuity along decreasing sequences of locally bounded
potentials, quasicontinuity, the unrestricted comparison principle---whose
non-Archimedean analogues are \emph{not} available here, for the reason
displayed in \Cref{sec:ch16-four-properties}: the local smoothing that makes
the complex arguments work has no counterpart on a Berkovich space.  The
one-variable prototype of the whole chapter is Lelong's theory of
plurisubharmonic functions and positive currents \cite{Lel68}.

\paragraph{The psh hierarchy.}
\Cref{prop:ch16-hierarchy-status} and the warning following it correct a
frequently repeated misreading.  Eventual stationarity of the filtration is
proved by a cardinality argument in \cite[Lemma~8.1.4]{CLD25} and is used in
the proofs of \cite[Propositions~8.1.5, 8.1.11 and 8.2.1]{CLD25}; the
literal sentence of \cite[Remark~8.1.6]{CLD25} denying it is superseded by
those authors' own lemma, and the same remark records the intended content,
namely that no known example lies outside the first stage.  CLD's
\cite[Definition~8.1.13]{CLD25} names that stage.  The open question is
therefore early collapse, and it is stated in that form in
\Cref{prob:frontier-cld-hierarchy}.  \Cref{rem:ch16-collapse-reduction} and
\Cref{rem:ch16-collapse-euclidean} are elementary observations included for
orientation only: the first characterizes collapse at each fixed stage by
one closure property but deliberately leaves intermediate finite stages
open; the second is the standard radial mollification argument, whose
ingredients are proved in
\Cref{prop:ch09-radial-mollification}.  Neither is a theorem of this book,
and no priority is claimed for either.

The finite real-affine-star theorem
\Cref{thm:ch16-collapse-finite-tropical-curves} is instead a book-local
first-draft result.  Its exact frozen statement and proof passed the direct
production Rethlas endpoint under the local archive label
\nolinkurl{20260827T103835Z_c8131b22d800}, using
\texttt{gpt-5.6-sol} at \texttt{xhigh}, with final verdict
\texttt{correct}, zero critical errors, and zero gaps.  The endpoint assigned
no run or item identifier, and none is inferred.  The precise distinction
between this zero-sedentarity ambient-induced model and other tropical
boundary structures is \Cref{rem:ch16-collapse-scope-firewall}.

\paragraph{Harmonic tropicalizations and the curve tests.}
The harmonic-chart theory used in
\Cref{sec:ch16-harmonic-charts} is \cite[\S7.8]{CLD25}; its enlargement, in
which harmonic tropicalizations become charts in their own right and the
resulting Dolbeault theory is compared with the CLD one, is
\cite{GJR26}, taken up in \Cref{sec:harmonic-trop}.  Curve-test psh and
pluriharmonic functions, which are a different notion and are never
identified with the CLD sheaf without a theorem, are
\cite{GR25Classical,GR25PL}; the comparison with semipositive PL functions
is \Cref{thm:gr-pl-comparison}.  The delta-form calculus of
\cite{GK17,GK19}, in which polyhedral curvature and tropical intersection
become part of the coefficient algebra, is developed in
\Cref{chap:delta-forms}; \Cref{thm:ch16-skeletal-product} is the
Bedford--Taylor shadow of that calculus, and
\Cref{rem:ch16-skeletal-gsmooth} is exactly the term that delta-forms are
designed to absorb.

\paragraph{Harmonic ambiguity.}
The source-localization lemma and canonical-collar bridge
\Cref{lem:p21p3-gf-legal-source-localization,%
prop:p21p3-gf-virtual-annulus-bridge}, the rigid-point zero-atom estimates
\Cref{lem:p21p3-gfc-rigid-zero-atom,%
lem:p21p3-componentwise-normal-rigid-zero-atom}, and the classifications
\Cref{prop:p21p3-gfc-smooth-rigidity,prop:p21p3-curve-rigidity,%
prop:p21p3-componentwise-rigid-exceptions} are
book-level deductions from the cited CLD, Ducros, and
Gubler--Jell--Rabinoff inputs; no novelty or priority claim is made.  They
classify positive closed currents of bidegree $(0,0)$ on the stated
one-dimensional analytic spaces over an arbitrary complete nontrivially
valued ground field and prove Bedford--Taylor independence for continuous
psh pairs.

The exact live bridge and reduced source files have respective SHA-256
prefixes \texttt{f3202ef13495}, \texttt{ff8225ed6206},
\texttt{c9f27245f560}, and \texttt{01ab07fdaced}.  Together they passed
production Rethlas run
\nolinkurl{20260828T042906Z_f2a1c1b9f2a3}, using
\texttt{gpt-5.6-sol} at \texttt{xhigh}, with final verdict
\texttt{correct}, zero critical errors, zero gaps, seventeen sequential
statement records, nineteen primary-source records, and no failed check.

The normal-curve rigid-exception extension
\Cref{prop:p21p3-rigid-exception-extension} is a further book-level
deduction.  Its exact statement and proof have SHA-256 prefixes
\texttt{cf53c4101b6f} and \texttt{cf94877ec4b2}.  Submitted after the exact
rig-smooth dependency, the materialized statement and proof have prefixes
\texttt{cc0c48725114} and \texttt{fd794304c18d}; they passed production
Rethlas run \nolinkurl{20260829T021315Z_cc0c48725114}, using
\texttt{gpt-5.6-sol} at \texttt{xhigh}, with final verdict
\texttt{correct}, zero critical errors, zero gaps, fourteen sequential
statement records, twenty-six primary-source records, and no failed check.
The proof uses graph connectedness, unibranchness of rigid points on a
normal curve, and the vanishing of smooth $(1,1)$-test germs at such points;
it does not promote the result to arbitrary non-rig-smooth normalizations.

The nilreduction and nonreduced extension in
\Cref{prop:p21p3-curve-rigidity-nonreduced} form a further book-level
deduction.  The exact live statement and proof have SHA-256 prefixes
\texttt{e7b9889f6c06} and \texttt{25fe49a46b1c}.  Installed after the exact
bridge and reduced prefix, these bytes passed the HTTP-archived production
Rethlas closure run
\nolinkurl{20260828T051601Z_db92555a72e1}, again using
\texttt{gpt-5.6-sol} at \texttt{xhigh}, with final verdict
\texttt{correct}, zero critical errors, zero gaps, twenty-three sequential
statement records, twenty-five primary-source records, and no failed check.
The independently checked combined nonreduced packet, with statement/proof
SHA-256 prefixes \texttt{dfdaa04a1cea} and \texttt{f8eb8968f9ae}, had
already passed run \nolinkurl{20260828T044008Z_dfdaa04a1cea} with the same
verdict and zero errors or gaps.  The earlier service-only run
\nolinkurl{20260828T045515Z_db92555a72e1} is corroborating provenance, not
the controlling gate.  The request, response, safe-header archive, and
persistent service evidence retain the full hashes and verification records.

The local normal rigid-point estimate
\Cref{lem:p21p3-componentwise-normal-rigid-zero-atom} has exact statement
and proof SHA-256 prefixes \texttt{469024d04e8e} and
\texttt{661100a4b25e}.  It passed production Rethlas run
\nolinkurl{20260829T060048Z_469024d04e8e}, item
\nolinkurl{pi_0d4235a44db2b33950f8dd55}, using
\texttt{gpt-5.6-sol} at \texttt{xhigh}, with final verdict
\texttt{correct}, zero critical errors, and zero gaps.  The componentwise
extension \Cref{prop:p21p3-componentwise-rigid-exceptions} has exact
statement and proof SHA-256 prefixes \texttt{5b03b18f2556} and
\texttt{3893a89ffd59}.  Its complete dependency materialization has statement
and proof prefixes \texttt{2bf673687e93} and \texttt{a27c40fbb132}; it
passed production run \nolinkurl{20260829T061505Z_2bf673687e93}, item
\nolinkurl{pi_ee24d64b88f0c4e60ad1acbd}, with the same model and scrutiny,
final \texttt{correct}, zero critical errors, and zero gaps.

An earlier materialization, production run
\nolinkurl{20260829T060846Z_9e819244c341}, item
\nolinkurl{pi_5c78fec3148f11d4fd5a15cb}, returned final
\texttt{wrong}, one critical error, and zero gaps because it omitted the
transitive canonical-collar bridge required by the imported reduced proof.
That noncontrolling packaging failure is retained in the archive; the exact
componentwise theorem bytes were not changed to conceal it.  The successful
run above supplied the missing exact dependency closure.

\Cref{rem:p21p3-curve-scope,rem:p21p3-rigid-exception-scope,%
rem:p21p3-nr-scope,rem:p213-normal-rigid-zero-atom-scope,%
rem:p21p3-componentwise-residual} record the exact scope.
In particular, no algebraic-closedness hypothesis is imposed, but the
ground field remains nontrivially valued and the goodness, separatedness,
strictness, boundarylessness, and dimension hypotheses are not relaxed.  In
the broadest one-dimensional result, each normalized reduced irreducible
component must be rig-smooth away from a locally finite set of rigid points.
Every Bedford--Taylor conclusion here concerns continuous psh pairs; the
literal possibly discontinuous predicate of
\cite[Definition~8.6.2]{CLD25} is not asserted.

\Cref{thm:ch16-pluriharmonic-criterion},
\Cref{prop:ch16-cell-rigidity} and \Cref{ex:ch16-cross-fan} are the
audited statements \Cref{thm:verified-cld-criterion},
\Cref{prop:verified-cell-rigidity} and \Cref{prop:verified-cross-fan} of
\Cref{chap:verified-extensions}, reproved here because this is where they
are used.  The cross-fan is a counterexample on an abstract calibrated fan;
the weighted configuration and pushed defect are realized analytically in
\Cref{prop:verified-cross-analytic-realization}, while
\Cref{rem:verified-cross-upstairs-firewall} isolates the stronger
upstairs-equality question that remains open.  The
distinction between the three harmonicity conditions tabulated in
\Cref{sec:ch16-reading-bt} follows \cite[Remark~7.8.3]{CLD25} for the first
two and \cite{GR25Classical} for the third.

\paragraph{The convex and real-geometric side.}
The identification of the top self-product of a convex potential with $n!$
times the Alexandrov measure is fixed in
\Cref{prop:ch09-ma-coordinates} and is used in
\Cref{sec:ch16-comparison} and \Cref{sec:ch16-worked}; the corresponding
supercurrent formalism, in which convex Bedford--Taylor products and mixed
volumes are developed intrinsically, is \cite{Ber26Convex}, with the
convention firewall recorded in \Cref{chap:cld-currents}.  The comparison of
real and non-Archimedean Monge--Amp\`ere operators in the toric case is
\cite{Vil20}, and the global invariant tropical version is
\cite{BGJKK25}.  Villena's comparison and the toric normalization of
\Cref{chap:model-calculations} agree, and
\Cref{ex:ch16-toric-normalization} records the factor.

\paragraph{What is new here, and what is not.}
The canonical-collar source-localization and branch-separation argument, the
rigid-point zero-atom estimates, the arbitrary-field dimension-one
classification and continuous-pair Bedford--Taylor consequence of
\Cref{prop:p21p3-curve-rigidity}, and its nilpotent-thickening extension
\Cref{prop:p21p3-curve-rigidity-nonreduced}, together with the normal-curve
rigid-exception extension
\Cref{prop:p21p3-rigid-exception-extension} and the reducible and nonreduced
componentwise extension
\Cref{prop:p21p3-componentwise-rigid-exceptions}, are book-level first-draft
deductions assembled from the cited inputs; no novelty or priority claim is
made.  The finite real-affine-star collapse of
\Cref{thm:ch16-collapse-finite-tropical-curves} is likewise a first-draft
result.  The remaining theorem families in this chapter come from the cited
sources.
The organizational
choices that differ from CLD are the four listed above: the sheaf-closure
principle as a named lemma, the explicit CLN constant, the single
harmonic-ambiguity section, and the delimitation table of
\Cref{sec:ch16-delimitation-table}, which exists because
\Cref{part:bfj,part:bj,part:comparisons} depend on not conflating the local
product with the global finite-energy ones.  The worked computations of
\Cref{sec:ch16-worked} are elementary but are carried to numbers on purpose:
every mass in this book that is later compared with an intersection number
is first checked in one of them.
